\beginnumbering
\eledsection*{Adhyāya 1}
\pstart \ledsidenote{\passref{2.1.1}}\skt{अथातो वात\-व्या\-धि\-नि\-दानं व्याख्यास्यामः ॥}\pend

\pstart\versepar \ledsidenote{\passref{2.1.3}}\edtext{\skt{धन्व\-न्त\-रि\-न्ध\-र्म्म\-भृ\-ता\-म्व\-रि\-ष्ठ\-म\-मृ\-तो\-द्भवं}}{\lemma{\skt{धन्वन्तरिन्}}\Afootnote{A H; \skt{धन्वन्तरिर्}\ N}}\skt{ । }\pend
\pstart\versepar \skt{चर\-णा\-वु\-प\-सं\-गृह्य सुश्रुतः परिपृच्छति ॥}\pend

\pstart\versepar \ledsidenote{\passref{2.1.4}}\skt{वायोः प्रकृतिभूतस्य व्यापन्नस्य च }\edtext{\skt{भूपते}}{\Afootnote{N H; \skt{कोपनैः ।}\ A}}\skt{ । }\pend
\pstart\versepar \skt{स्थानं }\edtext{\skt{रोगविभागञ्च}}{\Afootnote{N H; \skt{कर्म च रोगांश्च}\ A}}\skt{ }\edtext{\skt{वदस्व}}{\Afootnote{N H; \skt{वद मे}\ A}}\skt{ }\edtext{\skt{वदताम्वर}}{\lemma{\skt{वर}}\Afootnote{A; \skt{वरः ॥}\ N H}}\skt{ ॥}\pend

\pstart\versepar \ledsidenote{\passref{2.1.5}}\skt{तस्य तद्वचनं }\edtext{\skt{श्रुत्वा}}{\Afootnote{N H; \skt{प्रा}\ A}}\skt{ }\edtext{\skt{जगाद}}{\Afootnote{N H; \skt{ब्रवीद्भि}\ A}}\skt{ }\edtext{\skt{भगवानृषिः}}{\Afootnote{H; \skt{षजं वरः ।}\ A; \skt{भगवांनृषिः ।}\ N}}\skt{ । }\pend
\pstart\versepar \skt{स्वयंभूरेष }\edtext{\skt{भग\-वा\-न्वा\-यु\-रि\-त्य\-भि\-श\-ब्दितः}}{\lemma{\skt{भगवान्वायुर्}}\Afootnote{A H; \skt{भगवां वायुर्}\ N}}\skt{ ॥}\pend

\pstart\versepar \ledsidenote{\passref{2.1.6}}\edtext{\skt{स्वात\-न्त्र्या\-न्नि\-त्य\-भा\-वाच्च}}{\lemma{\skt{नित्यभावाच्च}}\Afootnote{A N; \skt{नित्यभवाच्च}\ H}}\skt{ सर्व्व\-ग\-त्वा\-त्त\-थैव च । }\pend
\pstart\versepar \skt{सर्वेषामेव सर्वात्मा सर्व\-लो\-क\-न\-म\-स्कृतः ॥}\pend

\pstart\versepar \ledsidenote{\passref{2.1.7}}\skt{स्थित्यु\-त्प\-त्ति\-वि\-ना\-शेषु भूतानामेष कारणं । }\pend
\pstart\versepar \skt{अव्यक्तो व्यक्तकर्म्मा च }\edtext{\skt{शीतो}}{\Afootnote{N H; \skt{रूक्षः}\ A}}\skt{ }\edtext{\skt{रूक्षो}}{\Afootnote{H; \skt{शीतो}\ A; \skt{रोक्षो}\ N}}\skt{ }\edtext{\skt{लघुश्चरः}}{\lemma{\skt{चरः}}\Afootnote{N H; \skt{खरः ॥}\ A}}\skt{ ॥}\pend

\pstart\versepar \ledsidenote{\passref{2.1.8}}\edtext{\skt{तिर्यग्गो}}{\Afootnote{A H; \skt{तियग्गो}\ N}}\skt{ द्विगुणश्चैव }\edtext{\skt{रजोबहुल}}{\Afootnote{A; \skt{रसवहुल}\ N; \skt{रजोवहुल}\ H}}\skt{ एव च । }\pend
\pstart\versepar \skt{अचिन्त्यवीर्यो दोषाणां नेता }\edtext{\skt{रोग\-स\-मू\-ह\-राट्}}{\Afootnote{A H; \skt{रोग\-स\-मू\-ह\-राट ॥}\ N}}\skt{ ॥}\pend

\pstart\versepar \ledsidenote{\passref{2.1.9}}\edtext{\skt{आशुचारी}}{\Afootnote{H; \skt{आशुकारी}\ A N\textsuperscript{pc}; \skt{आशुकीकारी}\ N\textsuperscript{ac}}}\skt{ }\edtext{\skt{मुहुश्चारी}}{\Afootnote{A N; \skt{मुहूश्चारी}\ H}}\skt{ पक्वा\-धा\-न\-गु\-दा\-लयः । }\pend
\pstart\versepar \skt{देहे }\edtext{\skt{विचरतस्तस्य}}{\lemma{\skt{विचरतस्}}\Afootnote{A N H\textsuperscript{pc}; \skt{विचरस्}\ H\textsuperscript{ac}}}\skt{ लक्षणानि }\edtext{\skt{निबोध}}{\Afootnote{A; \skt{निवोध}\ N H}}\skt{ मे ॥}\pend

\pstart\versepar \ledsidenote{\passref{2.1.10}}\edtext{\edtext{\skt{इन्द्रि\-या\-र्थो\-प\-स\-म्प\-त्ति\-र्दो\-ष\-धा\-त्व\-ग्नि\-सा\-म्यतां}}{\lemma{\skt{इन्द्रि\-या\-र्थो\-प\-स\-म्प\-त्तिर्}}\Afootnote{N H; \skt{इन्द्रि\-या\-र्थो\-प\-सं\-प्राप्तिं}\ A}}}{\lemma{\skt{दोष\-धा\-त्व\-ग्नि\-सा\-म्यतां}}\Afootnote{N H; \skt{दोष\-धा\-त्व\-ग्न्य\-वै\-कृ\-तम् ।}\ A}}\skt{ । }\pend
\pstart\versepar \skt{}\edtext{\skt{क्रिया\-णा\-मा\-नु\-लो\-म्यञ्च}}{\lemma{\skt{आनुलोम्यञ्च}}\Afootnote{A H; \skt{अनुलोम्यञ्च}\ N}}\skt{ कुर्या\-द्वा\-यु\-र\-दू\-षितः ॥}\pend

\pstart\versepar \ledsidenote{\passref{2.1.11}}\edtext{\skt{यथाग्नि}}{\Afootnote{H; \skt{यथाऽग्निः}\ A; \skt{यथाग्निः}\ N}}\skt{ पञ्चधा भिन्नो }\edtext{\skt{नाम\-स्था\-ना\-त्म\-क\-र्म्मभिः}}{\Afootnote{N H; \skt{नाम\-स्था\-न\-क्रि\-या\-मयैः ।}\ A}}\skt{ । }\pend
\pstart\versepar \skt{}\edtext{\edtext{\edtext{\skt{भिन्नोऽनिलस्तथा}}{\lemma{\skt{भिन्नो}}\Afootnote{A H; \skt{भिन्नो}\ N}}}{\lemma{\skt{ऽनिलस्}}\Afootnote{A H; \skt{निलस्}\ N}}\skt{ ह्येको}}{\lemma{\skt{तथा ह्य्}}\Afootnote{H; \skt{तस्था ह्य्}\ A N}}\skt{ }\edtext{\skt{नाम\-स्था\-न\-क्रि\-या\-मयैः}}{\Afootnote{A H; \skt{नाम\-स्था\-नो\-क्रि\-या\-मयैः ॥}\ N}}\skt{ ॥}\pend

\pstart\versepar \ledsidenote{\passref{2.1.12}}\edtext{\skt{प्राणोदानः}}{\Afootnote{N H; \skt{प्राणोदानौ}\ A}}\skt{ समानश्च }\edtext{\edtext{\edtext{\skt{व्यानोऽ\-पा\-न\-स्त\-थैव}}{\lemma{\skt{व्यानो}}\Afootnote{A; \skt{व्यानो}\ N H}}}{\lemma{\skt{ऽपानस्}}\Afootnote{\skt{चापा}\ A; \skt{पानस्}\ N H}}\skt{ च}}{\lemma{\skt{तथैव च}}\Afootnote{N H; \skt{न एव च ।}\ A}}\skt{ । }\pend
\pstart\versepar \skt{स्थानस्था }\edtext{\skt{मारुताः}}{\Afootnote{N H; \skt{मारुताः {\symfont `}}\ A}}\skt{ }\edtext{\skt{पञ्च}}{\Afootnote{N H; \emph{om.}\ A}}\skt{ यापयन्ति शरीरिणम् ॥}\pend

\pstart\versepar \ledsidenote{\passref{2.1.13}}\edtext{\skt{योऽनिलो}}{\Afootnote{H; \skt{यो वायुर्}\ A; \skt{यो निलो}\ N}}\skt{ वक्त्रसंचारी स प्राणो नाम देहधृक् । }\pend
\pstart\versepar \skt{}\edtext{\skt{सोऽन्न\-म्प्र\-वे\-श\-य\-त्यन्तः}}{\lemma{\skt{ऽन्नम्}}\Afootnote{A H; \skt{नम्}\ N}}\skt{ }\edtext{\skt{प्राणा\-श्चा\-प्य\-व\-ल\-म्बते}}{\lemma{\skt{अवलम्बते}}\Afootnote{A; \skt{अवलम्वते ॥}\ N H}}\skt{ ॥}\edtext{}{\lemma{\skt{देहधृक् । सो}}\Afootnote{A H; \skt{देहधृक् । सोन्}\ N}}\pend

\pstart\versepar \ledsidenote{\passref{2.1.14}}\skt{प्रायशः }\edtext{\skt{कुरुते}}{\Afootnote{A H; \skt{कुरुतेश्}\ N}}\skt{ }\edtext{\skt{चापि}}{\Afootnote{N H; \skt{दुष्टो}\ A}}\skt{ हिक्का\-श्वा\-सा\-दि\-का\-न्ग\-दान् । }\pend
\pstart\versepar \skt{उदानो }\edtext{\skt{नाम }\edtext{\edtext{\skt{योऽह्यूर्ध्वमुपैति}}{\lemma{\skt{ऽह्य्}}\Afootnote{\emph{om.}\ A; \skt{ह्य्}\ N; \skt{प्य्}\ H}}}{\lemma{\skt{ऊर्ध्वम्}}\Afootnote{N H; \skt{तूर्ध्वम्}\ A}}}{\lemma{\skt{नाम यो}}\Afootnote{A H; \skt{नाम योर्}\ N}}\skt{ पवनोत्तमः ॥}\edtext{}{\lemma{\skt{हिक्का\-श्वा\-सा\-दि\-का\-न्ग\-दान् । उदानो}}\Afootnote{A H; \skt{हिक्काश्वासादिकां गदान् । उदानो}\ N}}\pend

\pstart\versepar \ledsidenote{\passref{2.1.15}}\skt{तेन }\edtext{\skt{भाषि\-त\-गी\-ता\-दि\-वि\-शे\-षोऽ\-भि\-प्र\-व\-र्त्तते}}{\lemma{\skt{भाषि\-त\-गी\-ता\-दि\-वि\-शेषो}}\Afootnote{A N; \skt{भाषि\-त\-गी\-ता\-दि\-र्व्वि\-शेषो}\ H}}\skt{ । }\pend
\pstart\versepar \skt{}\edtext{\skt{ऊर्ध्व\-ज\-त्रु\-ग\-ता\-न्रो\-गा\-न्क\-रोति च}}{\Afootnote{A H; \skt{ऊर्ध्व\-ज\-त्रु\-ग\-ता\-न्रो\-गा\-ङ्क\-रोति च}\ N}}\skt{ विशेषतः ॥}\pend

\pstart\versepar \ledsidenote{\passref{2.1.16}}\edtext{\skt{आम\-प\-क्वा\-श\-य\-चरः}}{\Afootnote{A N; \skt{आम\-प\-क्वा\-स\-य\-चरः}\ H}}\skt{ समा\-नोऽ\-ग्नि\-स\-हा\-य\-वान् । }\pend
\pstart\versepar \skt{अन्नम्पचति तज्जाश्च }\edtext{\skt{भावा\-न्प्र\-वि\-वि\-नक्ति}}{\Afootnote{H; \skt{विशे\-षा\-न्वि\-वि\-नक्ति}\ A; \skt{भावा\-म्प्र\-वि\-वि\-नक्ति}\ N}}\skt{ }\edtext{\skt{सः}}{\Afootnote{N H; \skt{हि ॥}\ A}}\skt{ ॥}\edtext{}{\lemma{\skt{ऽग्नि\-स\-हा\-य\-वान् । अन्नम्}}\Afootnote{H; \skt{वह्निसङ्गतः । सोऽन्नं}\ A; \skt{ग्निसहायवान् ॥ अन्नम्}\ N}}\pend

\pstart\versepar \ledsidenote{\passref{2.1.17}}\edtext{\skt{गुल्मा\-ग्नि\-सं\-गा\-ती\-सा\-रा\-न्प्रा\-य\-शश्च}}{\Afootnote{N H; \skt{गुल्मा\-ग्नि\-सा\-दा\-ती\-सा\-र\-प्र\-भृती}\ A}}\skt{ }\edtext{\skt{करोति हि}}{\Afootnote{N H; \skt{न्कुरुते ग}\ A}}\skt{ ॥ }\pend
\pstart\versepar \skt{}\edtext{\skt{कृत्स्नदेहचरो}}{\Afootnote{N H\textsuperscript{pc}; \skt{दान् । कृत्स्नदेहचरो}\ A; \skt{कृत्स्न\-दे\-ह\-दे\-ह\-चरो}\ H\textsuperscript{ac}}}\skt{ व्यानो }\edtext{\skt{रस\-स\-म्व\-ह\-नो\-द्यतः}}{\Afootnote{A N; \skt{रस\-स\-म्वा\-ह\-नो\-द्यतः ।}\ H}}\skt{ । ॥}\pend

\pstart\versepar \ledsidenote{\passref{2.1.18}}\skt{स्वेदा\-सृ\-क्स्रा\-व\-ण\-श्चापि }\edtext{\skt{सर्व्वथा}}{\Afootnote{N H; \skt{पञ्चधा}\ A}}\skt{ चेष्टयत्यपि ॥ }\pend
\pstart\versepar \skt{क्रुद्धश्च कुरुते रोगान्प्रायशः सर्वदेहगान् ॥}\pend

\pstart\versepar \ledsidenote{\passref{2.1.19}}\skt{पक्वा\-धा\-ना\-ल\-योऽ\-पानः काले }\edtext{\skt{कर्षति}}{\Afootnote{A N\textsuperscript{pc} H; \skt{कर्षयति}\ N\textsuperscript{ac}}}\skt{ चाप्यधः । }\pend
\pstart\versepar \skt{}\edtext{\skt{वात\-मू\-त्र\-पु\-री\-षाणि}}{\Afootnote{N H; \skt{समीरणः शकृन्मूत्रं}\ A}}\skt{ शुक्र\-ग\-र्भा\-र्त्त\-वानि च ॥}\pend

\pstart\versepar \ledsidenote{\passref{2.1.20}}\skt{क्रुद्धश्च कुरुते रोगा\-न्घो\-रा\-न्ब\-स्ति\-गु\-दा\-श्र\-यान् ॥ }\pend
\pstart\versepar \skt{शुक्रदोषाः }\edtext{\skt{प्रमेहाश्च}}{\Afootnote{N H; \skt{प्रमेहास्तु}\ A}}\skt{ }\edtext{\skt{व्याना\-पा\-न\-प्र\-दो\-षजाः}}{\Afootnote{H; \skt{व्याना\-पा\-न\-प्र\-को\-पजाः ॥}\ A; \skt{व्याना\-पा\-न\-प्र\-दो\-षजा॥}\ N}}\skt{ ।}\edtext{}{\lemma{\skt{रोगा\-न्घो\-रा\-न्ब\-स्ति\-गु\-दा\-श्र\-यान् ॥ शुक्रदोषाः}}\Afootnote{\skt{रोगा\-न्घो\-रा\-न्ब\-स्ति\-गु\-दा\-श्र\-यान् । शुक्रदोष}\ A; \skt{रोगां घोरां वस्ति\-गु\-दा\-श्र\-यान् । शुक्रदोषा}\ N; \skt{रोगान्घोरान्वस्ति गुदाश्रयान् ॥ शुक्रदोषाः}\ H}}\pend

\pstart \ledsidenote{\passref{2.1.21ab}}\skt{युग\-प\-त्कु\-पि\-ता\-श्चापि देहं भिन्द्यु\-र\-सं\-श\-यम् ॥}\pend

\pstart \ledsidenote{\passref{2.1.21ab.1}}\edtext{\skt{॥ ० ॥}}{\Afootnote{\skt{॥}\ N; \skt{॥}\ H}}\pend

\pstart \ledsidenote{\passref{2.1.21cd}}\edtext{\skt{अतः}}{\Afootnote{N H; \skt{अत}\ A}}\skt{ }\edtext{\skt{सर्वान्प्रवक्ष्यामि}}{\Afootnote{H; \skt{ऊर्ध्वं प्रवक्ष्यामि}\ A; \skt{सर्वां प्रवाक्ष्यामि}\ N}}\skt{ }\edtext{\skt{नाना\-स्था\-ना\-न्त\-रा\-श्रि\-तान्}}{\Afootnote{H; \skt{नाना\-स्था\-ना\-न्त\-रा\-श्रितः ॥}\ A; \skt{नाना\-स्थि\-ना\-न्त\-रा\-स्रितां ।}\ N\textsuperscript{ac}; \skt{नाना\-स्था\-ना\-न्त\-रा\-स्रितां ।}\ N\textsuperscript{pc}}}\skt{ ।}\pend

\pstart\versepar \ledsidenote{\passref{2.1.22}}\edtext{\skt{बहुशः}}{\Afootnote{A; \skt{बहुसः}\ N; \skt{वहुशः}\ H}}\skt{ कुपितो }\edtext{\skt{वायु\-र्व्वि\-का\-रा\-न्कु\-रुते हि}}{\lemma{\skt{व्विकारान्कुरुते हि}}\Afootnote{A H; \skt{विकारं कुरुते हि}\ N}}\skt{ यान् ॥ }\pend
\pstart\versepar \skt{}\edtext{\skt{वायुरामाशये}}{\lemma{\skt{आमाशये}}\Afootnote{A; \skt{आमासये}\ N H}}\skt{ क्रुद्धः }\edtext{\skt{कुर्याच्छर्द्यादिकां}}{\Afootnote{N; \skt{छर्द्यादीन्कुरुते}\ A; \skt{कुर्याच्छर्द्यादिका}\ H}}\skt{ }\edtext{\skt{गदान्}}{\Afootnote{A N; \skt{न्गदान् ।}\ H}}\skt{ ।}\pend

\pstart\versepar \ledsidenote{\passref{2.1.23}}\edtext{\skt{मोहं}}{\Afootnote{A H; \skt{मोहाम्}\ N}}\skt{ }\edtext{\edtext{\skt{मूर्च्छाम्पिपासां}}{\lemma{\skt{मूर्च्छाम्}}\Afootnote{A H; \skt{मूर्च्छा}\ N}}\skt{ च}}{\lemma{\skt{पिपासां च}}\Afootnote{A H; \skt{पिपासास्च}\ N}}\skt{ }\edtext{\skt{हृद्ग्र\-ह\-म्पा\-र्श्व\-वे\-दनां}}{\lemma{\skt{पार्श्ववेदनां}}\Afootnote{A H; \skt{पार्श्ववेदनं ॥}\ N}}\skt{ ॥ }\pend
\pstart\versepar \skt{}\edtext{\skt{पक्वा\-स\-य\-स्थोऽ\-न्त्र\-कूजं}}{\lemma{\skt{पक्वासयस्थो}}\Afootnote{N H; \skt{पक्वाशयस्थो}\ A}}\skt{ }\edtext{\skt{शूलाध्मानौ}}{\Afootnote{H; \skt{शूलं नाभौ}\ A; \skt{शूलां ध्मानो}\ N}}\skt{ करोति च ।}\pend

\pstart\versepar \ledsidenote{\passref{2.1.24}}\skt{कृच्छ्र\-मू\-त्र\-पु\-री\-ष\-त्व\-मा\-ना\-ह\-न्त्रि\-क\-वे\-दनां ॥ }\pend
\pstart\versepar \skt{}\edtext{\skt{श्रोत्रा\-दि\-ष्वि\-न्द्रि\-य\-वधं}}{\lemma{\skt{श्रोत्रादिष्व्}}\Afootnote{A; \skt{स्रोत्रादिष्व्}\ N H}}\skt{ कुर्यात्क्रुद्धः समीरणः ।}\pend

\pstart\versepar \ledsidenote{\passref{2.1.25abc}}\skt{वैवर्ण्ण्यं स्फुरणं }\edtext{\skt{रूक्षं}}{\Afootnote{N H; \skt{रौक्ष्यं}\ A}}\skt{ }\edtext{\skt{सुप्तिं}}{\Afootnote{A H; \skt{सुप्ति}\ N}}\skt{ चुमुचुमायनं ॥ }\pend
\pstart\versepar \skt{त्वक्स्थो निस्तोदनं }\edtext{\skt{ग्रन्थीन्}}{\Afootnote{H; \skt{कुर्यात्}\ A; \skt{ग्रन्थी}\ N}}\pend

\pstart\versepar \ledsidenote{\passref{2.1.26bcd}}\edtext{\skt{सरु\-जा\-न्मां\-स\-सं\-श्रितः}}{\Afootnote{\skt{सशू\-ला\-न्मां\-स\-सं\-श्रितः ।}\ A; \skt{संरु\-जा\-म्मां\-स\-सं\-स्रितः ।}\ N; \skt{सरु\-जा\-न्मां\-स\-सं\-श्रितः । …}\ H}}\skt{ । }\pend
\pstart\versepar \skt{तथा }\edtext{\skt{मेदाश्रितः}}{\Afootnote{H; \skt{मेदःश्रितः}\ A; \skt{मेदस्रितः}\ N}}\skt{ }\edtext{\edtext{\skt{कुर्या\-द्ग्र\-न्थी\-न्म\-न्द\-रु\-जोऽ\-व्र\-णान्}}{\lemma{\skt{कुर्या\-द्ग्र\-न्थी\-न्म\-न्द\-रुजो}}\Afootnote{A H; \skt{कुर्याद्ग्रन्थी मन्दरुजो}\ N}}}{\lemma{\skt{ऽव्रणान्}}\Afootnote{A; \skt{व्रणान्}\ N H}}\skt{ ॥}\pend

\pstart\versepar \ledsidenote{\passref{2.1.27}}\edtext{\skt{कुर्याद्सिरागतः}}{\Afootnote{N H; \skt{कुर्यात्सिरागतः}\ A}}\skt{ शूलं सिरा\-कु\-ञ्च\-न\-पू\-रणं । }\pend
\pstart\versepar \skt{स्नायुप्राप्तः }\edtext{\skt{स्नायुजालं}}{\Afootnote{N H; \emph{om.}\ A}}\skt{ }\edtext{\edtext{\edtext{\skt{स्तम्भ\-य\-त्या\-क्षि\-प\-त्यपि}}{\lemma{\skt{स्तम्भयत्य्}}\Afootnote{N H; \skt{स्तम्भकम्पौ शूलम्}\ A}}}{\lemma{\skt{आक्षिपत्य्}}\Afootnote{N H; \skt{आक्षेपणं}\ A}}}{\lemma{\skt{अपि}}\Afootnote{N H; \skt{तथा ॥}\ A}}\skt{ ॥}\pend

\pstart\versepar \ledsidenote{\passref{2.1.28}}\edtext{\skt{हन्ति}}{\Afootnote{A H; \skt{हस्ति}\ N}}\skt{ सन्धिगतः }\edtext{\skt{संधीन्शूलशोफौ}}{\Afootnote{A H; \skt{सन्धी शूलशोफौ}\ N}}\skt{ करोति च । }\pend
\pstart\versepar \skt{}\edtext{\skt{अस्थिभेदः}}{\Afootnote{H; \skt{अस्थिशोषं}\ A; \skt{अस्थिभेद}\ N}}\skt{ }\edtext{\skt{प्रशोषश्च}}{\Afootnote{N H; \skt{प्रभेदं च}\ A}}\skt{ }\edtext{\skt{कुर्याच्छूलञ्च}}{\Afootnote{A H; \skt{कुर्याच्छूलश्च}\ N}}\skt{ }\edtext{\skt{तत्कृतः}}{\Afootnote{N H; \skt{तच्छ्रितः ॥}\ A}}\skt{ ॥}\pend

\pstart\versepar \ledsidenote{\passref{2.1.29}}\skt{तथा मज्जगते रुक्च न कदा\-चि\-त्प्र\-शा\-म्यति । }\pend
\pstart\versepar \skt{}\edtext{\edtext{\skt{अप्र\-वृ\-त्ति\-म्प्र\-वृ\-त्तिम्वा}}{\lemma{\skt{अप्रवृत्तिम्}}\Afootnote{N H; \skt{अप्रवृत्तिः}\ A}}}{\lemma{\skt{प्रवृत्तिम्वा}}\Afootnote{N H; \skt{प्रवृत्तिर्वा}\ A}}\skt{ }\edtext{\skt{विकृतां}}{\Afootnote{N H; \skt{विकृता}\ A}}\skt{ }\edtext{\skt{शुक्रगोनिलः}}{\Afootnote{N H; \skt{शुक्रगेऽनिले ॥}\ A}}\skt{ ॥}\pend

\pstart\versepar \ledsidenote{\passref{2.1.30}}\edtext{\skt{हस्त\-पा\-द\-शि\-रो\-धा\-तू\-न्तथा}}{\lemma{\skt{हस्त\-पा\-द\-शि\-रो\-धा\-तून्}}\Afootnote{A H; \skt{हस्त\-पा\-द\-सि\-रो\-धातू}\ N}}\skt{ सञ्चरति क्रमात् । }\pend
\pstart\versepar \skt{व्याप्नुयाद्वाखिलं देहम्वायुः सर्व्वगतो नृणाम् ॥}\pend

\pstart\versepar \ledsidenote{\passref{2.1.31}}\edtext{\skt{स्तम्भ\-ना\-क्षे\-प\-ण\-स्वा\-प\-शो\-फ\-शू\-लांश्च}}{\Afootnote{H; \skt{स्तम्भ\-ना\-क्षे\-प\-ण\-स्वा\-प\-शो\-फ\-शू\-लानि}\ A; \skt{स्तम्भ\-णा\-क्षे\-प\-ण\-स्वा\-प\-शो\-फ\-शू\-लांश्च}\ N}}\skt{ }\edtext{\skt{सर्वशः}}{\Afootnote{N H; \skt{सर्वगः ।}\ A}}\skt{ । }\pend
\pstart\versepar \skt{}\edtext{\skt{स्थानेषूक्तेषु}}{\Afootnote{A H; \skt{स्थानेसूक्तेषु}\ N}}\skt{ सम्मिश्रः }\edtext{\skt{सम्मिश्राः}}{\Afootnote{A H; \skt{सम्मिश्रा}\ N}}\skt{ कुरुते }\edtext{\skt{रुजाः}}{\Afootnote{N H; \skt{रुजः ॥}\ A}}\skt{ ॥}\pend

\pstart \ledsidenote{\passref{2.1.32ab}}\edtext{\skt{कुर्या\-द\-व\-य\-व\-स्थश्च}}{\Afootnote{N H; \skt{कुर्या\-द\-व\-य\-व\-प्राप्तो}\ A}}\skt{ }\edtext{\skt{मातरिश्वा}}{\Afootnote{N H; \skt{मारुतस्त्व}\ A}}\skt{ }\edtext{\skt{गदान्बहून्}}{\Afootnote{\skt{मितान्गदान् ।}\ A; \skt{गदां वहून्।}\ N; \skt{गदान्वहून् ।}\ H}}\skt{ ।}\pend

\pstart \ledsidenote{\passref{2.1.34cd}}\skt{प्राणे पित्तावृते }\edtext{\skt{च्छर्द्दि\-र्द्दा\-ह\-श्चै\-वो\-प\-जा\-यते}}{\lemma{\skt{द्दाहश्}}\Afootnote{A N\textsuperscript{pc} H; \skt{दाहआश्}\ N\textsuperscript{ac}}}\skt{ ॥}\pend

\pstart\versepar \ledsidenote{\passref{2.1.35}}\edtext{\skt{दौर्बल्यं}}{\Afootnote{A; \skt{दौर्वल्यं}\ N H}}\skt{ सदनन्तन्द्रा }\edtext{\skt{वैरस्यञ्च}}{\Afootnote{N H; \skt{वैवर्ण्यं च}\ A}}\skt{ कफावृते । }\pend
\pstart\versepar \skt{उदाने }\edtext{\skt{पित्तयुक्ते तु}}{\Afootnote{N H; \skt{पित्तसंयुक्ते}\ A}}\skt{ }\edtext{\skt{मोह\-मू\-र्च्छा\-भ्र\-म\-क्लमाः}}{\Afootnote{H; \skt{मूर्च्छा\-दा\-ह\-भ्र\-म\-क्लमाः ॥}\ A; \skt{मोह\-मू\-र्च्छा\-भ्र\-म\-क्लमा॥}\ N}}\skt{ ॥}\pend

\pstart\versepar \ledsidenote{\passref{2.1.36}}\edtext{\skt{अस्वेदहर्षो}}{\Afootnote{N H; \skt{अस्वेदहर्षौ}\ A}}\skt{ }\edtext{\skt{मन्दाग्निः}}{\Afootnote{N H; \skt{मन्दोऽग्निः}\ A}}\skt{ }\edtext{\skt{शीतता च}}{\Afootnote{N\textsuperscript{pc}; \skt{शीतस्तम्भौ}\ A; \emph{om.}\ N\textsuperscript{ac}; \skt{शीतता}\ H}}\skt{ }\edtext{\skt{कफावृते}}{\Afootnote{A N\textsuperscript{pc} H; \emph{om.}\ N\textsuperscript{ac}}}\skt{ । }\pend
\pstart\versepar \skt{}\edtext{\edtext{\skt{स्वेद\-दा\-हो\-ष्ण\-मूर्च्छाः स्युः}}{\Afootnote{\skt{स्वेद\-दा\-हो\-ष्ण\-मूर्च्छा स्युः}\ N; \skt{स्वेद\-दा\-हो\-ष्ण्य\-मूर्च्छाः स्युः}\ H}}\skt{ समाने }\edtext{\skt{पित्तसंयुते}}{\Afootnote{N H; \skt{पित्तसंयुक्ते स्वेद\-दा\-हौ\-ष्ण्य\-मू\-र्च्छ\-नम् ॥}\ A}}}{\lemma{\skt{स्वेद\-दा\-हो\-ष्ण\-मूर्च्छाः… स}}\Afootnote{\emph{om.}\ A}}\skt{ ॥}\pend

\pstart\versepar \ledsidenote{\passref{2.1.37}}\edtext{\skt{कफेन}}{\Afootnote{N H; \skt{कफाधि}\ A}}\skt{ }\edtext{\skt{संगे}}{\Afootnote{H; \skt{कं च}\ A; \skt{संसंगो}\ N}}\skt{ }\edtext{\skt{विण्मूत्रे}}{\Afootnote{H; \skt{विण्मूत्रं}\ A; \skt{विट्मूत्रे}\ N}}\skt{ }\edtext{\skt{गात्रहर्षश्च}}{\Afootnote{N; \skt{रोमहर्षः क}\ A; \skt{गात्रे हर्षश्च}\ H}}\skt{ }\edtext{\skt{जायते}}{\Afootnote{N H; \skt{फावृते ।}\ A}}\skt{ । }\pend
\pstart\versepar \skt{अपाने }\edtext{\skt{पित्तयुक्ते तु}}{\Afootnote{N H; \skt{पित्तसंयुक्ते}\ A}}\skt{ }\edtext{\skt{दाहोष्णो}}{\Afootnote{H; \skt{दाहौष्ण्ये}\ A; \skt{दाहौष्णो}\ N}}\skt{ }\edtext{\skt{रक्तमूत्रता}}{\Afootnote{N H; \skt{स्यादसृग्दरः ॥}\ A}}\skt{ ॥}\pend

\pstart\versepar \ledsidenote{\passref{2.1.38}}\skt{अधः\-का\-य\-गु\-रु\-त्वञ्च }\edtext{\skt{शीतता च}}{\Afootnote{N H; \skt{तस्मिन्नेव}\ A}}\skt{ कफावृते । }\pend
\pstart\versepar \skt{व्याने पित्तावृते दाहो }\edtext{\skt{गात्रविक्षेपणः}}{\Afootnote{N H; \skt{गात्रविक्षेपणं}\ A}}\skt{ क्लमः ॥}\pend

\pstart \ledsidenote{\passref{2.1.39.1}}\edtext{\edtext{\skt{स्तम्भ\-नो\-द्द\-ण्ड\-क\-श्चापि}}{\lemma{\skt{स्तम्भ\-नो\-द्द\-ण्ड\-कश्}}\Afootnote{N\textsuperscript{pc} H\textsuperscript{pc}; \emph{om.}\ N\textsuperscript{ac}; \emph{om.}\ H\textsuperscript{ac}}}}{\lemma{\skt{चापि}}\Afootnote{N\textsuperscript{pc} H\textsuperscript{pc}; \emph{om.}\ N\textsuperscript{ac} H\textsuperscript{ac}}}\skt{ }\edtext{\skt{शोथशूलं}}{\Afootnote{N\textsuperscript{pc} H\textsuperscript{pc}; \emph{om.}\ N\textsuperscript{ac} H\textsuperscript{ac}}}\skt{ }\edtext{\skt{कफावृते}}{\Afootnote{N\textsuperscript{pc} H\textsuperscript{pc}; \emph{om.}\ N\textsuperscript{ac} H\textsuperscript{ac}}}\skt{ ॥}\pend

\pstart\versepar \ledsidenote{\passref{2.1.40}}\skt{प्रायशः सुकुमाराणां मिथ्या\-हा\-र\-वि\-हा\-रिणां ॥ }\pend
\pstart\versepar \skt{}\edtext{\skt{दोषा\-ध्व\-म\-द्य\-प्र\-म\-दा\-व्या\-या\-मैश्च}}{\Afootnote{N H; \skt{रोगा\-ध्व\-प्र\-म\-दा\-म\-द्य\-व्या\-या\-मैश्चा}\ A}}\skt{ }\edtext{\skt{प्रपीडनात्}}{\Afootnote{\skt{तिपीडनात् ॥}\ A; \skt{प्रपीडिताः ।}\ N; \skt{प्रपीडितान् ॥}\ H}}\skt{ ॥}\pend

\pstart\versepar \ledsidenote{\passref{2.1.41}}\edtext{\skt{ऋतु\-दे\-श\-वि\-प\-र्या\-सा\-द\-सा\-त्म्या\-नाञ्च}}{\Afootnote{N H; \skt{ऋतु\-सा\-त्म्य\-वि\-प\-र्या\-सा\-त्स्ने\-हा\-दीनां च}\ A}}\skt{ भोजनात्। }\pend
\pstart\versepar \skt{स्थूल\-स्या\-व्या\-या\-म\-वतो वात\-र\-क्त\-म्प्र\-कु\-प्यति ॥}\edtext{}{\lemma{\skt{भोजनात्। स्थूल\-स्या\-व्या\-या\-म\-वतो}}\Afootnote{N H; \skt{विभ्रमात् । अव्यवाये तथा स्थूले}\ A}}\pend

\pstart\versepar \ledsidenote{\passref{2.1.42}}\edtext{\skt{हस्त्य\-श्वो\-ष्ट्रै\-र्ग\-च्छ\-तोऽ\-न्यैश्च}}{\lemma{\skt{ऽन्यैश्च}}\Afootnote{A H; \skt{न्यैश्च}\ N}}\skt{ वायुः }\pend
\pstart\versepar \skt{}\edtext{\skt{कोपं}}{\Afootnote{A H; \skt{कोपो}\ N}}\skt{ }\edtext{\skt{यायात्कारणैः}}{\Afootnote{N H; \skt{यातः कारणैः}\ A}}\skt{ }\edtext{\skt{सेवितैश्च}}{\Afootnote{N H; \skt{सेवितैः स्वैः ।}\ A}}\skt{ । }\pend
\pstart\versepar \skt{}\edtext{\skt{तीक्ष्णो\-ष्णा\-म्ल\-क्षा\-र\-शा\-का\-दि\-भोज्यैः}}{\Afootnote{A H; \skt{तीक्ष्णो\-ष्णा\-म्व्ल\-क्षा\-र\-शा\-का\-दि\-भौ\-ज्यैस्}\ N}}\skt{ }\pend
\pstart\versepar \skt{सन्तापाद्यैर्भूय }\edtext{\skt{आसेवितैश्च}}{\lemma{\skt{आसेवितैश्}}\Afootnote{N H; \skt{सा सेवितैश्}\ A}}\skt{ ॥}\pend

\pstart\versepar \ledsidenote{\passref{2.1.43abc}}\edtext{\skt{शीघ्रं}}{\Afootnote{N H; \skt{क्षिप्रं}\ A}}\skt{ }\edtext{\edtext{\skt{रक्त\-म्वि\-द्र\-व\-त्याशु}}{\lemma{\skt{विद्रवत्य्}}\Afootnote{N H; \skt{दुष्टिम्}\ A}}}{\lemma{\skt{आशु}}\Afootnote{N H; \skt{आयाति}\ A}}\skt{ तच्च }\pend
\pstart\versepar \skt{वायोर्मार्ग्गं }\edtext{\skt{संरुणद्ध्याशु}}{\lemma{\skt{संरुणद्ध्य्}}\Afootnote{A N; \skt{संरुणध्ध्य्}\ H}}\skt{ यातः । }\pend
\pstart\versepar \skt{}\edtext{\skt{क्रुद्धोऽ\-त्य\-र्थ\-म्मा\-र्ग\-रो\-धा\-द्वि\-प\-न्नस्}}{\lemma{\skt{मार्ग\-रो\-धा\-द्वि\-प\-न्नस्}}\Afootnote{N; \skt{मार्गरोधात्}\ A; \skt{मार्ग्ग\-रो\-धा\-द्वि\-प\-न्नस् …}\ H}}\skt{ ॥}\pend

\pstart\versepar \ledsidenote{\passref{2.1.44}}\skt{तत्सं\-पृ\-क्त\-म्वा\-युना दूषितेन ॥ }\pend
\pstart\versepar \skt{}\edtext{\skt{तत्प्रा\-ब\-ल्या\-दु\-च्यते}}{\Afootnote{A; \skt{तत्प्रा\-व\-ल्या\-दु\-च्यते}\ N H}}\skt{ वातरक्तं॥ }\pend
\pstart\versepar \skt{}\edtext{\skt{तद्वत्पित्तं}}{\Afootnote{A H; \skt{तद्वत्पित्त}\ N}}\skt{ }\edtext{\skt{दूषि\-ते\-ना\-सृ\-जाक्तं}}{\Afootnote{N; \skt{दूषि\-ते\-ना\-सृ\-जाक्तं श्लेष्मा दुष्टो दूषि\-ते\-ना\-सृ\-जाक्तः ॥}\ A; \skt{दूषि\-ते\-ना\-सृ\-जाक्तं । …}\ H}}\skt{ ।}\pend

\pstart\versepar \ledsidenote{\passref{2.1.45}}\edtext{\skt{अस्प\-र्शे\-च्छा\-तो\-द\-भे\-द\-प्र\-शोषो}}{\Afootnote{H; \skt{स्पर्शोद्विग्नौ तोद\-भे\-द\-प्र\-शोष}\ A; \skt{अस्प\-र्शे\-च्छा\-तो\-द\-भे\-द\-प्र\-शोषौ}\ N}}\skt{ }\pend
\pstart\versepar \skt{}\edtext{\skt{स्वापोपेतो}}{\Afootnote{H; \skt{स्वापोपेतौ}\ A N}}\skt{ वातरक्तेन पादौ ॥ }\pend
\pstart\versepar \skt{पित्तासृग्भ्यां }\edtext{\edtext{\skt{भव\-त\-स्तू\-ग्र\-दा\-हाव्}}{\lemma{\skt{भवतस्}}\Afootnote{N H; \emph{om.}\ A}}}{\lemma{\skt{तूग्रदाहाव्}}\Afootnote{N H; \skt{उग्रदाहौ}\ A}}\skt{ }\pend
\pstart\versepar \skt{}\edtext{\skt{अत्यर्थोष्णौ}}{\Afootnote{N H; \skt{भवे\-ता\-म\-त्य\-र्थोष्णौ}\ A}}\skt{ रक्तशोफौ मृदू च ।}\pend

\pstart\versepar \ledsidenote{\passref{2.1.46}}\skt{कण्डूमन्तौ }\edtext{\skt{श्वेतशीतौ}}{\Afootnote{A; \skt{स्वेतशीतौ}\ N H}}\skt{ }\edtext{\skt{सशोथौ}}{\Afootnote{N H; \skt{सशोफौ}\ A}}\skt{ }\pend
\pstart\versepar \skt{}\edtext{\skt{पीनस्तब्धौ}}{\Afootnote{A H; \skt{पीनस्तव्धौ}\ N}}\skt{ }\edtext{\skt{श्लेष्मदुष्टे तु}}{\Afootnote{A H; \skt{श्लेष्मदुष्टौ तु}\ N\textsuperscript{ac}; \skt{श्लेष्मदुष्टो तु}\ N\textsuperscript{pc}}}\skt{ }\edtext{\skt{रक्ते}}{\Afootnote{A H; \skt{रक्तौ । ।}\ N\textsuperscript{ac}; \skt{रक्तो । ।}\ N\textsuperscript{pc}}}\skt{ ॥ }\pend
\pstart\versepar \skt{}\edtext{\skt{सर्व्वैर्दुष्टे}}{\lemma{\skt{सर्व्वैर्}}\Afootnote{A; \skt{सर्व्वे}\ N; \skt{सर्व्व}\ H}}\skt{ शोणिते चापि दोषाः}\skt{ }\pend
\pstart\versepar \skt{स्वं स्वं }\edtext{\skt{लिङ्ग\-म्पा\-द\-यो\-र्द्द\-र्श\-यन्ति}}{\lemma{\skt{लिङ्गम्}}\Afootnote{N H; \skt{रूपं}\ A}}\skt{ ॥}\pend

\pstart\versepar \ledsidenote{\passref{2.1.48}}\skt{पाद\-यो\-र्मू\-ल\-मा\-स्थाय कदा\-चि\-द्ध\-स्त\-यो\-रपि । }\pend
\pstart\versepar \skt{}\edtext{\edtext{\skt{आखोर्व्विषमिव}}{\lemma{\skt{आखोर्}}\Afootnote{H; \skt{आस्वोर्}\ A; \skt{आखा}\ N}}}{\lemma{\skt{व्विषम्}}\Afootnote{A H; \skt{द्विषम्}\ N}}\skt{ }\edtext{\skt{क्रुद्धं}}{\Afootnote{A; \skt{क्रुद्धस्}\ N H}}\skt{ }\edtext{\skt{तद्दे\-ह\-म\-नु\-स\-र्प्पति}}{\lemma{\skt{तद्देहम्}}\Afootnote{A H; \skt{त देहम्}\ N}}\skt{ ॥}\pend

\pstart \ledsidenote{\passref{2.1.49ab}}\edtext{\skt{आजानुस्फुटितं}}{\Afootnote{A H; \skt{अजान्तुस्फुटितं}\ N}}\skt{ यच्च }\edtext{\skt{प्रभि\-न्न\-म्प्र\-स्रु\-तञ्च}}{\lemma{\skt{प्रस्रुतञ्च}}\Afootnote{A N; \skt{प्रश्रुतञ्च}\ H}}\skt{ यत् ।}\pend

\pstart \ledsidenote{\passref{2.1.50ab}}\edtext{\edtext{\edtext{\skt{वात\-र\-क्त\-म\-सा\-ध्य\-न्त\-द्याप्यं}}{\lemma{\skt{वातरक्तम्}}\Afootnote{N H; \skt{शोणितं तद}\ A}}}{\lemma{\skt{असाध्यन्}}\Afootnote{N H; \skt{साध्यं}\ A}}}{\lemma{\skt{तद्याप्यं}}\Afootnote{H; \skt{स्याद्याप्यं}\ A; \skt{तद्याप्य}\ N}}\skt{ }\edtext{\skt{सम्व\-त्स\-रो\-त्थि\-तम्}}{\Afootnote{A H; \skt{न्सम्वत्सरोमितः ॥}\ N\textsuperscript{ac}; \skt{न्सम्वत्सरोमितं ॥}\ N\textsuperscript{pc}}}\skt{ ॥}\pend

\pstart \ledsidenote{\passref{2.1.50cd}}\edtext{\skt{यदा तु}}{\Afootnote{A N; \skt{… यदा तु}\ H}}\skt{ }\edtext{\skt{धमनीः}}{\Afootnote{A H; \skt{धर्मनीः}\ N}}\skt{ सर्व्वाः }\edtext{\edtext{\skt{कुपितोऽभ्येति}}{\lemma{\skt{कुपितो}}\Afootnote{A H; \skt{कुपिते}\ N}}}{\lemma{\skt{ऽभ्येति}}\Afootnote{A; \skt{भ्येति}\ N H}}\skt{ मारुतः ।}\pend

\pstart\versepar \ledsidenote{\passref{2.1.51}}\skt{तदाक्षिपत्याशु }\edtext{\skt{मुहु\-र्मु\-हु\-र्दे\-ह\-म्मु\-हु\-श्चरः}}{\lemma{\skt{मुहुर्}}\Afootnote{A H; \skt{मुहु}\ N}}\skt{ ॥ }\pend
\pstart\versepar \skt{}\edtext{\skt{मुहु\-र्म्मु\-हु\-स्त्वा\-क्षि\-प\-णा\-दा\-क्षे\-पक}}{\lemma{\skt{त्वाक्षि\-प\-णा\-दा\-क्षे\-पक}}\Afootnote{N H; \skt{तदा\-क्षे\-पा\-दा\-क्षे\-पक}\ A}}\skt{ इति स्मृतः ।}\pend

\pstart\versepar \ledsidenote{\passref{2.1.52}}\edtext{\edtext{\edtext{\skt{यतोऽ\-प\-ता\-म्य\-तेऽ\-त्य\-र्थ\-मतो}}{\lemma{\skt{ऽपताम्यते}}\Afootnote{N; \skt{यं ताम्यते}\ H}}}{\lemma{\skt{ऽत्यर्थम्}}\Afootnote{H; \skt{त्यर्थम्}\ N}}\skt{ }\edtext{\skt{ज्ञेयोऽपतानकः}}{\lemma{\skt{ऽपतानकः}}\Afootnote{N H; \skt{ऽप\-ता\-न\-क\-संज्ञो यः}\ A}}}{\lemma{\skt{… ज्ञेयो}}\Afootnote{\emph{om.}\ A}}\skt{ ॥ }\pend
\pstart\versepar \skt{}\edtext{\skt{कफान्वितो}}{\Afootnote{H; \skt{पात\-य\-त्य\-न्त\-राऽ\-न्तरा । कफान्वितो}\ A; \skt{कफन्वितो}\ N}}\skt{ भृश\-म्वा\-यु\-स्ता\-स्वेव यदि तिष्ठति ।}\pend

\pstart \ledsidenote{\passref{2.1.53ab}}\edtext{\skt{स दण्ड\-व\-त्स्त\-म्भ\-यति स तु}}{\Afootnote{N H; \skt{स दण्ड\-व\-त्स्त\-म्भ\-यति कृच्छ्रो}\ A}}\skt{ }\edtext{\skt{दण्डापतानकः}}{\Afootnote{A N; \skt{दण्डापतानकः ॥ …}\ H}}\skt{ ॥}\pend

\pstart\versepar \ledsidenote{\passref{2.1.54}}\edtext{\skt{धनु\-स्तु\-ल्य\-न्न\-मे\-द्यस्तु}}{\lemma{\skt{धनुस्तुल्यन्}}\Afootnote{A H; \skt{धनुतुल्यन्}\ N}}\skt{ स }\edtext{\skt{धनुः\-स्त\-म्भ\-सं\-ज्ञितः}}{\Afootnote{H; \skt{धनुः\-स्त\-म्भ\-सं\-ज्ञकः ।}\ A; \skt{धनु\-त\-म्भ\-स\-ञ्ज्ञितः ।}\ N}}\skt{ । }\pend
\pstart\versepar \skt{}\edtext{\skt{अङ्गु\-ली\-गु\-ल्फ\-ज\-ठ\-र\-हृ\-द्व\-क्षो\-ग\-ल\-सं\-श्रितः}}{\Afootnote{A H; \skt{अङ्गु\-ली\-गु\-ल्फ\-ज\-ठ\-र\-हृ\-द्व\-क्षो\-ग\-ल\-सं\-श्रितंः ॥}\ N}}\skt{ ॥}\pend

\pstart\versepar \ledsidenote{\passref{2.1.55}}\edtext{\edtext{\skt{स्नायु\-प्र\-ता\-न\-म\-निलो}}{\lemma{\skt{स्नायुप्रतानम्}}\Afootnote{A H; \skt{स्नायुप्रतान}\ N}}}{\lemma{\skt{अनिलो}}\Afootnote{A H; \skt{निनिलो}\ N\textsuperscript{ac}; \skt{ननिलो}\ N\textsuperscript{pc}}}\skt{ यदाक्षिपति वेगवान् । }\pend
\pstart\versepar \skt{}\edtext{\skt{विष्ट\-ब्धा\-क्ष\-स्त\-ब्ध\-ह\-नु\-र्भ\-ग्न\-पार्श्वः}}{\lemma{\skt{भग्नपार्श्वः}}\Afootnote{A N; \skt{व्भग्नपार्श्वः}\ H}}\skt{ }\edtext{\skt{कफं}}{\Afootnote{A N; \skt{कफा}\ H}}\skt{ }\edtext{\skt{वमन्}}{\Afootnote{A N; \skt{न्वमन् ॥}\ H}}\skt{ ॥}\edtext{}{\lemma{\skt{वेगवान् । विष्ट\-ब्धा\-क्ष\-स्त\-ब्ध\-ह\-नुर्}}\Afootnote{\skt{वेगवान् । विष्टब्धाक्षः स्तब्धहनुर्}\ A; \skt{वेगवान् । विष्ट\-ब्धा\-क्ष\-स्त\-व्ध\-ह\-नुर्}\ N; \skt{वेगवान् । विष्ट\-व्धा\-क्ष\-स्त\-व्ध\-ह\-नुर्}\ H}}\pend

\pstart\versepar \ledsidenote{\passref{2.1.56}}\skt{अभ्य\-न्त\-र\-न्ध\-नु\-रिव यदा }\edtext{\skt{नाम्यति}}{\Afootnote{N H; \skt{नमति}\ A}}\skt{ मानवः । }\pend
\pstart\versepar \skt{}\edtext{\skt{तदा }\edtext{\skt{सोऽभ्यन्तरायामं}}{\lemma{\skt{ऽभ्यन्तरायामं}}\Afootnote{H; \skt{स्याभ्यन्तरायामं}\ A; \skt{भ्यन्तरायांमं}\ N}}}{\lemma{\skt{तदा सो}}\Afootnote{N H; \skt{तदाऽ}\ A}}\skt{ कुरुते मारुतो }\edtext{\skt{बली}}{\Afootnote{A; \skt{वली ॥}\ N H}}\skt{ ॥}\pend

\pstart\versepar \ledsidenote{\passref{2.1.57}}\edtext{\skt{बाह्य\-स्ना\-यु\-प्र\-ता\-नस्थो}}{\Afootnote{A; \skt{वाह्य\-स्ना\-यु\-प्र\-ता\-नस्थो}\ N H}}\skt{ }\edtext{\skt{बाह्या\-या\-म\-ङ्क\-रोति}}{\lemma{\skt{बाह्यायामङ्}}\Afootnote{A; \skt{वाह्यायामङ्}\ N H}}\skt{ च । }\pend
\pstart\versepar \skt{}\edtext{\skt{तमसाध्यं}}{\lemma{\skt{असाध्यं}}\Afootnote{N H; \skt{मसाध्यं}\ A}}\skt{ }\edtext{\skt{बुधाः}}{\Afootnote{A; \skt{वुधाः}\ N H}}\skt{ }\edtext{\skt{प्राहु\-र्व्व\-क्षः\-क\-ट्यू\-रु\-भ\-ञ्ज\-नम्}}{\lemma{\skt{व्वक्षः\-क\-ट्यू\-रु\-भ\-ञ्ज\-नम्}}\Afootnote{H; \skt{वक्षः\-क\-ठ्यू\-रु\-भ\-ञ्ज\-नम् ॥}\ A; \skt{वक्षः\-क\-यू\-रु\-भ\-ञ्जन ॥}\ N}}\skt{ ॥}\edtext{}{\lemma{\skt{करोति च । तम्}}\Afootnote{N H; \skt{करोति च । त}\ A}}\pend

\pstart\versepar \ledsidenote{\passref{2.1.58}}\skt{कफपित्तान्वितो वायुर्व्वायुरेव च }\edtext{\skt{केवलः।}}{\lemma{\skt{केवलः}}\Afootnote{A N; \skt{केवलंः।}\ H\textsuperscript{ac}; \skt{केवलं।}\ H\textsuperscript{pc}}}\skt{ }\pend
\pstart\versepar \skt{}\edtext{\skt{कुर्यादाक्षेपकं}}{\Afootnote{A H; \skt{कुर्या\-त्वा\-क्षे\-प\-कम्}\ N}}\skt{ त्वन्यं चतु\-र्थ\-म\-भि\-घा\-त\-जम् ॥}\pend

\pstart\versepar \ledsidenote{\passref{2.1.59}}\skt{गर्ब्भ\-पा\-त\-नि\-मि\-त्तश्च शोणि\-ता\-ति\-स्र\-वाच्च यः । }\pend
\pstart\versepar \skt{अभि\-घा\-त\-नि\-मि\-त्तश्च न सिध्य\-त्य\-प\-ता\-नकः ॥}\pend

\pstart\versepar \ledsidenote{\passref{2.1.60}}\edtext{\skt{अधो\-ग\-मा\-श्चो\-र्ध्व\-गाश्च}}{\lemma{\skt{चोर्ध्वगाश्च}}\Afootnote{H; \skt{स}\ A; \skt{चोर्ध्वगश्च}\ N}}\skt{ }\edtext{\edtext{\skt{तिर्यग्गाश्चानिलो}}{\lemma{\skt{तिर्यग्गाश्}}\Afootnote{H; \skt{तिर्यग्गा}\ A; \skt{तिर्यगाश्}\ N}}}{\lemma{\skt{चानिलो}}\Afootnote{N H; \skt{धमनीर्}\ A}}\skt{ }\edtext{\skt{बली}}{\Afootnote{\skt{ऊर्ध्वदेहगाः ।}\ A; \skt{वली ।}\ N H}}\skt{ । }\pend
\pstart\versepar \skt{}\edtext{\skt{यदा\-त्य\-र्थ\-म्प्र\-कु\-पितो}}{\lemma{\skt{यदात्यर्थम्}}\Afootnote{N H; \skt{यदा}\ A}}\skt{ }\edtext{\skt{धमनीः}}{\Afootnote{N H; \skt{ऽत्यर्थं मातरिश्वा}\ A}}\skt{ }\edtext{\skt{प्रतिपद्यते}}{\Afootnote{N H; \skt{प्रपद्यते ॥}\ A}}\skt{ ॥}\pend

\pstart\versepar \ledsidenote{\passref{2.1.61}}\skt{तदा\-न्य\-त\-र\-प\-क्षस्य सन्धि\-ब\-न्धा\-न्वि\-मो\-क्ष\-यन् । }\pend
\pstart\versepar \skt{हन्ति }\edtext{\skt{पक्षन्तमाहुश्च}}{\lemma{\skt{आहुश्च}}\Afootnote{N H; \skt{आहुर्हि}\ A}}\skt{ }\edtext{\skt{पक्षा\-घा\-त\-म्भि\-ष\-ग्वराः}}{\lemma{\skt{भिषग्वराः}}\Afootnote{A N\textsuperscript{pc} H; \skt{भिषग्वाराः ॥}\ N\textsuperscript{ac}}}\skt{ ॥}\edtext{}{\lemma{\skt{सन्धि\-ब\-न्धा\-न्वि\-मो\-क्ष\-यन् । हन्ति}}\Afootnote{A; \skt{सन्धि\-व\-न्धा\-न्य\-मो\-क्ष\-यन् । हन्ति}\ N; \skt{सन्धि\-व\-न्धा\-न्वि\-मो\-क्ष\-यन् । हन्ति}\ H}}\pend

\pstart\versepar \ledsidenote{\passref{2.1.62}}\edtext{\skt{तस्य}}{\Afootnote{N H; \skt{यस्य}\ A}}\skt{ कृत्स्नं शरी\-रा\-र्द्ध\-म\-क\-र्म्म\-ण्य\-म\-चे\-तनं । }\pend
\pstart\versepar \skt{ततः पतत्यसून्वापि त्यज\-त्य\-नि\-ल\-पी\-डितः ॥}\pend

\pstart \ledsidenote{\passref{2.1.32cd}}\edtext{\skt{दाह\-स\-न्ता\-प\-मूर्च्छा स्युर्व्वायौ}}{\lemma{\skt{दाह\-स\-न्ता\-प\-मूर्च्छा स्युर्}}\Afootnote{N; \skt{दाह\-सं\-ता\-प\-मूर्च्छाः स्युर्}\ A; \skt{दाहः सन्तापमूर्च्छा स्युर्}\ H}}\skt{ पित्तसमन्विते ।}\pend

\pstart \ledsidenote{\passref{2.1.33ab}}\edtext{\skt{शैत्य\-शो\-फ\-गु\-रु\-त्वञ्च}}{\Afootnote{H; \skt{शैत्य\-शो\-फ\-गु\-रु\-त्वानि}\ A; \skt{सैत्य\-शो\-फ\-गु\-रु\-त्वञ्च}\ N}}\skt{ तस्मिन्नेव }\edtext{\skt{कफान्विते}}{\Afootnote{N H; \skt{कफावृते ।}\ A}}\skt{ ॥}\pend

\pstart\versepar \ledsidenote{\passref{2.1.63}}\edtext{\skt{शुद्ध\-वा\-ता\-ह\-त\-म्पक्षं}}{\lemma{\skt{शुद्ध\-वा\-ता\-ह\-तम्}}\Afootnote{N H; \skt{शुद्धवातहतं}\ A}}\skt{ }\edtext{\skt{कृच्छ्रं}}{\Afootnote{N H; \skt{कृच्छ्र}\ A}}\skt{ साध्यतमम्विदुः ।}\skt{ }\pend
\pstart\versepar \skt{साध्यमन्येन संसृष्टमसाध्यं क्षयहेतुकम् ॥}\pend

\pstart\versepar \ledsidenote{\passref{2.1.68cdef}}\skt{उच्चै\-र्व्व्या\-ह\-र\-तोऽ\-त्यर्थं }\edtext{\skt{खादतः}}{\Afootnote{N H; \skt{स्वादतः}\ A}}\skt{ कठिनानि वा । }\pend
\pstart\versepar \skt{हसतो जृम्भतो }\edtext{\skt{भारा\-च्छ\-य\-ना\-द्वि\-ष\-मा\-दपि}}{\Afootnote{N H; \skt{भारा\-द्वि\-ष\-मा\-च्छ\-य\-ना\-दपि ॥}\ A}}\skt{ ॥}\pend

\pstart \ledsidenote{\passref{2.1.69cd}}\edtext{\skt{अर्द्दयित्वानिलो}}{\Afootnote{A N; \skt{… अर्द्दयित्वानिलो}\ H}}\skt{ }\edtext{\skt{वक्त्र\-म\-र्द्दि\-त\-ञ्ज\-न\-य\-त्यतः}}{\lemma{\skt{वक्त्रम्}}\Afootnote{A H; \skt{वक्रम्}\ N}}\skt{ ॥}\pend

\pstart\versepar \ledsidenote{\passref{2.1.70}}\skt{वक्रीभवति वक्त्रार्द्धङ्ग्रीवा }\edtext{\skt{चाप्युपवर्त्तते}}{\lemma{\skt{उपवर्त्तते}}\Afootnote{H; \skt{अपवर्तते}\ A N}}\skt{ ॥ }\pend
\pstart\versepar \skt{शिरश्चलति }\edtext{\skt{वाक्भंगो}}{\Afootnote{N; \skt{वाक्सङ्गो}\ A; \skt{वाक्भंगो}\ H}}\skt{ नेत्रादीनां च वैकृतं ।}\pend

\pstart \ledsidenote{\passref{2.1.71ab}}\edtext{\skt{ग्रीवा\-चि\-बु\-क\-द\-न्तानां}}{\Afootnote{A; \skt{ग्रीवा\-चि\-वु\-क\-द\-न्तानां}\ N H}}\skt{ }\edtext{\skt{तस्मिन्पार्श्वे च}}{\Afootnote{N H; \skt{तस्मिन्पार्श्वे तु}\ A}}\skt{ वेदना ॥}\pend

\pstart \ledsidenote{\passref{2.1.72cd}}\skt{तमर्दितमिति }\edtext{\skt{प्राहुर्व्व्याधिं}}{\lemma{\skt{व्व्याधिं}}\Afootnote{A H; \emph{om.}\ N}}\skt{ }\edtext{\skt{व्याधिविचक्षणाः}}{\Afootnote{N H; \skt{व्याधिविशारदाः ॥}\ A}}\skt{ ॥}\pend

\pstart\versepar \ledsidenote{\passref{2.1.73}}\skt{क्षीण\-स्या\-नि\-मि\-षा\-क्षस्य प्रस\-क्ता\-व्य\-क्त\-भा\-षिणः ॥ }\pend
\pstart\versepar \skt{न }\edtext{\skt{सिध्य\-त्य\-र्दि\-त\-म्बाढं}}{\lemma{\skt{बाढं}}\Afootnote{A; \skt{वाढं}\ N H}}\skt{ त्रिव\-र्ष\-म्वे\-प\-नस्य च ॥}\edtext{}{\lemma{\skt{प्रस\-क्ता\-व्य\-क्त\-भा\-षिणः ॥ न}}\Afootnote{N; \skt{प्रसक्तं सक्तभाषिणः । न}\ A; \skt{प्रश\-क्ता\-व्य\-क्त\-भा\-षिणः ॥ न}\ H}}\pend

\pstart\versepar \ledsidenote{\passref{2.1.74}}\edtext{\skt{पार्ष्णि\-प्र\-त्य\-ङ्गु\-लीनां या}}{\Afootnote{H; \skt{पार्ष्णि\-प्र\-त्य\-ङ्गु\-लीनां तु}\ A; \skt{पार्ष्णि\-प्र\-त्य\-ङ्गु\-लीना याङ्}\ N}}\skt{ कण्डरा }\edtext{\skt{सानिलार्द्दिता}}{\Afootnote{H; \skt{याऽनिलार्दिताः ।}\ A; \skt{यानिलार्दिता ॥}\ N}}\skt{ । }\pend
\pstart\versepar \skt{सक्थ्नः }\edtext{\skt{क्षेपन्न}}{\Afootnote{N H; \skt{क्षेपं नि}\ A}}\skt{ }\edtext{\skt{गृह्णाति}}{\Afootnote{N H; \skt{गृह्णीया}\ A}}\skt{ }\edtext{\skt{गृध्रसीति च सा}}{\Afootnote{N H; \skt{द्गृध्रसीति हि सा}\ A}}\skt{ }\edtext{\skt{स्मृता}}{\Afootnote{A H; \skt{स्मृताः ।}\ N}}\skt{ ॥}\pend

\pstart\versepar \ledsidenote{\passref{2.1.75}}\edtext{\skt{तल\-म्प्र\-त्यं\-गु\-ली\-नान्तु}}{\lemma{\skt{तलम्}}\Afootnote{N H; \skt{तल}\ A}}\skt{ कण्डरा }\edtext{\skt{वाहुपृष्ठतः}}{\Afootnote{N H; \skt{बाहुपृष्ठतः ।}\ A}}\skt{ । }\pend
\pstart\versepar \skt{}\edtext{\skt{बाह्वोः}}{\Afootnote{A; \skt{वाह्वोः}\ N H}}\skt{ कर्म्मक्षयकरी }\edtext{\skt{विश्वञ्चीतीह}}{\Afootnote{H; \skt{विश्वाचीति हि}\ A; \skt{विश्वञ्जती हि}\ N\textsuperscript{ac}; \skt{विश्वञ्चीती हि}\ N\textsuperscript{pc}}}\skt{ }\edtext{\skt{सोच्यते}}{\Afootnote{N; \skt{सा स्मृता ॥}\ A; \skt{चोच्यते ॥}\ H}}\skt{ ॥}\pend

\pstart\versepar \ledsidenote{\passref{2.1.76}}\edtext{\skt{वातशोणितजः}}{\Afootnote{A N; \skt{वातः शोणितजः}\ H}}\skt{ शोफो जानुमध्ये महारुजः । }\pend
\pstart\versepar \skt{}\edtext{\skt{ज्ञेयः}}{\Afootnote{N H; \skt{शिरः}\ A}}\skt{ }\edtext{\skt{क्रोष्टुकशीर्षेति}}{\Afootnote{N H; \skt{क्रोष्टुकपूर्वं तु}\ A}}\skt{ स्थूलः }\edtext{\skt{क्रोष्टु\-क\-शी\-र्ष\-वत्}}{\Afootnote{N H; \skt{क्रोष्टु\-क\-मू\-र्ध\-वत् ॥}\ A}}\skt{ ॥}\pend

\pstart\versepar \ledsidenote{\passref{2.1.77}}\skt{वायुः }\edtext{\skt{कट्याश्रितः}}{\Afootnote{N H; \skt{कट्यां स्थितः}\ A}}\skt{ सक्थ्नः }\edtext{\skt{कण्ड\-रा\-मा\-क्षि\-पे\-द्यदा}}{\lemma{\skt{कण्डराम्}}\Afootnote{A H; \skt{कण्ठराम्}\ N}}\skt{ । }\pend
\pstart\versepar \skt{}\edtext{\skt{खञ्जस्तदा}}{\lemma{\skt{खञ्जस्}}\Afootnote{N H; \skt{स्वञ्जस्}\ A}}\skt{ भवेज्जन्तुः पङ्गुः }\edtext{\skt{सक्थ्नो\-र्द्व\-यो\-र्व्य\-धात्}}{\lemma{\skt{व्यधात्}}\Afootnote{N H; \skt{वधात् ॥}\ A}}\skt{ ॥}\pend

\pstart\versepar \ledsidenote{\passref{2.1.78}}\edtext{\edtext{\skt{प्रकामम्वेपते}}{\lemma{\skt{प्रकामम्}}\Afootnote{H; \skt{प्रक्राम}\ A; \skt{प्रकाम}\ N}}}{\lemma{\skt{वेपते}}\Afootnote{H; \skt{न्वेपते}\ A N}}\skt{ }\edtext{\skt{जन्तुः}}{\Afootnote{N H; \skt{यस्तु}\ A}}\skt{ }\edtext{\skt{खञ्जन्निव च}}{\Afootnote{H; \skt{स्वञ्जन्निव च}\ A; \skt{खंजन्निव च}\ N}}\skt{ गच्छति । }\pend
\pstart\versepar \skt{}\edtext{\edtext{\skt{कला\-य\-ख\-ञ्ज\-न्त\-म्वि\-द्या\-न्मु\-क्त\-स\-न्धि\-प्र\-ब\-न्ध\-नम्}}{\lemma{\skt{कला\-य\-ख\-ञ्ज\-न्तम्}}\Afootnote{N H; \skt{कलायस्वञ्जं तं}\ A}}}{\lemma{\skt{विद्या\-न्मु\-क्त\-स\-न्धि\-प्र\-ब\-न्ध\-नम्}}\Afootnote{N; \skt{विद्या\-न्मु\-क्त\-स\-न्धि\-प्र\-ब\-न्ध\-नम् ॥}\ A; \skt{विद्या\-न्मु\-क्त\-स\-न्धि\-प्र\-व\-न्धनं ॥}\ H}}\skt{ ॥}\pend

\pstart\versepar \ledsidenote{\passref{2.1.79}}\skt{न्यस्ते तु }\edtext{\skt{विषमे}}{\Afootnote{N H; \skt{विषमं(मे)}\ A}}\skt{ पादे }\edtext{\skt{रुजं}}{\Afootnote{N H; \skt{रुजः}\ A}}\skt{ कुर्यात्समीरणः । }\pend
\pstart\versepar \skt{}\edtext{\skt{वातकण्टक}}{\Afootnote{A H; \skt{वातकण्ठकः}\ N}}\skt{ इत्येष }\edtext{\skt{जायते}}{\Afootnote{N H; \skt{विज्ञेयः ख}\ A}}\skt{ }\edtext{\skt{खलुकाश्रयः}}{\Afootnote{N H; \skt{डु(ल)काश्रितः ॥}\ A}}\skt{ ॥}\pend

\pstart\versepar \ledsidenote{\passref{2.1.80}}\skt{पादयोः कुरुते }\edtext{\edtext{\skt{दाह\-म्पि\-त्ता\-सृ\-क्स\-हि\-तोऽ\-निलः}}{\lemma{\skt{पित्तासृक्सहितो}}\Afootnote{A H; \skt{पित्तसृक्सहितो}\ N}}}{\lemma{\skt{ऽनिलः}}\Afootnote{A H; \skt{निलः ॥}\ N}}\skt{ ॥ }\pend
\pstart\versepar \skt{}\edtext{\skt{विशे\-ष\-त\-श्च\-ङ्क्र\-मतः}}{\lemma{\skt{चङ्क्रमतः}}\Afootnote{N H; \skt{चङ्क्रमणा}\ A}}\skt{ }\edtext{\skt{पाद\-दा\-ह\-न्त\-मा\-दि\-शेत्}}{\lemma{\skt{पाददाहन्तम्}}\Afootnote{N H; \skt{त्पाददाहं तम्}\ A}}\skt{ ॥}\pend

\pstart\versepar \ledsidenote{\passref{2.1.81}}\edtext{\skt{हृष्यते}}{\Afootnote{N H; \skt{हृष्यतश्}\ A}}\skt{ चरणौ यस्य }\edtext{\skt{भवेयाताञ्च}}{\Afootnote{N H; \skt{भवतश्च प्र}\ A}}\skt{ }\edtext{\skt{सुप्तकौ}}{\Afootnote{N H; \skt{सुप्तव}\ A}}\skt{ । }\pend
\pstart\versepar \skt{}\edtext{\skt{पादहर्षः स}}{\Afootnote{N H; \skt{त् । पादहर्षः स}\ A}}\skt{ विज्ञेयः कफ\-वा\-त\-प्र\-को\-पजः ॥}\pend

\pstart\versepar \ledsidenote{\passref{2.1.82}}\skt{अंसदेशस्थितो वायुः }\edtext{\edtext{\skt{शोष\-य\-त्यं\-स\-ब\-न्धनं}}{\lemma{\skt{शोषयत्य्}}\Afootnote{N H; \skt{शोषयित्वां}\ A}}}{\lemma{\skt{अंसबन्धनं}}\Afootnote{\skt{ऽसबन्धनम् ।}\ A; \skt{अन्सवन्धनं ॥}\ N; \skt{अंगवन्धनं ।}\ H}}\skt{ । }\pend
\pstart\versepar \skt{सिराश्चाकुञ्च्य तत्रस्थो }\edtext{\skt{जन\-य\-त्य\-व\-बा\-हु\-कम्}}{\lemma{\skt{अवबाहुकम्}}\Afootnote{A; \skt{अववाहकं ।}\ N; \skt{अववाहुकं ॥}\ H}}\skt{ ॥}\pend

\pstart\versepar \ledsidenote{\passref{2.1.83}}\skt{यदा }\edtext{\skt{शब्दवहम्वायुः}}{\lemma{\skt{वायुः}}\Afootnote{N H; \skt{स्रोतो वायुर्}\ A}}\skt{ }\edtext{\skt{श्रोत्र}}{\Afootnote{H; \emph{om.}\ A; \skt{स्रोत्र}\ N}}\skt{ आवृत्य तिष्ठति । }\pend
\pstart\versepar \skt{}\edtext{\skt{शुद्धः}}{\Afootnote{A H; \skt{शुद्ध}\ N}}\skt{ श्लेष्मान्वितो }\edtext{\skt{वापि}}{\Afootnote{N H; \skt{वाऽपि}\ A}}\skt{ }\edtext{\skt{बाधिर्यन्तेन}}{\lemma{\skt{बाधिर्यन्}}\Afootnote{A; \skt{वाधिर्यन्}\ N H}}\skt{ जायते ॥}\pend

\pstart\versepar \ledsidenote{\passref{2.1.85}}\skt{आवृत्य }\edtext{\skt{वायुः}}{\Afootnote{H; \emph{om.}\ A; \skt{वायुम्}\ N}}\skt{ सकफो }\edtext{\skt{धमनीः}}{\Afootnote{N H; \skt{वायुर्धमनीः}\ A}}\skt{ }\edtext{\skt{शब्दवाहिनीः}}{\Afootnote{N H; \skt{शब्दवाहीनीः ।}\ A}}\skt{ । }\pend
\pstart\versepar \skt{}\edtext{\edtext{\edtext{\edtext{\skt{नरा\-ङ्क\-रो\-त्य\-क्रि\-य\-का\-म्मू\-क\-मि\-न्मि\-र\-ग\-द्ग\-दाम्}}{\lemma{\skt{नराङ्}}\Afootnote{N H; \skt{नरा}\ A}}}{\lemma{\skt{करोत्य्}}\Afootnote{N H; \skt{न्करोत्य्}\ A}}}{\lemma{\skt{अक्रियकाम्}}\Afootnote{H; \skt{अक्रियका}\ A; \skt{अक्रियायकाम्}\ N}}}{\lemma{\skt{मूक\-मि\-न्मि\-र\-ग\-द्ग\-दाम्}}\Afootnote{N\textsuperscript{pc}; \skt{न्मूक\-मि\-न्मि\-ण\-ग\-द्ग\-दान् ॥}\ A; \skt{मूकमिन्मिर द्गदां ।}\ N\textsuperscript{ac}; \skt{मूक\-मि\-र्म्मि\-र\-ग\-द्गदां ॥}\ H}}\skt{ ॥}\pend

\pstart\versepar \ledsidenote{\passref{2.1.84}}\edtext{\skt{शिरोग्रीवहनूः}}{\Afootnote{\skt{हनु\-श\-ङ्ख\-शि\-रोग्री}\ A; \skt{शिरो\-ग्री\-व\-ह\-तुनुः}\ N\textsuperscript{ac}; \skt{शिरोग्रीवहनुः}\ N\textsuperscript{pc} H}}\skt{ }\edtext{\skt{शंखौ}}{\Afootnote{H; \skt{वं}\ A; \skt{संखै}\ N\textsuperscript{ac}; \skt{संखैः}\ N\textsuperscript{pc}}}\skt{ यस्य }\edtext{\skt{भिन्दन्निवानिलः}}{\Afootnote{A H; \skt{भिन्दति वानिलः ।}\ N}}\skt{ ।}\skt{ }\pend
\pstart\versepar \skt{कर्ण्णयोः कुरुते शूलं }\edtext{\skt{कर्ण्णशूलः स}}{\Afootnote{H; \skt{कर्णशूलं त}\ A; \skt{कर्ण्णशूल स}\ N}}\skt{ }\edtext{\skt{उच्यते}}{\Afootnote{N H; \skt{दुच्यते ॥}\ A}}\skt{ ॥}\pend

\pstart\versepar \ledsidenote{\passref{2.1.86}}\skt{अधो या }\edtext{\skt{वेदना}}{\Afootnote{N H; \skt{वेदनां}\ A}}\skt{ याति }\edtext{\skt{वर्च्चो\-मू\-त्रा\-श\-यो\-त्थिता}}{\Afootnote{A N; \skt{वर्च्चो\-मू\-त्रा\-स\-यो\-त्थिता ॥}\ H}}\skt{ ॥ }\pend
\pstart\versepar \skt{भिन्दतीव }\edtext{\skt{गुदोपस्थौ सा}}{\Afootnote{H; \skt{गुदोपस्थं सा}\ A; \skt{गुदोपस्थैसा}\ N}}\skt{ }\edtext{\skt{तूनीनाम}}{\Afootnote{N H; \skt{तूनीत्य}\ A}}\skt{ }\edtext{\skt{नामतः}}{\Afootnote{N H; \skt{भिधीयते ॥}\ A}}\skt{ ॥}\pend

\pstart\versepar \ledsidenote{\passref{2.1.87}}\skt{गुदोपस्थोत्थिता सैव }\edtext{\edtext{\skt{प्रति\-लो\-म\-म्प्र\-धा\-विता}}{\lemma{\skt{प्रतिलोमम्}}\Afootnote{N H; \skt{प्रतिलोम}\ A}}}{\lemma{\skt{प्रधाविता}}\Afootnote{N H; \skt{विसर्पिणी ।}\ A}}\skt{ । }\pend
\pstart\versepar \skt{वेगैः }\edtext{\skt{पक्वाशयं}}{\Afootnote{A N; \skt{पक्वासयं}\ H}}\skt{ }\edtext{\skt{याव\-त्प्र\-ति\-तू\-नीति सा}}{\Afootnote{N H; \skt{याति प्रतितूनीति सा}\ A}}\skt{ }\edtext{\skt{स्मृता}}{\Afootnote{A H; \skt{स्मृताः ।}\ N}}\skt{ ॥}\pend

\pstart\versepar \ledsidenote{\passref{2.1.88}}\edtext{\edtext{\skt{साटो\-प\-म\-त्यु\-ग्र\-रु\-ज\-मा\-ध्मा\-न\-मु\-द\-र\-म्भृशं}}{\lemma{\skt{आध्मानम्}}\Afootnote{N H; \skt{आध्मातम्}\ A}}}{\lemma{\skt{भृशं}}\Afootnote{A H; \skt{भृषं ॥}\ N}}\skt{ । }\pend
\pstart\versepar \skt{}\edtext{\edtext{\skt{आध्मा\-न\-न्त\-म्वि\-जा\-नी\-या\-द्घो\-र\-म्वा\-त\-नि\-रो\-ध\-जम्}}{\lemma{\skt{आध्मानन्तम्}}\Afootnote{N H; \skt{आध्मानमि}\ A}}}{\lemma{\skt{विजा\-नी\-या\-द्घो\-रम्}}\Afootnote{N H; \skt{ति जानीयाद्घोरं}\ A}}\skt{ ॥}\pend

\pstart\versepar \ledsidenote{\passref{2.1.89}}\edtext{\skt{विमु\-क्त\-पा\-र्श्व\-हृ\-द\-य\-न्त\-दे\-वा\-मा\-श\-यो\-त्थितं}}{\lemma{\skt{तदे\-वा\-मा\-श\-यो\-त्थितं}}\Afootnote{A N; \skt{तदे\-वा\-मा\-स\-यो\-त्थितं ।}\ H}}\skt{ । }\pend
\pstart\versepar \skt{प्रत्या\-ध्मा\-न\-म्वि\-जा\-नी\-या\-त्क\-फ\-व्या\-कु\-लि\-ता\-नि\-लम्}\pend

\pstart\versepar \ledsidenote{\passref{2.1.90}}\skt{नाभे\-र\-ध\-स्ता\-त्सं\-जातः संचारी यदि वाचलः ॥}\skt{ }\pend
\pstart\versepar \skt{}\edtext{\skt{अष्ठीलावद्घनं}}{\Afootnote{A H; \skt{अष्ठीला घन}\ N\textsuperscript{ac}; \skt{अष्ठीलावं घन}\ N\textsuperscript{pc}}}\skt{ }\edtext{\edtext{\skt{ग्रन्थि\-मू\-र्ध्व\-मा\-य\-त\-मु\-न्नतं}}{\lemma{\skt{ग्रन्थिम्}}\Afootnote{A H; \skt{ग्रन्थि}\ N}}}{\lemma{\skt{ऊर्ध्वम्}}\Afootnote{A N; \skt{ऊद्ध्वम्}\ H}}\skt{ । }\pend
\pstart\versepar \skt{}\edtext{\skt{वाता\-ष्ठी\-ला\-म्वि\-जा\-नी\-या\-द्ब\-हि\-र्म्मा\-र्ग्गा\-व\-रो\-धि\-नीम्}}{\lemma{\skt{विजा\-नी\-या\-द्ब\-हि\-र्म्मा\-र्ग्गा\-व\-रो\-धि\-नीम्}}\Afootnote{N; \skt{विजा\-नी\-या\-द्ब\-हि\-र्मा\-र्ग\-व\-रो\-धि\-नीम् ॥}\ A; \skt{विजा\-नी\-या\-द्व\-हि\-र्म्मा\-र्ग्गा\-व\-रो\-धिनीं ॥}\ H}}\skt{ ॥}\edtext{}{\lemma{\skt{… अ}}\Afootnote{\emph{om.}\ A}}\edtext{}{\lemma{\skt{… अ}}\Afootnote{\emph{om.}\ H}}\pend

\pstart\versepar \ledsidenote{\passref{2.1.91}}\edtext{\skt{एतामेव}}{\lemma{\skt{एताम्}}\Afootnote{N H; \skt{एनाम्}\ A}}\skt{ }\edtext{\skt{रुजायुक्तां}}{\Afootnote{A N; \skt{रुजायुक्तं}\ H}}\skt{ }\edtext{\skt{वात\-वि\-ण्मू\-त्र\-रो\-धिनीं}}{\Afootnote{A H; \skt{वात\-वि\-ट्मू\-त्र\-रो\-धिनीं ॥}\ N}}\skt{ । }\pend
\pstart\versepar \skt{}\edtext{\skt{प्रत्यष्ठीलामिति}}{\lemma{\skt{प्रत्यष्ठीलाम्}}\Afootnote{A H; \skt{प्रत्यष्ठीलम्}\ N}}\skt{ वदेज्जठरे तिर्यगुत्थिताम् ॥}\pend

\pstart \ledsidenote{\passref{2.1.trailer}}\edtext{\skt{इति}}{\Afootnote{H; \skt{इति सुश्रु\-त\-सं\-हि\-तायां निदानस्थाने}\ A; \skt{इति ॥ निदाने ॥ १ ॥ ॥ {\symfont ❈} ॥}\ N}}\skt{ }\edtext{\skt{वात\-व्या\-धि\-नि\-दानं}}{\Afootnote{A H; \emph{om.}\ N}}\skt{ }\edtext{\skt{प्रथमः}}{\Afootnote{H; \skt{नाम प्रथमोऽध्यायः ॥१ ॥}\ A; \emph{om.}\ N}}\skt{ ॥ १ ॥}\pend

\endnumbering
\beginnumbering
\eledsection*{Adhyāya 2}
\pstart \ledsidenote{\passref{2.2.1}}\edtext{\edtext{\skt{अथा\-तोऽ\-र्श\-सा\-न्नि\-दानं}}{\lemma{\skt{अथातो}}\Afootnote{A H; \skt{अथातो}\ N}}}{\lemma{\skt{ऽर्शसान्निदानं}}\Afootnote{A H; \skt{र्शसान्निदां}\ N}}\skt{ व्याख्यास्यामः ॥}\pend

\pstart\versepar \ledsidenote{\passref{2.2.3}}\edtext{\skt{षतर्शांसि}}{\Afootnote{\skt{षडर्शांसि}\ A H; \skt{षडर्षांसि}\ N}}\skt{ भवन्ति ॥ }\pend
\pstart\versepar \skt{}\edtext{\skt{वात\-पि\-त्तो\-क\-फ\-शो\-णि\-त\-स\-न्नि\-पातैः}}{\Afootnote{H; \skt{वात\-पि\-त्त\-क\-फ\-शो\-णि\-त\-स\-न्नि\-पातैः}\ A; \skt{वात\-पि\-त्त\-क\-फ\-शो\-णि\-त\-स\-न्नि\-पा\-तस्}\ N}}\skt{ सहजानि }\edtext{\skt{चेति}}{\Afootnote{N H; \skt{चेत् ॥}\ A}}\skt{ ॥}\pend

\pstart \ledsidenote{\passref{2.2.4}}\skt{तत्रानात्मवतां यथोक्तैः प्र }\edtext{\skt{कोपनैः}}{\Afootnote{N H; \skt{कोपणैर्}\ A}}\skt{ }\edtext{\skt{प्रकुपिताः}}{\Afootnote{N H; \skt{विरु\-द्धा\-ध्य\-श\-न\-स्त्री\-प्र\-स\-ङ्गो\-त्क\-टु\-का\-स\-न\-पृ\-ष्ठ\-या\-न\-व्वे\-ग\-वि\-धा\-र\-णा\-दि\-भि\-र्वि\-शेषैः प्रकुपिता}\ A}}\skt{ }\edtext{\edtext{\skt{दोषाश्चैः}}{\lemma{\skt{दोषाश्}}\Afootnote{N H; \skt{दोषा}\ A}}\skt{ कैकशः}}{\lemma{\skt{चैः कैकशः}}\Afootnote{\skt{एकशो}\ A; \skt{एकैकशः}\ N; \skt{चैकैकशः ।}\ H}}\skt{ । }\edtext{\skt{द्वन्द्वशः}}{\Afootnote{H; \skt{द्विशः}\ A; \skt{शः}\ N}}\skt{ }\edtext{\skt{सम स्ताः}}{\Afootnote{A N\textsuperscript{pc} H; \skt{स स्ताः}\ N\textsuperscript{ac}}}\skt{ शोणितसहिता वा यथोक्तं }\edtext{\skt{प्रसृताः}}{\Afootnote{N H; \skt{प्रसृ}\ A}}\skt{ । }\edtext{\skt{विशेषतो}}{\Afootnote{H; \skt{ताः}\ A; \skt{विषेसतो}\ N}}\skt{ }\edtext{\skt{मन्दाग्नेः}}{\Afootnote{H; \emph{om.}\ A; \skt{मन्दाग्नंः}\ N}}\skt{ }\edtext{\edtext{\skt{प्रधा\-न\-ध\-म\-नी\-र\-नु\-प्र\-पद्या}}{\lemma{\skt{प्रधा\-न\-ध\-म\-नीर्}}\Afootnote{N H; \skt{प्रधानधमनी}\ A}}\skt{ धो}}{\lemma{\skt{अनुप्रपद्या धो}}\Afootnote{H; \skt{रनुप्रपद्याधो}\ A; \skt{अनुप्रपध्याधो}\ N}}\skt{ गत्वा }\edtext{\edtext{\skt{गुदमागम्य}}{\lemma{\skt{गुदम्}}\Afootnote{H; \skt{गुद}\ A N}}}{\lemma{\skt{आगम्य}}\Afootnote{H; \skt{मागम्य}\ A N}}\skt{ प्रदूष्य }\edtext{\edtext{\edtext{\skt{गुद\-व\-ली\-र्म्मा\-स\-म्प्र\-रो\-हं\-क\-न्दा\-ञ्ज\-न\-यन्ति}}{\lemma{\skt{गुद\-व\-ली\-र्म्मा\-सम्}}\Afootnote{\skt{गुदवलीर्मांस}\ A; \skt{गुदावलीर्मांस}\ N; \skt{गुद\-व\-ली\-र्म्मां\-सम्}\ H}}}{\lemma{\skt{प्ररोहंकन्दाञ्}}\Afootnote{H\textsuperscript{ac}; \skt{प्ररो\-हा\-ञ्ज\-न\-यन्ति विशेषतो मन्दाग्नेः, तथा तृण\-का\-ष्ठो\-प\-ल\-लो\-ष्ठ\-व\-स्त्रा\-दिभिः शीतो\-द\-क\-सं\-स्प\-र्श\-नाद्वा कन्दाः}\ A; \skt{प्ररोहा जन}\ N; \skt{प्ररो}\ H\textsuperscript{pc}}}}{\lemma{\skt{जनयन्ति}}\Afootnote{H\textsuperscript{ac}; \skt{परि\-वृ\-द्धि\-मा\-सा\-द\-यन्ति,}\ A; \skt{यन्ति}\ N; \skt{हं कन्दाञ्जनयन्ति}\ H\textsuperscript{pc}}}\skt{ }\edtext{\skt{तान्य\-र्शां\-सी\-त्या\-च\-क्षते}}{\lemma{\skt{अर्शांसीत्य्}}\Afootnote{A H; \skt{अर्षासीत्य्}\ N}}\skt{ ॥}\pend

\pstart \ledsidenote{\passref{2.2.5}}\skt{तत्र }\edtext{\edtext{\skt{स्थूला\-न्त्र\-प्र\-ति\-ब\-द्ध\-म\-र्ध\-प\-ञ्च\-मा\-ङ\-गुलं}}{\lemma{\skt{स्थूला\-न्त्र\-प्र\-ति\-ब\-द्धम्}}\Afootnote{A; \skt{स्थूला\-न्त्र\-प्र\-ति\-व\-द्धम्}\ H}}}{\lemma{\skt{अर्ध\-प\-ञ्च\-मा\-ङ\-गुलं}}\Afootnote{H; \skt{अर्धपञ्चाङ्गुलं}\ A}}\skt{ }\edtext{\edtext{\edtext{\edtext{\skt{गुद\-मा\-हु\-स्त\-स्मि\-न\-व\-ल\-य\-स्तिस्रो}}{\lemma{\skt{गुदम्}}\Afootnote{H; \skt{गुद}\ A}}}{\lemma{\skt{आहुस्}}\Afootnote{H; \skt{माहुः,}\ A}}}{\lemma{\skt{तस्मि\-न\-व\-ल\-यस्}}\Afootnote{\skt{तस्मिन्वलयस्}\ A H}}}{\lemma{\skt{तिस्रो}}\Afootnote{A; \skt{तिस्रो}\ H}}\skt{ }\edtext{\skt{ध्यध्द\-ङ\-गु\-ला\-न्त\-र\-स\-म्भूताः}}{\Afootnote{\skt{ऽध्या\-ङ्गु\-ला\-न्त\-र\-सं\-भूताः}\ A; \skt{द्ध्यर्धा\-ङ\-गु\-ला\-न्त\-र\-स\-म्भूताः ।}\ H}}\skt{ । }\edtext{\skt{प्रवाहिणी}}{\Afootnote{H; \skt{प्रवाहणी}\ A}}\skt{ }\edtext{\skt{विसृजनी}}{\Afootnote{H; \skt{विसर्जनी}\ A}}\skt{ संवरणी }\edtext{\skt{चेति}}{\Afootnote{H; \skt{चेति चतु\-र\-ङ्गु\-ला\-यताः; सर्वा\-स्ति\-र्य\-गे\-का\-ङ्गु\-लो\-च्छ्रिताः ॥}\ A}}\skt{ ।}\pend

\pstart \ledsidenote{\passref{2.2.6ef}}\skt{रोमान्तेभ्यो }\edtext{\skt{यवाध्यर्धो}}{\Afootnote{A N; \skt{यवाद्ध्यर्थे}\ H}}\skt{ }\edtext{\skt{गुदौष्ठः}}{\Afootnote{H; \skt{गुदौष्ठः परिकीर्तितः ॥}\ A; \skt{गुदौष्ठ ।}\ N}}\skt{ ।}\pend

\pstart \ledsidenote{\passref{2.2.7}}\skt{प्रथमा तु }\edtext{\skt{गुदौ\-ष्ठा\-द\-ङ\-गु\-ल\-मात्रे}}{\Afootnote{H; \skt{गुदौ\-ष्ठा\-द\-ङ्गु\-ल\-मात्रे}\ A N}}\pend

\pstart \ledsidenote{\passref{2.2.8}}\skt{तेषां तु }\edtext{\skt{भविष्येतां}}{\Afootnote{\skt{भविष्यतां}\ A N H}}\skt{ पूर्वरूपाणि । }\edtext{\skt{अन्नेन}}{\Afootnote{N H; \skt{अन्नेऽ}\ A}}\skt{ श्रद्धा }\edtext{\edtext{\skt{कृच्छ्रा\-त्प\-क्ति\-र\-म्ली\-कान्न}}{\lemma{\skt{कृच्छ्रात्पक्तिर्}}\Afootnote{H; \skt{कृच्छ्रात्पक्ति}\ A; \skt{कृच्छात्पक्ति}\ N}}}{\lemma{\skt{अम्लीकान्न}}\Afootnote{H; \skt{रम्लीका प}\ A; \emph{om.}\ N}}\skt{ }\edtext{\skt{विजृम्भेः}}{\Afootnote{\skt{रिदाहो विष्टम्भः}\ A; \skt{का विष्टम्भः ॥}\ N; \emph{om.}\ H}}\skt{ }\edtext{\skt{सक्ति\-स\-द\-न\-मा\-चोपः}}{\Afootnote{\skt{पिपासा सक्थि\-स\-द\-न\-मा\-टोपः}\ A; \skt{सक्थि\-स\-द\-न\-मा\-टोपः}\ N; \skt{विष्टम्भः सक्थि\-स\-द\-न\-मा\-टोपः}\ H}}\skt{ }\edtext{\edtext{\skt{कार्श्य\-मु\-द्गा\-र\-बा\-हु\-ल्य\-मक्ष्णोः}}{\lemma{\skt{कार्श्य\-मु\-द्गा\-र\-बा\-हु\-ल्यम्}}\Afootnote{\skt{कार्श्य\-मु\-द्गा\-र\-बा\-हुल्य}\ A; \skt{कार्श्य\-मु\-द्गा\-र\-वा\-हु\-ल्यम्}\ N; \skt{कार्श्य मुद्गा\-र\-बा\-हु\-ल्यम्}\ H}}}{\lemma{\skt{अक्ष्णोः}}\Afootnote{H; \skt{मक्ष्णोः}\ A; \skt{अक्ष्णो}\ N}}\skt{ }\edtext{\skt{श्वयथुर्न्त्र दपरिकर्तनं}}{\Afootnote{\skt{श्वय\-थु\-र\-न्त्र\-कू\-जनं गुद\-प\-रि\-क\-र्तन}\ A; \skt{श्वय\-थु\-र्गु\-द\-प\-रि\-क\-र्त्तनं}\ N H}}\skt{ }\edtext{\skt{कासश्वासः}}{\Afootnote{N H; \skt{माशङ्का}\ A}}\skt{ }\edtext{\skt{पाण्डु\-रो\-ग\-ग्र\-ह\-णी\-दोषः}}{\Afootnote{N H; \skt{पाण्डु\-रो\-ग\-ग्र\-ह\-णी\-दोष}\ A}}\skt{ }\edtext{\skt{कदाचित्तदा}}{\Afootnote{\skt{शोषाणां कासश्वासौ बलहानिर्म्र}\ A; \skt{कदाचित्चन्द्रा वलहा}\ N; \skt{कदाचित्तदा}\ H}}\skt{ }\edtext{\skt{बल\-हा\-नि\-रि\-न्द्रि\-ययैः दौर्बल्यम्}}{\Afootnote{\skt{मस्तन्द्रा निद्रे\-न्द्रि\-य\-दौ\-र्बल्यं च ॥}\ A; \skt{निरि\-न्द्रि\-य\-दौ\-र्वल्यं चेति ।}\ N; \skt{बल\-हा\-नि\-रि\-न्द्रि\-य\-यै\-र्व्व\-त्वा\-ञ्चेति।}\ H}}\skt{ ।}\pend

\pstart \ledsidenote{\passref{2.2.9}}\edtext{\skt{जाते श्वेतान्ये व}}{\Afootnote{\skt{जातेष्वेतान्येव}\ A H; \skt{जातेष्वेताने}\ N}}\skt{ }\edtext{\skt{लक्षणानि}}{\Afootnote{A; \skt{व लिङ्गानि}\ N; \skt{लिंगानि}\ H}}\skt{ प्रव्यक्ततराणि भवन्ति ॥}\pend

\pstart \ledsidenote{\passref{2.2.10}}\skt{तत्र }\edtext{\skt{मारुतात्परिशु ष्कारु\-ण\-वि\-व\-र्णानि}}{\Afootnote{A; \skt{मारु\-ता\-त्प\-रि\-शु\-ष्का\-रु\-ण\-व\-र्णानि}\ N H}}\skt{ }\edtext{\skt{विष\-म\-म\-ध्य\-मानि}}{\Afootnote{H; \skt{विषममध्यानि}\ A; \skt{विष ध्यामानि}\ N\textsuperscript{ac}; \skt{विष म ध्यामानि}\ N\textsuperscript{pc}}}\skt{ }\edtext{\skt{कद\-म्ब\-पु\-ष्प\-तु\-ण्डि\-के\-री\-का\-ना\-डी\-मु\-ख\-सू\-ची\-मु\-खा\-कृ\-तीनि}}{\Afootnote{H; \skt{कद\-म्ब\-पु\-ष्प\-तु\-ण्डि\-के\-री\-ना\-डी\-मु\-कु\-ल\-सू\-ची\-मु\-खा\-कृ\-तीनि च}\ A; \skt{कद\-म्व\-पु\-ष्प\-तु\-ण्डि\-के\-रि\-का\-ना\-डी\-मु\-ख\-सू\-ची\-मु\-खा\-कृ\-तीनि}\ N}}\skt{ }\edtext{\skt{तैरुपद्रुतः}}{\lemma{\skt{तैर्}}\Afootnote{N H; \skt{भवन्ति; तैर्}\ A}}\skt{ सशूलं संह\-त\-मु\-प\-वे\-श्यते । }\edtext{\skt{कटी\-प्र\-ष्ठ\-गु\-द\-म\-देषु}}{\Afootnote{\skt{कटी\-पृ\-ष्ठ\-पा\-र्श्व\-मे\-ढ्र\-गु\-द\-मा\-भि\-प्र\-दे\-शेषु}\ A; \skt{कटी\-पृ\-ष्ठ\-गु\-द\-मे\-ढ्रेषु}\ N H}}\skt{ चास्य वेदना }\edtext{\skt{गुल्मा\-स्थी\-ला\-प्ली\-हो\-द\-राणि}}{\Afootnote{\skt{भवन्ति, गुल्मा\-ष्ठी\-ला\-प्ली\-हो\-द\-राणि}\ A; \skt{गुल्मा\-ष्ठी\-ला\-प्ली\-हो\-द\-राणि}\ N H}}\skt{ चास्य }\edtext{\edtext{\skt{तन्निमित्तान्येव}}{\lemma{\skt{तन्निमित्तान्य्}}\Afootnote{A H; \skt{तन्निमित्ता}\ N}}}{\lemma{\skt{एव}}\Afootnote{A H; \emph{om.}\ N}}\skt{ }\edtext{\skt{कृष्ण\-न\-ख\-न\-य\-न\-व\-द\-न\-मू\-त्र\-पूरी षस्श्च}}{\Afootnote{\skt{भवन्ति, कृष्ण\-त्व\-ङ्न\-ख\-न\-य\-न\-द\-श\-न\-व\-द\-न\-मू\-त्र\-पु\-री\-षश्च}\ A; \skt{नेव कृष्ण\-न\-ख\-न\-य\-न\-व\-द\-न\-मू\-त्र\-पु\-री\-ष\-व\-र्ण्णश्च}\ N; \skt{कृष्ण\-न\-ख\-न\-य\-न\-व\-द\-न\-मू\-त्र\-पूरी षवर्ण्णैश्च}\ H}}\skt{ पुरुषो भवति ॥}\pend

\pstart \ledsidenote{\passref{2.2.11}}\edtext{\skt{पित्ता\-न्तु\-नी\-ला\-ग्राणि}}{\Afootnote{H; \skt{पित्तान्नीलाग्राणि}\ A; \skt{पिनीलाग्राणि}\ N}}\skt{ }\edtext{\skt{तनूनि}}{\Afootnote{A H; \skt{मनूनि}\ N}}\skt{ }\edtext{\skt{विस र्पानि}}{\Afootnote{H; \skt{विसर्पीणि}\ A N}}\skt{ }\edtext{\skt{पीत्तावभासानि}}{\Afootnote{\skt{पीतावभासानि}\ A N H}}\skt{ यकृत्प्रकाशानि }\edtext{\skt{शुक\-जि\-ह्वा\-सं\-स्थ\-नानि}}{\Afootnote{\skt{शुक\-जि\-ह्वा\-सं\-स्था\-नानि}\ A N H}}\skt{ यवमध्यानि }\edtext{\skt{जलौ\-का\-व\-क्तु\-सदृ शानि}}{\Afootnote{\skt{जलौ\-को\-व\-क्त्र\-स\-दृ\-शानि}\ A; \skt{जलौकोवक्र सदृशानि}\ N; \skt{जलौ\-का\-व\-क्त्र\-स\-दृ\-शानि}\ H}}\skt{ }\edtext{\skt{प्रक्लिन्नानि}}{\Afootnote{N H; \skt{प्रक्लिन्नानि च}\ A}}\skt{ भवन्ति । }\edtext{\skt{तैरुपदृतः}}{\lemma{\skt{उपदृतः}}\Afootnote{H; \skt{उपद्रुतः}\ A N}}\skt{ सदाहं }\edtext{\skt{सदृषितमवि }\edtext{\skt{पक्व\-म\-ति\-सा\-र्यते}}{\lemma{\skt{अतिसार्यते}}\Afootnote{N H; \skt{तिसार्यते,}\ A}}}{\lemma{\skt{सदृषितमवि पक्वम्}}\Afootnote{\skt{सरुधिरम}\ A; \skt{सरु\-धि\-र\-म\-वि\-प\-क्वम्}\ N H}}\skt{ । }\edtext{\edtext{\skt{दाह\-ज्व\-र\-पि\-पा\-सा\-श्चा\-स्यो\-प\-द्रवा}}{\lemma{\skt{दाह\-ज्व\-र\-पि\-पा\-साश्}}\Afootnote{H; \skt{ज्वर\-दा\-ह\-पि\-पासा}\ A; \skt{दाह\-ज्व\-र\-पि\-पासा}\ N}}}{\lemma{\skt{चास्योपद्रवा}}\Afootnote{H; \skt{मूर्च्छा\-श्चा\-स्यो\-प\-द्रवा}\ A; \skt{मूर्छ्योपद्रवा}\ N}}\skt{ भवन्ति । }\edtext{\edtext{\skt{तन्निमित्तान्येव}}{\lemma{\skt{तन्निमित्तान्य्}}\Afootnote{H; \emph{om.}\ A; \skt{तन्निसिन्य्}\ N}}}{\lemma{\skt{एव}}\Afootnote{N H; \emph{om.}\ A}}\skt{ }\edtext{\edtext{\skt{पीत\-न\-ख\-न\-य\-न\-व\-द\-न\-मू\-त्र\-पु\-री\-षश्च}}{\lemma{\skt{पीत\-न\-ख\-न\-य\-न\-व\-द\-न\-मू\-त्र\-पु\-री\-षश्}}\Afootnote{\skt{पीत\-त्व\-ङ्न\-ख\-न\-य\-न\-द\-श\-न\-व\-द\-न\-मू\-त्र\-पु\-री\-षश्}\ A; \skt{पीत\-न\-ख\-न\-य\-न\-व\-द\-न\-मू\-त्र\-पु\-रीष}\ N; \skt{पीत\-न\-ख\-न\-न\-य\-न\-व\-द\-न\-मू\-त्र\-पु\-री\-षश्}\ H}}}{\lemma{\skt{च}}\Afootnote{A H; \skt{र्वर्ण्ण}\ N}}\skt{ पुरुषो भवति ।}\pend

\pstart \ledsidenote{\passref{2.2.12}}\edtext{\skt{श्लेष्मजानि}}{\Afootnote{N H; \skt{श्लेष्मजानि श्वेतानि}\ A}}\skt{ महामूलानि स्थिराणि वृत्तानि स्निग्धानि }\edtext{\skt{पाण्डुनि}}{\Afootnote{\skt{पाण्डू}\ A; \skt{पांडूनि}\ N H}}\skt{ }\edtext{\skt{शुक्लावभासानि}}{\Afootnote{H; \skt{नि}\ A; \skt{मुलावतानि}\ N}}\skt{ }\edtext{\skt{करी\-र\-प\-न\-सा\-स्थि\-गो\-स्त\-ना\-का\-नाणि न}}{\Afootnote{H; \skt{करी\-र\-प\-न\-सा\-स्थि\-गो\-स्त\-ना\-का\-राणि, न}\ A; \skt{करी\-र\-प\-ण\-सा\-स्थि\-गो\-स्त\-ना\-का\-राणि । न}\ N}}\skt{ भिद्यन्ते न }\edtext{\skt{स्रवन्ति च}}{\Afootnote{N H; \skt{स्रवन्ति}\ A}}\skt{ }\edtext{\skt{कण्डुबहुलानि च}}{\Afootnote{\skt{कण्डूबहुलानि च}\ A H; \skt{कण्डवहुलानि च}\ N}}\skt{ भवन्ति । }\edtext{\skt{तैरुपदृतः}}{\lemma{\skt{उपदृतः}}\Afootnote{H; \skt{उपद्रुतः}\ A N}}\skt{ }\edtext{\edtext{\edtext{\skt{सश्ले\-ष्म\-क\-म\-न\-ल्प\-मा\-मं\-वसा}}{\lemma{\skt{सश्लेष्मकम्}}\Afootnote{N H; \skt{सश्लेष्माण}\ A}}}{\lemma{\skt{अनल्पम्}}\Afootnote{H; \skt{मनल्पं}\ A; \skt{अनल्प}\ N}}\skt{ मेदः । }\edtext{\skt{प्रका\-श\-म\-ति\-सा\-र्यते}}{\lemma{\skt{अतिसार्यते}}\Afootnote{H; \skt{मतिसार्यते}\ A N}}}{\lemma{\skt{आमंवसा मेदः । प्रकाशम्}}\Afootnote{H; \skt{मांस\-धा\-व\-न\-प्र\-काश}\ A; \skt{मांस\-व\-सा\-मे\-दा\-प्र\-काश}\ N}}\skt{ । }\edtext{\skt{सोभू\-शी\-त\-ज्व\-रा\-रो\-च\-का\-वि\-पा\-क\-शिरो गौरवगातु}}{\Afootnote{\skt{शोफ\-शी\-त\-ज्व\-रा\-रो\-च\-का\-वि\-पा\-क\-शि\-रो\-गौर}\ A; \skt{शोफशीत ज्वरा\-रो\-च\-का\-वि\-पा\-क\-शि\-रो\-गौ\-र\-व\-ग्रात्र}\ N; \skt{सोफ\-शी\-त\-ज्व\-रा\-रो\-च\-का\-वि\-पा\-क\-शिरो गौरवगात्र}\ H}}\skt{ }\edtext{\skt{सदनानि}}{\Afootnote{N H; \skt{वाणि}\ A}}\skt{ }\edtext{\skt{यास्य}}{\Afootnote{\skt{चास्य}\ A N H}}\skt{ तन्निमित्तान्येव }\edtext{\skt{शुक्ल\-न\-ख\-न\-य\-न\-व\-द\-न\-मू\-त्र\-पु\-री\-ष\-व\-ण्णश्च}}{\Afootnote{\skt{भवन्ति, शुक्ल\-त्व\-ङ्ग\-ख\-न\-य\-न\-द\-श\-न\-व\-द\-न\-मू\-त्र\-पु\-री\-षश्च}\ A; \skt{शुक्ले न नख\-न\-य\-न\-व\-द\-न\-मू\-त्र\-पु\-री\-ष\-व\-र्ण्णश्च}\ N; \skt{शुक्ल\-न\-ख\-न\-य\-न\-व\-द\-न\-मू\-त्र\-पु\-री\-ष\-व\-र्ण्णश्च}\ H}}\skt{ पुरुषो भवति ।}\pend

\pstart \ledsidenote{\passref{2.2.13}}\skt{रक्तजानि }\edtext{\skt{तुन्यत्रोध प्ररोहविद्रुम}}{\Afootnote{H; \skt{न्यग्रो\-ध\-प्र\-रो\-ह\-वि\-द्रुम}\ A; \skt{तु न्यग्रो\-ध\-प्र\-मे\-ह\-वि\-द्रुम}\ N}}\skt{ }\edtext{\skt{काक\-ण\-न्ति\-का\-फ\-ल\-स\-दृ\-शानि}}{\Afootnote{A H; \skt{काक\-न\-न्ति\-क\-फ\-स\-दृशाः}\ N}}\skt{ पित्तलक्षणानि }\edtext{\skt{चेति}}{\Afootnote{N H; \skt{च,}\ A}}\skt{ । यदा }\edtext{\skt{तु}}{\Afootnote{N H; \skt{ऽव}\ A}}\skt{ }\edtext{\skt{गाढ\-पु\-री\-ष\-प्र\-पी\-डि\-तानि}}{\Afootnote{N H; \skt{गाढ\-पु\-री\-ष\-पी\-डि\-तानि}\ A}}\skt{ भवन्ति । }\edtext{\skt{तदात्यर्थं}}{\Afootnote{\skt{तदाऽत्यर्थं}\ A; \skt{तदात्यर्थन्}\ N; \skt{तद आत्यर्थं}\ H}}\skt{ }\edtext{\edtext{\skt{दुष्ट\-म\-न\-ल्प\-म\-सृ\-क्स\-हसा}}{\lemma{\skt{अनल्पम्}}\Afootnote{A; \skt{उश्ण}\ N; \skt{उष्णम्}\ H}}}{\lemma{\skt{असृक्सहसा}}\Afootnote{A H; \skt{मसृक्सहसा}\ N}}\skt{ विसृजन्ति । तस्य चा ति प्रवृत्तौ शोणिताति }\edtext{\edtext{\skt{योगो\-प\-द्र\-वा\-श्चास्य}}{\lemma{\skt{योगोपद्रवाश्}}\Afootnote{H; \skt{योगोपद्र}\ A N}}}{\lemma{\skt{चास्य}}\Afootnote{H; \skt{वा}\ A N}}\skt{ भवन्ति ॥}\pend

\pstart \ledsidenote{\passref{2.2.14}}\skt{सन्निपातजानि }\edtext{\skt{सर्व\-ल\-क्ष\-ण\-यु\-क्तानि}}{\Afootnote{N H; \skt{सर्व\-दो\-ष\-ल\-क्ष\-ण\-यु\-क्तानि ॥}\ A}}\skt{ ।}\pend

\pstart \ledsidenote{\passref{2.2.15}}\skt{सहजानि }\edtext{\skt{दुष्ट\-शु\-क्र\-शो\-णि\-त\-नि\-मि\-त्तानि}}{\Afootnote{N H; \skt{दुष्ट\-शो\-णि\-त\-शु\-क्र\-नि\-मि\-त्तानि,}\ A}}\skt{ । तेषां दोषत एव }\edtext{\skt{प्रसादनं}}{\Afootnote{H; \skt{प्रसाधनं}\ A; \skt{प्रसादन}\ N}}\skt{ }\edtext{\skt{कर्तव्यं}}{\Afootnote{A H; \skt{न्कर्तव्य}\ N}}\skt{ । }\edtext{\edtext{\skt{विशेषतश्चात्र}}{\lemma{\skt{विशेषतश्}}\Afootnote{A H; \skt{विशेषश्}\ N}}}{\lemma{\skt{चात्र}}\Afootnote{N H; \skt{चैतानि}\ A}}\skt{ दुर्दर्शनानि }\edtext{\skt{परुषारुण}}{\Afootnote{H; \skt{परुषाणि}\ A; \skt{परुणाषाण}\ N}}\skt{ }\edtext{\skt{पाणुनि}}{\Afootnote{\skt{पांसूनि}\ A; \skt{पाण्डूनि}\ N; \skt{पाण्दूनि}\ H}}\skt{ }\edtext{\edtext{\skt{दारु\-णा\-न्य\-न्त\-र्मु\-खानि}}{\lemma{\skt{दारुणान्य्}}\Afootnote{A H; \skt{दारुणा}\ N}}\skt{ । }\edtext{\skt{तैरुपविद्रुतः}}{\lemma{\skt{उपविद्रुतः}}\Afootnote{H; \skt{उपद्रुतः}\ A N}}}{\lemma{\skt{अन्तर्मुखानि । तैर्}}\Afootnote{A H; \skt{न्यन्तर्मुखा नि । तैर्}\ N}}\skt{ }\edtext{\edtext{\skt{कृषोऽ\-ल्प\-भु\-क्सि\-रा\-स\-न्तत}}{\lemma{\skt{कृषो}}\Afootnote{H; \skt{कृशो}\ A; \skt{कृशो}\ N}}}{\lemma{\skt{ऽल्प\-भु\-क्सि\-रा\-स\-न्तत}}\Afootnote{A H; \skt{ल्पभु\-क्शि\-रा\-स\-न्नत}\ N}}\skt{ }\edtext{\edtext{\skt{गात्रोऽल्पप्रजः}}{\lemma{\skt{गात्रो}}\Afootnote{A H; \skt{गात्रा}\ N}}}{\lemma{\skt{ऽल्पप्रजः}}\Afootnote{A H; \skt{ल्पप्रजः ।}\ N}}\skt{ }\edtext{\skt{क्षिणरेतश्च}}{\Afootnote{H; \skt{क्षीणरेताः क्षा}\ A; \skt{क्षीणरेताश्}\ N}}\skt{ }\edtext{\skt{भवति}}{\Afootnote{H; \skt{मस्वरः क्रोध\-नोऽ\-ल्पा\-ग्नि\-प्राणः परमलसश्च, तथा घ्राण\-शि\-रोऽ\-क्षि\-ना\-सा\-श्र\-व\-ण\-रोगी, सत\-त\-म\-न्त्र\-कू\-जा\-टो\-प\-हृ\-द\-यो\-प\-ले\-पा\-रो\-च\-क\-प्र\-भृ\-तिभिः पीड्यते ॥}\ A; \skt{च भवति ॥}\ N}}\skt{ ॥}\pend

\pstart\versepar \ledsidenote{\passref{2.2.16}}\edtext{\skt{भवति}}{\Afootnote{A H; \emph{om.}\ N}}\skt{ }\edtext{\skt{चात्र}}{\Afootnote{A H; \skt{श्लोक ॥}\ N}}\skt{ । }\pend
\pstart\versepar \skt{}\edtext{\edtext{\skt{बाह्य\-म\-ध्य\-म\-यो\-र्वल्योः}}{\lemma{\skt{बाह्य\-म\-ध्य\-म\-योर्}}\Afootnote{H; \skt{बाह्यमध्य}\ A; \skt{वाह्यमध्यमयो}\ N}}}{\lemma{\skt{वल्योः}}\Afootnote{H; \skt{वलि}\ A; \skt{वल्यो}\ N}}\skt{ }\edtext{\skt{प्रति\-कु\-र्या\-द्भि\-ष\-ग्वरः}}{\Afootnote{\skt{स्थानां प्रति\-कु\-र्या\-द्भि\-ष\-ग्वरः ।}\ A; \skt{प्रति\-कु\-र्या\-द्भि\-ष\-ग्वरः ।}\ N; \skt{प्रति\-कु\-र्या\-द्भि\-ष\-क्वरः ।}\ H}}\skt{ । }\pend
\pstart\versepar \skt{}\edtext{\edtext{\skt{वल्या\-म\-भ्य\-न्त\-रायां}}{\lemma{\skt{वल्याम्}}\Afootnote{N H; \emph{om.}\ A}}\skt{ तु}}{\lemma{\skt{अभ्यन्तरायां तु}}\Afootnote{H; \skt{अन्त\-र्व\-लि\-स\-मु\-त्थानां}\ A; \skt{अभ्यन्तरो यान्तु}\ N}}\skt{ }\edtext{\skt{प्रत्या\-ख्या\-या\-च\-रे\-त्क्रियां}}{\Afootnote{A H; \skt{प्रत्या\-ख्या\-या\-व\-चे\-न्क्रिया ॥}\ N}}\pend

\pstart \ledsidenote{\passref{2.2.17}}\edtext{\skt{प्रकुपितास्तु}}{\Afootnote{A H; \skt{प्रकुपिता}\ N}}\skt{ }\edtext{\skt{दोषा}}{\Afootnote{A H; \skt{स्त्य्दोषाः}\ N}}\skt{ }\edtext{\skt{मेढ्र\-म\-भि\-प्र\-पन्ना}}{\Afootnote{\skt{मेढ्र\-म\-भि\-प्र\-पन्ना}\ A H; \skt{मेढ्रु\-म\-भि\-प्र\-पन्ना}\ N}}\skt{ }\edtext{\skt{मांसशोणिते}}{\Afootnote{A H; \skt{मांसशोणित}\ N}}\skt{ }\edtext{\skt{प्रदूष्य}}{\Afootnote{A H; \skt{प्रदुष्य}\ N}}\skt{ }\edtext{\skt{कण्डूं}}{\Afootnote{A H; \skt{कण्डू}\ N}}\skt{ }\edtext{\skt{जनयन्ति}}{\Afootnote{A H; \skt{न्जनयन्ति}\ N}}\skt{ }\edtext{\skt{ततः}}{\Afootnote{A; \skt{तस्मि}\ N H}}\skt{ }\edtext{\edtext{\skt{कण्डूयनात्क्षतं}}{\Afootnote{A; \skt{न्कण्डू}\ N H}}\skt{ }\edtext{\skt{समुपजायते}}{\Afootnote{A; \skt{क्ष}\ N}}\skt{ }\edtext{\skt{तस्मिंश्च}}{\Afootnote{A; \skt{त}\ N; \skt{ते}\ H}}\skt{ }\edtext{\skt{क्षते}}{\Afootnote{A H; \skt{न्क्षते}\ N}}}{\lemma{\skt{यनात्… समुपजायते}}\Afootnote{\emph{om.}\ H}}\skt{ }\edtext{\skt{दुष्टमांसजाः}}{\Afootnote{A H; \skt{दुष्टं मास}\ N}}\skt{ प्ररोहाः }\edtext{\skt{पिच्छि\-ल\-ल\-रु\-धि\-र\-स्रा\-विणो}}{\Afootnote{A; \skt{पिच्छि\-ल\-रु\-धि\-र\-स्रा\-विणो}\ N H}}\skt{ जायन्ते }\edtext{\edtext{\edtext{\skt{कूर्च\-कि\-नोऽ\-भ्य\-न्त\-र\-मु\-प\-रि\-ष्टाद्वा}}{\lemma{\skt{कूर्चकिनो}}\Afootnote{A H; \skt{कुर्चकिनो}\ N}}}{\lemma{\skt{ऽभ्यन्तरम्}}\Afootnote{A H; \skt{भ्यन्तरम्}\ N}}\skt{ ते तु}}{\lemma{\skt{उपरिष्टाद्वा ते तु}}\Afootnote{H; \skt{उपरिष्टाद्वा, ते तु}\ A; \skt{उपपरिष्ठाद्वा । ते तु}\ N}}\skt{ }\edtext{\skt{शेफो}}{\Afootnote{A; \skt{सहं}\ N; \skt{शेफम्}\ H}}\skt{ }\edtext{\skt{विना\-श\-य\-न्त्यु\-प\-घ्नन्ति च}}{\Afootnote{\skt{विना\-श\-य\-न्त्यु\-प\-घ्नन्ति च}\ A; \skt{विना\-श\-य\-न्त्यु\-प\-घ्नन्ति}\ N; \skt{विना\-श\-य\-न्त्यु\-प\-घ्नन्ति}\ H}}\skt{ पुंस्त्वं }\edtext{\skt{योनि\-म\-भि\-प्र\-पन्नाः}}{\Afootnote{\skt{योनि\-म\-भि\-प्र\-पन्नाः}\ A; \skt{योनि\-म\-भि\-प्र\-पन्}\ N H}}\skt{ }\edtext{\edtext{\skt{सुकु\-मा\-रा\-न्दु\-र्ग\-न्धा\-न्पि\-च्छि\-ल\-रु\-धि\-र\-स्रा\-वि\-ण\-श्छ\-त्रा\-का\-रा\-न्क\-री\-रा\-ञ्ज\-न\-यन्ति}}{\lemma{\skt{सुकुमारान्}}\Afootnote{A; \skt{ना}\ N H}}}{\lemma{\skt{दुर्ग\-न्धा\-न्पि\-च्छि\-ल\-रु\-धि\-र\-स्रा\-वि\-ण\-श्छ\-त्रा\-का\-रा\-न्क\-री\-राञ्}}\Afootnote{\skt{दुर्ग\-न्धा\-न्पि\-च्छि\-ल\-रु\-धि\-र\-स्रा\-वि\-ण\-श्छ\-त्रा\-का\-रा\-न्क\-री\-राञ्}\ A; \skt{दुर्गन्धि पिच्छिलास्राविणः च्छत्राकारां प्ररोहां}\ N H}}\skt{ ते तु }\edtext{\edtext{\skt{योनि\-मु\-प\-घ्न\-न्त्या\-र्तवं}}{\lemma{\skt{उपघ्नन्त्य्}}\Afootnote{A H; \skt{उपघ्नन्}\ N}}\skt{ च}}{\lemma{\skt{आर्तवं च}}\Afootnote{A H; \skt{त्यार्तवश्च ।}\ N}}\skt{ । नाभि\-म\-भि\-प्र\-पन्नाः }\edtext{\edtext{\skt{सुकु\-मा\-रा\-न्दु\-र्ग\-न्धा\-न्पि\-च्छि\-ला\-न्ग\-ण्डू\-प\-द\-मु\-ख\-स\-दृ\-शा\-न्क\-री\-रा\-ञ्ज\-न\-यन्ति}}{\lemma{\skt{सुकुमारान्}}\Afootnote{A; \skt{सुकुमारा}\ N H}}}{\lemma{\skt{दुर्ग\-न्धा\-न्पि\-च्छि\-ला\-न्ग\-ण्डू\-प\-द\-मु\-ख\-स\-दृ\-शा\-न्क\-री\-राञ्}}\Afootnote{A; \skt{णां पिच्छिलं गण्डू\-प\-द\-मु\-ख\-स\-दृशां करीरां}\ N; \skt{न्पिच्छि\-ला\-न्ग\-ण्डू\-प\-द\-मु\-ख\-स\-दृशां करीराञ्}\ H}}\skt{ त एवोर्ध्वमागताः }\edtext{\skt{श्रोत्रा\-क्षि\-घ्रा\-ण\-व\-द\-ने\-ष्व\-र्शां\-स्यु\-प\-नि\-र्व\-र्त\-यन्ति}}{\lemma{\skt{उप\-नि\-र्व\-र्त\-यन्ति}}\Afootnote{A; \skt{उपजनयन्ति}\ N H}}\skt{ तत्र कर्णजेषु }\edtext{\skt{बाधिर्यं}}{\Afootnote{H; \skt{बाधिर्यं}\ A; \skt{वाधिर्यं}\ N}}\skt{ }\edtext{\skt{शूलं}}{\Afootnote{A H; \skt{शूलः}\ N}}\skt{ }\edtext{\skt{पूतिकर्णता च}}{\Afootnote{A; \skt{पूतिकर्ण्णताश्च ।}\ N; \skt{पूतिकर्णताञ्च}\ H}}\skt{ नेत्रजेषु वर्त्मा वरोधो वेदना स्रावो }\edtext{\skt{दर्शननाशश्च,}}{\lemma{\skt{दर्शननाशश्च}}\Afootnote{A N; \skt{दर्श\-न\-वि\-ना\-शश्च ॥}\ H}}\skt{ }\edtext{\skt{घ्राणजेषु}}{\Afootnote{A H; \skt{ग्राणजेषु}\ N}}\skt{ }\edtext{\edtext{\skt{प्रति\-श्या\-योऽ\-ति\-मात्रं}}{\lemma{\skt{प्रतिश्यायो}}\Afootnote{A H; \skt{प्रतिश्यायो}\ N}}}{\lemma{\skt{ऽतिमात्रं}}\Afootnote{A H; \skt{तिमात्रं}\ N}}\skt{ क्षवथुः }\edtext{\skt{कृच्छ्रोच्छ्वासता}}{\Afootnote{A; \skt{कृच्छ्रोच्छ्वासित ।}\ N; \skt{कृच्छ्रोच्छ्वासिता}\ H}}\skt{ }\edtext{\skt{पूतिनस्यं}}{\Afootnote{A; \skt{पूतिनास्य}\ N H}}\skt{ }\edtext{\skt{सानु\-ना\-सि\-क\-वा\-क्यत्वं}}{\Afootnote{A; \skt{शिरो तु नासिक वाक्यश्}\ N; \skt{शिरो\-रु\-क्सा\-नु\-ना\-सि\-क\-वा\-क्यञ्}\ H}}\skt{ }\edtext{\skt{शिरोदुःखं च}}{\Afootnote{A; \skt{च भवति ॥}\ N H}}\skt{ }\edtext{\skt{वक्त्रजेषु}}{\Afootnote{A H; \skt{वक्रेषु}\ N}}\skt{ }\edtext{\edtext{\edtext{\skt{कण्ठौ\-ष्ठ\-ता\-लू\-ना\-म\-न्य\-त\-म\-स्मिं\-स्तै\-र्ग\-द्ग\-द\-वा\-क्यता}}{\lemma{\skt{कण्ठौ\-ष्ठ\-ता\-लू\-नाम्}}\Afootnote{A; \skt{कण्ठाषु तालु}\ N; \skt{कण्ठोष्ठतालूष्व्}\ H}}}{\lemma{\skt{अन्यतम्}}\Afootnote{A; \skt{ष्ठन्यत}\ N; \skt{अन्यत}\ H}}}{\lemma{\skt{अस्मिंस्तैर्}}\Afootnote{A; \skt{मेषु तैर्}\ N H}}\skt{ }\edtext{\skt{रसाज्ञानं}}{\Afootnote{A; \skt{रसावरोण}\ N; \skt{रसावरोधो}\ H}}\skt{ मुखरोगाश्च भवन्ति ॥}\pend

\pstart \ledsidenote{\passref{2.2.18}}\skt{व्यानस्तु प्रकुपितः श्लेष्माणं परिगृह्य }\edtext{\skt{बहिः}}{\Afootnote{A; \skt{वहि}\ N; \skt{बहि}\ H}}\skt{ }\edtext{\skt{स्थिराणि}}{\Afootnote{A; \skt{र्ध्-अस्थिरं}\ N; \skt{र्द्धा स्थिराङ्}\ H}}\skt{ }\edtext{\skt{कीलवदर्शांसि}}{\Afootnote{A; \skt{कीलवदर्शान्ति}\ N; \skt{कीलवदर्शोऽभि}\ H}}\skt{ }\edtext{\skt{निर्वर्तयति}}{\Afootnote{A N; \skt{निवर्त्तयति}\ H}}\skt{ }\edtext{\skt{तानि}}{\Afootnote{A; \skt{तश्}\ N; \skt{तञ्}\ H}}\skt{ }\edtext{\edtext{\skt{चर्म\-की\-ला\-न्य\-र्शां\-सी\-त्या\-च\-क्षते}}{\lemma{\skt{चर्मकीलान्य्}}\Afootnote{A; \skt{चर्मकीलसि}\ N; \skt{चर्मकी}\ H}}}{\lemma{\skt{अर्शांसीत्य्}}\Afootnote{A; \skt{त्य्}\ N; \skt{लमीत्य्}\ H}}\skt{ ॥}\pend

\pstart\versepar \ledsidenote{\passref{2.2.19}}\edtext{\skt{भवन्ति}}{\Afootnote{A H; \emph{om.}\ N}}\skt{ }\edtext{\skt{चात्र}}{\Afootnote{A H; \emph{om.}\ N}}\skt{ }\pend
\pstart\versepar \skt{}\edtext{\skt{तस्य}}{\Afootnote{N H; \skt{तेषु}\ A}}\skt{ }\edtext{\edtext{\skt{तोदोऽथ}}{\lemma{\skt{तोदो}}\Afootnote{H; \skt{कीले}\ A; \skt{तो}\ N}}}{\lemma{\skt{ऽथ}}\Afootnote{H; \skt{षु}\ A; \skt{दोष}\ N}}\skt{ }\edtext{\skt{पारूष्यं}}{\Afootnote{H; \skt{निस्तोदो}\ A; \skt{रुम्}\ N}}\skt{ मारु\-ते\-नो\-प\-जा\-यते । }\pend
\pstart\versepar \skt{श्लेष्मणा तु सवर्ण्णत्वं }\edtext{\skt{ग्रथितत्वञ्च}}{\Afootnote{H; \skt{ग्रन्थित्वं च वि}\ A; \skt{ग्रंथितत्वश्च}\ N}}\skt{ निर्दिशेत् ॥}\pend

\pstart\versepar \ledsidenote{\passref{2.2.20}}\edtext{\skt{पित्त\-शो\-णि\-त\-ज\-रौक्ष्यं}}{\Afootnote{N H; \skt{पित्तशोणितजं रौक्ष्यं}\ A}}\skt{ }\edtext{\skt{कृष्णत्वं}}{\Afootnote{A; \skt{कृष्णता}\ N H}}\skt{ }\edtext{\skt{श्लक्ष्णता}}{\Afootnote{A N H\textsuperscript{pc}; \skt{श्लक्ष्नता}\ H\textsuperscript{ac}}}\skt{ तथा । }\pend
\pstart\versepar \skt{}\edtext{\skt{समु\-दी\-र्ण\-ख\-रत्वं च}}{\Afootnote{A H; \skt{समु\-दी\-र्ण्ण\-ख\-र\-त्वश्च}\ N}}\skt{ चर्मकीलस्य }\edtext{\skt{लक्षणम्॥}}{\lemma{\skt{लक्षणम्}}\Afootnote{A H; \skt{लक्षणां ॥}\ N}}\pend

\pstart\versepar \ledsidenote{\passref{2.2.21}}\edtext{\skt{अर्शसां}}{\Afootnote{A; \skt{दुर्नास्नां}\ N; \skt{दुर्न्नाम्नां}\ H}}\skt{ लक्षणं }\edtext{\skt{व्यासादुक्तं}}{\Afootnote{A; \skt{व्यासादुक्तं}\ N H}}\skt{ }\edtext{\edtext{\skt{सामान्यतस्तु}}{\lemma{\skt{सामान्यतस्}}\Afootnote{A; \skt{यत्च समासतः}\ N H}}}{\lemma{\skt{तु}}\Afootnote{A; \skt{तत्सर्व्वं प्राग\-भि\-हि\-ता\-त्सा\-ध\-यते चिकि}\ N; \skt{सन्निपात समुत्थानि सहजानि च वर्ज्ज}\ H}}\skt{ यत् । }\pend
\pstart\versepar \skt{तत्सर्वं }\edtext{\skt{प्राग्वि\-नि\-र्दि\-ष्टा\-त्सा\-ध\-ये\-द्भि\-षजां}}{\Afootnote{A; \emph{om.}\ N H}}\skt{ }\edtext{\skt{वरः}}{\Afootnote{A; \emph{om.}\ N H}}\skt{ ॥}\edtext{}{\lemma{\skt{यत् । तत्सर्वं}}\Afootnote{A; \skt{त्सकः ॥}\ N; \skt{येत् ।}\ H}}\pend

\pstart\versepar \ledsidenote{\passref{2.2.22}}\edtext{\skt{अर्शःसु}}{\Afootnote{A N; \emph{om.}\ H}}\skt{ }\edtext{\skt{दृश्यते}}{\Afootnote{A N; \emph{om.}\ H}}\skt{ }\edtext{\skt{रूपं}}{\Afootnote{A N; \emph{om.}\ H}}\skt{ }\edtext{\skt{यदा}}{\Afootnote{A N; \emph{om.}\ H}}\skt{ }\edtext{\skt{दोषद्वयस्य तु}}{\Afootnote{A N; \emph{om.}\ H}}\skt{ । }\pend
\pstart\versepar \skt{}\edtext{\skt{संसर्गं तं}}{\Afootnote{A N; \emph{om.}\ H}}\skt{ }\edtext{\skt{विजा\-नी\-या\-त्सं\-सर्गः स च}}{\Afootnote{A; \skt{विजा\-नी\-या\-त्सं\-सर्गः षड्वि}\ N; \emph{om.}\ H}}\skt{ }\edtext{\skt{षड्विधिः}}{\Afootnote{A; \skt{धसः ॥}\ N; \emph{om.}\ H}}\skt{ ॥}\pend

\pstart\versepar \ledsidenote{\passref{2.2.23}}\edtext{\skt{त्रिदो\-षा\-ण्य\-ल्प\-लि\-ङ्गानि}}{\Afootnote{\skt{त्रिदो\-षा\-ण्य\-ल्प\-लि\-ङ्गानि}\ A; \emph{om.}\ H}}\skt{ }\edtext{\skt{याप्यानि तु}}{\Afootnote{A; \emph{om.}\ H}}\skt{ विनिर्दिशेत् । }\pend
\pstart\versepar \skt{द्वन्द्वजानि }\edtext{\skt{द्वितीयायां}}{\Afootnote{A; \emph{om.}\ H}}\skt{ }\edtext{\skt{वलौ}}{\Afootnote{A; \emph{om.}\ H}}\skt{ }\edtext{\edtext{\skt{यान्याश्रितानि}}{\lemma{\skt{यान्य्}}\Afootnote{A; \emph{om.}\ H}}}{\lemma{\skt{आश्रितानि}}\Afootnote{A; \emph{om.}\ H}}\skt{ }\edtext{\skt{च}}{\Afootnote{A; \emph{om.}\ H}}\skt{ ॥}\edtext{}{\lemma{\skt{विनिर्दिशेत् । द्वन्द्वजानि}}\Afootnote{A; \emph{om.}\ H}}\pend

\pstart\versepar \ledsidenote{\passref{2.2.24}}\edtext{\skt{कृच्छ्रसाध्यानि}}{\Afootnote{A; \emph{om.}\ H}}\skt{ }\edtext{\edtext{\skt{तान्याहुः}}{\lemma{\skt{तान्य्}}\Afootnote{A; \emph{om.}\ H}}}{\lemma{\skt{आहुः}}\Afootnote{A; \emph{om.}\ H}}\skt{ }\edtext{\skt{परि\-सं\-व\-त्स\-राणि च}}{\Afootnote{A; \emph{om.}\ H}}\skt{ । }\pend
\pstart\versepar \skt{}\edtext{\skt{सन्नि\-पा\-त\-स\-मु\-त्थानि}}{\Afootnote{A N; \emph{om.}\ H}}\skt{ }\edtext{\skt{सहजानि तु}}{\Afootnote{A N; \emph{om.}\ H}}\skt{ }\edtext{\skt{वर्जयेत्}}{\Afootnote{A; \skt{वर्जयेत् । वल्यः सर्व्वा स्यू येषां हि दुर्ना\-म\-भि\-रु\-प\-द्रुताः ॥}\ N; \emph{om.}\ H}}\skt{ ॥}\edtext{}{\lemma{\skt{… सन्}}\Afootnote{\emph{om.}\ N}}\pend

\pstart\versepar \ledsidenote{\passref{2.2.25}}\edtext{\skt{सर्वाः }\edtext{\skt{स्युर्वलयो}}{\lemma{\skt{वलयो}}\Afootnote{A; \skt{र्व्वाश्च ये}\ H}}}{\lemma{\skt{सर्वाः स्युर्}}\Afootnote{A; \skt{वल्यः स}\ H}}\skt{ }\edtext{\skt{येषां}}{\Afootnote{A; \skt{षां हि}\ H}}\skt{ दुर्ना\-म\-भि\-रु\-प\-द्रुताः । }\pend
\pstart\versepar \skt{तैस्तु }\edtext{\skt{प्रतिहतो}}{\Afootnote{A H; \skt{प्रतिहितो}\ N}}\skt{ }\edtext{\edtext{\skt{वायुरपानः}}{\lemma{\skt{वायुर्}}\Afootnote{A H; \skt{वायु}\ N}}}{\lemma{\skt{अपानः}}\Afootnote{A H; \skt{रपानस्}\ N}}\skt{ सन्निवर्तते ॥}\edtext{}{\lemma{\skt{… तैस्}}\Afootnote{\emph{om.}\ N}}\pend

\pstart \ledsidenote{\passref{2.2.26}}\skt{ततो व्यानेन सङ्गम्य }\edtext{\skt{ज्योतिर्मृद्गाति}}{\lemma{\skt{मृद्गाति}}\Afootnote{A; \emph{om.}\ N H}}\skt{ }\edtext{\skt{देहिनाम्}}{\Afootnote{A; \skt{अस्य प्रमर्द्दतीति}\ N H}}\skt{ ॥}\pend

\pstart \ledsidenote{\passref{2.2.trailer}}\edtext{\edtext{\skt{इति सुश्रु\-त\-सं\-हि\-तायां }\edtext{\skt{निदा\-न\-स्था\-नेऽ\-र्शो\-नि\-दानं}}{\lemma{\skt{ऽर्शोनिदानं}}\Afootnote{A; \skt{अर्शोनिदा}\ H}}}{\lemma{\skt{…ऽ}}\Afootnote{\emph{om.}\ H}}\skt{ }\edtext{\skt{नाम}}{\Afootnote{A; \skt{ने}\ H}}\skt{ }\edtext{\skt{द्वितीयोऽध्यायः}}{\lemma{\skt{ऽध्यायः}}\Afootnote{A; \skt{॥ २ ॥ {\symfont ❈} ॥}\ N; \emph{om.}\ H}}}{\Afootnote{\emph{om.}\ N}}\skt{ ॥२ ॥}\pend

\endnumbering
\beginnumbering
\eledsection*{Adhyāya 3}
\pstart \ledsidenote{\passref{2.3.1}}\skt{अथातो श्मरीणान्निदानं }\edtext{\skt{व्याख्यास्यामः}}{\Afootnote{A H; \emph{om.}\ N}}\skt{ ॥}\pend

\pstart \ledsidenote{\passref{2.3.8}}\skt{तत्र श्लेष्माश्मरी }\edtext{\skt{श्लेष्म\-ल\-म\-न्न\-म\-भ्य\-स\-तोऽ\-त्य\-र्थ\-मु\-प\-लि\-प्याधः}}{\lemma{\skt{श्लेष्म\-ल\-म\-न्न\-म\-भ्य\-सतो}}\Afootnote{\skt{श्लेष्म\-ल\-म\-न्न\-म\-भ्य\-व\-ह\-रतो}\ A; \skt{श्लेष्म\-ल\-म\-न्न\-म\-भ्य\-शतो}\ N; \skt{श्लेष्म\-ल\-म\-न्न\-म\-भ्य\-सतो}\ H}}\skt{ }\edtext{\skt{परिवृद्धिं}}{\Afootnote{N H; \skt{परिवृद्धिं परिवृद्धिं}\ A}}\skt{ प्राप्य बस्ति\-मु\-ख\-म\-धि\-ष्ठाय }\edtext{\skt{श्रोतो}}{\Afootnote{H; \skt{स्रोतो}\ A; \skt{स्रोत्रो}\ N}}\skt{ }\edtext{\skt{निरुपणद्धि}}{\Afootnote{\skt{निरुणद्धि}\ A N H}}\skt{ । तस्य }\edtext{\skt{मूत्र\-प्र\-ति\-घा\-ता\-द्दा\-ल्यते}}{\Afootnote{A H; \skt{मूत्र\-प्र\-ति\-घा\-ता\-द्दा\-ल\-भ्यते}\ N}}\skt{ भिद्यते }\edtext{\skt{निस्तुद्यत}}{\Afootnote{A H; \skt{निश्तुद्यत}\ N}}\skt{ इव च }\edtext{\skt{बस्तिगुरुः}}{\Afootnote{N; \skt{बस्तिर्गुरुः}\ A; \skt{बवस्तिगुरुश्}\ H}}\skt{ }\edtext{\skt{च}}{\Afootnote{N H; \skt{शीतश्च}\ A}}\skt{ भवति । }\edtext{\skt{ऊष्मरी}}{\Afootnote{\skt{अश्मरी}\ A N H}}\skt{ चात्र श्वेता स्निग्धा महती }\edtext{\skt{कुक्कु\-टा\-ण्ड\-प्र\-तीमा}}{\Afootnote{H; \skt{कुक्कु\-टा\-ण्ड\-प्र\-ती\-काशा म}\ A; \skt{कुर्कु\-टा\-ण्ड\-प्र\-तीमा}\ N}}\skt{ }\edtext{\skt{मधुवर्णा भा}}{\Afootnote{H; \skt{धूकपुष्पवर्णा वा}\ A; \skt{सधुवर्ण्णा भा}\ N}}\skt{ भवति । तां }\edtext{\skt{श्लैष्मिकीमिति}}{\Afootnote{\skt{श्लैष्मिकीमिति}\ A; \skt{श्लेष्माश्मरीमिति}\ N H}}\skt{ }\edtext{\skt{विद्यात}}{\Afootnote{\skt{विद्यात्}\ A N H}}\skt{ ।}\pend

\pstart \ledsidenote{\passref{2.3.9}}\edtext{\skt{पित्तयुक्तस्तु}}{\Afootnote{N H; \skt{पित्तयुतस्तु}\ A}}\skt{ श्लेष्मा }\edtext{\edtext{\skt{संघा\-त\-मु\-प\-गम्य}}{\lemma{\skt{संघातम्}}\Afootnote{A N; \skt{संघात}\ H}}}{\lemma{\skt{उपगम्य}}\Afootnote{A N; \skt{मुपगम्य}\ H}}\skt{ }\edtext{\skt{यथोक्तं}}{\Afootnote{N H; \skt{यथोक्तां}\ A}}\skt{ परिवृद्धिम्प्राप्य }\edtext{\skt{बस्तिमुख}}{\Afootnote{\skt{बस्तिमुखम्}\ A N H}}\skt{ }\edtext{\skt{मधिष्ठाय}}{\Afootnote{\skt{अधिष्ठाय}\ A N H}}\skt{ }\edtext{\skt{स्रोतो}}{\Afootnote{A H; \skt{स्रोत्रो}\ N}}\skt{ निरुणद्धि । तस्य }\edtext{\skt{मूत्र\-प्र\-ति\-घा\-ता\-दू\-ष्यते}}{\Afootnote{\skt{मूत्रप्रतीघाता दूष्यते}\ A; \skt{मूत्रपतीघाता दूष्यते}\ N; \skt{मूत्र\-प्र\-ति\-घा\-ता\-दू\-ष्यते}\ H}}\skt{ }\edtext{\skt{चूष्यते}}{\Afootnote{N H; \skt{चूष्यते दह्यते}\ A}}\skt{ }\edtext{\skt{पच्यते}}{\Afootnote{H; \skt{पच्यत}\ A N}}\skt{ इव }\edtext{\skt{बस्तिरुष्णश्च}}{\Afootnote{\skt{बस्ति\-रु\-ष्ण\-वा\-तश्च}\ A; \skt{बस्तिरुष्णश्च}\ N H}}\skt{ भवति । अश्मरी चात्र सरक्ता पीतावभासा }\edtext{\skt{कृष्णा}}{\Afootnote{A; \skt{कृष्ण}\ N H}}\skt{ }\edtext{\skt{भल्ला\-त\-का\-प्र\-तिमा}}{\Afootnote{\skt{भल्ला\-त\-का\-स्थि\-प्र\-तिमा}\ A; \skt{भल्ला\-त\-क\-प्र\-तिमा}\ N H}}\skt{ }\edtext{\skt{मधुव भा}}{\Afootnote{\skt{मधुवर्णा वा}\ A; \skt{मधुवर्ण्णाभा}\ N H}}\skt{ भवति । तां }\edtext{\skt{पैत्तिकीमिति}}{\lemma{\skt{पैत्तिकीम्}}\Afootnote{A H; \skt{पतिकीम्}\ N}}\skt{ }\edtext{\skt{विद्यात}}{\Afootnote{\skt{विद्यात् ॥}\ A N H}}\skt{ ॥}\pend

\pstart \ledsidenote{\passref{2.3.10}}\edtext{\skt{वातयुक्तस्तु}}{\Afootnote{N H; \skt{वातयुतस्तु}\ A}}\skt{ }\edtext{\skt{श्लेष्मा\-सं\-घा\-त\-मु\-प\-गम्य}}{\Afootnote{\skt{श्लेष्मा संघा\-त\-मु\-प\-गम्य}\ A N H}}\skt{ }\edtext{\skt{यथोक्तं}}{\Afootnote{H; \skt{यथोक्तां}\ A N}}\skt{ परिवृद्धिम्प्राप्य }\edtext{\skt{बस्ति\-मु\-ख\-म\-धि\-ष्ठाय}}{\Afootnote{\skt{बस्ति\-मु\-ख\-म\-धि\-ष्ठाय}\ A N H}}\skt{ }\edtext{\skt{स्रोतो}}{\Afootnote{A H; \skt{स्रोत्रो}\ N}}\skt{ निरुणद्धि । तस्य मूत्र\-प्र\-ती\-घा\-ता\-त्तीव्रा वेदना भवति । }\edtext{\edtext{\skt{तयाऽत्यर्थं}}{\lemma{\skt{तया}}\Afootnote{H; \skt{तदा}\ A N}}}{\lemma{\skt{ऽत्यर्थं}}\Afootnote{A H; \skt{त्यर्थम्}\ N}}\skt{ }\edtext{\skt{पीडमानो}}{\Afootnote{H; \skt{पीड्यमानो}\ A N}}\skt{ }\edtext{\skt{दन्तान्खादन्ति}}{\Afootnote{\skt{दन्तान्खादति}\ A H; \skt{दन्तां खादति}\ N}}\skt{ नाभिम्पीडयति मेढ्रं }\edtext{\skt{प्रमृद्गाति}}{\Afootnote{A; \skt{प्रमृद्गा}\ N; \skt{प्रमृद्नाति}\ H}}\skt{ }\edtext{\skt{गर्द्धयति}}{\Afootnote{H; \skt{पायुं स्पृशति, विशर्धते,}\ A; \skt{ति गधयति}\ N}}\skt{ विदहति । }\edtext{\skt{वात\-मू\-त्र\-पु\-री\-ष\-कृ\-च्छ्र्ताद}}{\Afootnote{\skt{वात\-मू\-त्र\-पु\-री\-षाणि कृच्छ्रेण चास्य मेहतो}\ A; \skt{वात\-मू\-त्र\-पु\-रीष क्रच्छ्रता च}\ N; \skt{वात\-मू\-त्र\-पु\-री\-ष\-कृच्छ्र्ता च}\ H}}\skt{ }\edtext{\skt{भवति}}{\Afootnote{N H; \skt{निःसरन्ति;}\ A}}\skt{ । अश्मरी चात्र }\edtext{\skt{श्यावा}}{\Afootnote{A; \skt{श्या}\ N H}}\skt{ }\edtext{\skt{विषमा}}{\Afootnote{\skt{परुषा विषमा}\ A; \skt{साद्विषमा}\ N; \skt{मा विषमा}\ H}}\skt{ खरा }\edtext{\skt{परुषा}}{\Afootnote{H; \emph{om.}\ A; \skt{रु}\ N}}\skt{ }\edtext{\skt{कद\-म्ब\-पु\-ष्प\-व\-त्क\-ण्ट\-का\-चिता}}{\Afootnote{A; \skt{कद\-म्ब\-पु\-ष्प\-व\-त्क\-ण्ठ\-क\-चिता}\ N; \skt{कद\-म्ब\-पु\-ष्प\-व\-त्क\-ण्ट\-क\-चिता}\ H}}\skt{ भवति । तां वातिकीमिति विद्यात् ॥}\pend

\pstart \ledsidenote{\passref{2.3.11}}\edtext{\skt{प्राये\-णै\-ता\-स्ति\-स्रोऽ\-श्मर्यो}}{\lemma{\skt{प्रायेणैतास्}}\Afootnote{A; \skt{प्रापयेणैतास्}\ N; \skt{प्रायेण तास्}\ H}}\skt{ }\edtext{\skt{भवन्ति}}{\Afootnote{N H; \emph{om.}\ A}}\skt{ । }\edtext{\skt{दिवा\-स्व\-प्न\-स\-मा\-स\-ना\-ध्या\-स\-न\-शी\-त\-स्नि\-ग्ध\-म\-धु\-रा\-हा\-र\-प्रि\-य\-त्वा\-द्बा\-लानां}}{\Afootnote{\skt{दिवा\-स्व\-प्न\-स\-म\-श\-ना\-ध्य\-श\-न\-शी\-त\-स्नि\-ग्ध\-गु\-रु\-म\-धु\-रा\-हा\-र\-प्रि\-य\-त्वा\-द्वि\-शे\-षेण बालानां}\ A; \skt{दिवा\-स्व\-प्न\-स\-म\-श\-ना\-ध्य\-स\-र\-शी\-त\-स्नि\-ग्ध\-म\-धु\-रा\-हा\-र\-प्रि\-य\-त्वाद्बा}\ N; \skt{दिवा\-स्व\-प्न\-स\-म\-स\-ना\-ध्य\-स\-न\-शी\-त\-स्नि\-ग्ध\-म\-धु\-रा\-हा\-र\-प्रि\-य\-त्वाद्वा}\ H}}\skt{ }\edtext{\skt{तेषां}}{\Afootnote{\skt{भवन्ति; तेषाम्}\ A; \skt{लानान्तेषाम्}\ N H}}\skt{ }\edtext{\skt{एवा\-व\-ल्प\-ब\-स्ति\-का\-य\-त्वा\-द\-मां\-सो\-प\-व\-याच्च}}{\Afootnote{\skt{एवा\-ल्प\-ब\-स्ति\-का\-य\-त्वा\-द\-नु\-प\-चि\-त\-मां\-स\-त्वाच्च बस्तेः}\ A; \skt{एव\-चा\-ल्प\-ब\-स्ति\-का\-य\-त्वा\-द\-मां\-सो\-प\-च\-याच्च}\ N H}}\skt{ }\edtext{\skt{सुख\-ग्र\-ह\-णा\-ह\-र\-णाश्च}}{\Afootnote{N H; \skt{सुख\-ग्र\-ह\-णा\-ह\-रणा}\ A}}\skt{ भवन्ति । महतान्तु शुक्राश्मरी शुक्रनिमित्ता भवति ॥}\pend

\pstart \ledsidenote{\passref{2.3.12}}\edtext{\edtext{\edtext{\skt{मैथु\-न\-वि\-घा\-ता\-च्छ\-क्र\-मु\-प\-स्थि\-त\-म\-नि\-र्ग\-च्छ\-न्वि\-मा\-र्ग\-म\-नि\-लो\-वि\-गृ\-ष्क\-वृ\-ष\-ण\-यो\-र\-न्तरे}}{\lemma{\skt{मैथु\-न\-वि\-घा\-ता\-च्छ\-क्रम्}}\Afootnote{H; \skt{मैथु\-न\-वि\-घा\-ता\-द\-ति\-मै\-थु\-नाद्वा शुक्रं}\ A; \skt{मैथुनद्भिघाता शुक्रम्}\ N}}}{\lemma{\skt{उपस्थितम्}}\Afootnote{N H; \skt{चलितम्}\ A}}}{\lemma{\skt{अनि\-र्ग\-च्छ\-न्वि\-मा\-र्ग\-म\-नि\-लो\-वि\-गृ\-ष्क\-वृ\-ष\-ण\-योर्}}\Afootnote{\skt{अनि\-र्ग\-च्छ\-द्वि\-मा\-र्ग\-ग\-म\-ना\-द\-नि\-लोऽ\-भितः संगृह्य मेढ्र\-वृ\-ष\-ण\-योर्}\ A; \skt{अनिर्गच्छतन्ति मार्गगमनिलो विगृह्य वृषणयोर्}\ N; \skt{अनि\-र्ग\-च्छ\-न्वि\-मा\-र्ग\-म\-नि\-लो\-वि\-गृ\-ष्क\-वृ\-ष\-ण\-योर्}\ H}}\skt{ संहरति संहृत्य }\edtext{\skt{चोपशोषयति}}{\Afootnote{A H; \skt{चोपसोषयति ।}\ N}}\skt{ }\edtext{\edtext{\skt{सामू\-त्र\-मा\-र्ग\-मा\-वृत्य}}{\lemma{\skt{सामूत्रमार्गम्}}\Afootnote{A H; \skt{स मूत्रमार्गम्}\ N}}}{\lemma{\skt{आवृत्य}}\Afootnote{N H; \skt{आवृणोति,}\ A}}\skt{ मूत्रकृच्छ्रं }\edtext{\skt{बस्ति}}{\Afootnote{A H; \skt{बस्तिर्}\ N}}\skt{ }\edtext{\skt{शिरसि}}{\Afootnote{H; \emph{om.}\ A; \skt{असि}\ N}}\skt{ वेदनां }\edtext{\skt{शेपसि}}{\Afootnote{H; \skt{वृष}\ A; \skt{एर्शाम्}\ N}}\skt{ }\edtext{\skt{मुष्कयोश्च}}{\Afootnote{H; \skt{णयोश्च}\ A; \skt{उष्णयोश्च}\ N}}\skt{ }\edtext{\skt{श्वय\-थु\-मु\-त्पा\-द\-यति}}{\lemma{\skt{उत्पादयति}}\Afootnote{H; \skt{आपादयति,}\ A; \skt{उयोदयति}\ N}}\skt{ }\edtext{\skt{निपीडितामात्रे च}}{\Afootnote{\skt{पीडितमात्रे च}\ A; \skt{निपीडितमात्रे च}\ N H}}\skt{ }\edtext{\skt{तस्मिन्नवकाशे}}{\Afootnote{H; \skt{तस्मिन्नेव प्रवेशे}\ A; \skt{तस्मिंन व काश}\ N}}\skt{ प्रवि\-ल\-य\-मा\-प\-द्यते । तां शुक्राश्मरीमिति विद्यात् ॥}\pend

\pstart \ledsidenote{\passref{2.3.7}}\edtext{\skt{अथ}}{\Afootnote{A H; \emph{om.}\ N}}\skt{ }\edtext{\skt{जातासु}}{\Afootnote{A H; \emph{om.}\ N}}\skt{ }\edtext{\skt{वेदना}}{\Afootnote{H; \skt{नाभि\-ब\-स्ति\-से\-व\-नी\-मे\-ह\-ने\-ष्व\-न्य\-त\-म\-स्मि\-न्मे\-हतो वेदना}\ A; \skt{अथ तार्शु वेदना}\ N}}\skt{ }\edtext{\skt{सुमूत्रधारासु}}{\Afootnote{H; \skt{मूत्रधारा}\ A; \skt{मूत्रधारासु}\ N}}\skt{ }\edtext{\skt{सङ्गः स}}{\Afootnote{A H; \skt{शंगः श}\ N}}\skt{ }\edtext{\skt{रुधि\-र\-मू\-त्र\-वि\-कि\-रणः}}{\Afootnote{\skt{रुधिरमूत्रता मूत्रविकिरणं गोमे\-द\-क\-प्र\-का\-श\-म\-त्या\-विलं ससिकतं विसृजति; धाव\-न\-ल\-ङ्घ\-न\-प्ल\-व\-न\-पृ\-ष्ठ\-या\-नो\-ष्णा\-ध्व\-ग\-म\-नै\-श्चास्य वेदना भवन्ति ॥}\ A; \skt{रूधि\-र\-मू\-त्र\-वि\-कि\-रञ्चा ॥}\ N; \skt{रुधि\-र\-मू\-त्र\-वि\-कि\-र\-णञ्च ॥}\ H}}\skt{ ॥}\pend

\pstart\versepar \ledsidenote{\passref{2.3.18}}\edtext{\skt{नाभि\-पृ\-ष्ठ\-क\-टी\-मु\-ष्क\-गु\-द\-वं\-क्ष\-ण\-शे\-प\-साम्}}{\Afootnote{\skt{नाभि\-पृ\-ष्ठ\-क\-टू\-मु\-ष्क\-गु\-द\-व\-ङ्क्ष\-ण\-शे\-फ\-साम् ।}\ A; \skt{स्लोक नाभि\-पृ\-ष्ठ\-क\-टी\-मु\-ष्क\-गु\-द\-व\-क्षं\-ण\-से\-पसां ।}\ N; \skt{भवति चात्र नाभि\-पृ\-ष्ठ\-क\-टी\-मु\-ष्क\-गु\-द\-वं\-क्ष\-ण\-शे\-पसां ।}\ H}}\skt{ । }\pend
\pstart\versepar \skt{}\edtext{\edtext{\skt{नक्र\-द्वा\-र\-स्त\-नु\-त्व\-क्को\-मध्ये}}{\lemma{\skt{नक्रद्वारस्}}\Afootnote{\skt{एकवारस्}\ A; \skt{एकद्वारस्}\ N H}}}{\lemma{\skt{तनुत्वक्कोमध्ये}}\Afootnote{\skt{तनुत्वक्को मध्ये}\ A H; \skt{तन्तत्वक्को मध्यै}\ N}}\skt{ बस्तिरधोमुखः।}\pend

\pstart\versepar \ledsidenote{\passref{2.3.20}}\edtext{\skt{अलबुरिव}}{\lemma{\skt{अलबुर्}}\Afootnote{\skt{अलाब्वा}\ A; \skt{अलाबुर्}\ N H}}\skt{ रूपेण }\edtext{\edtext{\skt{सिरा\-स्ना\-यु\-भि\-रा\-वृतः}}{\lemma{\skt{सिरास्नायुभिर्}}\Afootnote{N H; \skt{सिरास्नायुप}\ A}}}{\lemma{\skt{आवृतः}}\Afootnote{N\textsuperscript{pc} H; \skt{रिग्रहः ।}\ A; \skt{आवृतआः ।}\ N\textsuperscript{ac}}}\skt{ । }\pend
\pstart\versepar \skt{}\edtext{\skt{मूत्रा\-स\-यो\-म\-ला\-धारः}}{\Afootnote{\skt{मूत्राशयो मलाधारः}\ A; \skt{मूत्राशयो मलाधार}\ N; \skt{मूत्रासयो मलाधारः}\ H}}\skt{ }\edtext{\skt{प्राणा\-य\-त\-न\-मु\-त्त\-मम्}}{\lemma{\skt{उत्तमम्}}\Afootnote{A H; \skt{उत्तमः ॥}\ N}}\skt{ ।}\pend

\pstart \ledsidenote{\passref{2.3.22cd}}\edtext{\skt{नाडी\-भि\-रु\-प\-नी\-तस्य}}{\lemma{\skt{उपनीतस्य}}\Afootnote{A H; \skt{उचनीतस्य}\ N}}\skt{ }\edtext{\skt{मूत्र\-स्या\-मा\-श\-या\-न्त\-रात}}{\Afootnote{\skt{मूत्र\-स्या\-मा\-श\-या\-न्त\-रात् ॥}\ A; \skt{मूत्र\-स्या\-मा\-ष\-या\-न्त\-रात् ।}\ N; \skt{मूत्र\-स्या\-स\-या\-न्त\-रात् ।}\ H\textsuperscript{ac}; \skt{मूत्र\-स्या\-मा\-स\-या\-न्त\-रात् ।}\ H\textsuperscript{pc}}}\skt{ ।}\pend

\pstart\versepar \ledsidenote{\passref{2.3.23}}\skt{जाग्रतः }\edtext{\skt{स्वप्नो}}{\Afootnote{\skt{स्वपतश्}\ A H; \skt{स्वतो}\ N}}\skt{ }\edtext{\skt{वापि स}}{\Afootnote{N H; \skt{चैव स}\ A}}\skt{ }\edtext{\skt{निष्यन्दे न}}{\Afootnote{N H; \skt{निःस्यन्दे न}\ A}}\skt{ पूर्यते ॥ }\pend
\pstart\versepar \skt{}\edtext{\skt{आमुखात्सलिले}}{\Afootnote{A H; \skt{आमुखासलिले}\ N}}\skt{ न्यस्तः पार्श्वेभ्यः }\edtext{\skt{प्रतिपूर्यते}}{\Afootnote{N H; \skt{पूर्यते नवः ॥}\ A}}\skt{ ॥}\pend

\pstart\versepar \ledsidenote{\passref{2.3.24}}\edtext{\skt{नवे}}{\Afootnote{N H; \emph{om.}\ A}}\skt{ }\edtext{\skt{घटे}}{\Afootnote{H; \skt{घटो}\ A N}}\skt{ यथा }\edtext{\skt{वद्धि\-ब\-स्ति\-मू\-त्रस्य}}{\Afootnote{H; \skt{तथा विद्धि बस्तिर्मूत्रेण}\ A; \skt{विद्धि बस्तिर्मूत्रस्य}\ N}}\skt{ पूर्यते । }\pend
\pstart\versepar \skt{}\edtext{\skt{एतेनैव तु}}{\Afootnote{N H; \skt{एवमेव प्र}\ A}}\skt{ }\edtext{\skt{कल्पेन}}{\Afootnote{N H; \skt{वेशेन}\ A}}\skt{ }\edtext{\skt{वातः}}{\Afootnote{A H; \skt{वात}\ N}}\skt{ पित्तं कफोऽपि वा ।}\pend

\pstart\versepar \ledsidenote{\passref{2.3.25}}\edtext{\skt{मूत्र\-यु\-क्त\-मु\-प\-स्ने\-हा\-न्प्र\-विष्य}}{\lemma{\skt{उप\-स्ने\-हा\-न्प्र\-विष्य}}\Afootnote{\skt{उप\-स्ने\-हा\-त्प्र\-विश्य}\ A; \skt{उपस्नेहा प्रविश्य}\ N; \skt{उप\-स्ने\-हा\-त्प्र\-विष्य}\ H}}\skt{ कुरुतेऽश्मरीन् ॥ }\pend
\pstart\versepar \skt{अप्सु }\edtext{\skt{स्वच्छास्वपि}}{\lemma{\skt{स्वच्छास्व्}}\Afootnote{N H; \skt{स्वच्छा(स्था)स्व्}\ A}}\skt{ यथा }\edtext{\skt{निषिक्तासु}}{\Afootnote{A H; \skt{निसिक्तासु}\ N}}\skt{ नवे }\edtext{\skt{मणौ}}{\Afootnote{H; \skt{घटे ॥}\ A; \skt{शृणौ भवेत्}\ N}}\skt{ ।}\edtext{}{\lemma{\skt{ऽश्मरीन् ॥ अप्सु}}\Afootnote{H; \skt{ऽश्मरीम् । अप्सु}\ A; \skt{श्मरीन ॥ अप्सु}\ N}}\pend

\pstart\versepar \ledsidenote{\passref{2.3.26}}\edtext{\skt{भवे\-त्का\-ला\-न्त\-रा\-त्पङ्कः}}{\Afootnote{H; \skt{कालान्तरेण पङ्}\ A; \skt{काला\-न्त\-रा\-त्प\-ङ्कस्}\ N}}\skt{ }\edtext{\skt{तव्द\-द\-श्म\-रि\-स\-म्भवः}}{\Afootnote{\skt{कः स्याद\-श्म\-री\-सं\-भ\-वस्}\ A; \skt{तद्व\-द\-श्म\-री\-सं\-भवः ॥}\ N; \skt{तद्व\-द\-श्म\-रि\-स\-म्भवः ॥}\ H}}\skt{ ॥ }\pend
\pstart\versepar \skt{}\edtext{\skt{संहृत्यापो}}{\Afootnote{H; \skt{तथा । संहन्त्यापो}\ A; \skt{संहृत्यापो}\ N}}\skt{ यथा }\edtext{\skt{दिव्यां}}{\Afootnote{N H; \skt{दिव्या}\ A}}\skt{ }\edtext{\edtext{\skt{मारुतोऽग्निश्च}}{\lemma{\skt{मारुतो}}\Afootnote{A H; \skt{मारुते}\ N}}}{\lemma{\skt{ऽग्निश्च}}\Afootnote{A H; \skt{ग्निश्च}\ N}}\skt{ वैद्युतः ॥}\pend

\pstart \ledsidenote{\passref{2.3.27ab}}\skt{तद्वद्वलासं }\edtext{\edtext{\skt{बस्तिस्थमुष्मा}}{\lemma{\skt{बस्तिस्थम्}}\Afootnote{H; \skt{बस्तिस्थम्}\ A N}}}{\lemma{\skt{उष्मा}}\Afootnote{H; \skt{ऊष्मा}\ A; \skt{ऊष्मां}\ N}}\skt{ }\edtext{\skt{संहृत्य}}{\Afootnote{H; \skt{संहन्ति}\ A; \skt{सहृत्य}\ N}}\skt{ सानिलः ॥}\pend

\pstart \ledsidenote{\passref{2.3.14cd}}\edtext{\skt{सा भिन्न\-मू\-र्त्ति\-र्व्वा\-तेन}}{\lemma{\skt{सा भिन्नमूर्त्तिर्}}\Afootnote{A H; \skt{सा भिन्नमूर्ति}\ N}}\skt{ }\edtext{\skt{शर्क्क\-रे\-त्य\-भि\-धी\-यते}}{\lemma{\skt{शर्क्करेत्य्}}\Afootnote{A H; \skt{सर्करेभ्य्}\ N}}\skt{ ।}\pend

\pstart\versepar \ledsidenote{\passref{2.3.27cdef}}\skt{मारुते }\edtext{\skt{विगुणे}}{\Afootnote{H; \skt{प्रगुणे}\ A; \skt{विगुणै}\ N}}\skt{ बस्तौ }\edtext{\skt{मूत्रं}}{\Afootnote{A H; \skt{मूत्र}\ N}}\skt{ }\edtext{\skt{सम्यग्र}}{\Afootnote{\skt{सम्यक्प्र}\ A; \skt{सम्यङ्न}\ N; \skt{सम्यग्न}\ H}}\skt{ वर्त्तते । }\pend
\pstart\versepar \skt{विकारा विविधाश्चापि प्रतिलोमे भवन्ति }\edtext{\skt{हीति}}{\Afootnote{H; \skt{हि ॥}\ A; \skt{हीत्रि ॥}\ N}}\skt{ ॥}\pend

\pstart \ledsidenote{\passref{2.3.trailer}}\edtext{\skt{अश्मरीनिदाने}}{\Afootnote{H; \skt{इति सुश्रु\-त\-सं\-हि\-तायां निदा\-न\-स्था\-नेऽ\-श्म\-री\-नि\-दानं नाम}\ A}}\skt{ }\edtext{\skt{तृतीयः}}{\Afootnote{\skt{तृतीयोऽध्यायः ॥}\ A; \skt{तृतीयः ॥ ꣸ ॥}\ H}}\skt{ ।}\pend

\endnumbering
\beginnumbering
\eledsection*{Adhyāya 4}
\pstart \ledsidenote{\passref{2.4.1}}\skt{अथातो भग\-न्द\-रा\-णा\-न्नि\-दानं }\edtext{\skt{व्याख्यास्यामः}}{\Afootnote{A H; \emph{om.}\ N}}\skt{ ॥}\pend

\pstart \ledsidenote{\passref{2.4.3}}\edtext{\skt{वात\-पि\-त्त\-श्ले\-ष्म\-स\-न्नि\-पा\-ता\-ग\-न्तु\-क\-नि\-मित्ताः}}{\Afootnote{N H; \skt{वात\-पि\-त्त\-श्ले\-ष्म\-स\-न्नि\-पा\-ता\-ग\-न्तु\-नि\-मित्ताः}\ A}}\skt{ }\edtext{\skt{शत\-पो\-न\-को\-ष्ट्र\-ग्रीव परि\-स्रा\-वि\-श\-म्बू\-का\-व\-र्तो\-न्मा\-र्गिको}}{\Afootnote{\skt{शत\-पो\-न\-को\-ष्ट्र\-ग्री\-व\-प\-रि\-स्रा\-वि\-श\-म्बू\-का\-व\-र्तो\-न्मा\-र्गिणो}\ A; \skt{शत\-पो\-न\-को\-ष्ट्र\-ग्री\-व\-प\-रि\-स्रा\-वि\-श\-म्वू\-का\-व\-र्त्तो\-न्मा\-र्गिको}\ N H}}\skt{ यथासंख्यं पञ्च }\edtext{\skt{भगन्दरा}}{\Afootnote{A H; \skt{भगरा}\ N}}\skt{ भवन्ति । ते तु }\edtext{\skt{गुद\-ब\-स्ति\-भ\-ग\-प्र\-दे\-श\-दा\-र\-णा\-भ\-ग\-न्दरा}}{\Afootnote{\skt{भग\-गु\-द\-ब\-स्ति\-प्र\-दे\-श\-दा\-र\-णाच्च भगन्दरा}\ A; \skt{गुद\-व\-स्ति\-भ\-ग\-प्र\-दे\-श\-द\-र्प\-णा\-द्भ\-ग\-द\-न्दरा}\ N\textsuperscript{ac}; \skt{गुद\-व\-स्ति\-भ\-ग\-प्र\-दे\-श\-द\-र्प\-णा\-द्भ\-ग\-न्दरा}\ N\textsuperscript{pc}; \skt{गुद\-ब\-स्ति\-भ\-ग\-प्र\-दे\-श\-दा\-र\-णा\-द्भ\-ग\-न्दरा}\ H}}\skt{ }\edtext{\skt{इत्युच्यन्ते}}{\lemma{\skt{उच्यन्ते}}\Afootnote{A H; \skt{उच्यते ।}\ N}}\skt{ । }\edtext{\skt{अभिन्नास्तु}}{\Afootnote{N H; \skt{अभिन्नाः}\ A}}\skt{ पिडका भिन्नास्तु }\edtext{\skt{भगन्दराः}}{\Afootnote{A N\textsuperscript{pc} H; \skt{सभगन्दराः ।}\ N\textsuperscript{ac}}}\skt{ ॥}\pend

\pstart \ledsidenote{\passref{2.4.5}}\skt{तत्रा\-प\-थ्य\-से\-वि\-ना\-म्वायुः प्रकुपितः सन्निवृत्तः स्थिरीभूतो }\edtext{\skt{गुदमभिगतो}}{\lemma{\skt{अभिगतो}}\Afootnote{N H; \skt{अभि}\ A}}\skt{ }\edtext{\skt{गुद\-द्वा\-रा\-द\-ङ्गुले}}{\Afootnote{N H; \skt{तो}\ A}}\skt{ }\edtext{\skt{द्व्यङ्गुले}}{\Afootnote{N\textsuperscript{pc} H; \skt{ऽङ्गुले}\ A; \skt{द्व्यद्यङ्गुले}\ N\textsuperscript{ac}}}\skt{ }\edtext{\skt{अङ्गुले वा}}{\Afootnote{\skt{द्व्यङ्गुले वा}\ A; \skt{त्र्यङ्गुले वा}\ N H}}\skt{ }\edtext{\edtext{\skt{मास\-शो\-णि\-त\-म\-भि\-प्र\-दू\-ष्या\-रु\-ण\-वर्ण्णां}}{\lemma{\skt{मासशोणितम्}}\Afootnote{\skt{मांसशोणिते}\ A; \skt{मांसशोणितम्}\ N H}}}{\lemma{\skt{अभि\-प्र\-दू\-ष्या\-रु\-ण\-वर्ण्णां}}\Afootnote{N H; \skt{प्रदू\-ष्या\-रु\-ण\-वर्णां}\ A}}\skt{ पिड\-का\-ञ्ज\-न\-यति । }\edtext{\skt{सास्य\-तो\-दा\-दी\-न्वे\-द\-ना\-वि\-शे\-षा\-नु\-प\-ज\-न\-यति}}{\Afootnote{\skt{साऽस्य तोदा\-दी\-न्वे\-द\-ना\-वि\-शे\-षा\-ञ्ज\-न\-यति,}\ A; \skt{सास्य तोदा\-दी\-म्वे\-द\-ना\-वि\-शे\-षा\-नु\-प\-ज\-न\-यति ॥}\ N; \skt{साऽस्य तोदा\-दी\-न्वे\-द\-ना\-वि\-शे\-षा\-नु\-प\-ज\-न\-यति ।}\ H}}\skt{ । }\edtext{\skt{अप्र\-ति\-क्रि\-य\-माणा च}}{\Afootnote{A H; \skt{अप्र\-ति\-क्रि\-य\-माना च}\ N}}\skt{ पाकमुपैति । }\edtext{\skt{मूत्रा\-स\-या\-भ्या\-स\-ग\-त\-त्वाच्च}}{\Afootnote{H; \skt{मूत्रा\-श\-या\-भ्या\-स\-ग\-त\-त्वाच्च}\ A N}}\skt{ }\edtext{\skt{व्रणप्रक्लिन्नः}}{\Afootnote{N H; \skt{व्रणः प्रक्लिन्नः}\ A}}\skt{ }\edtext{\edtext{\skt{शत\-पो\-न\-क\-व\-द\-णु\-मु\-खै\-श्छि\-द्रै\-रा\-पू\-र्यते}}{\lemma{\skt{शत\-पो\-न\-क\-व\-द\-णु\-मु\-खैश्}}\Afootnote{A; \skt{शत\-पो\-न\-क\-व\-द\-नु\-मु\-खैश्}\ N; \skt{शत\-पो\-न\-क\-व\-द\-णु\-मु\-खैश्}\ H}}}{\lemma{\skt{आपूर्यते}}\Afootnote{A H; \skt{आपूते ।}\ N}}\skt{ । तानि }\edtext{\edtext{\edtext{\edtext{\skt{चच्छि\-द्रा\-ण्य\-ज\-स्रं\-भे\-ना\-नु\-वि\-द्ध\-मा\-स्रा\-व\-म्प्र\-व\-हन्ति}}{\lemma{\skt{चच्छिद्राण्य्}}\Afootnote{\skt{च छिद्राण्य्}\ A H; \skt{च च्छिद्रा}\ N}}}{\lemma{\skt{अज\-स्रं\-भे\-ना\-नु\-वि\-द्धम्}}\Afootnote{\skt{अजस्रमच्छं फेनानुविद्धम्}\ A; \skt{ण्यजस्रम्फे नानुवि}\ N; \skt{अजस्रं फेनानुविद्धम्}\ H}}}{\lemma{\skt{आस्रावम्}}\Afootnote{H; \skt{अधिकमास्रावं}\ A; \skt{द्धमास्रावम्}\ N}}}{\lemma{\skt{प्रवहन्ति}}\Afootnote{H; \skt{स्रवन्ति,}\ A; \skt{प्रवहति ।}\ N}}\skt{ । }\edtext{\edtext{\skt{व्रणश्छिद्यते}}{\lemma{\skt{व्रणश्}}\Afootnote{N H; \skt{व्रणश्च}\ A}}}{\lemma{\skt{छिद्यते}}\Afootnote{N H; \skt{ताड्यते}\ A}}\skt{ भिद्यते }\edtext{\skt{ताड्यते}}{\Afootnote{N H; \skt{छिद्यते}\ A}}\skt{ सूचीभिरिव }\edtext{\skt{निस्तुद्यते}}{\Afootnote{A H; \skt{निस्दुद्यते}\ N}}\skt{ }\edtext{\edtext{\skt{गुद\-श्चा\-व\-दी\-र्यते}}{\lemma{\skt{गुदश्}}\Afootnote{H; \skt{गुदं}\ A N}}}{\lemma{\skt{चावदीर्यते}}\Afootnote{N H; \skt{चावदीर्यते, उपेक्षिते च}\ A}}\skt{ }\edtext{\edtext{\edtext{\skt{वात\-मू\-त्र\-पु\-री\-ष\-रे\-त\-सा\-म\-प्या\-ग\-मश्च}}{\lemma{\skt{वात\-मू\-त्र\-पु\-री\-ष\-रे\-त\-साम्}}\Afootnote{A H; \skt{वात\-मू\-त्र\-पु\-री\-ष\-रे\-तसा}\ N}}}{\lemma{\skt{अप्य्}}\Afootnote{A H; \skt{म}\ N}}\skt{ }\edtext{\skt{तैश्छिद्रैर्भवति}}{\lemma{\skt{छिद्रैर्}}\Afootnote{N H; \skt{एव छिद्रैर्}\ A}}}{\lemma{\skt{आगमश्च तैश्}}\Afootnote{A H; \skt{द्यागमश्च तैश्}\ N}}\skt{ }\edtext{\skt{तम्भगन्दरं}}{\lemma{\skt{भगन्दरं}}\Afootnote{A N\textsuperscript{pc} H; \skt{भगन्दरंआ}\ N\textsuperscript{ac}}}\skt{ }\edtext{\skt{शत\-पो\-न\-क\-मि\-त्या\-च\-क्षते}}{\lemma{\skt{शतपोनकम्}}\Afootnote{A H; \skt{शतपोकम्}\ N}}\skt{ ॥}\edtext{}{\lemma{\skt{भवति तम्}}\Afootnote{A N; \skt{व्भवति तम्}\ H}}\pend

\pstart \ledsidenote{\passref{2.4.6}}\skt{पित्तं तु }\edtext{\edtext{\skt{कुपि\-त\-म\-नि\-ले\-नाधः}}{\lemma{\skt{कुपितम्}}\Afootnote{N H; \skt{प्रकुपितम्}\ A}}}{\lemma{\skt{अनिलेनाधः}}\Afootnote{A H; \skt{अनिलेराधः ।}\ N}}\skt{ प्रेरितं }\edtext{\skt{पूर्व्व\-व\-दे\-वा\-व\-स्थितं}}{\Afootnote{H; \skt{पूर्व\-व\-द\-व\-स्थितं}\ A; \skt{पूर्व\-व\-द\-वा\-व\-स्थितं}\ N}}\skt{ }\edtext{\edtext{\skt{रक्ता\-न्त\-न्वी\-मु\-च्छ्रि\-ता\-ग्रा\-मु\-ष्ट्र\-ग्री\-वा\-का\-रा\-म्पि\-ड\-का\-ञ्ज\-न\-यति}}{\lemma{\skt{रक्तान्}}\Afootnote{A H; \skt{रक्ता}\ N}}}{\lemma{\skt{तन्वी\-मु\-च्छ्रि\-ता\-ग्रा\-मु\-ष्ट्र\-ग्री\-वा\-का\-राम्}}\Afootnote{\skt{तन्वी\-मु\-च्छ्रि\-ता\-मु\-ष्ट्र\-ग्री\-वाकां}\ A; \skt{त्तन्वी\-मु\-च्छ्रि\-ता\-ग्रा\-मु\-ष्ठ\-ग्री\-वा\-का\-राम्}\ N; \skt{तन्वी\-मु\-च्छ्रि\-ता\-ग्रा\-मु\-ष्ट्र\-ग्री\-वा\-का\-राम्}\ H}}\skt{ । }\edtext{\skt{सास्य\-दा\-हा\-दी\-न्वे\-द\-ना\-वि\-शे\-षा\-नु\-प\-ज\-न\-यति}}{\Afootnote{N; \skt{साऽस्य चोषा\-दी\-न्वे\-द\-ना\-वि\-शे\-षा\-ञ्ज\-न\-यति;}\ A; \skt{साऽस्य दाहा\-दी\-न्वे\-द\-ना\-वि\-शे\-षा\-नु\-प\-ज\-न\-यति ।}\ H}}\skt{ । }\edtext{\skt{अप्र\-ति\-क्रि\-य\-माणा च}}{\Afootnote{A H; \skt{अप्र\-ति\-ति\-क्रि\-यआ माणा च}\ N\textsuperscript{ac}; \skt{अप्रतितिक्रिय माणा च}\ N\textsuperscript{pc}}}\skt{ पाकमुपैति । व्रण\-श्चा\-ग्नि\-क्षा\-रा\-भ्या\-मिव दह्यते }\edtext{\edtext{\skt{दुर्ग\-न्ध\-मु\-ष्ण\-मा\-स्रावं}}{\lemma{\skt{उष्णम्}}\Afootnote{A H; \skt{उष्ण}\ N}}}{\lemma{\skt{आस्रावं}}\Afootnote{A H; \skt{मास्रावं}\ N}}\skt{ }\edtext{\edtext{\skt{स्रव\-त्यु\-पे\-क्षि\-तश्च}}{\lemma{\skt{स्रवत्य्}}\Afootnote{A H; \skt{स्रव}\ N}}}{\lemma{\skt{उपेक्षितश्च}}\Afootnote{A H; \skt{त्युपेक्षितश्च}\ N}}\skt{ वात\-मू\-त्र\-पु\-री\-ष\-रे\-तांसि विसृजति तं भग\-न्द\-र\-मु\-ष्ट्र\-ग्री\-व\-मि\-त्या\-च\-क्षते ॥}\pend

\pstart \ledsidenote{\passref{2.4.7}}\edtext{\skt{श्लेष्मा}}{\Afootnote{H; \skt{श्लेष्मा तु}\ A; \skt{श्लेआष्मा}\ N}}\skt{ }\edtext{\skt{प्रकुपितः}}{\Afootnote{A H; \skt{प्रकुपितास्}\ N}}\skt{ }\edtext{\skt{समीरणेनाधः}}{\Afootnote{A H; \skt{समीरणेनाध}\ N}}\skt{ }\edtext{\skt{प्रेरितः}}{\Afootnote{A H; \skt{प्रेरिपूत}\ N\textsuperscript{ac}; \skt{प्रेरित}\ N\textsuperscript{pc}}}\skt{ }\edtext{\skt{पूर्व्व\-व\-दे\-वा\-व\-स्थितः}}{\Afootnote{H; \skt{पूर्व\-व\-द\-व\-स्थितः}\ A; \skt{पूर्ववदेव अवस्थितः ।}\ N}}\skt{ शुक्लावभासां }\edtext{\skt{स्थिरा\-प्कं\-ण्डू\-म\-ती\-म्पि\-ड\-का\-ञ्ज\-न\-यति}}{\lemma{\skt{स्थिरा\-प्कं\-ण्डू\-म\-ती\-म्पि\-ड\-काञ्}}\Afootnote{\skt{स्थिरां कण्डूमतीं पिडकां}\ A N H}}\skt{ । }\edtext{\skt{सास्य\-क\-ण्ड्वा\-दी\-न्वे\-द\-ना\-वि\-शे\-षा\-नु\-प\-ज\-न\-यति}}{\Afootnote{N; \skt{साऽस्य कण्ड्वा\-दी\-न्वे\-द\-ना\-वि\-शे\-षा\-ञ\-ज\-न\-यति,}\ A; \skt{साऽस्य कण्ड्वा\-दी\-न्वे\-द\-ना\-वि\-शे\-षा\-नु\-प\-ज\-न\-यति ।}\ H}}\skt{ । अप्र\-ति\-क्रि\-य\-माणा च पाकमुपैति । व्रणश्च कठिनः }\edtext{\skt{ससंरम्भः}}{\Afootnote{N H; \skt{संरम्भी}\ A}}\skt{ कण्डूप्रायः }\edtext{\skt{पिच्छिलमास्रावं}}{\Afootnote{H; \skt{पिच्छि\-ल\-म\-ज\-स्र\-मा\-स्रावं}\ A; \skt{पिच्छि\-ल\-म\-ज\-स्र\-मा\-स्रावं}\ N}}\skt{ स्रव\-त्यु\-पे\-क्षि\-तश्च वात\-मू\-त्र\-पु\-री\-ष\-रे\-तान्सि विसृजति }\edtext{\edtext{\skt{तम्भ\-ग\-न्द\-र\-म्प\-रि\-स्रा\-वि\-ण\-मि\-त्या\-च\-क्षते}}{\lemma{\skt{भगन्दरम्}}\Afootnote{A H; \skt{भगदरम्}\ N}}}{\lemma{\skt{परिस्राविणम्}}\Afootnote{H; \skt{परीस्राविणम्}\ A; \skt{परिस्रविणम्}\ N}}\skt{ ॥}\pend

\pstart \ledsidenote{\passref{2.4.8}}\skt{वायुः प्रकुपितः }\edtext{\skt{प्रकुपितौ तु}}{\Afootnote{N H; \skt{प्रकुपितौ}\ A}}\skt{ }\edtext{\skt{पित्तश्लेष्माणौ}}{\Afootnote{A H; \skt{पित्तश्लेष्माणो}\ N}}\skt{ }\edtext{\skt{परिगृह्याधो}}{\Afootnote{A H; \skt{परिगृह्य धो}\ N}}\skt{ गत्वा }\edtext{\skt{पूर्व्ववदेव}}{\Afootnote{H; \skt{पूर्ववदव}\ A; \skt{पूर्ववदेव अव}\ N}}\skt{ स्थितः }\edtext{\skt{पादा\-ङ्गु\-ष्ठ\-प्र\-माणां}}{\Afootnote{H; \skt{पादा\-ङ्गु\-ष्ठा\-ग्र\-प्र\-माणां}\ A; \skt{पादा\-ङ्गु\-ष्ठ\-प्र\-मा\-णआं}\ N}}\skt{ }\edtext{\edtext{\skt{सर्व्व\-लिं\-गा\-म्पि\-ड\-क्ता\-ञ्ज\-न\-यन्ति}}{\lemma{\skt{पिडक्ताञ्}}\Afootnote{H; \skt{पिडकां}\ A N}}\skt{ । }\edtext{\skt{तेऽस्य}}{\lemma{\skt{ऽस्य}}\Afootnote{A H; \skt{स्य}\ N}}}{\lemma{\skt{जनयन्ति । ते}}\Afootnote{N H; \skt{जनयति, सा}\ A}}\skt{ }\edtext{\edtext{\skt{तोद\-दा\-ह\-क\-ण्ड्वा\-दी\-न्वे\-द\-ना\-वि\-शे\-षा\-नु\-प\-ज\-न\-य\-न्त्य\-प्र\-ति\-क्रि\-य\-मा\-णाश्च}}{\lemma{\skt{तोद\-दा\-ह\-क\-ण्ड्वा\-दी\-न्वे\-द\-ना\-वि\-शे\-षा\-नु\-प\-ज\-न\-यन्त्य्}}\Afootnote{H; \skt{तोद\-दा\-ह\-क\-ण्ड्वा\-दी\-न्वे\-द\-ना\-वि\-शे\-षा\-ञ्ज\-न\-यति,}\ A; \skt{तोद\-दा\-ह\-क\-ण्ड्वा\-दी\-म्वे\-द\-ना\-वि\-शे\-षा\-नुप जनयत्य्}\ N}}}{\lemma{\skt{अप्र\-ति\-क्रि\-य\-मा\-णाश्च}}\Afootnote{H; \skt{अप्र\-ति\-क्रि\-य\-माणा च}\ A; \skt{अप्र\-ति\-क्रि\-या\-मा\-णाम्च}\ N}}\skt{ पाकं }\edtext{\skt{गच्छन्ति}}{\Afootnote{N H; \skt{उअपैति,}\ A}}\skt{ । व्रणश्च }\edtext{\skt{नाना\-वि\-ध\-व\-र्ण्ण\-वे\-दना}}{\Afootnote{H; \skt{ना}\ A; \skt{नाना\-वि\-ध\-व\-र्ण्ण\-वे\-दरो}\ N}}\skt{ }\edtext{\edtext{\skt{नाना\-वि\-ध\-व\-र्ण्ण\-मा\-स्रावं}}{\lemma{\skt{नाना\-वि\-ध\-व\-र्ण्णम्}}\Afootnote{H; \skt{नाविधवर्णम्}\ A; \skt{गनाविधवर्ण्ण}\ N}}}{\lemma{\skt{आस्रावं}}\Afootnote{A H; \skt{मास्रावं}\ N}}\skt{ स्रवति }\edtext{\skt{पूर्ण्ण\-न\-दी\-श\-म्बु\-का\-व\-र्त्त\-व\-च्चा\-त्रा\-भ्यु\-त्ति\-ष्ठति}}{\Afootnote{\skt{पूर्ण\-न\-दी\-श\-म्बू\-का\-व\-र्त\-व\-च्चात्र समुत्तिष्ठन्ति}\ A; \skt{पूर्ण्ण\-न\-दी\-श\-म्वू\-का\-व\-र्त्त\-व\-च्चा\-त्रा\-भ्यु\-त्ति\-ष्ठति}\ N; \skt{पूर्ण्ण\-न\-दी\-श\-म्बु\-का\-व\-र्त्त\-व\-च्चात्र अभ्युत्तिष्ठति}\ H}}\skt{ वेदनाविशेषाः। तं भगन्दरं }\edtext{\skt{शम्बु\-का\-व\-र्त्त\-मि\-त्या\-च\-क्षते}}{\lemma{\skt{शम्बुकावर्त्तम्}}\Afootnote{H; \skt{शम्बूकावर्तम्}\ A; \skt{सवूकावर्त्तम्}\ N}}\skt{ ॥}\pend

\pstart \ledsidenote{\passref{2.4.9}}\skt{मूढेन }\edtext{\skt{तु}}{\Afootnote{N H; \skt{मांसलु}\ A}}\skt{ }\edtext{\edtext{\edtext{\skt{सास्थि\-श\-ल्य\-म\-न्न\-म\-भ्य\-व\-हृतं}}{\lemma{\skt{सास्थिशल्यम्}}\Afootnote{N H; \skt{ब्धेन यदस्थिशल्यम्}\ A}}}{\lemma{\skt{अन्नम्}}\Afootnote{N; \skt{अन्नेन स}\ A; \skt{अन्नम}\ H}}}{\lemma{\skt{अभ्यवहृतं}}\Afootnote{N; \skt{हाभ्यवहृतं}\ A; \skt{भ्यवहृतं}\ H}}\skt{ }\edtext{\edtext{\skt{यदा\-गा\-ढ\-पु\-री\-षो\-न्मि\-श्र\-म\-पा\-ने\-नाधः}}{\lemma{\skt{यदा\-गा\-ढ\-पु\-री\-षो\-न्मि\-श्रम्}}\Afootnote{\skt{यदाऽ\-व\-गा\-ढ\-पु\-री\-षो\-न्मि\-श्रम्}\ A; \skt{यदा गाढ\-पु\-री\-षो\-न्मि\-श्रम्}\ N H}}}{\lemma{\skt{अपानेनाधः}}\Afootnote{A H; \skt{अपानेनाध}\ N}}\skt{ }\edtext{\edtext{\skt{प्रेरि\-त\-म\-स\-म्य\-गा\-गतं}}{\lemma{\skt{प्रेरितम्}}\Afootnote{A H; \skt{प्रेरित}\ N}}}{\lemma{\skt{असम्यगागतं}}\Afootnote{A H; \skt{मंसम्यगागतं}\ N}}\skt{ गुदं }\edtext{\skt{क्षिणोति}}{\Afootnote{N H; \skt{अपक्षिणोति}\ A}}\skt{ । तदा }\edtext{\skt{तत्र}}{\Afootnote{N H; \emph{om.}\ A}}\skt{ }\edtext{\skt{क्षत\-नि\-मि\-त्ता\-ङ्ग\-ति\-रु\-प\-जा\-यते}}{\lemma{\skt{क्षत\-नि\-मि\-त्ता\-ङ्ग\-तिर्}}\Afootnote{H; \skt{क्षतनिमित्तः कोथ}\ A; \skt{क्षत\-नि\-मि\-त्ता\-द्ग\-तिर्}\ N}}\skt{ ॥ }\edtext{\skt{तस्मि\textasciitilde{}श्चक्षते}}{\Afootnote{\skt{तस्मिंश्च क्षते}\ A H; \skt{तस्मिश्च क्षते}\ N}}\skt{ }\edtext{\skt{पूय\-रु\-धि\-रा\-व\-र्ण्णै\-मां\-स\-कोथे}}{\Afootnote{H; \skt{पूय\-रु\-धि\-रा\-व\-की\-र्ण\-मां\-स\-कोथे}\ A; \skt{पूय\-रु\-धि\-रा\-व\-दी\-र्ण्ण\-मां\-स\-कोथे}\ N}}\skt{ }\edtext{\skt{भूमाविव}}{\Afootnote{\skt{भूमाविव}\ A N H}}\skt{ जलप्रक्लिन्नायां क्रिमयः }\edtext{\skt{सञ्जायते । ते}}{\Afootnote{H; \skt{संजायन्ते ते}\ A N}}\skt{ भक्षयन्तो }\edtext{\skt{गुदमनेकधा}}{\Afootnote{\skt{गुदमनेकधा}\ A N H}}\skt{ }\edtext{\edtext{\skt{पार्श्व\-तोऽ\-व\-दा\-र\-यन्ति}}{\lemma{\skt{पार्श्वतो}}\Afootnote{A H; \skt{पार्श्वतो}\ N}}}{\lemma{\skt{ऽवदारयन्ति}}\Afootnote{H; \skt{दारयन्ति,}\ A; \skt{वदारयन्ति ।}\ N}}\skt{ }\edtext{\skt{तस्य तैश्च}}{\Afootnote{N H; \skt{तस्य तैर्मा}\ A}}\skt{ }\edtext{\edtext{\edtext{\edtext{\skt{कृमि\-कृ\-तै\-र्म्मा\-र्गै\-र्व्वा\-त\-मू\-त्र\-पु\-री\-ष\-रे\-त\-तां\-स्य\-भि\-प्र\-प\-द्यन्ते}}{\lemma{\skt{कृमिकृतैर्}}\Afootnote{H; \skt{र्गैः कृमिकृतैर्}\ A; \skt{क्रिमिकृतैर्}\ N}}}{\lemma{\skt{म्मार्गैर्}}\Afootnote{H; \emph{om.}\ A; \skt{मार्गै}\ N}}}{\lemma{\skt{व्वात\-मू\-त्र\-पु\-री\-ष\-रे\-त\-तांस्य्}}\Afootnote{H; \skt{वात\-मू\-त्र\-पु\-री\-ष\-रे\-तांस्य्}\ A N}}\skt{ । }\edtext{\skt{तम्भ\-ग\-न्द\-र\-मु\-न्मा\-र्गि\-ण\-मि\-त्या\-क्षते}}{\lemma{\skt{आक्षते}}\Afootnote{H; \skt{आचक्षते ॥}\ A N}}}{\lemma{\skt{अभिप्रपद्यन्ते । तम्}}\Afootnote{N H; \skt{अभिनिःसरन्तिः; तं}\ A}}\skt{ ॥}\pend

\pstart\versepar \ledsidenote{\passref{2.4.13}}\edtext{\skt{भवति}}{\Afootnote{H; \emph{om.}\ A N}}\skt{ }\edtext{\skt{चात्र}}{\Afootnote{H; \emph{om.}\ A N}}\skt{ । }\pend
\pstart\versepar \skt{घोराः }\edtext{\skt{साधयितुन्दुःखाः}}{\lemma{\skt{साधयितुन्}}\Afootnote{A H; \skt{साधयितुर्}\ N}}\skt{ सर्व्व एव भगन्दराः । }\pend
\pstart\versepar \skt{तेष्व\-सा\-ध्य\-स्त्रि\-दो\-षोत्थः क्षतजश्च }\edtext{\skt{भगन्दर}}{\Afootnote{N H; \skt{भगन्दरः ॥}\ A}}\skt{ }\edtext{\skt{इति}}{\Afootnote{H; \emph{om.}\ A N}}\skt{ ॥ ४ ॥}\pend

\pstart \ledsidenote{\passref{2.4.trailer}}\edtext{\skt{भग\-न्द\-र\-नि\-दाने}}{\Afootnote{H; \skt{इति सुश्रु\-त\-सं\-हि\-तायां निदानस्थाने भग\-न्द\-र\-नि\-दानं नाम}\ A; \skt{इति । . । . ओ ॥}\ N}}\skt{ }\edtext{\skt{चतुर्थः}}{\Afootnote{H; \skt{चतुर्थोऽध्यायः ॥ ४ ॥}\ A; \emph{om.}\ N}}\skt{ ॥}\pend

\endnumbering
\beginnumbering
\eledsection*{Adhyāya 5}
\pstart \ledsidenote{\passref{2.5.1}}\skt{अथातः }\edtext{\skt{कुष्ठ\-नि\-दा\-न\-म्व्या\-ख्या\-स्यामः}}{\lemma{\skt{व्याख्यास्यामः}}\Afootnote{A H; \emph{om.}\ N}}\skt{ ॥}\pend

\pstart \ledsidenote{\passref{2.5.3}}\edtext{\skt{मिथ्या\-हा\-रा\-चा\-रस्य}}{\Afootnote{N H; \skt{मिथ्या\-हा\-रा\-चा\-रस्य विशे\-षा\-द्गु\-रु\-वि\-रु\-द्धा\-सा\-त्म्या\-जी\-र्णा\-हि\-ता\-शिनः स्नेहपीतस्य वान्तस्य वा व्याया\-म\-ग्रा\-म्य\-ध\-र्म\-से\-विनो ग्राम्या\-नू\-पौ\-द\-क\-मां\-सानि वा पय\-साऽ\-भी\-क्ष्ण\-म\-श्नतो यो वा मज्ज\-त्य\-प्सू\-ष्मा\-भि\-तप्तः सहसाछर्दिं वा प्रतिहन्ति, तस्य}\ A}}\skt{ पित्तश्लेष्माणौ }\edtext{\skt{प्रकुपितौ}}{\Afootnote{A H; \skt{प्रकुपितो}\ N}}\skt{ परिगृह्यानिलः }\edtext{\skt{प्रवृद्धस्तिर्यगाः}}{\lemma{\skt{तिर्यगाः}}\Afootnote{\skt{तिर्यग्गाः}\ A H N}}\skt{ सिराः }\edtext{\skt{संप्रपद्य}}{\Afootnote{H; \skt{सर्प्रपद्य}\ A; \skt{सम्प्रपद्यते}\ N}}\skt{ }\edtext{\skt{समूद्धूय}}{\Afootnote{H; \skt{समुद्धूय}\ A N}}\skt{ }\edtext{\skt{बाह्यमार्गम्प्रति}}{\lemma{\skt{बाह्यमार्गम्}}\Afootnote{H; \skt{बाह्यं मार्गं}\ A; \skt{वाह्यं मार्गम्}\ N}}\skt{ सम\-न्ता\-द्वि\-क्षि\-पति । यत्र }\edtext{\skt{यत्र}}{\Afootnote{N H; \skt{यत्र च}\ A}}\skt{ }\edtext{\skt{विक्षिप्तो}}{\Afootnote{H; \skt{दोषो विक्षिप्तो}\ A; \skt{विक्षिप्ता}\ N}}\skt{ निश्चरति तत्र तत्र मण्डलानि }\edtext{\skt{भवन्ति}}{\Afootnote{H; \skt{प्रादुर्भवन्ति,}\ A; \emph{om.}\ N}}\skt{ । }\edtext{\skt{तत्र च}}{\Afootnote{H; \skt{एवं समुत्पन्नस्त्वचि दोषस्तत्र तत्र च}\ A; \skt{भवन्ति । तत्र च}\ N}}\skt{ }\edtext{\edtext{\edtext{\skt{परि\-वृ\-द्धि\-म्प्रा\-प्या\-प्र\-ति\-क्रि\-य\-मा\-णोऽ\-भ्य\-न्त\-र\-म\-नु\-प्राप्तो}}{\lemma{\skt{प्राप्या\-प्र\-ति\-क्रि\-य\-माणो}}\Afootnote{A H; \skt{प्राप्या\-प्र\-ति\-क्रि\-या\-माणो}\ N}}}{\lemma{\skt{ऽभ्यन्तरम्}}\Afootnote{A N; \skt{भ्यन्तरम्}\ H}}}{\lemma{\skt{अनुप्राप्तो}}\Afootnote{N H; \skt{प्रतिपद्यते}\ A}}\skt{ }\edtext{\skt{धातू\-न्वि\-दू\-ष\-यति}}{\Afootnote{N H; \skt{धातू\-न\-भि\-दू\-ष\-यन् ॥}\ A}}\skt{ ।}\pend

\pstart \ledsidenote{\passref{2.5.4}}\edtext{\skt{तस्य च}}{\Afootnote{N H; \skt{तस्य}\ A}}\skt{ पूर्व्वरूपाणि }\edtext{\skt{त्वक्पा\-रू\-ष्य\-म\-क\-स्मा\-द्रो\-म\-हर्षः}}{\lemma{\skt{त्वक्पारूष्यम्}}\Afootnote{H; \skt{त्वक्पारुष्यम्}\ A N}}\skt{ }\edtext{\skt{कण्डु\-स्वे\-द\-ब\-हुत्वं}}{\Afootnote{\skt{कण्डूः स्वेदबाहुल्यम्}\ A; \skt{कण्डू स्वेदवहु}\ N; \skt{कण्डू\-स्वे\-द\-ब\-हुत्वं}\ H}}\skt{ }\edtext{\edtext{\edtext{\skt{सुप्त\-त्व\-म\-ङ्गा\-ना\-म\-सृ\-क्कृ\-ष्णता}}{\lemma{\skt{सुप्तत्वम्}}\Afootnote{H; \skt{अस्वेदनं वाङ्गप्रदेशानां स्वापः क्षतविसर्प}\ A; \skt{त्वसुप्तत्वम्}\ N}}}{\lemma{\skt{अङ्गानाम्}}\Afootnote{N H; \skt{णम्}\ A}}}{\lemma{\skt{असृक्कृष्णता}}\Afootnote{N H; \skt{असृजः कृष्णता}\ A}}\skt{ }\edtext{\skt{च}}{\Afootnote{N H; \skt{चेति ॥}\ A}}\skt{ ॥}\pend

\pstart \ledsidenote{\passref{2.5.5}}\skt{तत्र सप्त महाकुष्ठानि । }\edtext{\skt{एका\-द\-श\-क्षु\-द्र\-कु\-ष्ठा\-न्ये\-व\-म\-ष्टा\-दश}}{\lemma{\skt{अष्टादश}}\Afootnote{A H; \skt{अष्ठाद}\ N}}\skt{ }\edtext{\skt{भवन्ति}}{\Afootnote{H; \skt{कुष्ठानि भवन्ति ।}\ A; \skt{शा भवन्ति ।}\ N}}\skt{ । }\edtext{\skt{ते वा}}{\Afootnote{\skt{तत्र}\ A; \skt{तेषाम्}\ N H}}\skt{ }\edtext{\edtext{\skt{महा\-कु\-ष्ठा\-न्य\-रु\-णो\-दु\-म्ब\-र\-ऋ\-ष्य\-जि\-ह्व\-क\-पा\-ल\-का\-क\-ण\-क\-प\-ण्डु\-री\-कानि}}{\lemma{\skt{महाकुष्ठान्य्}}\Afootnote{A H; \skt{महाकुष्ठाण्य्}\ N}}}{\lemma{\skt{अरु\-णो\-दु\-म्ब\-र\-ऋ\-ष्य\-जि\-ह्व\-क\-पा\-ल\-का\-क\-ण\-क\-प\-ण्डु\-री\-कानि}}\Afootnote{\skt{अरु\-णो\-दु\-म्ब\-र\-र्ष्य(\-र्क्ष)\-जि\-ह्व\-क\-पा\-ल\-का\-क\-ण\-क\-पु\-ण्ड\-रीक}\ A; \skt{अरु\-णो\-दु\-म्व\-र\-ऋ\-श्य\-जि\-ह्व\-क\-पा\-ल\-का\-क\-ण\-क\-पु\-ण्ड\-री\-कानि}\ N; \skt{अरु\-णो\-दु\-म्ब\-र\-ऋ\-ष्य\-जि\-ह्व\-क\-पा\-ल\-का\-क\-णक पुण्डरीकानि}\ H}}\skt{ }\edtext{\edtext{\skt{दर्दृकुष्ठञ्चेति}}{\lemma{\skt{दर्दृकुष्ठञ्}}\Afootnote{\skt{दद्रुकुष्ठा}\ A; \skt{दद्रुकुष्ठश्}\ N; \skt{दर्द्द्रूकुष्ठञ्}\ H}}}{\lemma{\skt{चेति}}\Afootnote{N H; \skt{नीति ।}\ A}}\skt{ }\edtext{\skt{क्षूद्रकुष्ठानि}}{\Afootnote{H; \skt{क्षुद्रकुष्ठान्य्}\ A N}}\skt{ }\edtext{\skt{तु}}{\Afootnote{N H; \skt{अपि}\ A}}\skt{ }\edtext{\edtext{\skt{स्थूला\-रु\-ष्क\-म\-ष्ठै\-क\-कु\-ष्ठ\-च\-र्म्म\-दलं}}{\lemma{\skt{स्थूलारुष्कम्}}\Afootnote{A; \skt{स्थूलारुक्ष}\ N; \skt{स्थूलारुष्क}\ H}}}{\lemma{\skt{अष्ठै\-क\-कु\-ष्ठ\-च\-र्म्म\-दलं}}\Afootnote{\skt{महा\-कु\-ष्ठ\-मे\-क\-कुष्ठं चर्मदलं}\ A; \skt{महाकुष्ठ कजु\-ष्ठ\-श्च\-र्म\-द\-लम्}\ N; \skt{महाकु ष्ठैक\-कु\-ष्ठ\-च\-र्म्म\-दलं}\ H}}\skt{ }\edtext{\skt{परिसर्प्पो}}{\Afootnote{H; \emph{om.}\ A; \skt{परिसर्पा}\ N}}\skt{ }\edtext{\skt{विसर्प्पसिध्मं}}{\Afootnote{H; \skt{विसर्पः परिसर्पः सिध्मं}\ A; \skt{विसर्पा विसर्पसिध्मम्}\ N}}\skt{ विचर्च्चिका }\edtext{\skt{किटिभं}}{\Afootnote{N; \skt{किटिभं(मं)}\ A; \skt{किटिमं}\ H}}\skt{ }\edtext{\skt{पामा च}}{\Afootnote{A; \skt{पामार}\ N H}}\skt{ कसा चेति ॥}\pend

\pstart \ledsidenote{\passref{2.5.6}}\skt{सर्व्वाणि कुष्ठानि सवातानि सपित्तानि सश्लेष्माणि }\edtext{\skt{सक्रिमाणि}}{\Afootnote{H; \skt{सक्रिमीणि च}\ A; \skt{सक्रिःमीणि}\ N}}\skt{ भवन्ति । }\edtext{\skt{उत्पन्नस्य तु}}{\Afootnote{H; \skt{उत्सन्नतस्तु}\ A; \skt{उत्पन्न}\ N}}\skt{ }\edtext{\edtext{\skt{ग्रह\-ण\-म\-भि\-भ\-वान्त}}{\lemma{\skt{ग्रहणम्}}\Afootnote{H; \skt{दोषग्रहणम्}\ A; \skt{स्य ग्रहणम्}\ N}}}{\lemma{\skt{अभिभवान्त}}\Afootnote{\skt{अभिभवात्}\ A N; \skt{अभि भवान्त्}\ H}}\skt{ ॥}\pend

\pstart \ledsidenote{\passref{2.5.7}}\edtext{\skt{त्र}}{\Afootnote{\skt{तत्र}\ A N H}}\skt{ वातेनारुणं । }\edtext{\skt{पित्तेनौदुम्बरं ।अ\-ष्य\-जि\-ह्व\-क\-पा\-ल\-का\-क\-ण\-कानि}}{\Afootnote{\skt{पित्ते\-नो\-दु\-म्ब\-र\-र्ष्य(\-र्क्ष)\-जि\-ह्व\-क\-पा\-ल\-का\-क\-ण\-कानि,}\ A; \skt{पित्ते\-नो\-दु\-म्व\-र\-ऋ\-ष्य\-जि\-ह्व\-क\-पा\-ल\-का\-क\-ण\-कानि ।}\ N; \skt{पित्तेनौदुम्बरं ।ऋ\-ष्य\-जि\-ह्व\-क\-पा\-ल\-का\-क\-ण\-कानि ।}\ H}}\skt{ । }\edtext{\skt{श्लेष्मणा}}{\Afootnote{N H; \skt{स्लेष्मणा}\ A}}\skt{ पुण्डरीकं }\edtext{\edtext{\edtext{\skt{ददृ\-कु\-ष्ठ\-ञ्चे\-त्येषां}}{\lemma{\skt{ददृकुष्ठञ्}}\Afootnote{H; \skt{दद्रुकुष्ठं}\ A; \skt{दद्रुकु}\ N}}}{\lemma{\skt{चेत्य्}}\Afootnote{A H; \emph{om.}\ N}}\skt{ तु}}{\lemma{\skt{एषां तु}}\Afootnote{H; \skt{तेषां म}\ A; \emph{om.}\ N}}\skt{ }\edtext{\skt{महत्वं}}{\Afootnote{H; \skt{हत्त्वं क्रिया\-गु\-रु\-त्व\-मु\-त्तरो}\ A; \skt{ष्ठण्चेत्येषान्तु महत्वं}\ N}}\skt{ }\edtext{\edtext{\skt{सर्व्व\-धा\-त्व\-नु\-सा\-रि\-त्वा\-द\-सा\-ध्ध\-त्व\-ञ्चेति}}{\lemma{\skt{सर्व्व\-धा\-त्व\-नु\-सा\-रि\-त्वा\-द\-सा\-ध्ध\-त्वञ्}}\Afootnote{\skt{त्तरं धात्व\-नु\-प्र\-वे\-शा\-द\-सा\-ध्यत्वं}\ A; \skt{सर्व\-धा\-त्व\-नु\-सा\-रि\-त्वा\-द\-सा\-ध्यत्वा}\ N; \skt{सर्व्व\-धा\-त्व\-नु\-सारि त्वादसाध्यत्वञ्}\ H}}}{\lemma{\skt{चेति}}\Afootnote{A H; \skt{च्चेति ।}\ N}}\skt{ ।}\pend

\pstart \ledsidenote{\passref{2.5.8}}\skt{तत्र }\edtext{\skt{वाते\-ना\-रू\-ण\-व\-र्ण्णानि}}{\Afootnote{H; \skt{वाते\-ना\-रु\-णा\-भानि}\ A; \skt{वाते\-ना\-रु\-ण\-व\-र्ण्णानि}\ N}}\skt{ तनूनि विसर्प्पीणि }\edtext{\skt{तोद\-स्वा\-प\-यु\-क्तानि}}{\Afootnote{N H; \skt{तोद\-भे\-द\-स्वा\-प\-यु\-क्तान्य्}\ A}}\skt{ }\edtext{\skt{भवन्ति}}{\Afootnote{N H; \skt{अरुणानि । पित्तेन}\ A}}\skt{ ॥ }\edtext{\edtext{\skt{पक्वो\-दु\-म्ब\-र\-व\-र्ण्णा\-न्यौ\-दु\-म्ब\-राणि}}{\lemma{\skt{पक्वो\-दु\-म्ब\-र\-व\-र्ण्णान्य्}}\Afootnote{H; \skt{पक्वो\-दु\-म्ब\-र\-फ\-ला\-कृ\-ति\-व\-र्णान्य्}\ A; \skt{पक्वो\-दु\-म्व\-र\-व\-र्ण्णा\-न्यो\-दुम्व}\ N}}}{\lemma{\skt{औदुम्बराणि}}\Afootnote{A H; \skt{राणि}\ N}}\skt{ । }\edtext{\skt{अष्य\-जि\-ह्वे\-व\-ख\-राणि}}{\Afootnote{\skt{ऋष्य(\-क्ष)\-जि\-ह्वा\-प्र\-का\-शानि खराणि}\ A; \skt{ऋष्यजिह्वेव खराणि}\ N H}}\skt{ }\edtext{\skt{अष्यजिह्वानि}}{\Afootnote{\skt{ऋष्य(क्ष)जिह्वानि,}\ A; \skt{ऋश्यजिह्वानि ।}\ N; \skt{ऋष्यजिह्वानि ॥}\ H}}\skt{ ॥ कृष्ण\-क\-पा\-लि\-का\-प्र\-का\-शानि }\edtext{\skt{कपालकुष्ठानि ।काकणन्ति कफ\-ल\-स\-दृ\-शानि}}{\Afootnote{N H; \skt{कपालकुष्ठानि, काक\-ण\-न्ति\-का\-फ\-ल\-स\-दृ\-शान्य्}\ A}}\skt{ अतीव }\edtext{\skt{रक्तानि}}{\Afootnote{N H; \skt{रक्त}\ A}}\skt{ }\edtext{\skt{पर्यन्ते च}}{\Afootnote{N; \emph{om.}\ A; \skt{पच्यन्ते च}\ H}}\skt{ }\edtext{\skt{कृष्णानि}}{\Afootnote{A H; \skt{कृष्णाणि}\ N}}\skt{ }\edtext{\skt{काकणन्ति कानि}}{\Afootnote{N H; \skt{काकणकानि;}\ A}}\skt{ तेषां }\edtext{\skt{चतु\-र्ण्णा\-म\-प्या\-षा\-चो\-ष\-प\-रि\-दा\-ह\-धू\-मा\-य\-नानि}}{\Afootnote{\skt{चतु\-र्णा\-म\-प्यो\-ष\-चो\-ष\-प\-रि\-दा\-ह\-धू\-मा\-य\-नानि}\ A; \skt{पञ्चानाप्योषा चोष\-प\-रि\-दा\-ह\-धू\-मा\-य\-नानि}\ N; \skt{चतु\-र्ण्णा\-म\-प्या\-षा\-चो\-ष\-प\-रि\-दा\-ह\-धू\-मा\-य\-नानि ।}\ H}}\skt{ । }\edtext{\skt{क्षिप्रोत्थान}}{\Afootnote{A H; \skt{क्षिपोत्थान}\ N}}\skt{ प्रपा\-क\-भे\-दि\-त्वानि }\edtext{\skt{च}}{\Afootnote{N H; \skt{क्रिमिजन्म च}\ A}}\skt{ }\edtext{\skt{सामान्यानि}}{\Afootnote{N H; \skt{सामान्यानि लिङ्गानि ।}\ A}}\skt{ । }\edtext{\skt{पुण्ड\-री\-क\-प्र\-का\-शानि}}{\Afootnote{H; \skt{श्लेष्मणा पुण्ड\-री\-क\-प\-त्र\-प्र\-का\-शानि}\ A; \skt{पुण्ड\-री\-क\-प्र\-का\-शान्य}\ N}}\skt{ }\edtext{\edtext{\skt{पुण्ड\-री\-का\-न्य\-त\-सी\-पु\-ष्प\-व\-र्ण्णा\-नि\-ता\-म्राणि}}{\lemma{\skt{पुण्डरीकान्य्}}\Afootnote{N H; \skt{पौण्डरीकाणि,}\ A}}\skt{ वा}}{\lemma{\skt{अत\-सी\-पु\-ष्प\-व\-र्ण्णा\-नि\-ता\-म्राणि वा}}\Afootnote{A H; \skt{तसी\-पु\-ष्प\-व\-र्ण्णानि ताम्राणि वा}\ N}}\skt{ विसर्प्पीणि पिडका }\edtext{\skt{भवन्ति}}{\Afootnote{H; \skt{वन्ति च}\ A; \skt{वन्ति}\ N}}\skt{ }\edtext{\skt{दर्दृकुष्ठानि च}}{\Afootnote{\skt{दद्रुकुष्ठानि}\ A N; \skt{दर्द्द्रूकुष्ठानि च}\ H}}\skt{ }\edtext{\edtext{\edtext{\skt{तयो\-र्द्व\-यो\-र\-प्यु\-त्म\-न्नता}}{\lemma{\skt{तयोर्}}\Afootnote{A N; \skt{तयो}\ H}}}{\lemma{\skt{द्वयोर्}}\Afootnote{A N; \skt{र्द्वयोर्}\ H}}}{\lemma{\skt{अप्युत्मन्नता}}\Afootnote{\skt{अप्युत्सन्नता}\ A H; \skt{अप्युत्सन्नता ।}\ N}}\skt{ । परिमण्डलता }\edtext{\edtext{\edtext{\skt{कण्डु\-श्चि\-रो\-त्था\-न\-त्व\-ञ्चेति}}{\lemma{\skt{कण्डुश्}}\Afootnote{\skt{कण्डूश्}\ A H; \skt{कण्डू}\ N}}}{\lemma{\skt{चिरोत्थानत्वञ्}}\Afootnote{A H; \skt{श्चिरोत्थानत्वञ्}\ N}}}{\lemma{\skt{चेति}}\Afootnote{H; \skt{चेति सामान्यानि रूपाणि ॥}\ A; \skt{चे ॥}\ N}}\skt{ ॥}\pend

\pstart \ledsidenote{\passref{2.5.9}}\skt{क्षुद्रकुष्ठान्यत ऊर्ध्वं वक्ष्यामः। }\edtext{\skt{अरुःकुष्ठे}}{\Afootnote{H; \emph{om.}\ A; \skt{अरुः । कुष्ठम्}\ N}}\skt{ }\edtext{\edtext{\skt{स्थूल\-मू\-ला\-न्य\-रुंषि}}{\lemma{\skt{स्थूलमूलान्य्}}\Afootnote{H; \skt{स्थूलानि}\ A; \skt{स्थूलमूला}\ N}}}{\lemma{\skt{अरुंषि}}\Afootnote{H; \emph{om.}\ A; \skt{न्यरूंषि}\ N}}\skt{ }\edtext{\skt{संजायंते}}{\Afootnote{N H; \emph{om.}\ A}}\skt{ }\edtext{\skt{सन्धिषु}}{\Afootnote{N H; \skt{सन्धिष्व}\ A}}\skt{ }\edtext{\skt{दारुणानि}}{\Afootnote{H; \skt{तिदारुणानि}\ A; \skt{दारुणा}\ N}}\skt{ । }\edtext{\skt{महाकुमष्ठे}}{\Afootnote{\skt{स्थूलारुषि स्युः कठिनान्यरूंषि । त्वक्को\-च\-भे\-द\-स्व\-प\-ना\-ङ्ग\-सादाः कुष्ठे म}\ A; \skt{नि । महाकुष्ठे}\ N; \skt{महाकुष्ठे}\ H}}\skt{ }\edtext{\skt{सर्व्वदेहे}}{\Afootnote{N H; \skt{हत्पूर्वयुते}\ A}}\skt{ भवन्ति }\edtext{\skt{त्वक्कोचः}}{\Afootnote{\emph{om.}\ A; \skt{त्वक्संकोच}\ N; \skt{त्वक्संकोचः}\ H}}\skt{ }\edtext{\skt{भेदाङ्गदाहाः॥}}{\lemma{\skt{भेदाङ्गदाहाः}}\Afootnote{\emph{om.}\ A; \skt{स्वेद\-भे\-दा\-ङ्ग\-दाहाः}\ N H}}\pend

\pstart \ledsidenote{\passref{2.5.10}}\edtext{\skt{कृत्स्ने}}{\Afootnote{H; \emph{om.}\ A; \skt{कृत्सु}\ N}}\skt{ }\edtext{\skt{देहे}}{\Afootnote{N H; \emph{om.}\ A}}\skt{ }\edtext{\skt{यस्य}}{\Afootnote{N H; \emph{om.}\ A}}\skt{ }\edtext{\edtext{\edtext{\skt{कृष्णोऽ\-रु\-णो\-वा\-त\-च्चै\-का\-रव्यं}}{\lemma{\skt{कृष्णो}}\Afootnote{H; \skt{कृष्णारुणं येन भवे}\ A; \skt{कृष्णारुणो वा तम्चै}\ N}}\skt{ कुष्ठमा}}{\lemma{\skt{ष्णो… चैकारव्यं}}\Afootnote{\emph{om.}\ N}}\skt{ सुकष्टं}}{\lemma{\skt{ऽरु\-णो\-वा\-त\-च्चै\-का\-रव्यं कुष्ठमा सुकष्टं}}\Afootnote{\skt{च्छरीरं तदेककुष्ठं प्रवदन्ति कुष्ठम् । स्युर्येन}\ A; \skt{ऽरुणो वा तच्चैकाख्यं कुष्ठमाहुः सुकष्ठं ।}\ H}}\skt{ । }\edtext{\skt{कण्डु\-चो\-षौ\-तो\-दहौ तु}}{\Afootnote{H; \skt{कण्डू\-व्य\-थ\-नौ\-ष\-चो\-षास्}\ A; \skt{कण्डूचोषौ तोददाहौ तु}\ N}}\skt{ }\edtext{\skt{यस्य}}{\Afootnote{N H; \skt{त}\ A}}\skt{ }\edtext{\skt{काले}}{\Afootnote{N H; \skt{ले}\ A}}\skt{ }\edtext{\skt{चास्मि\-न्च\-र्म्म\-द\-ल\-म्व\-दन्ति}}{\lemma{\skt{चास्मि\-न्च\-र्म्म\-द\-लम्}}\Afootnote{H; \skt{षु तच्चर्मदलं}\ A; \skt{च्चास्मिञ्चर्मदालं}\ N}}\skt{ ।}\pend

\pstart \ledsidenote{\passref{2.5.12ab}}\edtext{\edtext{\skt{यस्मि\-न्स्फो\-टा\-स्रा\-व\-व\-न्त\-स्त\-माह्नः}}{\lemma{\skt{यस्मि\-न्स्फो\-टा\-स्रा\-व\-व\-न्त\-स्तम्}}\Afootnote{H; \skt{शनैः शरीरे पिडकाः स्रवन्त्यः सर्पन्ति यास्तं परिसर्पम्}\ A; \skt{यस्मिन्स्फोटा स्राववन्त तम्}\ N}}}{\lemma{\skt{आह्नः}}\Afootnote{\skt{आहुः}\ A N H}}\skt{ }\edtext{\edtext{\skt{पारी\-स\-र्प्प\-न्त\-च्छनैः}}{\lemma{\skt{पारीसर्प्पन्}}\Afootnote{H; \emph{om.}\ A; \skt{परीसर्पन्}\ N}}}{\lemma{\skt{तच्छनैः}}\Afootnote{N H; \emph{om.}\ A}}\skt{ }\edtext{\skt{सर्प्पमाणं}}{\Afootnote{N H; \emph{om.}\ A}}\skt{ ।}\pend

\pstart \ledsidenote{\passref{2.5.11cd}}\edtext{\skt{वैसर्प्यः}}{\Afootnote{\skt{मूर्च्छा\-वि\-दा\-हा\-र\-ति\-तो\-द\-पा\-का\-म्कृत्वा विसर्पः}\ A; \skt{वैसर्प्प}\ N; \skt{वैसर्प्पः}\ H}}\skt{ }\edtext{\skt{स्यात्सर्व्वतः}}{\Afootnote{N H; \skt{स भवेद्वि}\ A}}\skt{ }\edtext{\skt{सर्प्यते तु}}{\Afootnote{H; \skt{कारः ॥}\ A; \skt{सर्पनन्तु}\ N}}\skt{ }\edtext{\skt{त्वग्रक्तदीन्व्याप्य}}{\Afootnote{H; \emph{om.}\ A; \skt{त्वग्रक्तादीन्व्याप्य}\ N}}\skt{ }\edtext{\skt{भावान्सुशीघ्रम्}}{\Afootnote{H; \emph{om.}\ A; \skt{भावां सुशीघ्रं ॥}\ N}}\skt{ ॥}\pend

\pstart \ledsidenote{\passref{2.5.12cd}}\edtext{\skt{स्वच्छ\-स्वे\-त\-प्क\-ण्डु\-म\-च्चापि}}{\Afootnote{\skt{कण्ड्वन्वितं श्वेतमपायि}\ A; \skt{स्वच्छ स्वेतं कण्डुमच्चापि}\ N; \skt{स्वच्छ\-स्वे\-त\-ङ्क\-ण्ण्डु\-म\-च्चापि}\ H}}\skt{ }\edtext{\skt{भिष्मं}}{\Afootnote{\skt{सि}\ A; \skt{सिध्मम्}\ N H}}\skt{ }\edtext{\skt{परिध्वंसिम्रा }\edtext{\skt{यसश्चोर्द्वकाये}}{\lemma{\skt{चोर्द्वकाये}}\Afootnote{H; \skt{ऊर्ध्वकाये ॥}\ A; \skt{चोर्ध्वकाये ।}\ N}}}{\lemma{\skt{परिध्वंसिम्रा यसश्}}\Afootnote{\skt{ध्म विद्यात्तनु प्रायश}\ A; \skt{परिध्वंसि प्रायशश्}\ N; \skt{परिध्वंसि प्रायसश्}\ H}}\skt{ ।}\pend

\pstart \ledsidenote{\passref{2.5.12.1}}\edtext{\skt{पाणौ}}{\Afootnote{H; \skt{पानी}\ N}}\skt{ पादौ दारुणौ यस्य }\edtext{\skt{रुक्षै\-शो\-वा\-क\-ण\-डुः\-सा\-वि\-चर्च्चिः}}{\Afootnote{\skt{रूक्षौ सोषा कण्डू सा विचर्च्चि}\ N; \skt{रूक्षै शोषाकणडूः सा विचर्च्चिः}\ H}}\skt{ }\edtext{\skt{प्रदिष्ठा}}{\Afootnote{\skt{प्रदिष्टा}\ N H}}\skt{ ।}\pend

\pstart\versepar \ledsidenote{\passref{2.5.12.2}}\edtext{\skt{वैपाद्याख्यः}}{\Afootnote{H; \skt{वैपाद्या वैशो}\ N}}\skt{ }\edtext{\skt{पाद\-यो\-श्चा\-व\-दीर्ण्णु}}{\lemma{\skt{चावदीर्ण्णु}}\Afootnote{H; \skt{चावदीर्ण्णा}\ N}}\skt{ }\edtext{\edtext{\skt{सैवा\-र्थ\-ज्ञे\-स्ती\-व्र\-क\-ण्डु\-स\-दाहः।}}{\lemma{\skt{सैवार्थज्ञेस्}}\Afootnote{\skt{सैवार्थज्ञैस्}\ N H}}}{\lemma{\skt{तीव्र\-क\-ण्डु\-स\-दाहः}}\Afootnote{H; \skt{तीव्रकण्डू सदोहाः ।}\ N}}\skt{ }\pend
\pstart\versepar \skt{कृष्णं }\edtext{\skt{कण्डुम्मण्लं}}{\Afootnote{\skt{स्निग्धम्मण्डलङ्}\ N; \skt{कण्डुमण्डलं}\ H}}\skt{ }\edtext{\skt{कण्डरञ्च}}{\Afootnote{\skt{कङ्डुरञ्च}\ N H}}\skt{ }\edtext{\skt{श्यामोपेतं}}{\Afootnote{H; \skt{श्यावोपेतङ्}\ N}}\skt{ }\edtext{\edtext{\skt{कैटभम्पथते}}{\lemma{\skt{कैटभम्}}\Afootnote{H; \skt{किटिभम्}\ N}}}{\lemma{\skt{पथते}}\Afootnote{\skt{पथ्यते}\ N; \skt{पठ्यते}\ H}}\skt{ तु ॥}\pend

\pstart\versepar \ledsidenote{\passref{2.5.12.3}}\edtext{\skt{शुक्षाव पैडकाः}}{\Afootnote{\skt{सूक्ष्मा वहुः पिडकाः}\ N; \skt{शुक्ष्मा वह्व्यः पैडकाः}\ H}}\skt{ }\edtext{\skt{आवचत्यः}}{\Afootnote{\skt{स्राववत्यः}\ N H}}\skt{ पामेत्युक्ताः }\edtext{\skt{कण्डु\-म\-त्यु\-ग्र\-दाहाः।}}{\lemma{\skt{कण्डु\-म\-त्यु\-ग्र\-दाहाः}}\Afootnote{\skt{कण्डू\-म\-त्य\-ग्रु\-दाहाः ।}\ N; \skt{कण्डु\-म\-त्यु\-ग्र\-दाहाः ।}\ H}}\skt{ }\pend
\pstart\versepar \skt{}\edtext{\edtext{\skt{सैवा\-स्थे\-दा\-स्ती\-व्र\-दा\-है\-रु\-पे\-ता\-क्ज्ञथाः}}{\lemma{\skt{सैवास्थेदास्}}\Afootnote{\skt{सैवास्फोटास्}\ N H}}}{\lemma{\skt{उपेताक्ज्ञथाः}}\Afootnote{\skt{उपेता ज्ञेया}\ N; \skt{उपेता ज्ञेयाः}\ H}}\skt{ पाण्यौ }\edtext{\skt{कण्डुरुग्राः}}{\Afootnote{\skt{कच्छरुग्राः}\ N; \skt{कण्डुरुग्राः}\ H}}\skt{ }\edtext{\skt{स्थिचौ}}{\Afootnote{\skt{स्थिराश्}\ N; \skt{स्फिचौ}\ H}}\skt{ च ॥}\pend

\pstart \ledsidenote{\passref{2.5.12.4}}\skt{कृत्स्ने देहे }\edtext{\skt{पैड\-काः\-क\-ण्डु\-मत्यो }\edtext{\skt{तैरा\-स्रा\-वा\-रा\-क\-से\-त्यु\-च्यते}}{\lemma{\skt{आस्रा\-वा\-रा\-क\-सेत्य्}}\Afootnote{H; \skt{रास्रा\-वा\-रा\-क\-सेत्य्}\ N}}}{\lemma{\skt{पैड\-काः\-क\-ण्डु\-मत्यो तैर्}}\Afootnote{H; \skt{पैडकाः कण्डूमत्यो नै}\ N}}\skt{ तु । }\edtext{\skt{तत्रा\-रु\-ष्कं\-रा\-क\-सं\-यच्च}}{\Afootnote{N; \skt{तत्रारुष्कं राकसं यच्च}\ H}}\skt{ }\edtext{\skt{सिध्मं}}{\Afootnote{H; \skt{सिध्म}\ N}}\skt{ }\edtext{\skt{कफाधिषदेक}}{\Afootnote{\skt{कंफाधिक्यादेक}\ N; \skt{कफाधिष्कदेक}\ H}}\skt{ कुष्ठम्महच्च । । }\edtext{\skt{पित्ते द्रेका\-त्पा\-रि\-स\-र्प्यन्तु}}{\Afootnote{\skt{पित्तो\-द्रे\-प\-रि\-स\-र्पन्ति}\ N; \skt{पित्तो\-द्रे\-का\-त्पा\-रि\-स\-र्पन्तु}\ H}}\skt{ }\edtext{\skt{विद्यादृष्टो}}{\Afootnote{\skt{विद्यादुष्ठो}\ N; \skt{दुष्टो}\ H}}\skt{ निष्ठः कुष्ठ\-व\-र्ग\-स्त्रि\-दोषः ।}\pend

\pstart \ledsidenote{\passref{2.5.17}}\skt{किलासमपि }\edtext{\skt{कुष्ठविकार}}{\Afootnote{N H; \skt{कुष्ठविकल्प}\ A}}\skt{ एव । }\edtext{\skt{तत्तु}}{\Afootnote{N H; \skt{त}\ A}}\skt{ }\edtext{\edtext{\skt{त्रिविधम्वा}}{\lemma{\skt{त्रिविधम्}}\Afootnote{N H; \skt{त्त्रिविधं}\ A}}\skt{ ते}}{\lemma{\skt{वा ते}}\Afootnote{H; \skt{वातेन}\ A N}}\skt{ पित्तेन }\edtext{\skt{श्लेमणेति}}{\Afootnote{H; \skt{श्लेष्मणा चेति ।}\ A N}}\skt{ । कुष्ठ\-कि\-ला\-स\-यो\-र\-न्तरं }\edtext{\skt{त्वगतमेव}}{\lemma{\skt{त्वगतम्}}\Afootnote{\skt{त्वग्गतम्}\ A N H}}\skt{ किला\-स\-म\-प\-रि\-स्रावि च । }\edtext{\skt{तद्वातेन}}{\Afootnote{N; \skt{तद्वातेन}\ A H}}\skt{ }\edtext{\skt{मण्ड\-ल\-म\-रु\-ण\-म्प\-रि\-ध्वंसि च}}{\Afootnote{\skt{मण्डलमरुणं परुषं परिध्वंसि च,}\ A; \skt{मण्डलमरुणं परिध्वंसि च}\ N H}}\skt{ । पित्तेन }\edtext{\skt{पद्म\-प\-त्र\-प्र\-ती\-का\-सं\-स\-प\-रि\-धाहं च}}{\Afootnote{H; \skt{पद्म\-प\-त्र\-प्र\-ती\-काशं सपरिदाहं च}\ A N}}\skt{ । श्लेष्मणा }\edtext{\skt{स्वतं}}{\Afootnote{\skt{श्वेतं}\ A N; \skt{स्वेतं}\ H}}\skt{ स्निग्धं }\edtext{\skt{कण्डुरञ्च}}{\Afootnote{H; \skt{बहलं कण्डूमच्च ।}\ A; \skt{कणअडूरञ्च ।}\ N}}\skt{ । }\edtext{\skt{तत्र}}{\Afootnote{N H; \skt{तेषु}\ A}}\skt{ }\edtext{\skt{सम्ब\-द्ध\-म\-ण्ड\-ल\-म\-न्त\-जातं}}{\Afootnote{H; \skt{संब\-द्ध\-म\-ण्ड\-ल\-म\-न्ते\-जातं}\ A; \skt{सम्व\-द्ध\-म\-ण्ड\-ल\-म\-न्ते\-जात}\ N}}\skt{ }\edtext{\skt{रक्तरोमाणञ्च}}{\Afootnote{H; \skt{रक्तरोम चा}\ A; \skt{रक्तरोमाणञ्चा}\ N}}\skt{ साध्य\-म\-ग्नि\-द\-ग्धञ्च}\pend

\pstart \ledsidenote{\passref{2.5.18}}\edtext{\skt{तत्र}}{\Afootnote{N H; \emph{om.}\ A}}\skt{ }\edtext{\skt{कुष्ठेषु}}{\Afootnote{N H; \skt{कुष्ठेषु तु}\ A}}\skt{ }\edtext{\skt{त्वक्षंकोच}}{\Afootnote{H; \skt{त्वक्संकोच}\ A N}}\skt{ }\edtext{\skt{स्वेद\-णो\-भ\-स्वा\-प\-भे\-द\-कौ\-ण्य\-स्व\-रो\-प\-घातो}}{\Afootnote{\skt{स्वाप\-स्वे\-द\-शो\-फ\-भे\-द\-कौ\-ण्य\-स्व\-रो\-प\-घाता}\ A; \skt{स्वेद\-शो\-फ\-स्वा\-प\-भे\-द\-कौ\-ण्य\-स्व\-रो\-प\-घातो}\ N H}}\skt{ वातेन । }\edtext{\edtext{\skt{पाका\-व\-द\-र\-ण\-म\-ङ्गु\-लि\-प\-तनं}}{\lemma{\skt{पाका\-व\-द\-र\-णम्}}\Afootnote{N H; \skt{पाका\-व\-द\-र\-णाङ्}\ A}}}{\lemma{\skt{अङ्गुलिपतनं}}\Afootnote{N H; \skt{गुलिपतन}\ A}}\skt{ }\edtext{\skt{कर्ण्ण\-ना\-सा\-भ\-ङ्गा\-क्षि\-रागाः}}{\Afootnote{N H; \skt{कर्ण\-ना\-सा\-भ\-ङ्गा\-क्षि\-राग}\ A}}\skt{ }\edtext{\skt{क्षिप्रोत्यत्तयः}}{\Afootnote{\skt{सत्त्वोत्पत्तयः}\ A; \skt{क्षिपोत्पतयः ।}\ N; \skt{क्षिप्रोत्पत्तयः}\ H}}\skt{ पित्तेन । }\edtext{\skt{कण्डु\-व\-र्ण्ण\-भे\-दः\-शो\-भा\-स्रावाः}}{\Afootnote{\skt{कण्डू\-व\-र्ण\-भे\-द\-शो\-फा\-स्रा\-व\-गौ\-र\-वाणि}\ A; \skt{कण्डूवर्ण्णभेदः शोफः स्रावाः}\ N; \skt{कण्डूवर्ण्णभेदः शोफास्रावाः}\ H}}\skt{ श्लेष्मणा ।}\pend

\pstart \ledsidenote{\passref{2.5.19}}\edtext{\skt{तत्रा\-ति\-ब\-ल\-प्र\-वृत्तं}}{\Afootnote{H; \skt{तत्रा\-दि\-ब\-ल\-प्र\-वृत्तं}\ A; \skt{तत्रा\-ति\-व\-ल\-प्र\-वृ\-त्तम्}\ N}}\skt{ }\edtext{\skt{पुण्डरीकं}}{\Afootnote{N H; \skt{पौण्डरीकं}\ A}}\skt{ }\edtext{\skt{काक\-ण\-ञ्चा\-सा\-ध्य\-मिति ॥ श्लोकौ}}{\lemma{\skt{इति ॥ श्लोकौ}}\Afootnote{\emph{om.}\ A; \skt{इति ॥}\ N H}}\skt{ ॥}\pend

\pstart\versepar \ledsidenote{\passref{2.5.20}}\edtext{\skt{श्लोकौ}}{\Afootnote{H; \skt{भवन्ति चात्र ---}\ A; \skt{भ वति चात्र}\ N}}\skt{ ॥ }\pend
\pstart\versepar \skt{यथा }\edtext{\edtext{\skt{वनस्पतिर्ज्जातः}}{\lemma{\skt{वनस्पतिर्}}\Afootnote{H; \skt{बनस्पति}\ A; \skt{वनर्स्पतिर्}\ N}}}{\lemma{\skt{ज्जातः}}\Afootnote{N H; \skt{र्जातः}\ A}}\skt{ प्राप्य }\edtext{\skt{कालं}}{\Afootnote{H; \skt{काल}\ A N}}\skt{ }\edtext{\skt{क्रमेण तु}}{\Afootnote{N H; \skt{प्रकर्षणम् ।}\ A}}\skt{ । }\pend
\pstart\versepar \skt{अन्त\-र्भू\-मि\-म्वि\-गा\-हेत }\edtext{\skt{मूलै\-र्वृ\-ष्टि\-वि\-व\-र्द्धितैः}}{\lemma{\skt{वृष्टिविवर्द्धितैः}}\Afootnote{A H; \skt{वृष्टिविवर्द्धितः ॥}\ N}}\skt{ ॥}\pend

\pstart\versepar \ledsidenote{\passref{2.5.21}}\skt{एवं }\edtext{\skt{कुष्ठं}}{\Afootnote{A H; \skt{कुष्ठ}\ N}}\skt{ }\edtext{\skt{समुत्पन्नं}}{\Afootnote{A H; \skt{समुत्पन्न}\ N}}\skt{ त्वचि कालप्रकर्षतः । }\pend
\pstart\versepar \skt{क्रमेण }\edtext{\skt{धातून्प्राप्नोति}}{\Afootnote{H; \skt{धातून्व्याप्नोति}\ A; \skt{धातूं प्राप्नोति}\ N}}\skt{ }\edtext{\skt{नर\-स्या\-प्र\-ति\-का\-रिणः}}{\Afootnote{A H; \skt{नर\-स्या\-प्र\-ति\-का\-रणः ॥}\ N}}\skt{ ।}\pend

\pstart \ledsidenote{\passref{2.5.22}}\edtext{\skt{त्वत्स्थे}}{\Afootnote{\skt{स्पर्शहानिः स्वेदनत्व}\ A; \skt{त्वक्स्थै}\ N; \skt{त्वक्स्थे}\ H}}\skt{ }\edtext{\edtext{\skt{वैवर्ण्ण्यमङ्गेषु}}{\lemma{\skt{वैवर्ण्ण्यम्}}\Afootnote{H; \emph{om.}\ A; \skt{वैवर्ण्णम्}\ N}}}{\lemma{\skt{अङ्गेषु}}\Afootnote{N H; \skt{मीष}\ A}}\skt{ }\edtext{\skt{कुष्ठे}}{\Afootnote{H; \skt{त्कण्}\ A; \skt{कुष्ठै}\ N}}\skt{ }\edtext{\skt{रौक्ष्यञ्च}}{\Afootnote{H; \skt{डूश्च}\ A; \skt{रौक्ष्यैश्च}\ N}}\skt{ }\edtext{\skt{जायते}}{\Afootnote{N H; \skt{जायते । वैवर्ण्यं रूक्षभावश्च कुष्ठे त्वचि समाश्रिते ॥}\ A}}\skt{ ।}\pend

\pstart\versepar \ledsidenote{\passref{2.5.23}}\edtext{\skt{त्वकस्रावो}}{\Afootnote{H; \skt{त्वक्स्वापो}\ A; \skt{त्वक्स्रवो}\ N}}\skt{ }\edtext{\skt{रोममहर्षञ्च}}{\Afootnote{H; \skt{रोमहर्षश्च}\ A N}}\skt{ }\edtext{\skt{स्वेदस्याति}}{\Afootnote{N H; \skt{स्वेदस्याभि}\ A}}\skt{ प्रवर्त्तनम् ॥ }\pend
\pstart\versepar \skt{}\edtext{\skt{कण्डु\-र्व्वि\-पु\-य\-क\-श्चैव}}{\lemma{\skt{कण्डु\-र्व्वि\-पु\-य\-कश्}}\Afootnote{\skt{कण्डू\-र्वि\-पू\-य\-कश्}\ A N H}}\skt{ कुष्ठे शोणितसंश्रिते ।}\pend

\pstart\versepar \ledsidenote{\passref{2.5.24}}\edtext{\skt{बाहुल्यं}}{\Afootnote{A H; \skt{वाहुल्यम्}\ N}}\skt{ वक्त्रशोषश्च }\edtext{\skt{कार्क्क\-ष्य\-म्पि\-ड\-को\-द्गमः}}{\Afootnote{H; \skt{कार्कश्यं पिडकोद्गमः ।}\ A; \skt{कार्कश्ये पिडकोद्गमः ॥}\ N}}\skt{ ॥ }\pend
\pstart\versepar \skt{तोदः }\edtext{\skt{स्फोटस्थिरत्वञ्च}}{\Afootnote{H; \skt{स्फोटः स्थिरत्वं च}\ A N}}\skt{ कुष्ठे मांससमाश्रिते ॥}\pend

\pstart\versepar \ledsidenote{\passref{2.5.27}}\edtext{\edtext{\edtext{\skt{कौण्य\-ङ्ग\-ति\-क्ष\-योऽ\-ङ्गा\-ना\-म्भेदः}}{\lemma{\skt{गतिक्षयो}}\Afootnote{A N; \skt{गतिक्षयोङ्}\ H}}}{\lemma{\skt{ऽङ्गानाम्}}\Afootnote{A N; \skt{गानाम्}\ H}}}{\lemma{\skt{भेदः}}\Afootnote{H; \skt{संभेदः}\ A; \skt{भेद}\ N}}\skt{ }\edtext{\skt{क्षतविसर्प्यणः}}{\Afootnote{\skt{क्षतसर्पणम् ।}\ A; \skt{क्षतविसर्प्पणं ॥}\ N; \skt{क्षतविसर्प्यणं ॥}\ H}}\skt{ ॥ }\pend
\pstart\versepar \skt{}\edtext{\skt{मेदः}}{\Afootnote{H; \skt{शुक्र}\ A; \skt{मेद}\ N}}\skt{ स्थानगते }\edtext{\skt{लिंगन्पूर्व्वोक्तानि}}{\Afootnote{\skt{लिङ्गं प्रागुक्तानि}\ A; \skt{लिङ्गे पूर्वोक्तानि}\ N; \skt{लिंगम्पूर्व्वोक्तानि}\ H}}\skt{ तथैव च ।}\pend

\pstart\versepar \ledsidenote{\passref{2.5.26}}\edtext{\skt{नासा\-क्षि\-भ\-ङ्गो\-रा\-गाश्च}}{\Afootnote{H; \skt{नासा\-भ\-ङ्गोऽ\-क्षि\-रा\-गश्च}\ A; \skt{नासाक्षितङ्गो रागश्च}\ N}}\skt{ क्षते च क्रिमिसम्भवः ॥ }\pend
\pstart\versepar \skt{}\edtext{\skt{स्वरोपघातश्च}}{\Afootnote{N H; \skt{भवे\-त्स्व\-रो\-प\-घा\-तश्च}\ A}}\skt{ }\edtext{\skt{भवेदस्थिमज्जा}}{\Afootnote{\skt{ह्यस्थिमज्ज}\ A; \skt{भवेदस्थिमज्ज}\ N; \skt{भवे𑑎𑑎𑑎𑑎𑑎𑑎𑑎𑑎𑑎दस्थिमज्जा}\ H}}\skt{ समाश्रिते ॥}\pend

\pstart\versepar \ledsidenote{\passref{2.5.28}}\edtext{\skt{स्त्रीपुंसोः}}{\Afootnote{N H; \skt{स्त्रीपुंसयोः}\ A}}\skt{ }\edtext{\skt{कुष्ठ\-वा\-ह्न\-ल्या\-द्दु\-ष्ट\-शो\-णि\-त\-शु\-क्रयोः}}{\Afootnote{\skt{कुष्ठ\-दो\-षा\-द्दु\-ष्ट\-शो\-णि\-त\-शु\-क्रयोः ।}\ A; \skt{कुष्ठ\-वा\-हु\-ल्या\-द्दु\-ष्ट\-शो\-णि\-त\-शु\-क्रयोः}\ N H}}\skt{ । }\pend
\pstart\versepar \skt{यदपत्यं }\edtext{\skt{भवेज्जातं}}{\Afootnote{N H; \skt{तयोर्जातं}\ A}}\skt{ }\edtext{\skt{ज्ञेयन्तद पि}}{\lemma{\skt{तद पि}}\Afootnote{H; \skt{तदपि}\ A N}}\skt{ }\edtext{\skt{कुष्ठिलं}}{\Afootnote{N H; \skt{कुष्ठितम् ॥}\ A}}\skt{ ।}\pend

\pstart\versepar \ledsidenote{\passref{2.5.29}}\skt{कुष्ठमात्मवतः साध्यं त्वग्र\-क्त\-पि\-शि\-ता\-श्रितं ॥ }\pend
\pstart\versepar \skt{मेदोगतं भवे\-द्या\-प्य\-म\-सा\-ध्य\-मत उत्तरं ।}\pend

\pstart\versepar \ledsidenote{\passref{2.5.30}}\skt{देव\-द्र\-व्य\-गु\-रू\-द्र\-व्य\-प\-र\-दा\-रा\-भि\-म\-र्ष\-णात् । }\pend
\pstart\versepar \skt{}\edtext{\skt{पाप्मा\-पा\-प\-कृ\-त\-मे\-त\-त्कु\-ष्ठ\-मि\-त्य\-भि\-श\-ब्दितं ।}}{\lemma{\skt{कृतम्… अभिशब्दितं}}\Afootnote{\emph{om.}\ A}}\edtext{}{\lemma{\skt{देव\-द्र\-व्य\-गु\-रू\-द्र\-व्य\-प\-र\-दा\-रा\-भि\-म\-र्ष\-णात् । पाप्मा\-पा\-प\-कृ\-तम्}}\Afootnote{H; \skt{ब्रह्म\-स्त्री\-स\-ज्ज\-न\-व\-ध\-प\-र\-स्व\-ह\-र\-णा\-दिभिः । कर्मभिः पाप}\ A; \skt{देव\-द्र\-व्य\-गु\-रु\-द्र\-व्य\-प\-र\-दा\-रा\-भि\-म\-र्श\-नात् ॥ पाप्मा\-पा\-प\-कृ\-ताम्}\ N}}\pend

\pstart \ledsidenote{\passref{2.5.31}}\edtext{\skt{म्रियते}}{\Afootnote{A H; \skt{मृयते}\ N}}\skt{ यदि कुष्ठेन }\edtext{\skt{पुनर्ज्जाते न}}{\Afootnote{H; \skt{पुनर्जातेऽपि}\ A; \skt{पुनर्जाते पि}\ N}}\skt{ }\edtext{\skt{गच्छति}}{\Afootnote{N; \skt{गच्छति । गातः कष्टरतो रोगो यथा कुष्ठं प्रकीर्तितम् ॥}\ A; \skt{गच्छति ॥ 𑑎𑑎𑑎𑑎𑑎𑑎𑑎𑑎𑑎𑑎𑑎𑑎𑑎𑑎𑑎𑑎𑑎𑑎𑑎}\ H}}\skt{ ॥}\pend

\pstart\versepar \ledsidenote{\passref{2.5.32}}\skt{आहाराचारयोः प्रोक्तामास्थाय }\edtext{\skt{महती}}{\Afootnote{\skt{महतीं}\ A N H}}\skt{ }\edtext{\skt{क्रिंयाम्}}{\Afootnote{H; \skt{क्रियाम्}\ A N}}\skt{ ॥ }\pend
\pstart\versepar \skt{}\edtext{\edtext{\skt{औष\-धी\-ना\-म्व्वि\-षि\-ष्टानां}}{\lemma{\skt{औषधीनाम्}}\Afootnote{A H; \skt{ओषधीनाम्}\ N}}}{\lemma{\skt{व्विषिष्टानां}}\Afootnote{\skt{विशिष्टानां}\ A; \skt{विशैष्टानान्}\ N; \skt{विषिष्टानां}\ H}}\skt{ तपसश्च निषेवणात् ॥ }\pend
\pstart\versepar \skt{यस्तेन मुच्यते }\edtext{\skt{जन्तुः}}{\Afootnote{N H; \skt{जन्तुः स}\ A}}\skt{ }\edtext{\edtext{\edtext{\skt{पुण्या\-ङ्ग\-ति\-म\-वा\-प्नु\-यात्}}{\lemma{\skt{पुण्याङ्}}\Afootnote{A H; \skt{पुण्या}\ N}}}{\lemma{\skt{गतिम्}}\Afootnote{A H; \skt{द्गतिम्}\ N}}}{\lemma{\skt{अवाप्नुयात्}}\Afootnote{N H; \skt{आप्नुयात् ॥}\ A}}\skt{ ।}\edtext{}{\lemma{\skt{निषेवणात् ॥ यस्}}\Afootnote{A H; \skt{निशेवणात् ॥ यस्}\ N}}\pend

\pstart \ledsidenote{\passref{2.5.33ab}}\edtext{\skt{प्रवा\-ता\-द्गा\-त्र\-सं\-स्प\-र्शा\-न्नि\-श्वा\-सा\-त्स\-ह\-भो\-ज\-नात्}}{\lemma{\skt{प्रवा\-ता\-द्गा\-त्र\-सं\-स्प\-र्शान्}}\Afootnote{H; \skt{प्रस\-ङ्गा\-द्गा\-त्र\-सं\-स्प\-र्शान्}\ A; \skt{प्रवात गात्रसम्पर्शा}\ N}}\skt{ ॥}\pend

\pstart\versepar \ledsidenote{\passref{2.5.34}}\skt{कुष्ठञ्ज्वरश्च शोषश्च }\edtext{\skt{नेत्राभिष्यन्द}}{\Afootnote{A H; \skt{नेत्राभिष्यन्दि}\ N}}\skt{ एव च । }\pend
\pstart\versepar \skt{औप\-स\-र्गि\-क\-रो\-गाश्च संक्रामन्ति }\edtext{\skt{नरात्नरम्}}{\Afootnote{N H; \skt{नरान्नरम् ॥}\ A}}\skt{ ॥}\pend

\pstart \ledsidenote{\passref{2.5.trailer}}\edtext{\skt{इति ५ कुष्ठनिदाने}}{\Afootnote{\skt{इति सुश्रु\-त\-सं\-हि\-तायां निदानस्थाने कुष्ठनिदानं नाम}\ A; \skt{इति ॥ ॥ ओ ॥}\ N; \skt{इति ॥ ज ॥ कुष्ठनिदाने}\ H}}\skt{ }\edtext{\skt{पञ्चमः}}{\Afootnote{H; \skt{पञ्चमोऽध्यायः ॥५ ॥}\ A; \emph{om.}\ N}}\skt{ ।}\pend

\endnumbering
\beginnumbering
\eledsection*{Adhyāya 6}
\pstart \ledsidenote{\passref{2.6.1}}\skt{अथातः }\edtext{\edtext{\edtext{\skt{प्रमे\-हा\-णा\-न्नि\-दा\-न\-म्व्या\-ख्या\-स्यामः}}{\lemma{\skt{प्रमेहाणान्}}\Afootnote{H; \skt{प्रमेह}\ A; \skt{प्रमेहानान्}\ N}}}{\lemma{\skt{निदानम्}}\Afootnote{A H; \skt{निदानाध्यायां}\ N}}}{\lemma{\skt{व्याख्यास्यामः}}\Afootnote{A H; \skt{व्यांव्याखाश्यामः ॥}\ N}}\skt{ ॥}\pend

\pstart \ledsidenote{\passref{2.6.3}}\edtext{\edtext{\skt{दिवा\-स्व\-प्न\-प्र\-स\-क्त\-म\-लसं}}{\lemma{\skt{दिवा\-स्व\-प्न\-प्र\-स\-क्तम्}}\Afootnote{H; \skt{दिवा\-स्व\-प्ना\-व्या\-या\-मा\-ल\-स्य\-प्र\-सक्तं}\ A; \skt{दिवा\-स्व\-प्न\-प्र\-स\-क्ता\-नाम्}\ N}}}{\lemma{\skt{अलसं}}\Afootnote{N H; \emph{om.}\ A}}\skt{ }\edtext{\skt{शीत\-स्नि\-ग्ध\-म\-धु\-र\-मे\-द्य\-द्र\-वा\-न्न\-पा\-न\-म्पु\-रु\-ष\-ञ्जा\-नी\-या\-त्प्र\-मेही}}{\lemma{\skt{शीत\-स्नि\-ग्ध\-म\-धु\-र\-मे\-द्य\-द्र\-वा\-न्न\-पा\-नम्}}\Afootnote{\skt{शीत\-स्नि\-ग्ध\-म\-धु\-र\-मे\-द्य\-द्र\-वा\-न्न\-पा\-न\-से\-विनं}\ A N; \skt{शीत\-स्नि\-ग्ध\-म\-धु\-र\-मे\-ध्य\-द्र\-वा\-न्न\-पा\-नम्}\ H}}\skt{ भविष्यतीति ।}\pend

\pstart \ledsidenote{\passref{2.6.4}}\edtext{\skt{तस्यै\-वं\-प्र\-वृ\-त्तस्य}}{\Afootnote{N H; \skt{तस्य चैवंप्रवृत्तस्या}\ A}}\skt{ }\edtext{\skt{यदा}}{\Afootnote{N H; \skt{परिपक्वा एव}\ A}}\skt{ }\edtext{\skt{वात\-पि\-त्त\-श्ले\-ष\-माणो}}{\Afootnote{H; \skt{वात\-पि\-त्त\-श्ले\-ष्माणो यदा मेदसा सहै\-क\-त्व\-मु\-पेत्य}\ A; \skt{वात\-पि\-त्त\-श्ले\-ष्माणो मेद\-श्चा\-प\-रि\-प\-क्वम्}\ N}}\skt{ }\edtext{\edtext{\skt{मेद\-श्चो\-प\-रि\-प\-क्वा\-धो\-गत्वा}}{\lemma{\skt{मेदश्}}\Afootnote{H; \skt{मूत्र}\ A; \emph{om.}\ N}}}{\lemma{\skt{चोप\-रि\-प\-क्वा\-धो\-गत्वा}}\Afootnote{\skt{वाहि\-स्रो\-तां\-स्य\-नु\-सृ\-त्याधो गत्वा}\ A; \skt{अधो गत्वा}\ N; \skt{चोपरिपक्वाधो गत्वा}\ H}}\skt{ }\edtext{\edtext{\skt{बस्ति\-मु\-ख\-मा\-श्रित्य}}{\lemma{\skt{बस्तिमुखम्}}\Afootnote{H; \skt{बस्तेर्मुख}\ A; \skt{वस्तिमुखम्}\ N}}}{\lemma{\skt{आश्रित्य}}\Afootnote{N H; \skt{माश्रित्य}\ A}}\skt{ निर्भिद्यन्ते तदा }\edtext{\skt{प्रमे\-हा\-न्ज\-न\-यन्ति}}{\Afootnote{H; \skt{प्रमे\-हा\-ञ्ज\-न\-यन्ति}\ A N}}\skt{ ।}\pend

\pstart \ledsidenote{\passref{2.6.5}}\edtext{\edtext{\skt{तेषा\-म्पू\-र्व्व\-रू\-पाणि}}{\lemma{\skt{तेषाम्}}\Afootnote{N H; \skt{तेषां तु}\ A}}}{\lemma{\skt{पूर्व्वरूपाणि}}\Afootnote{N H; \skt{पूर्वरू}\ A}}\skt{ }\edtext{\skt{पाणि\-पा\-द\-त\-ल\-दाहः}}{\Afootnote{H; \skt{पाणि हस्त\-पा\-द\-त\-ल\-दाहः}\ A; \skt{पाणि\-ह\-पा\-द\-त\-ल\-दाहः ॥}\ N}}\skt{ }\edtext{\skt{स्निग्धः\-चि\-क्क\-न\-गा\-त्रता}}{\Afootnote{\skt{स्निग्ध\-पि\-च्छि\-ल\-गु\-रुता}\ A; \skt{स्विन्नचिकूण}\ N; \skt{स्निग्ध\-चि\-क्क\-न\-गा\-त्रता}\ H}}\skt{ }\edtext{\skt{मधु\-र\-शु\-क्ल\-मू\-त्रता}}{\Afootnote{H; \skt{गात्राणां मधु\-र\-शु\-क्ल\-मू\-त्रता}\ A; \skt{ग्राततां मधु\-र\-सू\-शु\-क्ल\-मू\-त्रता}\ N}}\skt{ तन्द्रा }\edtext{\skt{च}}{\Afootnote{N H; \skt{सादः पिपासा दुर्गन्धश्च श्वास\-स्ता\-लु\-ग\-ल\-जि\-ह्वा\-द\-न्तेषु मलो\-त्प\-त्ति\-र्ज\-टि\-ली\-भावः केशानां वृद्धिश्च नखानाम् ॥}\ A}}\skt{ ।}\pend

\pstart \ledsidenote{\passref{2.6.6}}\skt{तत्रा\-वि\-ल\-प्र\-भू\-त\-मू\-त्र\-ल\-क्षणाः सर्व्व एव }\edtext{\skt{प्रमेहा}}{\Afootnote{A H; \skt{प्रमेही}\ N}}\skt{ भवन्ति ।}\pend

\pstart \ledsidenote{\passref{2.6.7}}\skt{सर्व्व एव }\edtext{\skt{सर्व्व\-दो\-ष\-स\-मु\-त्थाश्च}}{\Afootnote{N H; \skt{सर्व\-दो\-ष\-स\-मुत्थाः}\ A}}\skt{ सह पिडकाभिः।}\pend

\pstart \ledsidenote{\passref{2.6.8}}\skt{तत्र }\edtext{\edtext{\skt{कफा\-दु\-द\-के\-क्षु\-वा\-लि\-का\-सु\-रा\-सि\-क\-ता\-श\-नै\-र्ल्ल\-व\-ण\-पि\-ष्ट\-सा\-न्द्र\-शु\-क्र\-मेहाः}}{\lemma{\skt{कफा\-दु\-द\-के\-क्षु\-वा\-लि\-का\-सु\-रा\-सि\-क\-ता\-श\-नैर्}}\Afootnote{\skt{कफा\-दु\-द\-के\-क्षु\-वा\-लि\-का\-सु\-रा\-सि\-क\-ता\-श\-नैर्}\ A H; \skt{कफा\-दु\-द\-का\-क्षु\-व\-लि\-का\-सु\-र\-रा\-सि\-क\-ता\-श\-नैर्}\ N}}}{\lemma{\skt{ल्लव\-ण\-पि\-ष्ट\-सा\-न्द्र\-शु\-क्र\-मेहाः}}\Afootnote{N H; \skt{लव\-ण\-पि\-ष्ट\-सा\-न्द्र\-शुक्र}\ A}}\skt{ }\edtext{\edtext{\skt{फेनमेहश्चेति}}{\lemma{\skt{फेनमेहश्}}\Afootnote{H; \skt{फेनमेहा}\ A; \skt{फेनमेहाश्}\ N}}}{\lemma{\skt{चेति}}\Afootnote{N H; \emph{om.}\ A}}\skt{ दश ॥ }\edtext{\skt{पित्ता\-न्नी\-ल\-ह\-रि\-द्रा\-म्व्ल\-क्षा\-र\-म\-ञ्जि\-ष्ठा\-मेहाः}}{\Afootnote{H; \skt{साध्याः, दोषदूष्याणां समक्रियत्वात् ; पित्ता\-न्नी\-ल\-ह\-रि\-द्रा\-म्ल\-क्षा\-र\-म\-ञ्जिष्ठा}\ A; \skt{पित्ता\-न्नी\-ल\-ह\-रि\-द्रा\-म्ल\-क्षा\-र\-म\-ञ्जि\-ष्था\-मेहाः}\ N}}\skt{ }\edtext{\edtext{\skt{शोणि\-त\-मे\-ह\-श्चेति}}{\lemma{\skt{शोणितमेहश्}}\Afootnote{H; \skt{शोणितमेहाः}\ A N}}}{\lemma{\skt{चेति}}\Afootnote{N H; \skt{षड्याप्याः, दोषदूष्याणां विषमक्रि}\ A}}\skt{ }\edtext{\skt{षट् ॥ }\edtext{\skt{वाता\-त्स\-र्प्पि\-र्व्व\-सा\-क्षौ\-द्र\-मे\-हा\-ह\-स्ति\-मे\-ह\-श्चेति}}{\lemma{\skt{चेति}}\Afootnote{N H; \emph{om.}\ A}}}{\lemma{\skt{षट् ॥ वाता\-त्स\-र्प्पि\-र्व्व\-सा\-क्षौ\-द्र\-मे\-हा\-ह\-स्ति\-मे\-हश्}}\Afootnote{H; \skt{यत्वात् ; वाता\-त्स\-र्पि\-र्व\-सा\-क्षौ\-द्र\-ह\-स्ति\-मे\-हाश्}\ A; \skt{षड् ॥ वाता\-त्स\-र्पि\-र्व\-सा\-क्षौ\-द्र\-मे\-हहा हस्तिमेहश्}\ N}}\skt{ }\edtext{\skt{चत्वारोऽसाध्याः॥}}{\lemma{\skt{ऽसाध्याः}}\Afootnote{N H; \skt{ऽसाध्यतमाः, महा\-त्य\-यि\-क\-त्वात् ॥}\ A}}\pend

\pstart\versepar \ledsidenote{\passref{2.6.9}}\skt{तत्र वात\-पि\-त्त\-मे\-दो\-भि\-र\-न्वितः }\edtext{\skt{श्लेष्मा\-स्वा\-न्मे\-हा\-न्ज\-न\-यन्ति}}{\Afootnote{H; \skt{श्लेष्मा श्लेष्म\-प्र\-मे\-हा\-ञ्ज\-न\-यति,}\ A; \skt{श्लेष्मश्वां मो मे हांजनयति ।}\ N}}\skt{ । }\pend
\pstart\versepar \skt{वात\-क\-फ\-शो\-णि\-त\-मे\-दो\-भि\-र\-न्वि\-त\-म्पित्तं । }\edtext{\skt{कफ\-पि\-त्त\-व\-सा\-म\-ज्जा\-मे\-दो\-भि\-र\-न्वितो}}{\lemma{\skt{कफ\-पि\-त्त\-व\-सा\-म\-ज्जा\-मे\-दो\-भिर्}}\Afootnote{H; \skt{पित्तप्रमेहान्, कफ\-पि\-त्त\-व\-सा\-म\-ज्ज\-मे\-दो\-भिर्}\ A; \skt{पित्त\-व\-सा\-म\-ज्जा\-मे\-दो\-भिर्}\ N}}\skt{ }\edtext{\edtext{\skt{वायुरिति}}{\lemma{\skt{वायुर्}}\Afootnote{N H; \skt{वायु}\ A}}}{\lemma{\skt{इति}}\Afootnote{N H; \skt{र्वातप्रमेहान् ॥}\ A}}\skt{ ॥}\pend

\pstart \ledsidenote{\passref{2.6.10}}\skt{तत्र }\edtext{\edtext{\edtext{\edtext{\skt{स्वेत\-म\-वे\-द\-न\-मु\-द\-क\-तु\-ल्य\-मु\-द\-क\-मेही}}{\lemma{\skt{स्वेतम्}}\Afootnote{H; \skt{श्वेत}\ A; \skt{श्वेतम्}\ N}}}{\lemma{\skt{अवेदनम्}}\Afootnote{N H; \skt{मवेदन}\ A}}}{\lemma{\skt{उदकतुल्यम्}}\Afootnote{N H; \skt{मुदकसदृश}\ A}}}{\lemma{\skt{उदकमेही}}\Afootnote{N H; \skt{मुदकमेही}\ A}}\skt{ मेहति । }\edtext{\edtext{\skt{ईक्षु\-र\-स\-तु\-ल्य\-मि\-क्षु\-वा\-लि\-का\-मेही}}{\lemma{\skt{ईक्षु\-र\-स\-तु\-ल्यम्}}\Afootnote{H; \skt{इक्षुरसतुल्य}\ A; \skt{इक्षु\-र\-स\-तु\-ल्यम्}\ N}}\skt{ ।सुरातुल्यं}}{\lemma{\skt{इक्षु\-वा\-लि\-का\-मेही ।सुरातुल्यं}}\Afootnote{H; \skt{मिक्षु\-वा\-लि\-का\-मेही; सुरातुल्यं}\ A; \skt{इक्षु\-व\-लि\-का\-मेही ॥ सुरातुल्यं}\ N}}\skt{ सुरामेही }\edtext{\skt{सरुजां}}{\Afootnote{H; \skt{सरुजं}\ A N}}\skt{ सिकतानुविद्धं सिकतामेही ।शनैः सकफं }\edtext{\skt{सान्द्रं}}{\Afootnote{H; \skt{मृत्स्नं}\ A; \skt{सा}\ N}}\skt{ }\edtext{\edtext{\skt{शनैर्म्मेही}}{\lemma{\skt{शनैर्}}\Afootnote{A N; \skt{शनै}\ H}}}{\lemma{\skt{म्मेही}}\Afootnote{A; \skt{मही ॥}\ N; \skt{र्म्मेही ।}\ H}}\skt{ । }\edtext{\skt{विषदं}}{\Afootnote{H; \skt{विशदं}\ A N}}\skt{ लवणतुल्यं }\edtext{\skt{लवणमेही ।हृ\-ष्ट\-रो\-मा\-पि\-ष्ट\-तुल्यं}}{\Afootnote{N H; \skt{लवणमेही; हृष्टरोमा पिष्टरसतुल्यं}\ A}}\skt{ पिष्टमेही । आविलं सान्द्रं सान्द्रमेही । }\edtext{\skt{शुक्रतुल्यं}}{\Afootnote{A H; \skt{शुक्रतु.यं}\ N}}\skt{ शुक्रमेही ॥ }\edtext{\edtext{\skt{स्तोकस्तोकं}}{\lemma{\skt{स्तोकस्}}\Afootnote{\skt{स्तोकं}\ A; \skt{स्तोक}\ N H}}}{\lemma{\skt{तोकं}}\Afootnote{\skt{स्तोकं}\ A N H}}\skt{ }\edtext{\skt{सफेनञ्च}}{\Afootnote{N H; \skt{सफेनम}\ A}}\skt{ }\edtext{\skt{फेनमेही}}{\Afootnote{H; \skt{च्छं फेनमेही}\ A; \skt{फनमेही}\ N}}\skt{ }\edtext{\skt{मेहति}}{\Afootnote{A H; \skt{महति ॥}\ N}}\skt{ ॥}\pend

\pstart\versepar \ledsidenote{\passref{2.6.11}}\skt{अत }\edtext{\edtext{\skt{ऊर्ध्द्व\-म्पि\-त्त\-नि\-मि\-त्ता\-न्व\-क्ष्यामः}}{\lemma{\skt{ऊर्ध्द्वम्}}\Afootnote{H; \skt{ऊर्ध्वं}\ A N}}}{\lemma{\skt{पित्त\-नि\-मि\-त्ता\-न्व\-क्ष्यामः}}\Afootnote{A H; \skt{पित्त\-नि\-मि\-त्त\-म्व\-क्ष्यामः ॥}\ N}}\skt{ । }\pend
\pstart\versepar \skt{}\edtext{\edtext{\skt{सफे\-न\-म\-च्छ\-न्नी\-ल\-न्नी\-ल\-मेही}}{\lemma{\skt{सफेनम्}}\Afootnote{H; \skt{सफेन}\ A N}}}{\lemma{\skt{अच्छ\-न्नी\-ल\-न्नी\-ल\-मेही}}\Afootnote{H; \skt{मच्छं नीलं नीलमेही}\ A; \skt{मच्छन्नीलमही}\ N}}\skt{ मेहति । सदाहं }\edtext{\skt{कद्कं}}{\Afootnote{\emph{om.}\ A; \skt{ककं}\ N; \skt{कटुकं}\ H}}\skt{ हरिद्राभं हरिद्रामेही }\edtext{\skt{मेहति}}{\Afootnote{N H; \emph{om.}\ A}}\skt{ । }\edtext{\skt{अम्ल\-र\-स\-ग\-न्ध\-म\-म्ल\-मेही}}{\Afootnote{A H; \skt{अम्व्ल\-र\-स\-ग\-न्ध\-म\-म्व्ल\-मेही ।}\ N}}\skt{ । }\edtext{\skt{सृत\-क्षा\-र\-प्र\-तिमं}}{\Afootnote{H; \skt{स्रुत\-क्षा\-र\-प्र\-तिमं}\ A N}}\skt{ क्षारमेही । }\edtext{\skt{मञ्जि\-ष्ठो\-द\-क\-प्र\-काशं}}{\Afootnote{A H; \skt{मञ्जिष्ठदकं प्रकाश}\ N}}\skt{ }\edtext{\skt{मञ्जिष्ठमेही}}{\Afootnote{N H; \skt{मञ्जिष्ठामेही;}\ A}}\skt{ । }\edtext{\skt{शोणितमेही}}{\Afootnote{N H; \skt{शोणितप्रकाशं शोणितमेही}\ A}}\skt{ }\edtext{\skt{शोणितम्मेहति}}{\lemma{\skt{शोणितम्}}\Afootnote{N H; \emph{om.}\ A}}\skt{ ॥}\pend

\pstart \ledsidenote{\passref{2.6.12}}\edtext{\skt{वात\-नि\-मि\-त्ता\-न्व\-क्ष्यामः}}{\Afootnote{H; \skt{अत ऊर्ध्वं वात\-नि\-मि\-त्ता\-न्व\-क्ष्यामः}\ A; \skt{वात\-नि\-मि\-त्ता\-न्व\-क्षामः ।}\ N}}\skt{ ॥ सर्प्पिःप्रकाशं }\edtext{\skt{सर्प्पिर्म्मेही}}{\Afootnote{A H; \skt{सर्पिम्मेही}\ N}}\skt{ मेहति । वसाप्रकाशं वसामेही । }\edtext{\skt{क्षौद्र\-स\-व\-र्ण्णं\-क्षौ\-द्र\-मेही ।म\-त्त\-मा\-त\-ङ्ग\-व\-द\-न\-प्र\-धू\-रं\-ह\-स्ति\-मे\-हि\-मे\-हति}}{\Afootnote{H; \skt{क्षौद्ररसवर्णं क्षौद्रमेही; मत्त\-मा\-त\-ङ्ग\-व\-द\-नु\-प्र\-बन्धं हस्तिमेही मेहति ॥}\ A; \skt{क्षौद्र\-स\-व\-र्ण्ण\-क्षौ\-द्र\-मेही ॥ मत्त\-मा\-त\-ङ्ग\-व\-द\-न\-प्र\-वद्धं मधुरं हस्तिमेही मेहति ॥}\ N}}\skt{ ॥}\pend

\pstart\versepar \ledsidenote{\passref{2.6.13}}\edtext{\edtext{\skt{उपद्रवानत}}{\Afootnote{H; \skt{उपद्रवानते}\ N}}\skt{ ऊर्ध्वम्वक्ष्यामः ॥ }\edtext{\skt{मक्षि\-को\-प\-स\-र्प्प\-ण\-मा\-ल\-स्य\-मा\-स्यो\-प\-देहः}}{\Afootnote{H; \skt{मेद\-क्षि\-को\-प\-स\-र्प\-ण\-मा\-लस्य मास्योपदेहः ।}\ N}}}{\lemma{\skt{… म}}\Afootnote{\emph{om.}\ A}}\skt{ }\edtext{\skt{प्रतिश्यायः}}{\Afootnote{A H; \skt{प्रतिश्याय}\ N}}\skt{ । }\edtext{\edtext{\edtext{\skt{शैथि\-ल्य\-म\-रो\-च\-काऽ\-वि\-पाकौ}}{\lemma{\skt{शैथिल्यम्}}\Afootnote{\skt{शैथिल्या}\ A; \skt{सैथिल्यम्}\ N; \skt{शैथिल्य}\ H}}}{\lemma{\skt{अरोचका}}\Afootnote{N; \skt{रोचका}\ A; \skt{मरोचका}\ H}}}{\lemma{\skt{ऽविपाकौ}}\Afootnote{H; \skt{विपाकाः}\ A; \skt{विपाको}\ N}}\skt{ }\edtext{\edtext{\skt{कफ\-प्र\-से\-क\-श्छ\-र्द्यति}}{\lemma{\skt{कफप्रसेकश्}}\Afootnote{H; \skt{कफप्रसेक}\ A N}}}{\lemma{\skt{छर्द्यति}}\Afootnote{H; \skt{च्छर्दि}\ A; \skt{स्छर्द्यति}\ N}}\skt{ निद्राकास }\edtext{\skt{इति}}{\Afootnote{N H; \skt{श्वासाश्चेति}\ A}}\skt{ }\edtext{\skt{श्लेष्म\-जा\-ना\-मु\-प\-द्रवा}}{\lemma{\skt{उपद्रवा}}\Afootnote{A H; \skt{उपद्र}\ N}}\skt{ }\edtext{\skt{भवन्ति}}{\Afootnote{H; \emph{om.}\ A; \skt{न्भवन्ति ॥}\ N}}\skt{ ॥ }\edtext{\edtext{\skt{वृष\-ण\-यो\-र\-व\-द\-रणं}}{\lemma{\skt{वृषणयोर्}}\Afootnote{A N; \skt{वृषणयो}\ H}}}{\lemma{\skt{अवदरणं}}\Afootnote{A N; \skt{रवदरणं}\ H}}\skt{ }\edtext{\skt{बस्तिभेदो}}{\Afootnote{A H; \skt{वस्तिभेदो}\ N}}\skt{ }\edtext{\edtext{\skt{मेढ्रतोदोऽम्लीका}}{\lemma{\skt{मेढ्रतोदो}}\Afootnote{A H; \skt{मे}\ N}}\skt{ पिपासा ज्वरोऽतीसारो}}{\lemma{\skt{ऽम्लीका पिपासा ज्वरोऽतीसारो}}\Afootnote{H; \skt{हृदि शूल\-म\-म्ली\-का\-ज्व\-रा\-ती\-सा\-रा\-रो\-चका वमथुः परिधूपनं दाहो}\ A; \skt{तादा\-म्ली\-का\-पि\-पा\-सा\-ज्व\-रो\-ती\-सारो}\ N}}\skt{ }\edtext{\skt{मूर्च्छापाण्डुरोग}}{\Afootnote{N H; \skt{मूर्च्छा पिपासा निद्रानाशः पाण्डुरोगः}\ A}}\skt{ }\edtext{\skt{इति}}{\Afootnote{N H; \emph{om.}\ A}}\skt{ }\edtext{\edtext{\skt{पित्त\-जा\-ना\-मु\-प\-द्रवा}}{\lemma{\skt{पित्तजानाम्}}\Afootnote{H; \skt{पीतविण्मू}\ A; \skt{पित्तजा}\ N}}}{\lemma{\skt{उपद्रवा}}\Afootnote{H; \skt{त्रनेत्रत्वं}\ A; \emph{om.}\ N}}\skt{ }\edtext{\skt{भवन्ति}}{\Afootnote{H; \skt{चेति पैत्ति}\ A; \emph{om.}\ N}}\skt{ ॥ }\edtext{\skt{हृद्ग्रहौ}}{\Afootnote{H; \skt{कानां; हृद्ग्रहो}\ A; \skt{ना ॥ हृद्ग्रहो}\ N}}\skt{ }\edtext{\edtext{\skt{दौर्ब्ब\-ल्य\-म\-नि\-द्रा\-लम्भः}}{\lemma{\skt{दौर्ब्बल्यम्}}\Afootnote{H; \skt{लौल्य}\ A; \skt{दौर्वल्यम्}\ N}}}{\lemma{\skt{अनिद्रालम्भः}}\Afootnote{N H; \skt{मनिद्रा स्तम्भः}\ A}}\skt{ कम्पः }\edtext{\skt{शूलो}}{\Afootnote{\skt{शूलं}\ A; \skt{शुलो}\ N; \skt{शूलो}\ H}}\skt{ }\edtext{\skt{बध्दपुरीषत्वं}}{\Afootnote{H; \skt{बद्धपुरीषत्वं}\ A; \skt{वध्द\-पु\-री\-ष\-त्वम्}\ N}}\skt{ }\edtext{\skt{शोषः}}{\Afootnote{N H; \emph{om.}\ A}}\skt{ }\edtext{\skt{कासः}}{\Afootnote{N H; \emph{om.}\ A}}\skt{ }\edtext{\skt{श्वास}}{\Afootnote{H; \skt{चे}\ A; \skt{स्वास}\ N}}\skt{ }\edtext{\skt{इति}}{\Afootnote{N H; \skt{ति}\ A}}\skt{ }\edtext{\skt{वात\-जा\-ना\-मु\-प\-द्रवा}}{\lemma{\skt{उपद्रवा}}\Afootnote{N H; \emph{om.}\ A}}\skt{ }\edtext{\skt{भवन्ति}}{\Afootnote{N H; \emph{om.}\ A}}\skt{ ॥ }\pend
\pstart\versepar \skt{एवमेते विंशति प्रमेहाः सोपद्रवा व्याख्याताः ॥}\pend

\pstart \ledsidenote{\passref{2.6.14}}\skt{तत्र वसा\-मे\-दो\-भ्या\-म\-भि\-प\-न्न\-श\-री\-रस्य }\edtext{\edtext{\skt{दोषै\-र\-नु\-ग\-त\-धातोः}}{\lemma{\skt{दोषैर्}}\Afootnote{N H; \skt{त्रिभिर्दोषैश्}\ A}}}{\lemma{\skt{अनुगतधातोः}}\Afootnote{N H; \skt{चानुगतधातोः}\ A}}\skt{ प्रमेहिणो }\edtext{\skt{नवपिडकाः}}{\Afootnote{N H; \skt{दशपिडका}\ A}}\skt{ }\edtext{\skt{सञ्जायन्ते}}{\Afootnote{N H; \skt{जायन्ते ।}\ A}}\skt{ ॥ }\edtext{\skt{तद्यथा}}{\Afootnote{A H; \skt{तद्यत्वा}\ N}}\skt{ ॥ }\edtext{\skt{सराविका}}{\Afootnote{H; \skt{शराविका}\ A N}}\skt{ । }\edtext{\skt{सर्षषी}}{\Afootnote{H; \skt{सर्षपिका,}\ A; \skt{सर्षपी ॥}\ N}}\skt{ । कच्छपिका । जालिनी । }\edtext{\skt{पुत्त्रिणी}}{\Afootnote{N H; \skt{विनता, पुत्रिणी,}\ A}}\skt{ । मसूरिका । }\edtext{\skt{अलजी}}{\Afootnote{A H; \skt{अंजल जी}\ N}}\skt{ । विदारिका । विद्रधिका चेति ॥}\pend

\pstart\versepar \ledsidenote{\passref{2.6.15}}\edtext{\skt{अन्तोन्नता च}}{\Afootnote{H; \skt{शरावमात्रा}\ A; \skt{अन्तोन्नतान्च}\ N}}\skt{ }\edtext{\skt{तदृपा}}{\Afootnote{\skt{तद्रूपा}\ A N H}}\skt{ निम्नमध्या }\edtext{\skt{सराविका}}{\Afootnote{H; \skt{शराविका ।}\ A; \skt{शरीविकाः ।}\ N}}\skt{ । }\pend
\pstart\versepar \skt{}\edtext{\skt{गौर\-स\-र्ष\-प\-सं\-स्था\-ना\-त\-त्प्र\-माणा च}}{\Afootnote{A N; \skt{गौर\-स\-र्ष\-प\-सं\-स्थाना𑑎𑑎 तत्प्रमाणा च}\ H}}\skt{ }\edtext{\skt{सर्षषी}}{\Afootnote{N H; \skt{सार्षपी ॥}\ A}}\skt{ ॥}\pend

\pstart\versepar \ledsidenote{\passref{2.6.16}}\skt{सदाहा कूर्म्मसंस्थाना ज्ञेया कच्छपिका }\edtext{\skt{बुधैः}}{\Afootnote{A H; \skt{वुधैः ।}\ N}}\skt{ । }\pend
\pstart\versepar \skt{जालिनी तीव्रदाहा तु मांस\-जा\-ल\-स\-मा\-वृता ॥}\pend

\pstart \ledsidenote{\passref{2.6.17}}\edtext{\edtext{\skt{महत्यल्पचिता}}{\lemma{\skt{महत्य्}}\Afootnote{N H; \skt{महती पिडका नीला पिडका विनता स्मृता । मह}\ A}}}{\lemma{\skt{अल्पचिता}}\Afootnote{N H; \skt{त्यल्पाचिता}\ A}}\skt{ }\edtext{\skt{ज्ञेया}}{\Afootnote{A H; \skt{यो}\ N}}\skt{ पिडका }\edtext{\skt{चापि}}{\Afootnote{N H; \skt{सा तु}\ A}}\skt{ पुत्रिणी ।}\pend

\pstart\versepar \ledsidenote{\passref{2.6.18}}\edtext{\skt{मसू\-र\-सं\-स्था\-न\-समा}}{\Afootnote{H; \skt{मसू\-र\-स\-म\-सं\-स्थाना}\ A; \skt{मसुरसंस्थान समो}\ N}}\skt{ }\edtext{\skt{विज्ञेया तु}}{\Afootnote{H; \skt{ज्ञेया सा तु}\ A; \skt{विज्ञेयान्तु}\ N}}\skt{ मसूरिका ॥ }\pend
\pstart\versepar \skt{रक्ता }\edtext{\skt{शिता}}{\Afootnote{H; \skt{सिता}\ A N}}\skt{ }\edtext{\skt{स्फोटचिता}}{\Afootnote{H; \skt{स्फोटवती}\ A; \skt{स्फोटाचिता}\ N}}\skt{ दारुणा त्वलजी भवेत् ।}\pend

\pstart\versepar \ledsidenote{\passref{2.6.19}}\edtext{\skt{विदा\-री\-क\-न्द\-व\-द्वृत्ता}}{\Afootnote{A H; \skt{विदा\-री\-क\-न्द\-व\-द्धृत्ता}\ N}}\skt{ कठिना च विदारिका ॥ }\pend
\pstart\versepar \skt{}\edtext{\edtext{\skt{विद्र\-धे\-र्ल्ल\-क्ष\-णै\-र्युक्ता}}{\lemma{\skt{विद्रधेर्}}\Afootnote{N H; \skt{विद्रधे}\ A}}}{\lemma{\skt{ल्लक्षणैर्}}\Afootnote{N H; \skt{र्लक्षणैर्}\ A}}\skt{ ज्ञेया विद्रधिका }\edtext{\skt{बुधैः}}{\Afootnote{A H; \skt{वुधैः ।}\ N}}\skt{ ॥}\pend

\pstart \ledsidenote{\passref{2.6.19.1}}\edtext{\skt{ये यन्मया}}{\Afootnote{N; \skt{ये यन्मयाः}\ H}}\skt{ स्मृता }\edtext{\skt{मेहास्तेषां}}{\lemma{\skt{मेहास्}}\Afootnote{H; \skt{मेहा}\ N}}\skt{ }\edtext{\skt{एताम्तु}}{\Afootnote{H; \skt{मेहास्तु}\ N}}\skt{ }\edtext{\skt{तन्मया}}{\Afootnote{N; \skt{तन्मयाः ॥}\ H}}\skt{ ॥}\pend

\pstart\versepar \ledsidenote{\passref{2.6.20}}\skt{गुदे हृदि शिरस्यंसे पृष्ठे }\edtext{\skt{मर्म्मसु}}{\Afootnote{N H; \skt{मर्मणि}\ A}}\skt{ चोत्थिताः। }\pend
\pstart\versepar \skt{सोपद्रवा }\edtext{\skt{दुर्ब्बलाग्रैः}}{\Afootnote{\skt{दुर्बलाग्रेः}\ A; \skt{दुर्वलाग्नः}\ N; \skt{दुर्ब्बलाग्नैः}\ H}}\skt{ पिडकाः परिवर्ज्जयेत् ॥}\pend

\pstart\versepar \ledsidenote{\passref{2.6.21}}\edtext{\skt{कृत्स्नं}}{\Afootnote{A H; \skt{कृत्स्न}\ N}}\skt{ शरीरं निष्पीड्य }\edtext{\skt{मेदो\-म\-ज्जा\-व\-सा\-युतः।}}{\lemma{\skt{मेदो\-म\-ज्जा\-व\-सा\-युतः}}\Afootnote{H; \skt{मेदो\-म\-ज्ज\-व\-सा\-युतः ।}\ A; \skt{मेदो\-म\-ज्ज\-व\-सा\-युताः ॥}\ N}}\skt{ }\pend
\pstart\versepar \skt{}\edtext{\skt{अधः}}{\Afootnote{A H; \skt{अतः}\ N}}\skt{ प्रक्रमते }\edtext{\edtext{\skt{वायु\-स्ते\-ना\-सा\-ध्यास्तु}}{\lemma{\skt{वायुस्}}\Afootnote{N H; \skt{वयु}\ A}}}{\lemma{\skt{तेनासाध्यास्तु}}\Afootnote{N H; \skt{स्तेनासाध्यास्तु}\ A}}\skt{ वातजाः।}\pend

\pstart\versepar \ledsidenote{\passref{2.6.22}}\edtext{\skt{प्रमेहे}}{\Afootnote{H; \skt{प्रमेह}\ A N}}\skt{ }\edtext{\edtext{\edtext{\skt{पूर्व्व\-रू\-पा\-णा\-मा\-कृ\-ति\-र्यत्र}}{\lemma{\skt{पूर्व्वरूपाणाम्}}\Afootnote{N H; \skt{पूर्वरूपाणा}\ A}}}{\lemma{\skt{आकृतिर्}}\Afootnote{N H; \skt{माकृति}\ A}}}{\lemma{\skt{यत्र}}\Afootnote{N H; \skt{र्यत्र}\ A}}\skt{ दृश्यते । }\pend
\pstart\versepar \skt{}\edtext{\skt{किञ्चिदस्यधिकं}}{\Afootnote{\skt{किंचिच्चाप्यधिकं}\ A; \skt{किञ्चिदन्यधिकं}\ N; \skt{किञ्चिदत्यधिकं}\ H}}\skt{ मूत्रं तं }\edtext{\skt{प्रमेहीति}}{\Afootnote{N H; \skt{प्रमेहिणम्}\ A}}\skt{ }\edtext{\skt{निर्द्दिशेत्}}{\Afootnote{N H; \skt{आदिशेत् ॥}\ A}}\skt{ ।}\pend

\pstart \ledsidenote{\passref{2.6.23}}\edtext{\skt{कृच्छ्राण्यर्द्धानि}}{\lemma{\skt{कृच्छ्राण्य्}}\Afootnote{H; \skt{कृत्स्नान्य्}\ A; \skt{कृत्स्नाण्य्}\ N}}\skt{ वा यस्मि\-न्पू\-र्व्व\-रू\-पाणि मानवे । }\edtext{\skt{प्रवृत्तं}}{\Afootnote{H; \skt{प्रवृत्त}\ A; \skt{वृत्त}\ N}}\skt{ }\edtext{\skt{मूत्रमत्यर्थं तं}}{\lemma{\skt{अत्यर्थं तं}}\Afootnote{A H; \skt{अत्यर्थन्तम्}\ N}}\skt{ प्रमे\-हि\-ण\-मा\-दि\-शेत् ॥}\pend

\pstart\versepar \ledsidenote{\passref{2.6.24}}\edtext{\skt{पिडकापीडितं}}{\Afootnote{N H; \skt{पिडकापिडितं}\ A}}\skt{ }\edtext{\skt{गाढ\-मु\-प\-स्ट\-ष्ट\-मु\-प\-द्रवैः}}{\lemma{\skt{उपस्टष्टम्}}\Afootnote{\skt{उपसृष्टम्}\ A N H}}\skt{ । }\pend
\pstart\versepar \skt{}\edtext{\skt{मधु\-मे\-हि\-ण\-मा\-चष्टे}}{\lemma{\skt{मधुमेहिणम्}}\Afootnote{H; \skt{मधुमेहिनम्}\ A N}}\skt{ स }\edtext{\skt{चासाध्यतमः}}{\Afootnote{N H; \skt{चासाध्यः प्रकी}\ A}}\skt{ }\edtext{\skt{स्यातेः॥}}{\lemma{\skt{स्यातेः}}\Afootnote{\skt{र्तितः ॥}\ A; \skt{स्मृतः}\ N H}}\pend

\pstart\versepar \ledsidenote{\passref{2.6.25}}\skt{स चापि }\edtext{\skt{गमनात्स्थानं}}{\Afootnote{A H; \skt{गमनां स्थानं}\ N}}\skt{ }\edtext{\edtext{\skt{स्थाना\-दा\-स\-न\-मेव}}{\lemma{\skt{स्थानादासनम्}}\Afootnote{N H; \skt{स्थानदासनम्}\ A}}\skt{ च}}{\lemma{\skt{एव च}}\Afootnote{N H; \skt{इच्छति ।}\ A}}\skt{ । }\pend
\pstart\versepar \skt{आसनाद्वृणुते शय्यां }\edtext{\skt{शयने}}{\Afootnote{N H; \skt{शयना}\ A}}\skt{ }\edtext{\skt{स्वप्नमिच्छति}}{\lemma{\skt{स्वप्नम्}}\Afootnote{N; \skt{त्स्वप्नम्}\ A; \skt{स्वप्तुम्}\ H}}\skt{ ॥}\pend

\pstart \ledsidenote{\passref{2.6.26}}\skt{यथा }\edtext{\edtext{\edtext{\skt{शुक्ला\-दि\-व\-र्ण्णा\-ना\-म्प\-ञ्चा\-ना\-मु\-त्क\-र्षा\-प\-क\-र्ष\-कृ\-तेन}}{\lemma{\skt{शुक्ला\-दि\-व\-र्ण्णा\-नाम्}}\Afootnote{N H; \skt{हि वर्णानां}\ A}}}{\lemma{\skt{पञ्चानाम्}}\Afootnote{N H; \skt{पञ्चाना}\ A}}}{\lemma{\skt{उत्क\-र्षा\-प\-क\-र्ष\-कृ\-तेन}}\Afootnote{N H; \skt{मुत्क\-र्षा\-प\-क\-र्ष\-कृ\-तेन}\ A}}\skt{ संयो\-ग\-वि\-शे\-षेण }\edtext{\skt{स्वेत\-ब\-भृ\-क\-पि\-ल\-क\-पो\-त\-मे\-च\-का\-दीनां}}{\Afootnote{H; \skt{शब\-ल\-ब\-भ्रु\-क\-पि\-ल\-क\-पो\-त\-मे\-च\-का\-दीनां}\ A; \skt{स्वेत\-व\-अ\-भ्रू\-क\-पि\-ल\-क\-पो\-त\-मे\-च\-का\-दी\-नाम्}\ N}}\skt{ }\edtext{\edtext{\edtext{\skt{वर्ण्णा\-ना\-म\-ने\-के\-षा\-मु\-त्प\-त्ति\-र्भ\-वति}}{\lemma{\skt{वर्ण्णा\-ना\-म\-ने\-के\-षाम्}}\Afootnote{\skt{वर्णा\-ना\-म\-ने\-केषा}\ A; \skt{वर्ण्णानां अनेकषाम्}\ N; \skt{वर्ण्णा\-ना\-म\-ने\-के\-षाम्}\ H}}}{\lemma{\skt{उत्पत्तिर्}}\Afootnote{N H; \skt{मुत्पत्ति}\ A}}}{\lemma{\skt{भवति}}\Afootnote{N H; \skt{र्भवति,}\ A}}\skt{ । एवमेव }\edtext{\skt{दोष\-धा\-तु\-म\-ला\-हा\-र\-वि\-शे\-षे\-णो\-त्क\-र्षा\-प\-क\-र्ष\-कृ\-तेन}}{\Afootnote{A H; \skt{दोष\-धा\-तु\-म\-ला\-हा\-र\-वि\-शे\-षे\-णो\-त्क\-र्षा\-प\-क\-र्ष\-कृ\-तेम}\ N}}\skt{ संयो\-ग\-वि\-शे\-षेण }\edtext{\skt{प्रमे\-ह\-ना\-ना\-क\-र\-ण\-म्भ\-वति}}{\lemma{\skt{प्रमे\-ह\-ना\-ना\-क\-र\-णम्}}\Afootnote{N H; \skt{प्रमेहाणां नानाकरणं}\ A}}\skt{ ॥}\pend

\pstart\versepar \ledsidenote{\passref{2.6.27}}\edtext{\skt{भवति}}{\Afootnote{A H; \emph{om.}\ N}}\skt{ }\edtext{\skt{चात्र}}{\Afootnote{A H; \emph{om.}\ N}}\skt{ ॥ }\pend
\pstart\versepar \skt{सर्व्व एव }\edtext{\skt{प्रमेहास्तु}}{\Afootnote{N H; \skt{प्रमेहास्तु}\ A}}\skt{ }\edtext{\skt{काले\-ना\-प्र\-ति\-का\-रिणः}}{\Afootnote{H; \skt{काले\-ना\-प्र\-ति\-कु\-र्वतः ।}\ A; \skt{काले प्रतिकारिणो ।}\ N}}\skt{ }\pend
\pstart\versepar \skt{}\edtext{\skt{मधु\-मे\-ह\-त्व\-मा\-यान्ति}}{\lemma{\skt{मधुमेहत्वम्}}\Afootnote{A H; \skt{मधुमेहात्वम्}\ N}}\skt{ तदा }\edtext{\skt{साध्या}}{\Afootnote{N H; \skt{ऽसाध्या}\ A}}\skt{ }\edtext{\skt{भवन्ति च}}{\Afootnote{N; \skt{भवन्ति हि ॥}\ A; \skt{भवन्ति चेति ॥ ० ॥ ॥}\ H}}\skt{ ॥}\pend

\pstart \ledsidenote{\passref{2.6.trailer}}\edtext{\skt{इति}}{\Afootnote{\skt{इति सुश्रु\-त\-सं\-हि\-तायां निदानस्थाने}\ A; \skt{त्र ६ ॥ ० ॥}\ N; \emph{om.}\ H}}\skt{ }\edtext{\skt{प्रमेहनिदाने}}{\Afootnote{H; \skt{प्रमेहनिदानं नाम}\ A; \emph{om.}\ N}}\skt{ }\edtext{\skt{षष्ठः}}{\Afootnote{H; \skt{षष्ठोऽध्यायः ॥ ६ ॥}\ A; \emph{om.}\ N}}\skt{ ॥}\pend

\endnumbering
\beginnumbering
\eledsection*{Adhyāya 7}
\pstart \ledsidenote{\passref{2.7.1}}\skt{अथात }\edtext{\edtext{\skt{उद\-रा\-णा\-न्नि\-दा\-न\-म्व्या\-ख्या\-स्यामः}}{\lemma{\skt{निदानम्}}\Afootnote{A H; \skt{निदान व्याख्यास्यामः ॥}\ N}}}{\lemma{\skt{व्याख्यास्यामः}}\Afootnote{A H; \emph{om.}\ N}}\skt{ ॥}\pend

\pstart\versepar \ledsidenote{\passref{2.7.3}}\edtext{\edtext{\skt{धन्व\-न्त\-रि\-र्ध\-र्म्म\-भृ\-ता\-म्व\-रिष्ठो}}{\lemma{\skt{धन्वन्तरिर्}}\Afootnote{A H; \skt{धन्वन्तिरिर्}\ N}}}{\lemma{\skt{वरिष्ठो}}\Afootnote{A H; \skt{परिष्ठो}\ N}}\skt{ राज\-र्षि\-रि\-न्द्र\-प्र\-तिमो }\edtext{\skt{महात्मा}}{\Afootnote{N H; \skt{ऽभवद्यः ।}\ A}}\skt{ । }\pend
\pstart\versepar \skt{}\edtext{\edtext{\skt{ब्रह्म\-र्षि\-पु\-त्र\-म्वि\-न\-यो\-प\-पन्नं}}{\lemma{\skt{ब्रह्मर्षिपुत्रम्}}\Afootnote{A H; \skt{व्रह्मर्षिपुत्रम्}\ N}}}{\lemma{\skt{विनयोपपन्नं}}\Afootnote{A H; \skt{विनायोपपन्नं}\ N}}\skt{ }\edtext{\skt{शिष्यं}}{\Afootnote{A H; \skt{शिष्य}\ N}}\skt{ शुभं }\edtext{\skt{सुश्रु\-त\-म\-न्व\-शा\-सत्}}{\lemma{\skt{अन्वशासत्}}\Afootnote{N H; \skt{अन्वशात्सः ॥}\ A}}\skt{ ॥}\pend

\pstart\versepar \ledsidenote{\passref{2.7.4}}\skt{पृथ\-क्स\-म\-स्तै\-रपि }\edtext{\skt{दोषैः}}{\Afootnote{H\textsuperscript{ac}; \skt{चेह दोषैः}\ A; \skt{चापि दोषैः}\ N H\textsuperscript{pc}}}\skt{ }\edtext{\edtext{\skt{प्लीहो\-द\-र\-म्ब\-द्ध\-गु\-द\-व\-म्व\-दन्ति}}{\lemma{\skt{बद्धगुदवम्}}\Afootnote{\skt{बद्धगुदं}\ A H; \skt{वद्धगुदम्}\ N}}}{\lemma{\skt{वदन्ति}}\Afootnote{N H; \skt{तथैव ।}\ A}}\skt{ । }\pend
\pstart\versepar \skt{}\edtext{\skt{आगन्तुकं}}{\Afootnote{A H; \skt{आगन्तुक}\ N}}\skt{ }\edtext{\skt{सप्त\-म\-म\-ष्ट\-मन्तु}}{\lemma{\skt{अष्टमन्तु}}\Afootnote{N H; \skt{अष्टमं च}\ A}}\skt{ दकोदरं }\edtext{\skt{चैव}}{\Afootnote{N H; \skt{चेति}\ A}}\skt{ }\edtext{\skt{भवन्ति}}{\Afootnote{N H; \skt{वदन्ति}\ A}}\skt{ तानि ॥}\pend

\pstart\versepar \ledsidenote{\passref{2.7.5}}\edtext{\skt{सुदु\-र्ब्ब\-ला\-ग्ने\-र\-हि\-ता\-श\-नस्य}}{\lemma{\skt{सुदुर्ब्बलाग्नेर्}}\Afootnote{A H; \skt{सुदुर्वलाग्नेर्}\ N}}\skt{ }\pend
\pstart\versepar \skt{}\edtext{\skt{वृद्धिं}}{\Afootnote{N H; \skt{संशु\-ष्क\-पू\-त्य\-न्न\-नि\-षे\-व\-णाद्वा स्नेहा\-दि\-मि\-थ्या\-च\-र\-णाच्च जन्तोर्वृद्धिं}\ A}}\skt{ गताः कोष्ठ\-म\-भि\-प्र\-पन्नाः ॥}\pend

\pstart\versepar \ledsidenote{\passref{2.7.6}}\edtext{\skt{गुल्मा\-कृ\-ति\-व्य\-ञ्ज\-न\-ल\-क्ष\-णानि}}{\Afootnote{H; \skt{गुल्मा\-कृ\-ति\-व्य\-ञ्जि\-त\-ल\-क्ष\-णानि}\ A; \skt{गुल्माकृतिं व्यंज\-न\-ल\-क्ष\-णानि}\ N}}\skt{ कुर्व्वन्ति }\edtext{\skt{घोराण्युदराणि}}{\Afootnote{N H; \skt{घोराण्युदराणि}\ A}}\skt{ दोषाः ॥ }\pend
\pstart\versepar \skt{}\edtext{\skt{कोष्ठा\-दु\-प\-स्ने\-ह\-व\-द\-न्न\-सारो}}{\Afootnote{A H; \skt{कोष्ठा\-दु\-प\-स्ने\-ह\-व\-द\-न्त\-सारो}\ N}}\skt{ निःसृत्य }\edtext{\skt{दुष्टोऽ\-नि\-ल\-सं\-प्र\-युक्तः।}}{\lemma{\skt{ऽनि\-ल\-सं\-प्र\-युक्तः}}\Afootnote{H; \skt{ऽनि\-ल\-वे\-ग\-नुन्नः ॥}\ A; \skt{निलसंप्रयुक्ताः।}\ N}}\pend

\pstart \ledsidenote{\passref{2.7.7ab}}\skt{त्वचः }\edtext{\skt{समुन्नाम्य}}{\Afootnote{N H; \skt{समुन्नम्य}\ A}}\skt{ शनैः }\edtext{\skt{समस्ताद्धि}}{\Afootnote{N H; \skt{समन्ताद्वि}\ A}}\skt{ }\edtext{\skt{सर्प्यमाणो}}{\Afootnote{N H; \skt{वर्धमानो}\ A}}\skt{ जठरं करोति ॥}\pend

\pstart \ledsidenote{\passref{2.7.8cd}}\edtext{\skt{यद्गृ\-ह्य\-पृ\-ष्ठो\-द\-र\-पा\-र्श्व\-वस्ती}}{\Afootnote{H; \skt{संगृह्य पार्श्वो\-द\-र\-पृ\-ष्ठ\-ना\-भीर्}\ A; \skt{प्रगृ\-ह्य\-पृ\-ष्ठौ\-द\-र\-पा\-र्श्व\-वस्ती}\ N}}\skt{ }\edtext{\skt{न विवर्द्धते}}{\Afootnote{\skt{यद्वर्धते}\ A; \skt{न्विवर्धते}\ N H}}\skt{ कृष्ण\-सि\-रा\-व\-नद्धं ।}\pend

\pstart\versepar \ledsidenote{\passref{2.7.9}}\edtext{\skt{समूढवातं}}{\Afootnote{N H; \emph{om.}\ A}}\skt{ }\edtext{\skt{सरुजं}}{\Afootnote{N H; \skt{सशूलम्}\ A}}\skt{ }\edtext{\skt{सशब्दं}}{\Afootnote{H; \skt{आना\-ह\-व\-दु\-ग्र\-शब्दं}\ A; \skt{सशव्दं}\ N}}\skt{ सतो\-द\-भे\-द\-म्प\-व\-ना\-त्म\-क\-न्तत् ॥ }\pend
\pstart\versepar \skt{सचो\-ष\-कृ\-ष्णा\-ज्व\-र\-दा\-ह\-युक्तं }\edtext{\skt{पीताः}}{\Afootnote{N H; \skt{पीतं}\ A}}\skt{ सिरा }\edtext{\skt{यत्र च}}{\Afootnote{N H; \skt{भान्ति च}\ A}}\skt{ }\edtext{\skt{भान्ति}}{\Afootnote{N H; \skt{यत्र}\ A}}\skt{ }\edtext{\skt{पीताः}}{\Afootnote{A H; \skt{पीता ।}\ N}}\edtext{}{\lemma{\skt{तत् ॥ सचो\-ष\-कृ\-ष्णा\-ज्व\-र\-दा\-ह\-युक्तं}}\Afootnote{H; \skt{तत् । यच्चो\-ष\-तृ\-ष्णा\-ज्व\-र\-दा\-ह\-युक्तं}\ A; \skt{तत् ॥ सचो\-ष\-कृ\-ष्णा\-ज्व\-र\-दा\-ह\-युक्तां}\ N}}\pend

\pstart\versepar \ledsidenote{\passref{2.7.10}}\skt{पीता\-क्षि\-वि\-ण्मू\-त्र\-न\-खा\-न\-नस्य }\edtext{\skt{पित्तो\-द\-र\-न्तं\-त्व\-चि\-रा\-भि\-वृद्धिः}}{\Afootnote{H; \skt{पित्तोदरं तत्त्व\-चि\-रा\-भि\-वृद्धि ।}\ A; \skt{विदो\-त्तो\-द\-र\-न्त\-त्व\-चि\-रा\-भि\-वृद्धिः ॥}\ N}}\skt{ ॥ }\pend
\pstart\versepar \skt{यच्छीतलं शुक्ल\-सि\-रा\-व\-नद्धं गुरुस्थिरं शुक्ल\-न\-खा\-न\-नस्य ।}\pend

\pstart\versepar \ledsidenote{\passref{2.7.11}}\edtext{\skt{स्निग्ध\-म्म\-ह\-त्सा\-द\-न\-शो\-फ\-युक्तं}}{\lemma{\skt{मह\-त्सा\-द\-न\-शो\-फ\-युक्तं}}\Afootnote{H; \skt{महच्छोफयुतं}\ A; \skt{मह\-त्सी\-द\-न\-शो\-फ\-युक्तं}\ N}}\skt{ }\edtext{\edtext{\skt{चिरा\-भि\-वृ\-द्धि\-म्प्र\-थितं}}{\lemma{\skt{चिराभिवृद्धिम्}}\Afootnote{N H; \skt{ससादं कफोदरं तत्तु चिराभिवृ}\ A}}}{\lemma{\skt{प्रथितं}}\Afootnote{N H; \skt{द्धि ।}\ A}}\skt{ कफात्तत् ॥ }\pend
\pstart\versepar \skt{}\edtext{\edtext{\skt{स्त्रियोऽ\-न्न\-पा\-न\-न्न\-ख\-रो\-म\-मू\-त्रै\-र्व्वि\-डा\-र्त्त\-वै\-र्यु\-क्त\-म\-सा\-धु\-वृत्ताः।}}{\lemma{\skt{ऽन्न\-पा\-न\-न्न\-ख\-रो\-म\-मू\-त्रैर्}}\Afootnote{H; \skt{ऽन्नपानं नखरोममूत्र}\ A; \skt{न्नधा\-न\-न्न\-ख\-रो\-म\-मू\-त्रैर्}\ N}}}{\lemma{\skt{व्विडार्त्तवैर्}}\Afootnote{A H; \skt{व्विडार्त्तवै}\ N}}\edtext{}{\lemma{\skt{कफात्तत् ॥ स्त्रियो}}\Afootnote{N H; \skt{स्त्रियो}\ A}}\pend

\pstart\versepar \ledsidenote{\passref{2.7.12}}\skt{यस्मै }\edtext{\skt{प्रय\-च्छ\-न्त्य\-र\-यो\-ग\-रांश्च}}{\lemma{\skt{अरयोगरांश्च}}\Afootnote{\skt{अरयो गरांश्च}\ A H; \skt{अरसोगरांश्च}\ N}}\skt{ }\edtext{\skt{दुष्टा\-म्बु\-दू\-षी\-वि\-ष\-से\-व\-नाद्वा}}{\Afootnote{A H; \skt{दुष्टां वुदू\-षी\-वि\-ष\-से\-व\-नाद्वा ॥}\ N}}\skt{ ॥ }\pend
\pstart\versepar \skt{तेनाशु रक्तं कुपिताश्च दोषाः }\edtext{\skt{कुर्युः}}{\Afootnote{N H; \skt{कुर्वन्}\ A}}\skt{ }\edtext{\edtext{\skt{सुघो\-र\-ञ्ज\-ठ\-र\-न्त्रि\-लिङ्गं}}{\lemma{\skt{सुघोरञ्}}\Afootnote{H; \skt{ति घोरं}\ A; \skt{सुघौरञ्}\ N}}}{\lemma{\skt{त्रिलिङ्गं}}\Afootnote{A H; \skt{त्रिलङ्गं}\ N}}\skt{ ।}\pend

\pstart\versepar \ledsidenote{\passref{2.7.13}}\edtext{\skt{तच्छी\-त\-वा\-ता\-त\-प\-दु\-र्द्दि\-नेषु}}{\Afootnote{N H; \skt{तच्छी\-त\-वा\-ता\-भ्र\-स\-मु\-द्भ\-वेषु}\ A}}\skt{ विशेषतः कुप्यति दह्यते च ॥ }\pend
\pstart\versepar \skt{स चातुरो मूर्च्छति }\edtext{\skt{सम्प्रयुक्तः}}{\Afootnote{H; \skt{संप्रसक्तं}\ A; \skt{संप्रयुक्तं}\ N}}\skt{ पाण्डुः }\edtext{\skt{कृषः}}{\Afootnote{N H; \skt{कृशः}\ A}}\skt{ शुष्यति तृष्णया च ॥}\pend

\pstart\versepar \ledsidenote{\passref{2.7.14}}\edtext{\skt{दूष्योदरं}}{\Afootnote{N H; \skt{प्रकीर्तितं दूष्युदरं}\ A}}\skt{ }\edtext{\edtext{\skt{कीर्त्ति\-त\-मे\-त\-देवं}}{\lemma{\skt{कीर्त्तितम्}}\Afootnote{N H; \emph{om.}\ A}}}{\lemma{\skt{एतदेवं}}\Afootnote{H; \skt{तु घोरं}\ A; \skt{एव एवं}\ N}}\skt{ प्लीहोदरं कीर्तयतो }\edtext{\skt{निबोधः}}{\Afootnote{N H; \skt{निबोध ।}\ A}}\skt{ ॥ }\pend
\pstart\versepar \skt{विदा\-ह्य\-भि\-ष्य\-न्दि\-र\-तस्य जन्तोः }\edtext{\edtext{\skt{प्रदू\-ष्ट\-म\-त्य\-र्थ\-म\-सृ\-क्क\-फश्च}}{\lemma{\skt{प्रदूष्टम्}}\Afootnote{N H; \skt{प्रदुष्टम्}\ A}}}{\lemma{\skt{अत्य\-र्थ\-म\-सृ\-क्क\-फश्}}\Afootnote{\skt{अत्य\-र्थ\-म\-सृ\-क्क\-फश्}\ A; \skt{अत्य\-र्त्थ\-म\-सृ\-क्क\-फश्}\ N H}}\skt{ ।}\pend

\pstart\versepar \ledsidenote{\passref{2.7.15}}\edtext{\skt{प्लीहाभिवृद्धिं}}{\Afootnote{A H; \skt{प्लीहाभिवृ}\ N}}\skt{ }\edtext{\skt{कुरुते}}{\Afootnote{H; \skt{सततं करोति}\ A; \skt{द्धि कुरुते}\ N}}\skt{ }\edtext{\skt{प्लिहोत्थं}}{\Afootnote{N H; \skt{प्लीहो}\ A}}\skt{ }\edtext{\skt{प्रचक्षत}}{\Afootnote{H; \emph{om.}\ A; \skt{प्रचक्षते त}\ N}}\skt{ }\edtext{\edtext{\skt{तज्ज\-ठ\-र\-म्प्र\-वृद्धं}}{\lemma{\skt{तज्जठरम्}}\Afootnote{H; \skt{दरं त}\ A; \skt{ठज्जठरं}\ N}}}{\lemma{\skt{प्रवृद्धं}}\Afootnote{N H; \skt{त्प्रवदन्ति}\ A}}\skt{ ॥ }\pend
\pstart\versepar \skt{}\edtext{\skt{तद्वामपार्श्वे}}{\Afootnote{H; \skt{तज्ज्ञाः । वामे च पार्श्वे}\ A; \skt{तद्वामपार्श्व}\ N}}\skt{ परिवृद्धिमेति विशेषतः सीदति चातुरोऽत्र ।}\pend

\pstart\versepar \ledsidenote{\passref{2.7.16}}\edtext{\skt{मन्द\-ज्व\-रा\-ग्निः\-क\-फ\-पि\-त्त\-लिं\-गै\-रू\-प\-दृतः}}{\lemma{\skt{ऊपदृतः}}\Afootnote{N H; \skt{उपद्रुतः}\ A}}\skt{ }\edtext{\skt{क्षीण\-ब\-लोऽ\-ति\-पाण्डुः।}}{\lemma{\skt{क्षीणबलो}}\Afootnote{A H; \skt{क्षीणो वलो}\ N}}\skt{ }\pend
\pstart\versepar \skt{}\edtext{\skt{सव्ये तु}}{\Afootnote{N H; \skt{सव्येत}\ A}}\skt{ }\edtext{\skt{पार्श्वे}}{\Afootnote{N H; \skt{रस्मि}\ A}}\skt{ }\edtext{\skt{यकृतिप्रदुष्टे}}{\Afootnote{N H; \skt{न्यकृति प्रदुष्टे}\ A}}\skt{ ज्ञेयं }\edtext{\skt{यकृद्दाल्यदरं}}{\Afootnote{\skt{यकृद्दाल्युदरं}\ A H; \skt{वकृ\-द्दा\-ल्यु\-द\-रन्}\ N}}\skt{ }\edtext{\skt{तदेव}}{\Afootnote{A N; \skt{तदेव ।}\ H}}\skt{ ।}\pend

\pstart\versepar \ledsidenote{\passref{2.7.17}}\edtext{\edtext{\edtext{\edtext{\edtext{\skt{यस्या\-न्त\-म\-न्नै\-रू\-प\-ल\-पि\-भि\-र्व्वा\-बा\-ला\-श्म\-भि\-र्व्वा\-पि\-हितं}}{\lemma{\skt{यस्यान्तम्}}\Afootnote{H; \skt{यस्यान्त्र}\ A; \skt{यस्यान्त}\ N}}}{\lemma{\skt{अन्नैर्}}\Afootnote{H; \skt{मन्नैर्}\ A N}}}{\lemma{\skt{ऊपलपिभिर्}}\Afootnote{\skt{उपलेपिभिर्}\ A H; \skt{अपलेपिभिर्}\ N}}}{\lemma{\skt{व्वाबालाश्मभिर्}}\Afootnote{A H; \skt{व्वालाश्मभिर्}\ N}}}{\lemma{\skt{व्वापिहितं}}\Afootnote{N H; \skt{वा सहितैः}\ A}}\skt{ यथावत् ॥ }\pend
\pstart\versepar \skt{सञ्चीयते तत्र मलः }\edtext{\skt{सदोषः}}{\Afootnote{A H; \skt{प्रदोषः}\ N}}\skt{ }\edtext{\skt{शनैः}}{\Afootnote{N H; \emph{om.}\ A}}\skt{ }\edtext{\skt{शनैः}}{\Afootnote{N H; \emph{om.}\ A}}\skt{ }\edtext{\skt{शज्व\-र\-व\-क्तु\-नाड्यां}}{\Afootnote{\skt{क्रमेण नाड्यामिव संकरो हि ॥}\ A; \skt{शकेरवन्तु नान्याः ।}\ N; \skt{शङ्करवत्तु नाड्यां ।}\ H}}\skt{ ।}\edtext{}{\lemma{\skt{यथावत् ॥ सञ्चीयते}}\Afootnote{N H; \skt{पृथग्वा । संचीयते}\ A}}\pend

\pstart\versepar \ledsidenote{\passref{2.7.18}}\skt{निरुध्यते चास्य }\edtext{\skt{गुदं}}{\Afootnote{N H; \skt{गुदे}\ A}}\skt{ पुरीषं निरेति कृच्छ्रादपि }\edtext{\skt{चाल्पमल्पं}}{\Afootnote{A N; \skt{चाल्पमल्पं ॥}\ H}}\skt{ ॥ }\pend
\pstart\versepar \skt{}\edtext{\skt{हृन्नाभिमध्ये}}{\Afootnote{A H; \skt{हृन्नौभिमध्ये}\ N}}\skt{ परिवृद्धिमेति }\edtext{\edtext{\edtext{\skt{तस्यो\-द\-र\-म्व\-द्ध\-गु\-द\-म्व\-दन्ति}}{\lemma{\skt{तस्योदरम्}}\Afootnote{N H; \skt{त(य)च्चोदरं}\ A}}}{\lemma{\skt{वद्धगुदम्}}\Afootnote{N H; \skt{विट्समगन्}\ A}}}{\lemma{\skt{वदन्ति}}\Afootnote{N H; \skt{धिकं च ॥}\ A}}\skt{ ॥}\pend

\pstart \ledsidenote{\passref{2.7.19cd}}\edtext{\skt{शल्य\-न्त\-था\-नो\-प\-हितं}}{\lemma{\skt{तथानोपहितं}}\Afootnote{H; \skt{यदन्नोपहितं}\ A; \skt{तथान्नोपहितं}\ N}}\skt{ }\edtext{\edtext{\skt{गुदान्त्रम्भिनत्ति}}{\lemma{\skt{गुदान्त्रम्}}\Afootnote{N H; \skt{तदन्त्रं}\ A}}}{\lemma{\skt{भिनत्ति}}\Afootnote{A H; \skt{भिन्नन्ति}\ N}}\skt{ }\edtext{\edtext{\skt{वर्च्चा\-ग\-त\-म\-न्यतो}}{\lemma{\skt{वर्च्चागतम्}}\Afootnote{H; \skt{यस्यागतम्}\ A; \skt{वर्चो गतम्}\ N}}}{\lemma{\skt{अन्यतो}}\Afootnote{N H; \skt{अन्यथा}\ A}}\skt{ वा ॥}\pend

\pstart\versepar \ledsidenote{\passref{2.7.20}}\edtext{\skt{तस्मा\-त्सृ\-ता\-न्त्रा\-त्स\-लि\-ल\-प्र\-काशः}}{\Afootnote{H; \skt{तस्मा\-त्स्रु\-तोऽ\-न्त्रा\-त्स\-लि\-ल\-प्र\-काशः}\ A; \skt{तस्मा\-त्स्रु\-ता\-न्त्रा\-त्स\-लि\-ल\-प्र\-काशः}\ N}}\skt{ स्रावः }\edtext{\skt{स्रवेद्वै}}{\Afootnote{A; \skt{स्रवेद्वै}\ N H}}\skt{ गुदतस्तु भूयः। }\pend
\pstart\versepar \skt{}\edtext{\edtext{\skt{नाभे\-र\-ध\-श्चो\-द\-र\-मेति}}{\lemma{\skt{अधश्}}\Afootnote{A N; \skt{अध}\ H}}}{\lemma{\skt{चोदरम्}}\Afootnote{A N; \skt{श्चोदरम्}\ H}}\skt{ }\edtext{\skt{वृद्धिन्निस्तुद्यते}}{\lemma{\skt{वृद्धिन्}}\Afootnote{A H; \skt{वृद्धि}\ N}}\skt{ }\edtext{\skt{दाल्यति}}{\Afootnote{N H; \skt{ऽतीव विदह्यते}\ A}}\skt{ }\edtext{\skt{चातिस्रावं}}{\Afootnote{\skt{च ॥}\ A; \skt{च त्रिमात्रं ।}\ N; \skt{चाति मात्रं ।}\ H}}\skt{ ।}\pend

\pstart\versepar \ledsidenote{\passref{2.7.21}}\edtext{\skt{एतत्परिस्राव्यु }\edtext{\skt{रम्प्रविष्टं}}{\lemma{\skt{प्रविष्टं}}\Afootnote{N H; \skt{प्रदिष्टं}\ A}}}{\lemma{\skt{एतत्परिस्राव्यु रम्}}\Afootnote{H; \skt{एत\-त्प\-रि\-स्रा\-व्यु\-दरं}\ A; \skt{एत\-त्प\-रि\-स्रा\-व्यु\-दरं}\ N}}\skt{ दकोदरं कीर्त्तयतो निबोधः। }\pend
\pstart\versepar \skt{यः }\edtext{\edtext{\skt{स्नेह\-पी\-तोऽ\-प्य\-नु\-वा\-सितो}}{\lemma{\skt{ऽप्य्}}\Afootnote{H; \skt{ऽ}\ A N}}\skt{ वा}}{\lemma{\skt{अनुवासितो वा}}\Afootnote{H; \skt{प्यनुवासितो वा}\ A N}}\skt{ वान्तो }\edtext{\edtext{\skt{विरिक्तोऽप्यथवा}}{\lemma{\skt{ऽप्य्}}\Afootnote{N H; \skt{ऽ}\ A}}}{\lemma{\skt{अथवा}}\Afootnote{N H; \skt{प्यथवा}\ A}}\skt{ }\edtext{\skt{निचूढः॥}}{\lemma{\skt{निचूढः}}\Afootnote{N; \skt{निरूढः}\ A H}}\edtext{}{\lemma{\skt{निबोधः। यः}}\Afootnote{H; \skt{निबोध । यः}\ A; \skt{निवोधः । यः}\ N}}\pend

\pstart\versepar \ledsidenote{\passref{2.7.22}}\edtext{\skt{पिबेज्जलं}}{\Afootnote{A H; \skt{पिवेज्जलं}\ N}}\skt{ शीतलमाशु तस्य स्रोतांसि }\edtext{\skt{दूष्यन्ति हि}}{\Afootnote{N H; \skt{दुष्यन्ति हि}\ A}}\skt{ तद्वहानि । }\pend
\pstart\versepar \skt{}\edtext{\skt{स्नेहो\-प\-लि\-प्ते\-ष्व\-थ\-वापि}}{\lemma{\skt{अथवापि}}\Afootnote{N H; \skt{अथवाऽपि}\ A}}\skt{ तेषु दकोदरं }\edtext{\skt{पूर्व्व\-व\-द\-भ्यु\-पैति}}{\Afootnote{A; \skt{पूर्व्व\-व\-द\-भ्यु\-प्यैति ॥}\ N H}}\skt{ ॥}\pend

\pstart\versepar \ledsidenote{\passref{2.7.23}}\edtext{\skt{स्निग्ध\-म्म\-ह\-त्त\-त्प\-रि\-वृ\-त्त\-नाभिः}}{\lemma{\skt{मह\-त्त\-त्प\-रि\-वृ\-त्त\-नाभिः}}\Afootnote{N H; \skt{मह\-त्सं\-प\-रि\-वृ\-त्त\-नाभि}\ A}}\skt{ }\edtext{\edtext{\edtext{\skt{समा\-त\-त\-म्पू\-र्व्व\-मि\-वा\-म्बुना}}{\lemma{\skt{समाततम्}}\Afootnote{N H; \skt{भृशोन्नतं}\ A}}}{\lemma{\skt{पूर्व्वम्}}\Afootnote{H; \skt{पूर्णम्}\ A N}}\skt{ च}}{\lemma{\skt{इवाम्बुना च}}\Afootnote{A H; \skt{इवाम्वुना च ।}\ N}}\skt{ । }\pend
\pstart\versepar \skt{यथा }\edtext{\skt{दृश्यति}}{\Afootnote{N; \skt{दृतिः}\ A; \skt{दृति}\ H}}\skt{ क्षुभ्यति कम्पते च }\edtext{\skt{शब्दायते}}{\Afootnote{A H; \skt{शव्दायते}\ N}}\skt{ चापि दकोदरन्तत् ॥}\pend

\pstart\versepar \ledsidenote{\passref{2.7.24}}\edtext{\skt{आध्मनो}}{\Afootnote{N; \skt{आध्मानं}\ A; \skt{अध्वनो}\ H}}\skt{ गमने }\edtext{\edtext{\edtext{\skt{शक्ति\-र्द्दौ\-ब्ब\-ल्य\-न्दु\-र्ब्ब\-ला\-ग्निता}}{\lemma{\skt{शक्तिर्}}\Afootnote{H; \skt{ऽशक्तिर्}\ A; \skt{शक्ति}\ N}}}{\lemma{\skt{द्दौब्बल्यन्}}\Afootnote{\skt{दौर्बल्यं}\ A H; \skt{दौर्व्वल्य्दु}\ N}}}{\lemma{\skt{दुर्ब्बलाग्निता}}\Afootnote{A H; \skt{र्व्वलाग्निता}\ N}}\skt{ । }\pend
\pstart\versepar \skt{शोफः }\edtext{\skt{सदनमङ्गनां}}{\lemma{\skt{अङ्गनां}}\Afootnote{N; \skt{अङ्गानां}\ A H}}\skt{ }\edtext{\skt{संगो\-वा\-त\-पु\-री\-षयोः}}{\Afootnote{\skt{सङ्गो वातपुरीषयोः}\ A H; \skt{संगो\-वा\-त\-पु\-री\-षो\-शयोः ॥}\ N}}\skt{ ॥ }\pend
\pstart\versepar \skt{}\edtext{\skt{दाहतन्द्री च}}{\Afootnote{H; \skt{दाहस्तृष्णा च}\ A; \skt{दाहस्तन्द्री च}\ N}}\skt{ }\edtext{\skt{मूर्च्छा च}}{\Afootnote{N H; \skt{सर्वेषु}\ A}}\skt{ जठरेषु भवन्ति हि ।}\pend

\pstart\versepar \ledsidenote{\passref{2.7.24.1}}\edtext{\skt{छर्दि\-श्चै\-वा\-ति\-सा\-रश्च त्र्ष्णा}}{\lemma{\skt{चैवातिसारश्च त्र्ष्णा}}\Afootnote{N; \skt{चैवातिसारश्च तृष्णो}\ H}}\skt{ }\edtext{\skt{दौर्व्वल्यमेव}}{\lemma{\skt{दौर्व्वल्यम्}}\Afootnote{N; \skt{दौर्ब्बल्यम्}\ H}}\skt{ च । }\pend
\pstart\versepar \skt{}\edtext{\skt{मूत्रग्रहः}}{\Afootnote{N; \skt{मूद्रग्रहः}\ H}}\skt{ पार्श्वशूलं स्वर\-भे\-द\-प्र\-मू\-ढता । }\pend
\pstart\versepar \skt{हिक्काश्वासश्च }\edtext{\skt{कासश्च}}{\Afootnote{H; \skt{काशश्च}\ N}}\skt{ अरु\-चि\-श्चा\-प्यु\-प\-द्रवाः ॥}\pend

\pstart\versepar \ledsidenote{\passref{2.7.25}}\skt{अन्ते }\edtext{\skt{सलिलभावञ्च}}{\Afootnote{N H; \skt{सलिलभावं हि}\ A}}\skt{ }\edtext{\skt{भवन्ति}}{\Afootnote{N H; \skt{भजन्ते}\ A}}\skt{ जठराणि तु । }\pend
\pstart\versepar \skt{सर्वाणि }\edtext{\skt{परिपक्वानि}}{\Afootnote{N H; \skt{एव परीपाका}\ A}}\skt{ }\edtext{\edtext{\skt{तस्मा\-त्त\-म्प\-रि\-व\-र्ज्ज\-येत्}}{\lemma{\skt{तस्मात्तम्}}\Afootnote{N H; \skt{त्तदा ता}\ A}}}{\lemma{\skt{परिवर्ज्जयेत्}}\Afootnote{N H; \skt{नि विवर्जयेत् ॥}\ A}}\skt{ ॥}\pend

\pstart \ledsidenote{\passref{2.7.trailer}}\edtext{\edtext{\skt{इत्यु\-द\-र\-नि\-दाने}}{\lemma{\skt{इत्य्}}\Afootnote{H; \skt{इति सुश्रु\-त\-सं\-हि\-तायां निदानस्थाने}\ A; \skt{७ ॥}\ N}}}{\lemma{\skt{उदरनिदाने}}\Afootnote{H; \skt{उदरनिदानं नाम}\ A; \emph{om.}\ N}}\skt{ }\edtext{\skt{सप्तमः}}{\Afootnote{H; \skt{सप्तमोऽध्यायः ॥७ ॥}\ A; \emph{om.}\ N}}\skt{ ॥}\pend

\endnumbering
\beginnumbering
\eledsection*{Adhyāya 8}
\pstart \ledsidenote{\passref{2.8.1}}\skt{अथातो }\edtext{\edtext{\skt{मूढ\-ग\-र्भा\-णा\-न्नि\-दा\-न\-म्व्या\-ख्या\-स्यामः}}{\lemma{\skt{मूढगर्भाणान्}}\Afootnote{N H; \skt{मूढगर्भ}\ A}}}{\lemma{\skt{व्याख्यास्यामः}}\Afootnote{A H; \emph{om.}\ N}}\skt{ ॥}\pend

\pstart \ledsidenote{\passref{2.8.3}}\edtext{\edtext{\skt{ग्राम्य\-ध\-र्म्म\-या\-न\-वा\-ह\-ना\-ध्व\-ग\-म\-न\-प्र\-स्ख\-ल\-न\-प्र\-प\-त\-न\-पी\-ड\-ना\-भि\-घा\-त\-वि\-ष\-म\-श\-य\-नो\-प\-वा\-स\-वे\-गा\-भि\-घा\-ता\-ति\-रू\-क्ष\-क\-टु\-क\-ति\-क्त\-भो\-ज\-ना\-जी\-र्ण्ण\-ग\-र्भ\-शा\-त\-न\-प्र\-भृ\-ति\-भि\-र\-भि\-घा\-त\-वि\-शे\-षना}}{\lemma{\skt{ग्राम्य\-ध\-र्म्म\-या\-न\-वा\-ह\-ना\-ध्व\-ग\-म\-न\-प्र\-स्ख\-ल\-न\-प्र\-प\-त\-न\-पी\-ड\-ना\-भि\-घा\-त\-वि\-ष\-म\-श\-य\-नो\-प\-वा\-स\-वे\-गा\-भि\-घा\-ता\-ति\-रू\-क्ष\-क\-टु\-क\-ति\-क्त\-भो\-ज\-ना\-जी\-र्ण्ण\-ग\-र्भ\-शा\-त\-न\-प्र\-भृ\-ति\-भिर्}}\Afootnote{H; \skt{ग्राम्य\-ध\-र्म\-या\-न\-वा\-ह\-ना\-ध्व\-ग\-म\-न\-प्र\-स्ख\-ल\-न\-प्र\-प\-त\-न\-प्र\-पी\-डा\-न\-धा\-व\-ना\-भि\-घा\-त\-वि\-ष\-म\-श\-य\-ना\-स\-नो\-प\-वा\-स\-वे\-गा\-भि\-घा\-ता\-ति\-रू\-क्ष\-क\-टु\-ति\-क्त\-भो\-ज\-न\-शो\-का\-ति\-क्षा\-र\-से\-व\-ना\-ति\-सा\-र\-व\-म\-न\-वि\-रे\-च\-न\-प्रे\-ङ्खो\-ल\-ना\-जी\-र्ण\-ग\-र्भ\-शा\-त\-न\-प्र\-भृ\-ति\-भिर्}\ A; \skt{ग्राम्य\-ध\-र्म\-या\-न\-वा\-ह\-ना\-ध्व\-ग\-म\-न\-प्र\-स्ख\-ल\-न\-प्र\-प\-त\-न\-पी\-ड\-ना\-भि\-घा\-त\-वि\-ष\-म\-श\-य\-ना\-स\-नो\-प\-वा\-स\-वे\-गा\-भि\-घा\-ता\-ति\-रू\-क्षा\-क\-टु\-क\-ति\-क्त\-भो\-ज\-ना\-जी\-र्ण्ण\-ग\-र्भ\-शा\-त\-न\-प्र\-भृ\-ति\-भिर्}\ N}}}{\lemma{\skt{अभि\-घा\-त\-वि\-शे\-षना}}\Afootnote{H; \skt{विशे\-षै\-र्ब\-न्धा\-नान्म्}\ A; \skt{अभि\-घा\-त\-वि\-शे\-ष\-च\-त्व\-रान्य्}\ N}}\skt{ }\edtext{\skt{च्छिद्यन्ते}}{\Afootnote{H; \skt{उच्यते}\ A; \skt{उच्छिद्यन्ते}\ N}}\skt{ }\edtext{\skt{गर्भम्फलमिव}}{\lemma{\skt{गर्भम्}}\Afootnote{H; \skt{गर्भः,}\ A; \skt{गर्व्भम्}\ N}}\skt{ }\edtext{\skt{वृत्तबन्धनात् । सख\-लु\-मु\-क्त\-ब\-न्ध\-न\-त्वा\-द्ग\-र्भ\-श\-य्या\-म\-ति\-क्रम्य}}{\lemma{\skt{वृत्तबन्धनात् । सख\-लु\-मु\-क्त\-ब\-न्ध\-न\-त्वा\-द्ग\-र्भ\-श\-य्याम्}}\Afootnote{H; \skt{वृन्त\-ब\-न्ध\-ना\-द\-भि\-घा\-त\-वि\-शेषैः; स विमुक्तबन्धनो गर्भाशयम्}\ A; \skt{वृत्तवन्धनात् । स खलु मुक्त\-व\-न्ध\-न\-त्वा\-द्ग\-र्भ\-श\-य्याम्}\ N}}\skt{ }\edtext{\edtext{\skt{यकृ\-त्प्ली\-हा\-न्त्र\-वि\-व\-गै\-र\-व\-स्र\-न्स\-मानः}}{\lemma{\skt{यकृ\-त्प्ली\-हा\-न्त्र\-वि\-व\-गैर्}}\Afootnote{H; \skt{यकृ\-त्प्ली\-हा\-न्त्र\-वि\-व\-रैर्}\ A N}}}{\lemma{\skt{अवस्रन्समानः}}\Afootnote{A H; \skt{अवस्रस्तमानः}\ N}}\skt{ }\edtext{\skt{कोष्ठ}}{\Afootnote{A H; \skt{कोष्ठं}\ N}}\skt{ }\edtext{\skt{संक्षोभयति}}{\Afootnote{N H; \skt{संक्षो\-भ\-मा\-पा\-द\-यति,}\ A}}\skt{ । }\edtext{\skt{तस्याः}}{\Afootnote{N H; \skt{तस्या}\ A}}\skt{ }\edtext{\skt{कोष्ठ\-सं\-क्षो\-भा\-द्वा\-यु\-र\-पानो}}{\Afootnote{N; \skt{जठ\-र\-सं\-क्षो\-मा\-द्वा\-यु\-र\-पानो}\ A; \skt{कोष्ठ\-स\-क्षो\-मा\-द्वा\-यु\-र\-व्यंनो}\ H}}\skt{ }\edtext{\skt{मूढः}}{\Afootnote{A H; \skt{मूढं}\ N}}\skt{ }\edtext{\edtext{\skt{पार्श्व\-ब\-स्ति\-शि\-रौ\-द\-र\-यो\-नि\-शू\-ला\-ना\-ह\-मू\-त्र\-स\-द्गा\-ना\-पा\-द्य\-ग\-र्भ\-स्प्र\-च्या\-व\-यति}}{\lemma{\skt{पार्श्व\-ब\-स्ति\-शि\-रौ\-द\-र\-यो\-नि\-शू\-ला\-ना\-ह\-मू\-त्र\-स\-द्गा\-ना\-पा\-द्य\-ग\-र्भस्}}\Afootnote{H; \skt{पार्श्व\-ब\-स्ति\-शी\-र्षो\-द\-र\-यो\-नि\-शू\-ला\-ना\-ह\-मू\-त्र\-स\-ङ्गा\-ना\-म\-न्य\-त\-म\-मा\-पाद्य गर्भं}\ A; \skt{प्रार्श्ववस्तिशिर उद\-र\-यो\-नि\-शू\-ला\-ना\-ह\-मू\-त्र\-स\-द्ग\-ना\-षा\-द्य\-ग\-र्भस्}\ N}}}{\lemma{\skt{प्रच्यावयति}}\Afootnote{N H; \skt{च्यावयति}\ A}}\skt{ । तरुणं }\edtext{\skt{शोणितभावेन}}{\Afootnote{N H; \skt{शोणितस्रावेण; तमेव}\ A}}\skt{ }\edtext{\edtext{\skt{कदा\-चि\-द्वि\-वृ\-द्ध\-म\-स\-म्य\-गा\-ग\-त\-म\-प\-त्य\-प\-थ\-म\-नु\-प्रा\-प्त\-म\-नि\-र्ग\-च्छ\-न्त\-म\-पा\-न\-वै\-गु\-ण्य\-स\-म्मो\-हि\-त\-ङ्गर्भं}}{\lemma{\skt{अनिर्गच्छन्तम्}}\Afootnote{N H; \skt{अनिर}\ A}}}{\lemma{\skt{अपा\-न\-वै\-गु\-ण्य\-स\-म्मो\-हि\-तङ्}}\Afootnote{H; \skt{स्यमानं विगु\-णा\-पा\-न\-सं\-मो\-हितं}\ A; \skt{अपा\-न\-वै\-गु\-ण्या\-स\-मो\-हितं}\ N}}\skt{ }\edtext{\skt{मूढ\-ग\-र्भ\-मि\-त्या\-च\-क्षते}}{\lemma{\skt{आचक्षते}}\Afootnote{A N; \skt{आक्षते ॥}\ H}}\skt{ ॥}\pend

\pstart \ledsidenote{\passref{2.8.4}}\edtext{\skt{स}}{\Afootnote{N H; \skt{ततः}\ A}}\skt{ कीलः }\edtext{\skt{प्रतिखुरो}}{\Afootnote{A H; \skt{प्रतिखुरो}\ N}}\skt{ }\edtext{\skt{बीजकः}}{\Afootnote{A H; \skt{वीजकः}\ N}}\skt{ परिघ }\edtext{\skt{इति}}{\Afootnote{N H; \skt{इ}\ A}}\skt{ }\edtext{\skt{चतुर्व्विधो मूढगर्भो }\edtext{\skt{भवतीत्येके}}{\lemma{\skt{एके}}\Afootnote{N H; \emph{om.}\ A}}\skt{ भाषन्ते}}{\lemma{\skt{ति… भवतीत्य्}}\Afootnote{\emph{om.}\ A}}\skt{ ॥ य }\edtext{\skt{ऊर्ध्वं}}{\Afootnote{N H; \skt{ऊर्ध्व}\ A}}\skt{ }\edtext{\skt{शिरःपादाभ्यां}}{\Afootnote{N H; \skt{बाहुशिरःपादो यो}\ A}}\skt{ योनि\-मु\-ख\-न्नि\-रु\-णद्धि कील इव स कीलः । }\edtext{\skt{निःसृ\-त\-ह\-स्त\-पा\-द\-शि\-रः\-का\-य\-शक्तः स}}{\Afootnote{H; \skt{निःसृ\-त\-ह\-स्त\-पा\-द\-शिराः कायसङ्गी}\ A; \skt{निःसृ\-त\-ह\-स्त\-पा\-द\-शिरः । कायशक्तः}\ N}}\skt{ }\edtext{\skt{प्रतिखुरः । यस्तु}}{\Afootnote{N H; \skt{प्रतिखुरः; यो}\ A}}\skt{ निर्गच्छति }\edtext{\skt{शिरोभुजः स}}{\Afootnote{\skt{एकशिरोभुजः स}\ A; \skt{शिरो भुजः स}\ N H}}\skt{ }\edtext{\skt{बीजकः}}{\Afootnote{\skt{बीजकः; यस्तु}\ A; \skt{वीजकः}\ N H}}\skt{ । परिघ इव योनि\-मु\-ख\-मा\-वृत्य }\edtext{\skt{तिष्ठेत्स}}{\Afootnote{N H; \skt{तिष्ठति स}\ A}}\skt{ परिघ }\edtext{\skt{इति}}{\Afootnote{N H; \skt{इति चतुर्विधो भवतीत्य्}\ A}}\skt{ । }\edtext{\skt{तत्तु न}}{\Afootnote{N H; \skt{एके भाषन्ते । तत्तु न}\ A}}\skt{ }\edtext{\skt{सम्यक् ॥ स}}{\Afootnote{N H; \skt{सम्यक्; कस्मात्स}\ A}}\skt{ यदा }\edtext{\skt{विगु\-णा\-नि\-ल\-पी\-डि\-तोऽ\-प\-त्य\-प\-थ\-म\-ने\-कधा}}{\lemma{\skt{विगु\-णा\-नि\-ल\-पी\-डितो}}\Afootnote{H; \skt{विगु\-णा\-नि\-ल\-प्र\-पी\-डितो}\ A; \skt{विगु\-णा\-नि\-पी\-डितो}\ N}}\skt{ }\edtext{\skt{प्रतिपद्यते}}{\Afootnote{N H; \skt{प्रपद्यते}\ A}}\skt{ । तदा }\edtext{\skt{चतुःसंख्या}}{\Afootnote{N H; \skt{सङ्ख्या}\ A}}\skt{ हीयते ॥}\edtext{}{\lemma{\skt{भाषन्ते ॥ य}}\Afootnote{N H; \skt{तत्र}\ A}}\pend

\pstart \ledsidenote{\passref{2.8.5}}\edtext{\skt{तत्र }\edtext{\skt{कश्चिद्द्वाभ्यां}}{\lemma{\skt{चिद्द्वाभ्यां}}\Afootnote{N H; \skt{श्चिद्द्वाभ्यां}\ A}}}{\lemma{\skt{तत्र कश्}}\Afootnote{N H; \skt{तत्र क}\ A}}\skt{ सक्थिभ्यां }\edtext{\edtext{\skt{योनि\-म\-भि\-प्र\-प\-द्यते}}{\lemma{\skt{योनिम्}}\Afootnote{N H; \skt{योनि}\ A}}}{\lemma{\skt{अभिप्रपद्यते}}\Afootnote{N H; \skt{मुखं प्रतिपद्यते;}\ A}}\skt{ । }\edtext{\skt{कश्चि\-दा\-भु\-ग्नै\-क\-सक्थि}}{\Afootnote{H; \skt{कश्चि\-दा\-भु\-ग्नै\-क\-स\-क्थिर्}\ A N}}\skt{ । }\edtext{\skt{कश्चि\-दा\-भु\-ग्न\-शि\-रः\-स्फि\-दे\-शे\-न\-ति\-र्य\-गा\-गतः}}{\Afootnote{H; \skt{एकेन; कश्चि\-दा\-भु\-ग्न\-स\-क्थि\-श\-रीरः स्फिग्देशेन तिर्यगागतः;}\ A; \skt{कश्चि\-दा\-भु\-ग्न\-शिरः स्फिग्देशेन तिर्यगागतः ।}\ N}}\skt{ । }\edtext{\skt{कश्चि\-दु\-द\-र\-पा\-र्श्व\-पृ\-ष्ठा\-ना\-म\-न्य\-त\-मेन}}{\lemma{\skt{कश्चि\-दु\-द\-र\-पा\-र्श्व\-पृ\-ष्ठा\-नाम्}}\Afootnote{N H; \skt{कश्चि\-दु\-रः\-पा\-र्श्व\-पृ\-ष्ठा\-नाम्}\ A}}\skt{ }\edtext{\skt{योनि\-द्वा\-र\-म्पि\-धाय तिष्ठति}}{\lemma{\skt{पिधाय तिष्ठति}}\Afootnote{N H; \skt{पिधा\-या\-व\-ति\-ष्ठते; अन्तः\-पा\-र्श्वा\-प\-वृ\-त्त\-शिराः}\ A}}\skt{ । कश्चिदेकेन }\edtext{\skt{बाहुना}}{\Afootnote{A H; \skt{वाहुना}\ N}}\skt{ }\edtext{\skt{पार्श्वा\-प\-वृ\-त्त\-शिरा}}{\Afootnote{N H; \emph{om.}\ A}}\skt{ । }\edtext{\skt{कश्चि\-दा\-भु\-ग्न\-शि\-रा\-बा\-हु\-द्व\-येन}}{\Afootnote{A H; \skt{कश्चि\-दा\-भु\-ग्न\-शि\-रा\-वा\-हु\-द्व\-येन ।}\ N}}\skt{ । कश्चि\-दा\-भु\-ग्न\-मध्यो हस्त\-पा\-द\-शि\-रोभिः । कश्चिदेकेन }\edtext{\skt{सक्थ्ना\-यो\-नि\-म\-भि\-प्र\-प\-द्या\-प\-रेण}}{\Afootnote{\skt{सक्थ्ना योनिमुखं प्रति\-प\-द्य\-तेऽ\-प\-रेण}\ A; \skt{सक्थ्न\-का\-यो\-नि\-म\-भि\-प्र\-प\-द्या\-प\-रेण}\ N; \skt{सक्थ्ना\-यो\-नि\-म\-भि\-प्र\-प\-द्या\-प\-रेण}\ H}}\skt{ }\edtext{\edtext{\skt{पायु\-मि\-त्य\-ष्ट\-वि\-धा\-मू\-ढ\-ग\-र्भ\-ग\-ति\-रु\-द्दिष्टा}}{\lemma{\skt{अष्ट\-वि\-धा\-मू\-ढ\-ग\-र्भ\-ग\-तिर्}}\Afootnote{H; \skt{अष्टविधा मूढगर्भगति}\ A; \skt{अष्टविधो मूढ\-ग\-र्भा\-ग\-तिर्}\ N}}}{\lemma{\skt{उद्दिष्टा}}\Afootnote{H; \skt{रुद्दिष्टा}\ A; \skt{उद्दिष्टाः ।}\ N}}\skt{ }\edtext{\skt{समासेन}}{\Afootnote{A H; \skt{समासेनः}\ N}}\skt{ ।}\pend

\pstart \ledsidenote{\passref{2.8.6}}\skt{तत्र }\edtext{\skt{द्वावन्त्यौ}}{\Afootnote{\skt{द्वावन्त्या}\ A; \skt{द्वाभ्यां वंत्यौ}\ N; \skt{द्वावन्त्यौ}\ H}}\skt{ }\edtext{\skt{मूढ\-ग\-र्भा\-व\-साध्यौ}}{\Afootnote{N; \skt{वसाध्यौ}\ A; \skt{मूढगर्भावसा}\ H}}\skt{ }\edtext{\edtext{\skt{शेषेष्वपि}}{\lemma{\skt{शेषेष्व्}}\Afootnote{N; \skt{मूढगर्भौ, शेषा}\ A; \skt{ध्यौ शेषेष्व्}\ H}}}{\lemma{\skt{अपि}}\Afootnote{H; \skt{नपि}\ A; \skt{अपि तु}\ N}}\skt{ }\edtext{\skt{विप\-री\-ते\-न्द्रि\-या\-र्थ\-सा\-क्षे\-प\-क\-यो\-नि\-भ्रं\-स\-स\-म्व\-र\-ण\-स\-क्क\-ल\-श्वा\-स\-का\-स\-भ्र\-म\-नि\-पी\-डि\-तंश्च}}{\Afootnote{\skt{विप\-री\-ते\-न्द्रि\-या\-र्था\-क्षे\-प\-क\-यो\-नि\-भ्रं\-श\-स्ं\-व\-र\-ण\-म\-क्क\-ल्ल\-श्वा\-स\-का\-स\-भ्र\-म\-नि\-पी\-डिता}\ A; \skt{विप\-री\-ते\-न्द्रि\-या\-र्थ\-सा\-क्षे\-प\-क\-यो\-नि\-भ्र\-ञ्स\-स\-म्व\-र\-ण\-स\-क्क\-ल्ल\-श्वा\-स\-का\-स\-भ्र\-म\-नि\-पी\-डि\-तंश्च}\ N; \skt{विप\-री\-ते\-न्द्रि\-या\-र्थ\-सा\-क्षे\-प\-क\-यो\-नि\-भ्रं\-स\-स\-म्व\-र\-ण\-स\-क्क\-ल\-श्वा\-स\-का\-स\-भ्र\-म\-नि\-पी\-डि\-तॉंश्च}\ H}}\skt{ }\edtext{\skt{परिहरेत्}}{\Afootnote{N H; \skt{न्परिहरेत् ॥}\ A}}\skt{ ॥}\pend

\pstart\versepar \ledsidenote{\passref{2.8.7}}\edtext{\skt{भवन्ति}}{\Afootnote{A H; \skt{भ वन्ति चात्र श्लोकाः}\ N}}\skt{ }\edtext{\skt{चात्र}}{\Afootnote{A H; \emph{om.}\ N}}\skt{ }\edtext{\skt{श्लोकाः॥}}{\lemma{\skt{श्लोकाः}}\Afootnote{H; \emph{om.}\ A N}}\skt{ }\pend
\pstart\versepar \skt{कालस्य }\edtext{\skt{परिमाणेन}}{\Afootnote{H; \skt{परिणामेन}\ A; \skt{परिमाणौ न}\ N}}\skt{ }\edtext{\skt{मुक्त\-वृ\-त्ता\-द्य\-था\-फलं}}{\Afootnote{N; \skt{मुक्तं वृन्ताद्यथा पहम् ।}\ A; \skt{मुक्त\-वृ\-त्ता\-द्य\-था\-फलं ।}\ H}}\skt{ । }\pend
\pstart\versepar \skt{प्रपद्यते स्वभावेन नान्यथा }\edtext{\skt{पृथिवीतलम्}}{\Afootnote{N H; \skt{पतितुं ध्रुवम् ॥}\ A}}\skt{ ॥}\pend

\pstart\versepar \ledsidenote{\passref{2.8.8}}\skt{एवं कालप्रकर्षेण मुक्तो नाडीनि बन्धनात् । }\pend
\pstart\versepar \skt{गर्भाशयस्थो हि }\edtext{\skt{गर्भो}}{\Afootnote{H; \skt{यो गर्भो}\ A; \skt{जज गर्भो}\ N}}\skt{ जननाय प्रपद्यते ॥}\edtext{}{\lemma{\skt{बन्धनात् । गर्भाशयस्थो हि}}\Afootnote{\skt{बन्धनात् । गर्भाशयस्थो}\ A; \skt{वन्धनात् । गर्भाशयस्थो हि}\ N; \skt{वम्बन्धनात् । गर्भासयस्थो हि}\ H}}\pend

\pstart\versepar \ledsidenote{\passref{2.8.9}}\edtext{\skt{क्रिमि\-वा\-ता\-भि\-घा\-तैस्तु}}{\Afootnote{N H; \skt{कृमि\-वा\-ता\-भि\-घा\-तैस्तु}\ A}}\skt{ }\edtext{\skt{तदे\-वो\-प\-द्रु\-त\-म्फलं}}{\lemma{\skt{तदे\-वो\-प\-द्रु\-तम्}}\Afootnote{N H; \skt{तदेवोपद्रुतं}\ A}}\skt{ । }\pend
\pstart\versepar \skt{}\edtext{\skt{पतत्यकालेऽपि}}{\lemma{\skt{ऽपि}}\Afootnote{A; \skt{पि}\ N H}}\skt{ यथा तथा }\edtext{\skt{स्याद्गर्भविच्युतिः॥}}{\lemma{\skt{स्याद्गर्भविच्युतिः}}\Afootnote{A H; \skt{स्यां गर्भविच्युतिः ॥}\ N}}\pend

\pstart\versepar \ledsidenote{\passref{2.8.10}}\edtext{\skt{आचतुर्थान्ततो}}{\Afootnote{N H; \skt{आचतुर्थात्ततो}\ A}}\skt{ }\edtext{\skt{मासात्प्रसृते}}{\Afootnote{H; \skt{मासात्प्रस्रवे}\ A; \skt{मासात्प्रसूते}\ N}}\skt{ }\edtext{\skt{गर्भविच्युतः}}{\Afootnote{N H; \skt{द्गर्भविच्युतिः ।}\ A}}\skt{ । }\pend
\pstart\versepar \skt{ततः स्थिरशरीरस्य }\edtext{\skt{पादः}}{\Afootnote{N H; \skt{पातः}\ A}}\skt{ पञ्चमषष्ठयोः ॥}\pend

\pstart\versepar \ledsidenote{\passref{2.8.11}}\edtext{\skt{अप\-वि\-द्ध\-श\-रीरा तु}}{\Afootnote{N H; \skt{प्रविध्यति शिरो या तु}\ A}}\skt{ शीताङ्गी }\edtext{\skt{निरपत्रपा}}{\Afootnote{A H; \skt{निरपत्रया ।}\ N}}\skt{ । }\pend
\pstart\versepar \skt{नीलोद्धतसिरा }\edtext{\skt{हन्ति मा}}{\Afootnote{N H; \skt{हन्ति सा}\ A}}\skt{ }\edtext{\skt{गर्भन्स}}{\Afootnote{\skt{गर्भं स}\ A H; \skt{गर्भस}\ N}}\skt{ }\edtext{\skt{चताःपूनः}}{\Afootnote{H; \skt{च तां तथा ॥}\ A; \skt{चताम्पूनः ॥}\ N}}\skt{ ॥}\pend

\pstart\versepar \ledsidenote{\passref{2.8.12}}\edtext{\skt{गर्भा\-स्प\-न्द\-न\-मा\-वी\-ना\-म्प्र\-णाशः}}{\Afootnote{A; \skt{गर्भास्पंदन नोपीनां प्रणाशः}\ N; \skt{गर्भा\-स्प\-न्द\-न\-मा\-वी\-ना\-म्प्र\-णाशः}\ H}}\skt{ श्यावपाण्डुता । }\pend
\pstart\versepar \skt{}\edtext{\skt{भवे\-दु\-च्छ्वा\-स\-पू\-तित्वं}}{\Afootnote{N H; \skt{भव\-त्यु\-च्छ्वा\-स\-पू\-तित्वं}\ A}}\skt{ }\edtext{\edtext{\skt{शूलश्चान्तर्म्मृते}}{\lemma{\skt{शूलश्}}\Afootnote{N H; \skt{शूलं}\ A}}}{\lemma{\skt{चान्तर्म्मृते}}\Afootnote{H; \skt{चान्तर्भृते}\ A; \skt{चान्तमृंते}\ N}}\skt{ शिशौ ॥}\pend

\pstart\versepar \ledsidenote{\passref{2.8.13}}\edtext{\skt{मान\-सा\-ग\-न्तु\-भि\-र्म्मा\-तु\-रु\-प\-तापैः}}{\lemma{\skt{म्मातुरुपतापैः}}\Afootnote{\skt{मातुरुपतापैः}\ A N H}}\skt{ प्रपीडितः । }\pend
\pstart\versepar \skt{गर्भो व्यापद्यते कुक्षौ }\edtext{\skt{व्याधिभिश्च}}{\Afootnote{A H; \skt{व्याधिभि सु}\ N}}\skt{ }\edtext{\skt{निपीडितः}}{\Afootnote{N H; \skt{प्रपीडितः ॥}\ A}}\skt{ ॥}\pend

\pstart\versepar \ledsidenote{\passref{2.8.14}}\edtext{\skt{कुक्षौ}}{\Afootnote{H; \skt{बस्त}\ A; \skt{कुक्षे}\ N}}\skt{ }\edtext{\skt{मातुर्व्विपन्नाया}}{\Afootnote{N H; \skt{मारविपन्नायाः}\ A}}\skt{ }\edtext{\skt{गर्भः}}{\Afootnote{N H; \skt{कुक्षिः}\ A}}\skt{ प्रस्पन्दते }\edtext{\skt{स्त्रियाः}}{\Afootnote{N H; \skt{यदि । तत्क्षणा}\ A}}\skt{ । }\pend
\pstart\versepar \skt{}\edtext{\skt{जन्मकाले}}{\Afootnote{N H; \skt{ज्जन्मकाले}\ A}}\skt{ }\edtext{\skt{मुहुर्त्तात्तं}}{\Afootnote{N H; \skt{तं}\ A}}\skt{ }\edtext{\skt{पाट\-यि\-त्वो\-द्ध\-रे\-च्छि\-शुम्}}{\Afootnote{N H; \skt{पाट\-यि\-त्वो\-द्ध\-रे\-द्भि\-षक् ॥}\ A}}\skt{ ॥}\pend

\pstart\versepar \ledsidenote{\passref{2.8.14.1}}\skt{यदा सोऽन्तर्म्मृतो गर्भः }\edtext{\edtext{\skt{शुनो\-ब\-स्ति\-रि\-वा\-ततः}}{\lemma{\skt{शुनोबस्तिर्}}\Afootnote{H; \skt{शुनोवस्ति}\ N}}}{\lemma{\skt{इवाततः}}\Afootnote{H; \skt{रिवाततः ।}\ N}}\skt{ । }\pend
\pstart\versepar \skt{}\edtext{\skt{तेनावृता च}}{\Afootnote{H; \skt{तोनावृता च}\ N}}\skt{ }\edtext{\skt{नार्यस्तु}}{\Afootnote{H; \skt{नार्यास्तु}\ N}}\skt{ कुक्षिरुन्नह्यते भृशम् ॥}\pend

\pstart\versepar \ledsidenote{\passref{2.8.14.2}}\skt{उत्क्षिप्य इव चाङ्गानि }\edtext{\skt{मूत्रबस्तिश्च}}{\Afootnote{H; \skt{मूत्रवस्तिश्च}\ N}}\skt{ विद्यते । }\pend
\pstart\versepar \skt{क्लोम\-प्ली\-हा\-य\-कृ\-च्चैव }\edtext{\skt{फुफ्फुसं}}{\Afootnote{H; \skt{पुप्फुसं}\ N}}\skt{ }\edtext{\skt{हृदयन्तथा}}{\Afootnote{H; \skt{हृद्युयन्तथा ॥}\ N}}\skt{ ॥ }\pend
\pstart\versepar \skt{गर्भेण }\edtext{\skt{पीडिता }\edtext{\skt{ह्येत\-दू\-र्ध्व\-म्प्र\-क्रा\-मति}}{\lemma{\skt{एतदूर्ध्वम्}}\Afootnote{N; \skt{ह्येत दूर्ध्वम्}\ H}}}{\lemma{\skt{पीडिता ह्य्}}\Afootnote{N; \skt{पीडिता}\ H}}\skt{ स्त्रियाः । }\pend
\pstart\versepar \skt{सा सूयते मुह्यति च कृच्छ्रोच्छ्वासा च }\edtext{\skt{जायते}}{\Afootnote{H; \skt{ज्ञाते ॥}\ N}}\skt{ ॥}\edtext{}{\lemma{\skt{स्त्रियाः । सा}}\Afootnote{H; \skt{स्त्रिया । सा}\ N}}\pend

\pstart\versepar \ledsidenote{\passref{2.8.14.3}}\edtext{\skt{पूति\-ग\-न्ध्य\-न्त\-था\-स्वेदो}}{\Afootnote{H; \skt{पूतिगन्धस्तथा स्वेदो}\ N}}\skt{ }\edtext{\skt{जिह्वातालू च}}{\Afootnote{H; \skt{जिह्वातालुश्च}\ N}}\skt{ शुष्यति । }\pend
\pstart\versepar \skt{वेपते }\edtext{\skt{स्राम्यति}}{\Afootnote{H; \skt{भ्राम्यति}\ N}}\skt{ तथा }\edtext{\skt{जीवि\-त\-ञ्चो\-प\-रु\-ध्यते}}{\lemma{\skt{जीवितञ्}}\Afootnote{H; \skt{जीवतं}\ N}}\skt{ ॥ }\pend
\pstart\versepar \skt{}\edtext{\edtext{\skt{एतै\-र्ल्लि\-ङ्गै\-र्व्वि\-जा\-नी\-या\-न्मृ\-त\-ग\-र्भ\-ञ्चि\-कि\-त्सक}}{\lemma{\skt{ल्लिङ्गैर्}}\Afootnote{N; \skt{ल्लिङ्गौर्}\ H}}}{\lemma{\skt{व्विजा\-नी\-या\-न्मृ\-त\-ग\-र्भञ्}}\Afootnote{H; \skt{विज्ञा\-नी\-या\-न्मृ\-त\-ग\-र्भञ्}\ N}}\skt{ इति ॥}\pend

\pstart \ledsidenote{\passref{2.8.trailer}}\edtext{\edtext{\skt{मूढ\-ग\-र्भ\-नि\-दा\-नेऽ\-ष्टमः॥}}{\lemma{\skt{मूढ\-ग\-र्भ\-नि\-दाने}}\Afootnote{H; \skt{इति सुश्रु\-त\-सं\-हि\-तायां निदानस्थाने मूढ\-ग\-र्भ\-नि\-दानं ना}\ A; \skt{नि ॥ द्र ॥ ९ ॥}\ N}}}{\lemma{\skt{ऽष्टमः}}\Afootnote{H; \skt{माष्टमोऽध्यायः ॥ ८ ॥}\ A; \emph{om.}\ N}}\pend

\endnumbering
\beginnumbering
\eledsection*{Adhyāya 9}
\pstart \ledsidenote{\passref{2.9.1}}\skt{अथातो विद्रधीनां }\edtext{\skt{निदा\-न\-म्व्या\-ख्या\-स्यामः}}{\lemma{\skt{व्याख्यास्यामः}}\Afootnote{A H; \emph{om.}\ N}}\skt{ ।}\pend

\pstart\versepar \ledsidenote{\passref{2.9.3}}\skt{सर्व्वामरगुरुः श्रीमा\-न्नि\-मि\-त्ता\-न्त\-र\-भू\-मिपः । }\pend
\pstart\versepar \skt{}\edtext{\skt{शिष्यायोवाच}}{\Afootnote{A H; \skt{सिष्यायोवाच}\ N}}\skt{ निखि\-ल\-मि\-द\-म्वि\-द्र\-धि\-ल\-क्ष\-णम् ॥}\pend

\pstart\versepar \ledsidenote{\passref{2.9.4}}\edtext{\skt{त्वग्र\-क्त\-मां\-स\-मे\-दान्सि}}{\Afootnote{A H; \skt{त्वग्रक्तमांसं मेदांसि}\ N}}\skt{ प्रदूष्यास्थि समाश्रिताः । }\pend
\pstart\versepar \skt{दोषाः शोफं }\edtext{\skt{शनै\-र्घो\-र\-ञ्ज\-न\-य\-न्त्यु\-च्छ्रिता}}{\lemma{\skt{घोर\-ञ्ज\-न\-य\-न्त्यु\-च्छ्रिता}}\Afootnote{H; \skt{घोरं जनयन्त्युच्छ्रिता}\ A N}}\skt{ }\edtext{\skt{भृशम्}}{\Afootnote{A H; \skt{भृषम् ॥}\ N}}\skt{ ॥}\pend

\pstart\versepar \ledsidenote{\passref{2.9.5}}\edtext{\skt{महामूलं}}{\Afootnote{A H; \skt{महामूला}\ N}}\skt{ }\edtext{\skt{रुजावन्तं}}{\Afootnote{A H; \skt{रुजावन्त}\ N}}\skt{ }\edtext{\edtext{\skt{वृत्तम्वाप्यथवा}}{\lemma{\skt{वाप्य्}}\Afootnote{\skt{चा (वा)}\ A; \skt{चा प्य्}\ N; \skt{वा}\ H}}\skt{ यतं}}{\lemma{\skt{अथवा यतं}}\Afootnote{N; \skt{ऽप्य\-थ\-वा++\-आ\-य\-तम् ।}\ A; \skt{प्यथवायतं ।}\ H}}\skt{ । }\pend
\pstart\versepar \skt{}\edtext{\edtext{\skt{सविद्राधिरिति}}{\lemma{\skt{सविद्राधिर्}}\Afootnote{H; \skt{तमाहुर्विद्रधिं}\ A; \skt{स विद्रधिर्}\ N}}}{\lemma{\skt{इति}}\Afootnote{N H; \emph{om.}\ A}}\skt{ }\edtext{\skt{ख्याते}}{\Afootnote{H; \skt{धीरा}\ A; \skt{स्या ख्यातो}\ N}}\skt{ }\edtext{\skt{विज्ञेयः}}{\Afootnote{N H; \skt{विज्ञेयः स}\ A}}\skt{ }\edtext{\skt{षड्विधश्च}}{\Afootnote{N H; \skt{च षङ्विधः ॥}\ A}}\skt{ }\edtext{\skt{यः}}{\Afootnote{N H; \emph{om.}\ A}}\skt{ ॥}\pend

\pstart\versepar \ledsidenote{\passref{2.9.6}}\skt{पृथग्दोषैः समस्तैश्च क्षतेनाप्यसृजा तथा । }\pend
\pstart\versepar \skt{षण्णामपि हि तेषान्तु लक्षणं सम्प्रवक्ष्यते ॥}\pend

\pstart\versepar \ledsidenote{\passref{2.9.7}}\edtext{\edtext{\skt{कृष्णोऽरुणो}}{\lemma{\skt{कृष्णो}}\Afootnote{A H; \skt{कृष्णा}\ N}}\skt{ वा}}{\lemma{\skt{ऽरुणो वा}}\Afootnote{A; \skt{रुणो वा}\ N H}}\skt{ परुषो }\edtext{\skt{भृष\-म\-त्य\-र्थ\-वे\-दनः}}{\lemma{\skt{भृषम्}}\Afootnote{N H; \skt{भृशम्}\ A}}\skt{ । }\pend
\pstart\versepar \skt{}\edtext{\skt{क्षिप्रो\-त्था\-न\-प्र\-पा\-कश्च}}{\Afootnote{N H; \skt{चित्रो\-त्था\-न\-प्र\-पा\-कश्च}\ A}}\skt{ }\edtext{\edtext{\skt{विद्र\-धि\-र्व्वा\-त\-सं\-भवः}}{\lemma{\skt{विद्रधिर्}}\Afootnote{A N\textsuperscript{pc} H; \skt{विध्रद्रधिर्}\ N\textsuperscript{ac}}}}{\lemma{\skt{व्वातसंभवः}}\Afootnote{A H; \skt{वातलक्षणः ॥}\ N}}\skt{ ॥}\pend

\pstart\versepar \ledsidenote{\passref{2.9.8}}\edtext{\skt{पक्वो\-दु\-म्ब\-र\-सं\-काशः}}{\Afootnote{A H; \skt{पक्वो\-दु\-म्ब\-र\-श\-ङ्काशः}\ N}}\skt{ }\edtext{\skt{श्यावो वा}}{\Afootnote{A H; \skt{स्यावो वा}\ N}}\skt{ ज्वरदाहवान् । }\pend
\pstart\versepar \skt{क्षिप्रो\-त्था\-न\-प्र\-पा\-कश्च विद्रधिः पित्तसंभवः ॥}\pend

\pstart\versepar \ledsidenote{\passref{2.9.9}}\edtext{\skt{सरावसदृशः}}{\Afootnote{H; \skt{शरवसदृशः}\ A; \skt{शरावसदृशः}\ N}}\skt{ पाण्डुः }\edtext{\skt{शीतः}}{\Afootnote{A H; \skt{सीत}\ N}}\skt{ }\edtext{\skt{स्निग्धाल्पवेदनः}}{\Afootnote{H; \skt{स्तब्धोऽल्पवेदनः ।}\ A; \skt{शिन्ग्धो ल्पवेदनः ।}\ N}}\skt{ । }\pend
\pstart\versepar \skt{}\edtext{\skt{चिरो\-त्था\-न\-प्र\-पा\-कश्च}}{\Afootnote{A H; \skt{चिरो\-स्था\-न\-प्र\-पा\-कश्च}\ N}}\skt{ }\edtext{\skt{सकुण्डश्च}}{\Afootnote{H; \skt{सकण्डुश्च}\ A; \skt{सकण्डृश्च}\ N}}\skt{ }\edtext{\skt{कफात्मकः}}{\Afootnote{N H; \skt{कफोत्थितः ॥}\ A}}\skt{ ॥}\pend

\pstart\versepar \ledsidenote{\passref{2.9.10}}\edtext{\skt{तनु\-पी\-त\-सि\-ता\-श्चै\-षा\-मा\-स्रावाः}}{\lemma{\skt{तनु\-पी\-त\-सि\-ताश्}}\Afootnote{A H; \skt{तनु\-पी\-त\-शि\-ताश्}\ N}}\skt{ क्रमशः स्मृताः । }\pend
\pstart\versepar \skt{नाना\-व\-र्ण्ण\-रु\-जा\-स्रावो }\edtext{\skt{घाटालो}}{\Afootnote{A H; \skt{घाटाले}\ N}}\skt{ }\edtext{\skt{विषमो}}{\Afootnote{A H; \skt{विषवो}\ N}}\skt{ महान् ॥}\pend

\pstart\versepar \ledsidenote{\passref{2.9.11}}\skt{विषमम्पच्यते चापि विद्रधिः सान्निपातिकः । }\pend
\pstart\versepar \skt{तैस्तै\-र्भा\-वै\-र\-भि\-हते क्षते वा }\edtext{\skt{पथ्यसेविनः}}{\Afootnote{N H; \skt{ऽपथ्यसेविनः ॥}\ A}}\skt{ ॥}\pend

\pstart\versepar \ledsidenote{\passref{2.9.12}}\skt{क्षतोष्मा }\edtext{\skt{वायुर्व्विसृतः}}{\Afootnote{H; \skt{वायुविसृतः}\ A; \skt{चायुर्विसृतः ।}\ N}}\skt{ सर\-क्त\-म्पि\-त्त\-मी\-र\-येत् । }\pend
\pstart\versepar \skt{ज्वरस्तृष्णा च दाहश्च जायते }\edtext{\skt{चास्य}}{\Afootnote{H; \skt{तस्य}\ A; \skt{वा स्य}\ N}}\skt{ देहिनः ॥}\edtext{}{\lemma{\skt{ईरयेत् । ज्वरस्}}\Afootnote{A H; \skt{ईरयेत् । ज्वर}\ N}}\pend

\pstart\versepar \ledsidenote{\passref{2.9.13}}\edtext{\edtext{\edtext{\skt{आग\-न्तु\-र्व्वि\-द्र\-धि\-स्त्वेष}}{\lemma{\skt{आगन्तुर्}}\Afootnote{H; \skt{एष विद्रधिरागन्तुः}\ A; \skt{आगन्तु}\ N}}}{\lemma{\skt{व्विद्रधिस्}}\Afootnote{H; \emph{om.}\ A; \skt{विद्रधि}\ N}}}{\lemma{\skt{त्वेष}}\Afootnote{H; \emph{om.}\ A; \skt{ष}\ N}}\skt{ पित्त\-वि\-द्र\-धि\-ल\-क्षणः । }\pend
\pstart\versepar \skt{}\edtext{\skt{कृष्णस्फोटावृतः}}{\Afootnote{A H; \skt{कृष्णःस्फोटावृतः}\ N}}\skt{ }\edtext{\skt{श्याव\-स्ती\-व्र\-दा\-ह\-रु\-जा\-ज्वरः}}{\lemma{\skt{श्यावस्}}\Afootnote{A H; \skt{स्याव}\ N}}\skt{ ॥}\pend

\pstart\versepar \ledsidenote{\passref{2.9.14}}\edtext{\skt{पित्त\-वि\-द्र\-धि\-लि\-ङ्गश्च}}{\Afootnote{H; \skt{पित्त\-वि\-द्र\-धि\-लि\-ङ्गस्तु}\ A; \skt{पित्त\-वि\-द्र\-धि\-लि\-ङ्गस्य}\ N}}\skt{ रक्त\-वि\-द्र\-धि\-रु\-च्यते । }\pend
\pstart\versepar \skt{उक्ता }\edtext{\skt{विद्रधयो}}{\Afootnote{A H; \skt{विद्रधयो ह्य्}\ N}}\skt{ }\edtext{\skt{ह्येते}}{\Afootnote{A H; \skt{एते}\ N}}\skt{ }\edtext{\skt{तेष्वसाध्यस्तु}}{\lemma{\skt{असाध्यस्तु}}\Afootnote{A H; \skt{असाध्यश्च}\ N}}\skt{ सर्व्वजः ॥}\pend

\pstart \ledsidenote{\passref{2.9.15ab}}\edtext{\skt{अभ्यन्तरानत}}{\Afootnote{N H; \skt{आभ्य\-न्त\-रा\-न\-तस्}\ A}}\skt{ }\edtext{\edtext{\skt{ऊर्ध्व\-म्वि\-द्र\-धी\-न्स\-म्प्र\-व\-क्ष्यते}}{\lemma{\skt{ऊर्ध्वम्}}\Afootnote{N H; \skt{तूर्ध्वं}\ A}}}{\lemma{\skt{विद्र\-धी\-न्स\-म्प्र\-व\-क्ष्यते}}\Afootnote{H; \skt{विद्र\-धी\-न्प\-रि\-च\-क्षते ।}\ A; \skt{विद्रधी सं प्रवक्ष्यते ।}\ N}}\skt{ ।}\pend

\pstart \ledsidenote{\passref{2.9.16cd}}\skt{पृथक्संभूय वा दोषाः कुपिता }\edtext{\skt{गुल्मरूपिणम्}}{\Afootnote{A H; \skt{गुल्मरूपिणः ॥}\ N}}\skt{ ॥}\pend

\pstart\versepar \ledsidenote{\passref{2.9.17}}\skt{वल्मी\-क\-व\-त्स\-मु\-न्न\-द्ध\-मन्तः कुर्वन्ति विद्रधिं । }\pend
\pstart\versepar \skt{गुदे }\edtext{\skt{बस्तिमुखे}}{\Afootnote{A N\textsuperscript{pc} H; \skt{बैस्तिमुखे}\ N\textsuperscript{ac}}}\skt{ नाभ्यां कुक्षौ वंक्षणयोस्तथा ॥}\pend

\pstart\versepar \ledsidenote{\passref{2.9.18}}\skt{वृक्कयोः }\edtext{\skt{प्लीह्नियकृति}}{\Afootnote{N H; \skt{यकृति}\ A}}\skt{ }\edtext{\skt{हृदिवाक्लोम्नि}}{\Afootnote{N H; \skt{प्लीह्नि हृदये क्लोम्नि}\ A}}\skt{ }\edtext{\edtext{\skt{वाप्यथ}}{\lemma{\skt{वाप्य्}}\Afootnote{N H; \skt{वा}\ A}}}{\lemma{\skt{अथ}}\Afootnote{\skt{तथा ।}\ A; \skt{अच ।}\ N; \skt{एथ ।}\ H}}\skt{ । }\pend
\pstart\versepar \skt{}\edtext{\edtext{\skt{एषामुक्तानि}}{\lemma{\skt{एषाम्}}\Afootnote{N H; \skt{तेषां}\ A}}}{\lemma{\skt{उक्तानि}}\Afootnote{N H; \emph{om.}\ A}}\skt{ लिंगानि }\edtext{\skt{बाह्य\-वि\-द्र\-धि\-ल\-क्षणैः}}{\Afootnote{N; \skt{जानी\-या\-द्बा\-ह्य\-वि\-द्र\-धि\-ल\-क्षणैः ॥}\ A; \skt{बाह्य\-वि\-द्र\-धि\-ल\-क्षण्णैः ॥}\ H}}\skt{ ॥}\pend

\pstart\versepar \ledsidenote{\passref{2.9.19}}\edtext{\skt{आम\-प\-क्वे\-ष\-णी\-याच्च}}{\Afootnote{H; \skt{आम\-प\-क्वै\-ष\-णी\-याच्च}\ A; \skt{आम\-प\-क्व\-षं\-णी\-याच्च}\ N}}\skt{ पक्वा पक्व\-म्वि\-भा\-व\-येत् । }\pend
\pstart\versepar \skt{अधिष्ठान }\edtext{\skt{विशेषेण}}{\Afootnote{A H; \skt{विशेषण}\ N}}\skt{ लिंगं शृणु विशेषतः ॥}\edtext{}{\lemma{\skt{विभावयेत् । अधिष्ठान}}\Afootnote{N H; \skt{विनिर्दिशेत् । अधिष्ठान}\ A}}\pend

\pstart\versepar \ledsidenote{\passref{2.9.20}}\skt{गुदे }\edtext{\skt{वातनिरोधस्तु}}{\Afootnote{A H; \skt{गतनिरोधस्तु}\ N}}\skt{ }\edtext{\skt{कृच्छ्राल्पमूत्रता}}{\Afootnote{H; \skt{बस्तौ कृच्छ्राल्पमूत्रता}\ A N}}\skt{ । }\pend
\pstart\versepar \skt{नाभ्यां हिक्का }\edtext{\skt{तथाटोपः}}{\Afootnote{N H; \skt{तथा++आटोपः}\ A}}\skt{ कुक्षौ मारु\-त\-को\-प\-नम् ॥}\pend

\pstart\versepar \ledsidenote{\passref{2.9.21}}\edtext{\skt{कटी\-पृ\-ष्ट\-ग्र\-ह\-स्तीव्रो}}{\lemma{\skt{कटी\-पृ\-ष्ट\-ग्र\-हस्}}\Afootnote{H; \skt{कटी\-पृ\-ष्ठ\-ग्र\-हस्}\ A N}}\skt{ }\edtext{\skt{वंक्षणोत्थे तु}}{\Afootnote{A H; \skt{वङ्क्षणोच्छे तु}\ N}}\skt{ विद्रधौ । }\pend
\pstart\versepar \skt{}\edtext{\skt{वृक्कयोः}}{\Afootnote{A H; \skt{वृक्षयोः}\ N}}\skt{ }\edtext{\skt{पार्श्वसंकोचः}}{\Afootnote{A H; \skt{या स्वसंकोचः}\ N}}\skt{ }\edtext{\skt{प्लीह्न्यु\-च्छ्वा\-स\-रो\-ध\-नम्}}{\lemma{\skt{उच्छ्वा\-स\-रो\-ध\-नम्}}\Afootnote{\skt{उच्छ्वा\-सा\-व\-रो\-ध\-नम् ॥}\ A; \skt{उच्छ्वासनिरोधं ॥}\ N; \skt{उच्छ्वा\-स\-नि\-रो\-धनं ॥}\ H}}\skt{ ॥}\pend

\pstart\versepar \ledsidenote{\passref{2.9.22}}\skt{सर्व्वा\-ङ्ग\-प्र\-ग्र\-ह\-स्तीव्रो }\edtext{\skt{हृदिकासश्च}}{\Afootnote{H; \skt{हृदि शूलश्च}\ A; \skt{हृदि कामश्च}\ N}}\skt{ }\edtext{\skt{जायते}}{\Afootnote{N H; \skt{दारुणः ।}\ A}}\skt{ । }\pend
\pstart\versepar \skt{}\edtext{\skt{श्वासो वा}}{\Afootnote{H; \skt{श्वासो}\ A; \skt{श्वासौ वा}\ N}}\skt{ यकृति तृष्णा }\edtext{\skt{क्लोम्नि}}{\Afootnote{N H; \emph{om.}\ A}}\skt{ }\edtext{\edtext{\edtext{\skt{पेपीयतेऽप्यपः}}{\lemma{\skt{पेपीयते}}\Afootnote{H; \skt{च पिपासा क्लोमजे}\ A; \skt{पै पीयते प्येषः ॥}\ N}}}{\lemma{\skt{ऽप्य्}}\Afootnote{H; \skt{ऽधि}\ A; \emph{om.}\ N}}}{\lemma{\skt{अपः}}\Afootnote{H; \skt{का ॥}\ A; \emph{om.}\ N}}\skt{ ॥}\pend

\pstart\versepar \ledsidenote{\passref{2.9.23}}\skt{आमो वा }\edtext{\skt{यदि}}{\Afootnote{H; \skt{यदि वा}\ A N}}\skt{ }\edtext{\skt{पक्वो वा}}{\Afootnote{N H; \skt{पक्वो}\ A}}\skt{ महान्वा यदि }\edtext{\skt{वेतरः}}{\Afootnote{A H; \skt{वेतर ।}\ N}}\skt{ । }\pend
\pstart\versepar \skt{सर्व्वो }\edtext{\skt{मर्म्मोत्थितत्वास्तु}}{\Afootnote{H; \skt{मर्मोत्थितश्चापि}\ A; \skt{मर्मोत्थितश्चात्रै}\ N}}\skt{ विद्रधिः }\edtext{\skt{कष्ट}}{\Afootnote{A H; \skt{कष्ठ}\ N}}\skt{ उच्यते ॥}\pend

\pstart\versepar \ledsidenote{\passref{2.9.23.add}}\edtext{\skt{नाभे\-रू\-प\-रि\-जातो यो}}{\lemma{\skt{ऊपरिजातो यो}}\Afootnote{H; \skt{उपरिजातो ।यो}\ N}}\skt{ }\edtext{\skt{मर्म्माभ्यासे च}}{\Afootnote{H; \skt{मर्माख्यां से च}\ N}}\skt{ विद्रधिः । }\pend
\pstart\versepar \skt{यः पक्वः }\edtext{\skt{सोऽप्यसाध्यः}}{\lemma{\skt{ऽप्यसाध्यः}}\Afootnote{\skt{प्यसाध्यः}\ N H}}\skt{ }\edtext{\edtext{\skt{स्याद्य\-श्चो\-र्द्ध\-मु\-दर्वे}}{\lemma{\skt{चोर्द्धम्}}\Afootnote{H; \skt{चोत्व}\ N}}}{\lemma{\skt{उदर्वे}}\Afootnote{H; \skt{सुदुरै}\ N}}\skt{ तथा ॥}\pend

\pstart\versepar \ledsidenote{\passref{2.9.24}}\edtext{\skt{नाभेरूपरिजाः}}{\lemma{\skt{ऊपरिजाः}}\Afootnote{H; \skt{उपरिजाः}\ A N}}\skt{ पक्वा }\edtext{\skt{यान्त्युर्ध्वमितरे}}{\lemma{\skt{उर्ध्वम्}}\Afootnote{H; \skt{ऊर्ध्वम्}\ A N}}\skt{ त्वधः । }\pend
\pstart\versepar \skt{}\edtext{\skt{अधः}}{\Afootnote{N H; \skt{जीवत्यधो}\ A}}\skt{ }\edtext{\skt{सृतेषु}}{\Afootnote{H; \skt{निःस्रुतेषु}\ A; \skt{स्रुतेषु}\ N}}\skt{ }\edtext{\skt{जीवेत्तु}}{\Afootnote{H; \emph{om.}\ A; \skt{जीवन्तु}\ N}}\skt{ }\edtext{\skt{सृतेषूर्ध्वन्न}}{\Afootnote{H; \skt{स्रुतेषूर्ध्वं न}\ A; \skt{स्रुतेशूर्ध्वन्न}\ N}}\skt{ जीवति ॥}\pend

\pstart\versepar \ledsidenote{\passref{2.9.25}}\edtext{\skt{हृद्ब\-स्ति\-ना\-भि\-वर्ज्या ये}}{\Afootnote{N H; \skt{हृन्ना\-भि\-ब\-स्ति\-वर्ज्या ये}\ A}}\skt{ तेषु भिन्नेषु बाह्यतः । }\pend
\pstart\versepar \skt{}\edtext{\skt{जीवे\-त्क\-दा\-चि\-त्पु\-रुषो}}{\Afootnote{A H; \skt{जीवेत्कदाचि पुरुषो}\ N}}\skt{ नेतरेषु कदाचन ॥}\pend

\pstart\versepar \ledsidenote{\passref{2.9.26}}\edtext{\skt{स्त्रीणा\-म\-व\-प्र\-जा\-ता\-ना\-म्प्र\-जा\-ता\-ना\-न्तथा}}{\lemma{\skt{अव\-प्र\-जा\-ता\-नाम्}}\Afootnote{N H; \skt{अपप्रजातानां}\ A}}\skt{ }\edtext{\skt{हितैः}}{\Afootnote{N H; \skt{ऽहितैः ।}\ A}}\skt{ । }\pend
\pstart\versepar \skt{दाहज्वरकरो घोरो जायते रक्तविद्रधिः ॥}\pend

\pstart\versepar \ledsidenote{\passref{2.9.27}}\edtext{\edtext{\skt{स्त्रीणाम्मिथ्यो}}{\lemma{\skt{स्त्रीणाम्}}\Afootnote{N H; \skt{अ}\ A}}}{\lemma{\skt{मिथ्यो}}\Afootnote{H; \skt{पि स}\ A; \skt{मिथ्या}\ N}}\skt{ }\edtext{\edtext{\edtext{\skt{प्रजा\-ता\-ना\-म\-सृ\-क्क\-प्या\-दनिः}}{\lemma{\skt{प्रजातानाम्}}\Afootnote{H; \skt{म्यक्प्रजातानाम्}\ A; \skt{प्रजातानाम्}\ N}}}{\lemma{\skt{असृक्कप्य्}}\Afootnote{H; \skt{असृक्का}\ A; \skt{मसृक्का}\ N}}}{\lemma{\skt{आदनिः}}\Afootnote{H; \skt{यादनिः}\ A; \skt{यादनि}\ N}}\skt{ }\edtext{\skt{सृतां}}{\Afootnote{H; \skt{सृतम्}\ A N}}\skt{ । }\pend
\pstart\versepar \skt{}\edtext{\edtext{\skt{रक्त\-ज\-म्वि\-द्र\-धि\-ङ्कु\-र्या\-त्कुक्षौ}}{\lemma{\skt{रक्तजम्}}\Afootnote{A H; \skt{रक्तजां}\ N}}}{\lemma{\skt{कुर्यात्कुक्षौ}}\Afootnote{N H; \skt{कुर्यात्कुक्सौ}\ A}}\skt{ }\edtext{\skt{मक्कलसंज्ञितं}}{\Afootnote{N H; \skt{मक्क\-ल्ल\-सं\-ज्ञि\-तम् ॥}\ A}}\skt{ ।}\pend

\pstart\versepar \ledsidenote{\passref{2.9.28}}\edtext{\edtext{\skt{सप्ता\-हा\-त्नो\-प\-शा\-न्त\-श्चे\-त्ततः}}{\lemma{\skt{सप्ता\-हा\-त्नो\-प\-शा\-न्तश्}}\Afootnote{H; \skt{सप्ता\-हा\-न्नो\-प\-शा\-न्तश्}\ A; \skt{शप्ता\-हा\-न्नो\-प\-शा\-न्तश्}\ N}}}{\lemma{\skt{चेत्ततः}}\Afootnote{A H; \skt{चे}\ N}}\skt{ }\edtext{\skt{सम्प्रतिपच्यते}}{\Afootnote{H; \skt{ऽसौ संप्रपच्यते ।}\ A; \skt{त्त्तस्स\-म्प्र\-ति\-प\-द्यते ।}\ N}}\skt{ ॥ }\pend
\pstart\versepar \skt{}\edtext{\edtext{\skt{विशेषमथ}}{\lemma{\skt{विशेषम्}}\Afootnote{A N; \skt{विशेष}\ H}}}{\lemma{\skt{अथ}}\Afootnote{A N; \skt{मथ}\ H}}\skt{ वक्ष्यामि स्पष्ट\-म्वि\-द्र\-धि\-गु\-ल्मयोः ॥}\pend

\pstart \ledsidenote{\passref{2.9.29cd}}\edtext{\skt{कस्मात्न}}{\Afootnote{H; \skt{कस्मान्न}\ A N}}\skt{ पच्यते गुल्मो विद्रधिः पाकमेति च ।}\pend

\pstart \ledsidenote{\passref{2.9.30ab}}\skt{न निबन्धोऽस्ति गुल्मानाम्विद्रधिः }\edtext{\skt{सनिबन्धनः}}{\Afootnote{A H; \skt{सनिबन्धनं ॥}\ N}}\skt{ ॥}\pend

\pstart\versepar \ledsidenote{\passref{2.9.31}}\skt{विवरानुचरो }\edtext{\edtext{\skt{गुल्मम्त्वप्सु}}{\lemma{\skt{गुल्मम्}}\Afootnote{H; \skt{ग्रन्थिर्}\ A; \skt{गुल्मै ह्य्}\ N}}}{\lemma{\skt{त्वप्सु}}\Afootnote{H; \skt{अप्सु}\ A N}}\skt{ }\edtext{\skt{बुद्बुदकोपमः}}{\Afootnote{N H; \skt{बुद्बुदको यथा ।}\ A}}\skt{ । }\pend
\pstart\versepar \skt{एवम्प्रकारो गुल्मस्तु तस्मात्पाकन्न गच्छति ॥}\pend

\pstart \ledsidenote{\passref{2.9.32cd}}\edtext{\skt{माम्श\-शो\-णि\-त\-ही\-नस्य}}{\Afootnote{N; \skt{मांस\-शो\-णि\-त\-ही\-नत्वा}\ A}}\skt{ }\edtext{\skt{गुल्मो न}}{\Afootnote{\skt{द्गुल्मः पा}\ A; \skt{गुल्म्लो न}\ N}}\skt{ }\edtext{\skt{परिपच्यते}}{\Afootnote{N; \skt{कं न गच्छति ॥}\ A}}\skt{ ।}\pend

\pstart \ledsidenote{\passref{2.9.32ab}}\edtext{\skt{मांस\-शो\-णि\-त\-बा\-हु\-ल्या\-त्पा\-क\-ङ्ग\-च्छति}}{\lemma{\skt{मांस\-शो\-णि\-त\-बा\-हु\-ल्या\-त्पा\-कङ्}}\Afootnote{A N; \skt{मांस\-शो\-णि\-त\-वा\-हु\-ल्या\-त्पा\-कङ्}\ H}}\skt{ विद्रधिः ।}\pend

\pstart\versepar \ledsidenote{\passref{2.9.33}}\skt{गुल्मस्तिष्ठति }\edtext{\skt{दोषैः }\edtext{\edtext{\skt{स्वैर्व्वि\-द्र\-धि\-र्म्मा\-न्स\-शो\-णिते}}{\lemma{\skt{व्विद्रधिर्}}\Afootnote{A N H\textsuperscript{pc}; \skt{व्विद्रधिर्}\ H\textsuperscript{ac}}}}{\lemma{\skt{म्मान्सशोणिते}}\Afootnote{A H; \skt{माम्शसोणिते ।}\ N}}}{\lemma{\skt{दोषैः स्वैर्}}\Afootnote{N H; \skt{दोषे स्वे}\ A}}\skt{ ॥ }\pend
\pstart\versepar \skt{}\edtext{\skt{विद्रधिः}}{\Afootnote{A H; \skt{विद्रधिं}\ N}}\skt{ पच्यते तस्माद्गुल्मश्चापि न पच्यते ।}\pend

\pstart \ledsidenote{\passref{2.9.34ab}}\skt{हृन्नाभिबस्तिजः पक्वो वर्ज्यो यश्च त्रिदोषजः ॥}\pend

\pstart\versepar \ledsidenote{\passref{2.9.34.add}}\skt{विद्रधेः पूर्व्वरूपाणि }\edtext{\skt{विण्मू\-त्रा\-नि\-ल\-स\-ङ्ग्रहः}}{\Afootnote{H; \skt{विद्सू\-त्रा\-नि\-ल\-स\-द्ग्रहः ॥}\ N}}\skt{ । }\pend
\pstart\versepar \skt{}\edtext{\edtext{\skt{भ्रमो\-द्ग\-म\-र्द्दो\-वै\-रस्यं}}{\lemma{\skt{भ्रमो\-द्ग\-म\-र्द्दो\-वैर्}}\Afootnote{\skt{भ्रमो\-द्ग\-म\-र्दो\-वेर्}\ N; \skt{भ्रमोऽ\-द्ग\-म\-र्द्दो\-वैर्}\ H}}}{\lemma{\skt{अस्यं}}\Afootnote{H; \skt{अस्य}\ N}}\skt{ }\edtext{\skt{काण्ट\-व्य\-म्म\-धु\-चा\-स्यता}}{\Afootnote{H; \skt{काम्मधुरस्यता ।}\ N}}\skt{ ॥ }\pend
\pstart\versepar \skt{}\edtext{\edtext{\edtext{\skt{मुहु\-र्म्मु\-हु\-स्त\-था\-त्यर्थ}}{\lemma{\skt{मुहुर्}}\Afootnote{N; \skt{मुहु}\ H}}}{\lemma{\skt{म्मुहुस्}}\Afootnote{N; \skt{र्म्मुहुस्}\ H}}}{\lemma{\skt{तथात्यर्थ}}\Afootnote{H; \skt{तथात्यथ}\ N}}\skt{ यथा }\edtext{\edtext{\skt{स्थानम्भवेद्रुजा}}{\lemma{\skt{स्थानम्}}\Afootnote{H; \skt{स्ठानम्}\ N}}}{\lemma{\skt{भवेद्रुजा}}\Afootnote{N; \skt{उवेद्रुजा ।}\ H}}\skt{ । }\pend
\pstart\versepar \skt{विवर्द्धते ज्वरश्चास्य तृष्णा चास्य प्रवाधत इति ॥}\pend

\pstart \ledsidenote{\passref{2.9.trailer}}\edtext{\skt{॥ ॥वि\-द्र\-धि\-नि\-दाने}}{\Afootnote{H; \skt{इति सुश्रु\-त\-सं\-हि\-तायां निदानस्थाने विद्रधिनिदानं नाम}\ A; \skt{९}\ N}}\skt{ }\edtext{\skt{नवमः}}{\Afootnote{H; \skt{नवमोऽध्यायः ॥ ९ ॥}\ A; \emph{om.}\ N}}\skt{ ॥}\pend

\endnumbering
\beginnumbering
\eledsection*{Adhyāya 10}
\pstart \ledsidenote{\passref{2.10.1}}\skt{अथातो }\edtext{\edtext{\skt{विस\-र्प्प\-ना\-डी\-स्त\-न\-रो\-गा\-णा\-न्नि\-दा\-न\-म्व्या\-ख्या\-स्यामः}}{\lemma{\skt{विस\-र्प्प\-ना\-डी\-स्त\-न\-रो\-गा\-णा\-न्नि\-दा\-नम्}}\Afootnote{N H; \skt{विस\-र्प\-ना\-डी\-स्त\-न\-रो\-ग\-नि\-दानं}\ A}}}{\lemma{\skt{व्याख्यास्यामः}}\Afootnote{A H; \skt{व्या ॥}\ N}}\skt{ ॥}\pend

\pstart\versepar \ledsidenote{\passref{2.10.3}}\edtext{\skt{त्वग्मां\-स\-शो\-णि\-त\-गताः}}{\Afootnote{H; \skt{त्वङ्मां\-स\-शो\-णि\-त\-गताः}\ A N}}\skt{ }\edtext{\skt{कुपिताश्च}}{\Afootnote{N H; \skt{कुपितास्तु}\ A}}\skt{ दोषाः }\edtext{\edtext{\skt{सर्व्वा\-ङ्ग\-चा\-रि\-ण\-मिह}}{\lemma{\skt{सर्व्वा\-ङ्ग\-चा\-रि\-णम्}}\Afootnote{N H; \skt{सर्वा\-ङ्ग\-सा\-रि\-णम्}\ A}}}{\lemma{\skt{इह}}\Afootnote{N H; \skt{इहा}\ A}}\skt{ स्थितमात्मलिङ्गं । }\pend
\pstart\versepar \skt{कुर्व्वन्ति }\edtext{\skt{विस्तृ\-त\-म\-नु\-न्न\-त\-माशु}}{\lemma{\skt{विस्तृ\-त\-म\-नु\-न्न\-तम्}}\Afootnote{H; \skt{विस्तृ\-त\-म\-नु\-न्न\-तम्}\ A N}}\skt{ }\edtext{\skt{शोफन्तं}}{\Afootnote{A H; \skt{शोफं त}\ N}}\skt{ सर्व्वतो }\edtext{\skt{विसरणात्तु}}{\Afootnote{H; \skt{विसरणाच्च}\ A; \skt{विसरञा तु}\ N}}\skt{ }\edtext{\skt{विसर्प्पमाहुः}}{\lemma{\skt{आहुः}}\Afootnote{A N\textsuperscript{pc} H; \skt{आमहुः ।}\ N\textsuperscript{ac}}}\skt{ ॥}\pend

\pstart \ledsidenote{\passref{2.10.4}}\edtext{\skt{वातात्मको}}{\Afootnote{A H; \skt{वातात्मको}\ N}}\skt{ }\edtext{\skt{सितमृदुः}}{\Afootnote{A H; \skt{सितमृदुः}\ N}}\skt{ }\edtext{\edtext{\skt{परुषोऽङ्गमर्द्दः}}{\lemma{\skt{परुषो}}\Afootnote{A H; \skt{परुषोङ्}\ N}}}{\lemma{\skt{ऽङ्गमर्द्दः}}\Afootnote{H; \skt{ऽङ्गमर्द}\ A; \skt{गमर्द्द}\ N}}\skt{ }\edtext{\skt{संभे\-द\-तो\-द\-प\-व\-न\-ज्व\-र\-लि\-ङ्ग\-युक्तः}}{\Afootnote{N; \skt{संभे\-द\-तो\-द\-प\-व\-न\-ज्व\-र\-लि\-ङ्ग\-युक्तः । गण्डैर्यदा तु विष\-मै\-र\-ति\-दू\-षि\-त\-त्वा\-द्युक्तः स एव कथितः खलु वर्जनीयः ॥}\ A; \skt{संभे\-द\-तो\-द\-प\-व\-न\-ज्व\-र\-लि\-ङ्ग\-युक्तः । …}\ H}}\skt{ ।}\pend

\pstart \ledsidenote{\passref{2.10.5}}\skt{पित्तात्मको- }\edtext{\edtext{\skt{द्रूत\-ग\-ति\-र्ज्व\-र\-दा\-ह\-पाकः}}{\lemma{\skt{द्रूतगतिर्}}\Afootnote{H; \skt{द्रुतगति}\ A N}}}{\lemma{\skt{ज्वरदाहपाकः}}\Afootnote{H; \skt{र्ज्वरदाहपाक}\ A; \skt{र्ज्वरदाहपाकः}\ N}}\skt{ }\edtext{\skt{स्फोट\-प्र\-भे\-द\-ब\-हुल}}{\Afootnote{H; \skt{स्फोट\-प्र\-भे\-द\-ब\-हुलः}\ A; \skt{स्फोट\-प्र\-भे\-द\-व\-हुल}\ N}}\skt{ क्षतजप्रकाशः ॥ }\edtext{\edtext{\skt{सद्यः\-क्ष\-ता\-न्वि\-त\-मि\-वापि}}{\lemma{\skt{सद्यः\-क्ष\-ता\-न्वि\-तम्}}\Afootnote{H; \skt{दोष\-प्र\-वृ\-द्धि\-ह\-त\-मांस}\ A; \skt{सद्यः कृतान्वितम्}\ N}}}{\lemma{\skt{इवापि}}\Afootnote{N H; \emph{om.}\ A}}\skt{ }\edtext{\edtext{\skt{हि-तम्विहाति}}{\lemma{\skt{हि-तम्}}\Afootnote{N H; \skt{सिरो}\ A}}}{\lemma{\skt{विहाति}}\Afootnote{H; \skt{यदा स्या}\ A; \skt{विभाति ।}\ N}}\skt{ }\edtext{\skt{स्रोतो\-ज\-क\-र्द्द\-म\-व\-पुर्न्न}}{\Afootnote{H; \skt{त्स्रोतो\-ज\-क\-र्द\-म\-निभो न}\ A; \skt{स्रोतो\-ज\-क\-र्द\-म\-व\-पुर्ना}\ N}}\skt{ तदा स सिध्येत् ।}\pend

\pstart \ledsidenote{\passref{2.10.6}}\skt{श्लेष्मा-त्मकः सरति }\edtext{\skt{मन्द\-म\-शी\-घ्र\-पाकी}}{\Afootnote{H; \skt{मन्द\-म\-शी\-घ्र\-पाकः}\ A; \skt{मन्द\-म\-शी\-घ्र\-पाकी}\ N}}\skt{ स्निग्धः }\edtext{\skt{स्थिर\-प्र\-चु\-र\-क\-ण्ड\-र\-न\-न्य\-वर्ण्णः}}{\Afootnote{H; \skt{सित\-श्व\-य\-थु\-र\-ल्प\-रु\-गु\-ग्र\-हण्डुः ।}\ A; \skt{स्थिरः प्रचु\-र\-क\-ण्डु\-र\-न\-न्य\-वर्ण्णः ।}\ N}}\skt{ ॥ }\edtext{\edtext{\skt{सर्व्वा\-त्म\-क\-स्त्रि\-वि\-ध\-व\-र्ण्ण\-रु\-जा\-व\-गाढः}}{\lemma{\skt{सर्व्वात्मकस्}}\Afootnote{H; \skt{सर्वात्मक}\ A; \skt{सर्वात्कमस्}\ N}}}{\lemma{\skt{त्रिवि\-ध\-व\-र्ण्ण\-रु\-जा\-व\-गाढः}}\Afootnote{N H; \skt{स्त्रिवि\-ध\-व\-र्ण\-रु\-जोऽ\-व\-गाढः}\ A}}\skt{ पक्वो न }\edtext{\skt{सिध्यति स}}{\Afootnote{N H; \skt{सिध्यति च}\ A}}\skt{ }\edtext{\skt{मांस\-सि\-रा\-प्र\-णा\-शात्}}{\Afootnote{H; \skt{मांस\-सि\-रा\-प्र\-शा\-तात् ॥}\ A; \skt{मंशिइरा प्रणाशात् ॥}\ N\textsuperscript{ac}; \skt{मंशिरा प्रणाशात् ॥}\ N\textsuperscript{pc}}}\skt{ ।}\pend

\pstart \ledsidenote{\passref{2.10.7}}\edtext{\skt{सद्यः\-क्ष\-त\-म्व्र\-ण\-मु\-पेत्य}}{\lemma{\skt{सद्यः\-क्ष\-त\-म्व्र\-णम्}}\Afootnote{H; \skt{सद्यः\-क्ष\-त\-व्र\-णम्}\ A; \skt{सद्यः कृतव्रणम्}\ N}}\skt{ नरस्य }\edtext{\skt{रक्तं}}{\Afootnote{N H; \skt{पित्तं}\ A}}\skt{ }\edtext{\skt{पित्तञ्च}}{\Afootnote{H; \skt{रक्तं च}\ A; \skt{पतित्तञ्च}\ N}}\skt{ }\edtext{\skt{दोषबहुलस्य}}{\Afootnote{A H; \skt{दोषवहुलस्य}\ N}}\skt{ करोति }\edtext{\skt{शोफम्}}{\Afootnote{N H; \skt{शोहम् ।}\ A}}\skt{ ॥ श्यावं }\edtext{\edtext{\skt{सलो-\-हि\-त\-म\-ति\-ज्व\-र\-दा\-ह\-पाकः}}{\lemma{\skt{सलो-हितम्}}\Afootnote{H; \skt{सलोहित}\ A N}}}{\lemma{\skt{अति\-ज्व\-र\-दा\-ह\-पाकः}}\Afootnote{H; \skt{मति\-ज्व\-र\-दा\-ह\-पाकं}\ A N}}\skt{ स्फोटैः कुल\-त्थ\-स\-दृ\-शै\-र\-सि\-तैश्च कीर्ण्णम् ॥}\pend

\pstart \ledsidenote{\passref{2.10.8}}\skt{सिद्ध्यन्ति }\edtext{\skt{वात\-क\-फ\-पि\-त्त\-कृता}}{\Afootnote{A H; \skt{वात\-क\-फ\-पि\-त्त\-कृतो}\ N}}\skt{ विसर्प्पाः सर्व्वात्मकः }\edtext{\skt{क्षतकृतश्च न}}{\Afootnote{A H; \skt{क्षतकृतश्च नि}\ N}}\skt{ सिद्धिमेति ॥ }\edtext{\skt{पित्ता\-त्म\-को\-ञ्ज\-न\-व\-पुश्च}}{\Afootnote{H; \skt{पैत्तानिलावपि च}\ A; \skt{पित्ता\-त्मा\-कं\-ज\-न\-व\-पुश्च}\ N}}\skt{ }\edtext{\skt{यदा}}{\Afootnote{N H; \skt{दर्शि}\ A}}\skt{ }\edtext{\skt{तदास्यात्कृच्छ्राच्च}}{\Afootnote{H; \skt{तपूर्वलिङ्गौ सर्वे च}\ A; \skt{तदा स्यात्कृच्छ्रश्च}\ N}}\skt{ मर्म्मसु भवन्ति हि }\edtext{\skt{सव्व}}{\Afootnote{H; \emph{om.}\ A; \skt{सर्व्व}\ N}}\skt{ }\edtext{\skt{एव}}{\Afootnote{N H; \skt{कृच्छ्रसाध्याः ॥}\ A}}\skt{ ।}\pend

\pstart\versepar \ledsidenote{\passref{2.10.9}}\edtext{\skt{यः }\edtext{\skt{शोफमीक्षति}}{\lemma{\skt{ईक्षति}}\Afootnote{N H; \emph{om.}\ A}}}{\lemma{\skt{यः शोफम्}}\Afootnote{N H; \skt{शोफं}\ A}}\skt{ }\edtext{\edtext{\skt{सुप\-क्क्व\-म\-प\-क्क्व\-सञ्जी}}{\lemma{\skt{सुपक्क्वम्}}\Afootnote{H; \skt{न पक्वम्}\ A; \skt{सुपक्वम्}\ N}}\skt{ यो वा}}{\lemma{\skt{अपक्क्वसञ्जी यो वा}}\Afootnote{H; \skt{इति पक्वमुपेक्षते यो यो वा}\ A; \skt{अपक्वसंज्ञी यो वा}\ N}}\skt{ व्रणं }\edtext{\edtext{\skt{प्रचु\-र\-पू\-य\-म\-सा\-धु\-वृत्तः}}{\lemma{\skt{प्रचुरपूयम्}}\Afootnote{N H; \skt{प्रचुरपूय}\ A}}}{\lemma{\skt{असाधुवृत्तः}}\Afootnote{H; \skt{मसाधुवृत्तः ।}\ A; \skt{असाधुवृत्तं ॥}\ N}}\skt{ । }\pend
\pstart\versepar \skt{अभ्यन्तरं प्रविशति प्रविदार्य तस्य स्थानानि पूर्व्वविहितानि ततः स पूयः ।}\pend

\pstart \ledsidenote{\passref{2.10.10}}\edtext{\edtext{\edtext{\skt{तस्या\-ति\-मा\-त्र\-ग\-म\-ना\-द्ग\-ति\-रि\-त्य\-तस्तु}}{\lemma{\skt{तस्या\-ति\-मा\-त्र\-ग\-म\-ना\-द्ग\-तिर्}}\Afootnote{A H; \skt{तस्या\-ति\-मा\-त्र\-ग\-म\-ना\-द्गति}\ N}}}{\lemma{\skt{इत्य्}}\Afootnote{A H; \skt{रित्य्}\ N}}}{\lemma{\skt{अतस्तु}}\Afootnote{N H\textsuperscript{pc}; \skt{अतश्च}\ A; \skt{अस्तु}\ H\textsuperscript{ac}}}\skt{ }\edtext{\skt{नाडी च}}{\Afootnote{H; \skt{नाडीव}\ A N}}\skt{ }\edtext{\skt{यद्वहति}}{\Afootnote{N H; \skt{यद्वहति}\ A}}\skt{ }\edtext{\skt{तेन}}{\Afootnote{A H; \skt{तेनु}\ N}}\skt{ }\edtext{\skt{मता तु}}{\Afootnote{A H; \skt{मतानु}\ N}}\skt{ नाडी ॥ }\edtext{\edtext{\skt{दोषै\-स्त्रि\-बि\-र्ब्भ\-वति}}{\lemma{\skt{त्रिबिर्}}\Afootnote{A H; \skt{त्रिभिर्}\ N}}\skt{ या}}{\lemma{\skt{ब्भवति या}}\Afootnote{\skt{भवति सा}\ A N; \skt{भ्भवति या}\ H}}\skt{ पृथगेकशश्च संमूर्च्छितैरपि च }\edtext{\skt{शल्यनिमित्ततो}}{\Afootnote{A H; \skt{सल्यनिमित्ततो}\ N}}\skt{ }\edtext{\skt{न्या}}{\Afootnote{N H; \skt{ऽन्या ॥}\ A}}\skt{ ।}\pend

\pstart \ledsidenote{\passref{2.10.11}}\edtext{\skt{तत्रा\-नि\-ला\-त्प\-रु\-ष\-सू\-क्ष्म\-मुखी}}{\Afootnote{A H; \skt{तत्रा\-नि\-ला\-त्पा\-रु\-ष\-सू\-क्ष्मुखी}\ N}}\skt{ सशूला फेना\-नु\-वि\-द्ध\-म\-धिकं स्रवति }\edtext{\skt{क्षपासु}}{\Afootnote{\skt{क्षपायाम् । तृट्ता\-प\-तो\-द\-स\-द\-न\-ज्व\-र\-भे\-द\-हेतुः पीतं स्रव\-त्य\-धि\-क\-मु\-ष्ण\-म\-हःसु पित्तात् ॥}\ A; \skt{क्षपाम् । पित्ता तु तृट्ज्वरकरी परिदाहयुक्ता पीतं स्रव\-त्य\-धि\-क\-मु\-ष्ण\-महः सुचापि ॥}\ N; \skt{क्षपासु ॥ …}\ H}}\skt{ ॥}\pend

\pstart \ledsidenote{\passref{2.10.11a}}\skt{पित्तात्तु तृद्करकरी परिदाहयुक्ता पीतं स्रव\-त्य\-धि\-क\-मु\-ष्ण\-म\-हस्सु चापि ।}\pend

\pstart \ledsidenote{\passref{2.10.12}}\skt{ज्ञेया }\edtext{\skt{कफा\-द्ब\-हु\-घ\-ना\-र्ज्जु\-न\-पि\-च्छि\-लास्रा}}{\Afootnote{A H; \skt{कफा\-द्व\-हु\-घ\-ना\-र्ज\-न\-पि\-च्छि\-लास्रा}\ N}}\skt{ }\edtext{\skt{स्तस\-क\-ण्डु\-र\-रु\-जा\-र\-जनी प्रवृद्धाः}}{\Afootnote{\skt{रात्रिस्रुतिः स्तिमि\-त\-रु\-क्क\-ठिना सकण्डूः । दोष\-द्व\-या\-भि\-हि\-त\-ल\-क्ष\-ण\-द\-र्श\-नेन तिस्रो गती\-र्व्य\-ति\-क\-र\-प्र\-भ\-वास्तु विद्यात् ॥}\ A; \skt{स्तव्धाः सकण्डुररुजा रजनीप्रवृद्धाः ।}\ N; \skt{स्त सक\-ण्डु\-र\-रु-\-जा\-र\-जनी प्रवृद्धाः ॥ … ॥}\ H}}\skt{ ॥ ॥}\pend

\pstart \ledsidenote{\passref{2.10.13}}\edtext{\skt{दाहज्वरश्च}}{\Afootnote{H; \skt{दाहज्वरश्व}\ A N}}\skt{ सन\-मू\-र्च्छ\-न\-व\-क्त्र\-शोषा }\edtext{\skt{यस्या}}{\Afootnote{N H; \skt{यस्यां}\ A}}\skt{ }\edtext{\edtext{\skt{भव\-न्त्य\-भि\-हि\-तानि}}{\lemma{\skt{भवन्त्य्}}\Afootnote{H; \skt{भवन्}\ A; \skt{भव}\ N}}\skt{ च}}{\lemma{\skt{अभिहितानि च}}\Afootnote{H; \skt{त्यभिहितानि च}\ A N}}\skt{ लक्षणानि । }\edtext{\skt{तामा\-दि\-शे\-त्प\-व\-न\-पि\-त्त\-क\-फ\-प्र\-को\-पा\-द्घो\-रा\-ङ्ग\-ति\-त्त्व\-सु\-ह-\-रा\-मिव}}{\lemma{\skt{आदि\-शे\-त्प\-व\-न\-पि\-त्त\-क\-फ\-प्र\-को\-पा\-द्घो\-रा\-ङ्ग\-ति\-त्त्व\-सु\-ह-\-राम्}}\Afootnote{H; \skt{आदि\-शे\-त्प\-व\-न\-पि\-त्त\-क\-फ\-प्र\-को\-पा\-द्घो\-रा\-मसु क्षयकरीम्}\ A; \skt{आदि\-शे\-त्प\-व\-न\-पि\-त्त\-क\-फ\-प्र\-को\-पा\-द्घोरां गतिं त्वसुहराम्}\ N}}\skt{ कालरात्रिम् ॥}\pend

\pstart \ledsidenote{\passref{2.10.14}}\edtext{\edtext{\skt{नष्ट\-ङ्क\-थ\-ञ्चि\-द\-णु\-मा\-र्ग्ग\-मु\-दी\-रि\-तेषु}}{\lemma{\skt{कथ\-ञ्चि\-द\-णु\-मा\-र्ग्गम्}}\Afootnote{H; \skt{कथं\-चि\-द\-नु\-मार्ग}\ A; \skt{कथ\-ञ्चि\-द\-नु\-मा\-र्गम्}\ N}}}{\lemma{\skt{उदीरितेषु}}\Afootnote{N H; \skt{मुदीरितेषु}\ A}}\skt{ }\edtext{\skt{स्थानेषु}}{\Afootnote{A H; \skt{स्थाने}\ N}}\skt{ }\edtext{\edtext{\skt{शल्यमचिरेण}}{\lemma{\skt{शल्यम्}}\Afootnote{N H; \skt{शल्य}\ A}}}{\lemma{\skt{अचिरेण}}\Afootnote{N H; \skt{मचिरेण}\ A}}\skt{ गतिङ्करोति । सा फेनिलं }\edtext{\edtext{\edtext{\skt{मथि\-त\-म\-च्छ\-म\-सृ\-ग्वि\-मि\-श्र\-मुष्णं}}{\lemma{\skt{मथितमच्छम्}}\Afootnote{H; \skt{मथितमच्छ}\ A N}}}{\lemma{\skt{असृग्विमिश्रम्}}\Afootnote{H; \skt{मसृग्विमिश्र}\ A; \skt{मसृग्विनिश्र ।}\ N}}}{\lemma{\skt{उष्णं}}\Afootnote{H; \skt{मुष्णं}\ A; \skt{मुष्ण}\ N}}\skt{ }\edtext{\skt{करोति}}{\Afootnote{N H; \skt{स्रवेत}\ A}}\skt{ सहसा }\edtext{\skt{सु सरुजञ्च}}{\Afootnote{H; \skt{सरुजा च}\ A; \skt{सरुजश्च}\ N}}\skt{ नित्यम् ॥}\pend

\pstart\versepar \ledsidenote{\passref{2.10.15}}\edtext{\skt{यावन्त्यो}}{\Afootnote{N H; \skt{यावत्यो}\ A}}\skt{ }\edtext{\skt{गतयश्चैव}}{\lemma{\skt{चैव}}\Afootnote{N H; \skt{यैश्च}\ A}}\skt{ कारणैः }\edtext{\skt{संभवन्ति हि}}{\Afootnote{A H; \skt{सम्भवन्ति हिः ।}\ N}}\skt{ । }\pend
\pstart\versepar \skt{}\edtext{\skt{तावन्त्यः}}{\Afootnote{H; \skt{तावन्तः}\ A; \skt{तावन्त्य}\ N}}\skt{ }\edtext{\skt{स्तनरोगास्तु}}{\Afootnote{H; \skt{स्तनरोगाः स्युः}\ A; \skt{स्तनरो}\ N}}\skt{ }\edtext{\edtext{\skt{तैरेव}}{\lemma{\skt{तैर्}}\Afootnote{H; \skt{स्त्रीणां तैर्}\ A; \skt{गास्तु तैर्}\ N}}\skt{ च}}{\lemma{\skt{एव च}}\Afootnote{N H; \skt{एव}\ A}}\skt{ }\edtext{\skt{भवन्ति ह}}{\Afootnote{N H; \skt{हेतुभिः ॥}\ A}}\skt{ ॥}\pend

\pstart\versepar \ledsidenote{\passref{2.10.16}}\skt{धम\-न्य\-स्सं\-वृ\-त\-द्वाराः कन्यानां स्तनसंश्रिताः । }\pend
\pstart\versepar \skt{}\edtext{\skt{दोषा-\-ग\-ति\-त्वा\-त्तासां}}{\Afootnote{N H; \skt{दोषा\-वि\-स\-र\-णा\-त्तासां न}\ A}}\skt{ }\edtext{\skt{हि}}{\Afootnote{N H; \skt{भवन्ति}\ A}}\skt{ }\edtext{\skt{स्तन\-रो\-गा\-न\-स\-न्त्यतः}}{\Afootnote{N H; \skt{स्तनामयाः ॥}\ A}}\skt{ ॥}\pend

\pstart\versepar \ledsidenote{\passref{2.10.17}}\skt{तासामेव प्रजा\-ता\-ना\-ङ्ग\-र्ब्भि\-णी\-नाञ्च ताः पुनः । }\pend
\pstart\versepar \skt{स्वभावादेव विवृता जायन्ते }\edtext{\skt{संभतः}}{\Afootnote{H; \skt{संभवन्त्यतः}\ A N}}\skt{ ॥}\pend

\pstart\versepar \ledsidenote{\passref{2.10.18}}\skt{रसप्रसादो मधुरः पक्वा\-हा\-र\-नि\-मि\-त्तजः । }\pend
\pstart\versepar \skt{}\edtext{\skt{कृत्स्नाद्देहात्स्तनौ}}{\Afootnote{H; \skt{कृत्स्नदेहात्स्तनौ}\ A; \skt{कृत्स्नाद्देहात्स्तनौ}\ N}}\skt{ }\edtext{\skt{प्राप्य}}{\Afootnote{N H; \skt{प्राप्तः}\ A}}\skt{ स्तन्य\-मि\-त्य\-भि\-धी\-यते ॥}\pend

\pstart \ledsidenote{\passref{2.10.19ab}}\edtext{\skt{विषशस्तेष्वपि}}{\lemma{\skt{विषशस्तेष्व्}}\Afootnote{H\textsuperscript{pc}; \skt{विषस्तेष्व्}\ H\textsuperscript{ac}}}\skt{ गात्रेषु यथा शुक्रन्न दृश्यते ।}\pend

\pstart\versepar \ledsidenote{\passref{2.10.20}}\skt{तदेव }\edtext{\skt{चेष्टायुवतिं}}{\Afootnote{H; \skt{चेष्टयुवते}\ A; \skt{चेष्टायुवति}\ N}}\skt{ }\edtext{\skt{दर्श\-ना\-त्स्म\-र\-णा\-दपि}}{\Afootnote{N H; \skt{र्दर्श\-ना\-त्स्म\-र\-णा\-दपि ।}\ A}}\skt{ ॥ }\pend
\pstart\versepar \skt{}\edtext{\skt{संश\-ब्दा\-त्स्प\-र्श\-ना\-द्वापि}}{\Afootnote{N H; \skt{शब्द\-सं\-श्र\-व\-णा\-त्स्पर्शा}\ A}}\skt{ }\edtext{\skt{संहर्षाच्च}}{\Afootnote{N H; \skt{त्संहर्षाच्च}\ A}}\skt{ प्रवर्त्तते ।}\pend

\pstart\versepar \ledsidenote{\passref{2.10.21}}\edtext{\edtext{\skt{सुप्र\-स\-न्न\-म\-न\-श्चात्र}}{\lemma{\skt{सुप्र\-स\-न्न\-म\-नश्}}\Afootnote{N H; \skt{सुप्रसन्नं मनस्}\ A}}}{\lemma{\skt{चात्र}}\Afootnote{N H; \skt{तत्र}\ A}}\skt{ }\edtext{\skt{दर्शने}}{\Afootnote{N H; \skt{हर्षणे}\ A}}\skt{ हेतुरुच्यते ॥ }\pend
\pstart\versepar \skt{}\edtext{\skt{आहा\-र\-र\-स\-वी\-र्य\-त्वा\-देवं}}{\Afootnote{N H; \skt{आहा\-र\-र\-स\-यो\-नि\-त्वा\-देवं}\ A}}\skt{ स्तन्यमपि }\edtext{\skt{स्त्रियः}}{\Afootnote{H; \skt{स्त्रियाः}\ A N}}\skt{ ।}\pend

\pstart \ledsidenote{\passref{2.10.19cd}}\skt{सर्व्व\-दे\-हा\-श्रि\-त\-त्वाच्च शुक्र\-ल\-क्ष\-ण\-मु\-च्यते ॥}\pend

\pstart\versepar \ledsidenote{\passref{2.10.22}}\edtext{\skt{तदपत्यस्य}}{\Afootnote{H; \skt{तदेवापत्य}\ A; \skt{सर्व\-दे\-हा\-स्रि\-त\-त्वाच्च शुक्र\-ल\-क्ष\-ण\-मु\-च्यते ॥ तदपत्यस्य}\ N}}\skt{ संस्प\-र्शा\-द्द\-र्श\-ना\-त्स्म\-र\-णा\-दपि ॥ }\pend
\pstart\versepar \skt{ग्रहणाच्च }\edtext{\skt{शरीरेभ्यः}}{\Afootnote{H; \skt{शरीरस्य}\ A; \skt{शरीरस्यः}\ N}}\skt{ }\edtext{\skt{शुक्र\-व\-त्सं\-प्र\-व\-र्त्तते}}{\Afootnote{A H; \skt{शुक्र\-व\-त्स\-म\-प्र\-व\-र्त्तते ॥}\ N}}\skt{ ॥}\pend

\pstart\versepar \ledsidenote{\passref{2.10.23}}\skt{स्नेहो }\edtext{\edtext{\skt{निरन्तरस्तासां}}{\lemma{\skt{निरन्तरस्}}\Afootnote{A H; \skt{निर}\ N}}}{\lemma{\skt{तासां}}\Afootnote{H; \skt{तत्र}\ A; \skt{स्तासां}\ N}}\skt{ }\edtext{\skt{प्रस्नवे}}{\Afootnote{H; \skt{प्रस्रवे}\ A N}}\skt{ हेतुरुच्यते । }\pend
\pstart\versepar \skt{तत्क\-षा\-य\-म्भ\-वे\-द्वा\-ता\-त्क्षि\-प्तञ्च }\edtext{\skt{प्लवते म्भसि}}{\Afootnote{N H; \skt{प्लवतेऽम्भसि ॥}\ A}}\skt{ ॥}\pend

\pstart \ledsidenote{\passref{2.10.24}}\edtext{\skt{पित्तादम्लं}}{\Afootnote{A H; \skt{पित्तादम्व्लं}\ N}}\skt{ सकटुकं }\edtext{\skt{राज्यम्भसि च}}{\Afootnote{H; \skt{राज्योऽम्भसि च}\ A; \skt{राज्योम्भसि च}\ N}}\skt{ }\edtext{\skt{पीतिका}}{\Afootnote{H; \skt{पीतिकाः ।}\ A N}}\skt{ । }\edtext{\skt{श्लेष्मा\-द्घ\-न\-म्पि\-च्छि\-लञ्च}}{\lemma{\skt{श्लेष्माद्घनम्}}\Afootnote{H; \skt{कफाद्घनं}\ A; \skt{श्लेष्माघन}\ N}}\skt{ जले चाप्यवसीदति ॥ }\edtext{\skt{सर्व्वैर्दुष्टं}}{\lemma{\skt{दुष्टं}}\Afootnote{N H; \skt{दुष्टैः}\ A}}\skt{ }\edtext{\skt{सर्व्व\-लि\-ङ्ग\-म\-भि\-घा\-ताच्च}}{\lemma{\skt{अभिघाताच्च}}\Afootnote{A H; \skt{मभिघाताच्च}\ N}}\skt{ दु-ष्यति ।}\pend

\pstart \ledsidenote{\passref{2.10.26}}\edtext{\skt{सक्षीरो वा प्यदुग्धो}}{\lemma{\skt{सक्षीरो वा प्य्}}\Afootnote{H; \skt{सक्षीरौ वाऽ}\ A; \skt{सक्षीरौ वाप्य्}\ N}}\skt{ वा }\edtext{\skt{दोषः}}{\Afootnote{H; \emph{om.}\ A; \skt{दोषं}\ N}}\skt{ प्राप्य }\edtext{\skt{स्तनौ}}{\Afootnote{N H; \skt{दोषः स्तनौ}\ A}}\skt{ }\edtext{\skt{स्त्रियाः}}{\Afootnote{N H; \skt{स्त्रियाः ।}\ A}}\skt{ । }\edtext{\skt{प्रदूष्य}}{\Afootnote{N H; \skt{रक्तं मांसं च सन्दूष्य}\ A}}\skt{ }\edtext{\skt{मांसरुधिर}}{\Afootnote{H; \emph{om.}\ A; \skt{मांसरुधिरे}\ N}}\skt{ स्तनरोगाय कल्पते ।}\edtext{}{\lemma{\skt{अदुग्धो वा}}\Afootnote{H; \skt{प्यदुग्धौ वा}\ A; \skt{अदुग्धौ वा}\ N}}\pend

\pstart \ledsidenote{\passref{2.10.27}}\edtext{\edtext{\skt{षण्णामपि}}{\lemma{\skt{षण्णाम्}}\Afootnote{H; \skt{पञ्चानाम्}\ A; \skt{पञ्चाम्}\ N}}\skt{ च}}{\lemma{\skt{अपि च}}\Afootnote{N H; \skt{अपि}\ A}}\skt{ }\edtext{\skt{तेषां}}{\Afootnote{H; \skt{तेषां तु}\ A; \skt{तेषा}\ N}}\skt{ }\edtext{\skt{हि}}{\Afootnote{N H; \skt{हित्वा शो}\ A}}\skt{ }\edtext{\edtext{\edtext{\skt{रक्त\-ज\-म्वि\-द्र\-धि\-म्विना}}{\lemma{\skt{रक्तजम्}}\Afootnote{N H; \skt{णित}\ A}}}{\lemma{\skt{विद्रधिम्}}\Afootnote{A H; \skt{विद्रधि}\ N}}}{\lemma{\skt{विना}}\Afootnote{H; \emph{om.}\ A; \skt{वाना ॥}\ N}}\skt{ ॥ लक्षणानि समानानि }\edtext{\edtext{\skt{बाह्य\-वि\-द्र\-धि\-ल\-क्ष\-णै\-रिति}}{\lemma{\skt{बाह्य\-वि\-द्र\-धि\-ल\-क्ष\-णैर्}}\Afootnote{A H; \skt{वाह्य\-वि\-द्र\-धि\-ल\-क्ष\-णैर्}\ N}}}{\lemma{\skt{इति}}\Afootnote{H; \emph{om.}\ A N}}\skt{ ॥}\pend

\pstart \ledsidenote{\passref{2.10.27a}}\skt{प्रभ\-ञ्ज\-न\-क\-दु\-र्न्नामा भगदारणं ।}\pend

\pstart \ledsidenote{\passref{2.10.27b}}\skt{कुष्ठ\-प्र\-मे\-ह\-ज\-ढ\-ग\-र्ब्भञ्च विद्रधिम् ॥}\pend

\pstart \ledsidenote{\passref{2.10.27c}}\skt{विस\-र्प्प\-स्त\-न\-ना\-डीभिः पूर्यते दशकोपनः ॥}\pend

\pstart \ledsidenote{\passref{2.10.trailer}}\edtext{\skt{विस\-र्प्प\-ना\-डी\-स्त\-न\-रो\-ग\-नि\-दा\-न\-द\-शमः}}{\Afootnote{H; \skt{इति सुश्रु\-त\-सं\-हि\-तायां निदानस्थाने विस\-र्प\-ना\-डी\-स्त\-न\-रो\-ग\-नि\-दानं नाम दशमोऽध्यायः ॥१० ॥}\ A; \skt{इति ॥ निदाने ॥.. {\symfont ❈} ॥}\ N}}\skt{ ॥}\pend

\endnumbering
\beginnumbering
\eledsection*{Adhyāya 11}
\pstart \ledsidenote{\passref{2.11.1}}\skt{अथातो }\edtext{\edtext{\skt{ग्रन्थ्य\-प\-च्य\-र्व्वु\-द\-ग\-ल\-ग\-ण्डा\-ना\-न्नि\-दा\-न\-म्व्या\-ख्या\-स्यामः}}{\lemma{\skt{ग्रन्थ्य\-प\-च्य\-र्व्वु\-द\-ग\-ल\-ग\-ण्डा\-नान्}}\Afootnote{H; \skt{ग्रन्थ\-प\-च्य\-र्बु\-द\-ग\-ल\-ग\-ण्डानां}\ A; \skt{ग्रन्थ्य\-प\-च्य\-र्वु\-द\-ग\-ल\-ग\-ण्डान्}\ N}}}{\lemma{\skt{व्याख्यास्यामः}}\Afootnote{A H; \skt{व्याखाश्यामः ॥}\ N}}\skt{ ॥}\pend

\pstart \ledsidenote{\passref{2.11.3}}\skt{वातादयो }\edtext{\edtext{\skt{मान्स\-म\-सृ\-क्प्र\-दुष्टाः}}{\lemma{\skt{मान्सम्}}\Afootnote{A H; \skt{मांस}\ N}}}{\lemma{\skt{असृक्प्रदुष्टाः}}\Afootnote{H; \skt{असृक्च दुष्टाः}\ A; \skt{मसृक्दुष्टाः}\ N}}\skt{ सन्दूष्य }\edtext{\skt{सेदश्च}}{\Afootnote{H; \skt{मेदश्च}\ A N}}\skt{ }\edtext{\skt{तथा}}{\Afootnote{H; \skt{कफा}\ A; \skt{त}\ N}}\skt{ }\edtext{\skt{सिरांश्च}}{\Afootnote{N H; \skt{नुविद्धम् ।}\ A}}\skt{ वृत्तो\-न्न\-त\-म्वि\-ग्र\-थि\-तन्तु }\edtext{\skt{शोफं}}{\Afootnote{A H; \skt{शोफ}\ N}}\skt{ }\edtext{\skt{कुर्वन्त्यतो}}{\Afootnote{A H; \skt{न्कुर्वन्त्यतो}\ N}}\skt{ }\edtext{\edtext{\skt{ग्रन्थिरिति}}{\lemma{\skt{ग्रन्थिर्}}\Afootnote{N H; \skt{ग्रन्थि}\ A}}}{\lemma{\skt{इति}}\Afootnote{N H; \skt{रिति}\ A}}\skt{ प्रदिष्टः ॥}\pend

\pstart \ledsidenote{\passref{2.11.4}}\skt{आयम्यते }\edtext{\skt{वृञ्जति}}{\Afootnote{H; \skt{व्यथ्यत एति}\ A; \skt{वृति}\ N}}\skt{ }\edtext{\skt{तुद्यते च}}{\Afootnote{N H; \skt{तोदं}\ A}}\skt{ प्रत्यस्य ते }\edtext{\skt{साथ्स्यति}}{\Afootnote{H; \skt{कृत्यत}\ A; \skt{मथ्यति}\ N}}\skt{ }\edtext{\skt{भिद्यते च}}{\Afootnote{N H; \skt{एति भेदम् ।}\ A}}\skt{ । }\edtext{\skt{कृष्णो}}{\Afootnote{A H; \skt{कृष्णो}\ N}}\skt{ }\edtext{\skt{मृदु\-र्व्ब\-स्ति\-रि\-वा\-त\-तश्च}}{\lemma{\skt{मृदुर्व्बस्तिर्}}\Afootnote{\skt{ऽमृदुर्बस्तिर्}\ A; \skt{मृदुर्वस्तिर्}\ N; \skt{मृदुर्व्वस्तिर्}\ H}}\skt{ }\edtext{\skt{भिन्नः}}{\Afootnote{H; \skt{बिन्नः}\ A; \skt{भिन्न}\ N}}\skt{ }\edtext{\skt{स्रवे\-च्चा\-नि\-ल\-जो\-म्रे\-मच्छं}}{\Afootnote{H; \skt{स्रवे\-च्चा\-नि\-ल\-जोऽ\-स्र\-म\-च्छम् ॥}\ A; \skt{स्रवे\-च्चा\-नि\-ल\-जो\-स्र\-मच्छं ॥}\ N}}\skt{ ॥}\pend

\pstart \ledsidenote{\passref{2.11.5}}\skt{दन्दह्यते धूप्यति चूष्यते च पापच्यते प्रज्वलतीव चापि । }\edtext{\skt{रक्तः}}{\Afootnote{A H; \skt{रक्ता}\ N}}\skt{ }\edtext{\skt{सपित्तो}}{\Afootnote{H; \skt{सपीतो}\ A; \emph{om.}\ N}}\skt{ }\edtext{\edtext{\skt{प्यथवापि}}{\lemma{\skt{प्य्}}\Afootnote{H; \skt{ऽप्य}\ A; \skt{पि}\ N}}}{\lemma{\skt{अथवापि}}\Afootnote{H; \skt{थवाऽपि}\ A; \skt{त्तोप्यथवापि}\ N}}\skt{ }\edtext{\skt{पित्ताद्भिन्नः}}{\Afootnote{H; \skt{पित्ताद्भिन्नः}\ A N}}\skt{ }\edtext{\edtext{\skt{स्रवे\-द्दु\-ष्ण\-म\-तीव}}{\lemma{\skt{स्रवेद्दुष्णम्}}\Afootnote{\skt{स्रवेदुष्णम्}\ A N; \skt{स्रवेद्दु}\ H\textsuperscript{ac}; \skt{स्रवेद्दुष्ण}\ H\textsuperscript{pc}}}}{\lemma{\skt{अतीव}}\Afootnote{A N; \skt{ष्णमतीव}\ H\textsuperscript{ac}; \skt{दुष्णमतीव}\ H\textsuperscript{pc}}}\skt{ चास्रम् ॥}\pend

\pstart\versepar \ledsidenote{\passref{2.11.6}}\edtext{\skt{पीतो}}{\Afootnote{H; \skt{शीतोऽ}\ A; \skt{शीतो}\ N}}\skt{ }\edtext{\skt{विवर्णो}}{\Afootnote{A H; \skt{विवर्ण्णो}\ N}}\skt{ }\edtext{\skt{ल्परुजो}}{\Afootnote{H; \skt{ऽल्परुजो}\ A; \skt{ल्परुजो}\ N}}\skt{ }\edtext{\skt{तिकण्डूः}}{\Afootnote{A H; \skt{तिकण्डूः}\ N}}\skt{ }\edtext{\skt{पाषा\-ण\-व\-त्सं\-ह\-न\-नो\-प\-पन्नः}}{\Afootnote{A H; \skt{पाषा\-ण\-व\-त्सं\-ह\-न\-तो\-प\-प\-पन्न्नः ॥}\ N}}\skt{ }\pend
\pstart\versepar \skt{। }\edtext{\skt{चिराधि}}{\Afootnote{H; \skt{चिराभि}\ A N}}\skt{ वृद्धिश्च }\edtext{\skt{कफ\-प्र\-को\-पा\-द्भिन्नः}}{\Afootnote{A H; \skt{कफ\-प्र\-को\-पा\-द्भि\-न्नस्}\ N}}\skt{ स्रवे\-च्छु\-क्ल\-घ\-नञ्च पूयम् ॥}\pend

\pstart\versepar \ledsidenote{\passref{2.11.7}}\edtext{\skt{शरी\-र\-वृ\-द्धि\-क्ष\-य\-वृ\-द्धि\-हानिः}}{\Afootnote{A H; \skt{शरीरवृद्धिः क्षयवृद्धिहान्नीः}\ N}}\skt{ स्निग्धो }\edtext{\edtext{\skt{महा\-नू\-क\-ण्डु\-यु\-तोऽ\-रु\-जश्च}}{\lemma{\skt{महा\-नू\-क\-ण्डु\-युतो}}\Afootnote{H; \skt{महानल्परुजो}\ A; \skt{महाकण्डयु}\ N}}}{\lemma{\skt{ऽरुजश्च}}\Afootnote{H; \skt{ऽतिकण्डूः ।}\ A; \skt{तोरूजाश्च ।}\ N}}\skt{ । }\pend
\pstart\versepar \skt{मेदःकृतो }\edtext{\skt{दृश्यति}}{\Afootnote{H; \skt{गच्छति}\ A; \skt{दृश्यते}\ N}}\skt{ चात्र भिन्ने }\edtext{\skt{पिण्या\-क\-क\-ल्क\-प्र\-ति\-मन्तु}}{\Afootnote{H; \skt{पिण्या\-क\-स\-र्पिः\-प्र\-तिमं तु}\ A; \skt{पिण्या\-क\-ल्क\-प्र\-ति\-मन्तु}\ N}}\skt{ मेदः ॥}\pend

\pstart\versepar \ledsidenote{\passref{2.11.8}}\edtext{\skt{व्याया\-म\-जा\-तै\-र\-व\-लस्य तैस्तैराक्षिप्य}}{\lemma{\skt{अवलस्य तैस्तैर्}}\Afootnote{H; \skt{अबलस्य तैस्तैर्}\ A N}}\skt{ }\edtext{\skt{वायुर्हि}}{\Afootnote{A H; \skt{वायुर्हि}\ N}}\skt{ सिराप्रतानं । }\pend
\pstart\versepar \skt{}\edtext{\skt{सङ्कुञ्च्य}}{\Afootnote{H; \skt{संपीड्य}\ A; \skt{संच्य}\ N}}\skt{ }\edtext{\skt{सपिण्ड्य}}{\Afootnote{H; \skt{सङ्कोच्य}\ A; \skt{स}\ N}}\skt{ विशोष्य चापि }\edtext{\edtext{\edtext{\skt{ग्रन्थि\-ङ्क\-रो\-त्यु\-न्न\-त\-माशु}}{\lemma{\skt{ग्रन्थिङ्}}\Afootnote{A H; \skt{ग्रन्थि}\ N}}}{\lemma{\skt{करोत्य्}}\Afootnote{A H; \skt{करो}\ N}}}{\lemma{\skt{उन्नतम्}}\Afootnote{A H; \skt{त्युन्नतम्}\ N}}\skt{ वृत्तम् ॥}\pend

\pstart\versepar \ledsidenote{\passref{2.11.9}}\skt{ग्रन्थिः सिराजः स तु कृच्छ्रसाध्यो भवेद्यदि }\edtext{\skt{स्यात्स\-रु\-ज\-श्च\-लश्च}}{\Afootnote{H; \skt{स्यात्स\-रु\-ज\-श्च\-लश्च ।}\ A N}}\skt{ । }\pend
\pstart\versepar \skt{}\edtext{\skt{अरुक्स}}{\Afootnote{A H; \skt{अरु स}\ N}}\skt{ }\edtext{\skt{एवाप्यचलो}}{\Afootnote{N H; \skt{एवाप्यचलो}\ A}}\skt{ महांश्च मर्मो\-त्थि\-त\-श्चा\-पि\-वि\-व\-र्ज्ज\-नीयः ॥}\pend

\pstart\versepar \ledsidenote{\passref{2.11.10}}\edtext{\skt{हन्व\-स्थि\-क\-क्ष्या\-क्ष\-क\-वा\-हु\-स\-न्धि\-म\-न्या\-ग\-ले\-षू\-प\-चि\-तन्तु}}{\Afootnote{H; \skt{हन्व\-स्थि\-क\-क्षा\-क्ष\-क\-बा\-हु\-स\-न्धि\-म\-न्या\-ग\-ले\-षू\-प\-चितं तु}\ A; \skt{हृव\-स्थि\-क\-क्षा\-क्ष\-क\-वा\-हु\-स\-न्धि\-म\-न्या\-ग\-ले\-षू\-प\-चितं तु}\ N}}\skt{ मेदः । }\pend
\pstart\versepar \skt{}\edtext{\skt{ग्रन्थिं}}{\Afootnote{A H; \emph{om.}\ N}}\skt{ स्थिरं वृत्त\-म\-था\-य\-तम्वा }\edtext{\edtext{\edtext{\skt{स्निग्धङ्कफश्चास्य}}{\lemma{\skt{स्निग्धङ्}}\Afootnote{A H; \skt{स्निंग्ध}\ N}}}{\lemma{\skt{कफश्}}\Afootnote{A H; \skt{कफाश्}\ N}}}{\lemma{\skt{चास्य}}\Afootnote{H; \skt{चाल्प}\ A; \skt{चा}\ N}}\skt{ रुजङ्करोति ॥}\pend

\pstart \ledsidenote{\passref{2.11.11}}\edtext{\skt{तं }\edtext{\edtext{\edtext{\skt{ग्रन्थि\-भि\-स्त्वा\-म\-ल\-का\-स्थि\-मा\-त्रै\-र्म्म\-त्स्या\-ण्ड\-जा\-ल\-प्र\-ति\-मै\-स्त\-थान्यैः}}{\lemma{\skt{त्वाम\-ल\-का\-स्थि\-मा\-त्रैर्}}\Afootnote{A H; \skt{स्त्वाम\-ल\-का\-स्थि\-मा\-त्रैर्}\ N}}}{\lemma{\skt{म्मत्स्या\-ण्ड\-जा\-ल\-प्र\-ति\-मैस्}}\Afootnote{A H; \skt{मर्त्या\-ण्ड\-जा\-ल\-प्र\-ति\-मैस्}\ N}}}{\lemma{\skt{तथान्यैः}}\Afootnote{N H; \skt{तथाऽन्यैः ।}\ A}}}{\lemma{\skt{तं ग्रन्थिभिस्}}\Afootnote{A H; \skt{तैर्ग्रन्थिभि}\ N}}\skt{ }\edtext{\edtext{\skt{अन\-न्य\-व\-र्णै\-रु\-प\-ची\-य\-मानं}}{\lemma{\skt{अनन्यवर्णैर्}}\Afootnote{A H; \skt{अनन्यवर्णै}\ N}}}{\lemma{\skt{उपचीयमानं}}\Afootnote{A H; \skt{रूपचीयमानं}\ N}}\skt{ }\edtext{\skt{चय\-प्र\-क\-र्षा\-द\-पचिं}}{\Afootnote{H; \skt{चय\-प्र\-क\-र्षा\-द\-पचीं}\ A; \skt{चय\-प्र\-क\-र्षा\-द\-प\-चिम्}\ N}}\skt{ वदन्ति ॥}\pend

\pstart\versepar \ledsidenote{\passref{2.11.12}}\edtext{\edtext{\skt{कण्ड्व\-न्वि\-ता\-स्तेऽ\-ल्प\-रुजा}}{\lemma{\skt{कण्ड्वन्वितास्ते}}\Afootnote{H; \skt{कण्डूयुतास्ते}\ A; \skt{कण्दुम्वितास्ते}\ N}}}{\lemma{\skt{ऽल्परुजा}}\Afootnote{\skt{ऽल्परुजः}\ A; \skt{ल्परुज}\ N; \skt{ल्परुजा}\ H}}\skt{ प्रभिन्नाः स्रवन्ति नश्यन्ति भवन्ति चान्ये । }\pend
\pstart\versepar \skt{मेदःकफाभ्यां खलु }\edtext{\skt{रोग}}{\Afootnote{A H; \skt{रोग रोग}\ N}}\skt{ एष सुदुस्तरो }\edtext{\skt{वर्ष\-ग\-णा\-नु\-बन्धी}}{\Afootnote{A H; \skt{वर्षगण्{\symfont ¦}अनुवन्धि ॥}\ N}}\skt{ ॥}\pend

\pstart \ledsidenote{\passref{2.11.13}}\edtext{\skt{गात्रप्रदेशे}}{\Afootnote{A H; \skt{गात्रप्रदोशे}\ N}}\skt{ }\edtext{\skt{क्वचिदेव}}{\Afootnote{N H; \skt{क्वचिदेव}\ A}}\skt{ दोषाः संमूर्च्छिता }\edtext{\edtext{\skt{मांस\-म\-सृ\-क्प्र\-दूष्य}}{\lemma{\skt{मांसम्}}\Afootnote{A H; \skt{मांस}\ N}}}{\lemma{\skt{असृक्प्रदूष्य}}\Afootnote{H; \skt{अभिप्रदूष्य ।}\ A; \skt{सप्रष्य ।}\ N}}\skt{ । }\edtext{\skt{वृत्तं}}{\Afootnote{A H; \skt{वृत्त}\ N}}\skt{ }\edtext{\edtext{\skt{मृदु\-म्म\-न्द\-रु\-ज\-म्म\-हा\-न्त\-म\-न\-ल्प\-मू\-ल\-ञ्चि\-र\-वृ\-द्ध्य\-पा\-कम्}}{\lemma{\skt{मृदुम्}}\Afootnote{H; \skt{स्थिरं}\ A; \skt{न्मृद}\ N}}}{\lemma{\skt{मन्दरुजम्}}\Afootnote{A H; \skt{मन्दरूजम्}\ N}}\skt{ ॥}\pend

\pstart\versepar \ledsidenote{\passref{2.11.14}}\edtext{\skt{कुर्व्वन्ति}}{\Afootnote{A H; \skt{कुर्व}\ N}}\skt{ }\edtext{\edtext{\edtext{\skt{मांसो\-च्छ्र\-य\-म\-भ्य\-गा\-ध\-न्त\-म\-र्व्वुदं}}{\lemma{\skt{मांसोच्छ्रयम्}}\Afootnote{H; \skt{मांसोपच}\ A; \skt{मांसोयम्}\ N}}}{\lemma{\skt{अभ्यगाधन्तम्}}\Afootnote{N H; \skt{यं तु शोफं तम्}\ A}}}{\lemma{\skt{अर्व्वुदं}}\Afootnote{N H; \skt{अर्बुदं}\ A}}\skt{ शास्त्रविदो वदन्ति । }\pend
\pstart\versepar \skt{वातेन पित्तेन कफेन चापि }\edtext{\skt{रक्तेन}}{\Afootnote{A N H\textsuperscript{pc}; \skt{रक्ते}\ H\textsuperscript{ac}}}\skt{ मांसेन च मेदसा च ॥}\pend

\pstart\versepar \ledsidenote{\passref{2.11.15}}\edtext{\skt{यज्जायते}}{\Afootnote{H; \skt{तज्जायते}\ A; \skt{यज्जायते}\ N}}\skt{ तस्य च लक्षणानि ग्रन्थेः }\edtext{\skt{समाननि}}{\Afootnote{H; \skt{समानानि}\ A N}}\skt{ सदा भवन्ति । }\pend
\pstart\versepar \skt{}\edtext{\skt{दोषाः}}{\Afootnote{N; \skt{दोषः}\ A; \skt{दोषा}\ H}}\skt{ }\edtext{\skt{प्रदुष्टा}}{\Afootnote{H; \skt{प्रदुष्टो}\ A N}}\skt{ रुधिरं }\edtext{\skt{सिरासु}}{\Afootnote{N H; \skt{सिरास्तु}\ A}}\skt{ }\edtext{\skt{सकुञ्च्य}}{\Afootnote{H; \skt{संपीड्य}\ A; \skt{संकुञ्च्य}\ N}}\skt{ }\edtext{\skt{सम्पिण्ड्य}}{\Afootnote{H; \skt{सङ्कोच्य}\ A; \skt{सपिञ्}\ N}}\skt{ }\edtext{\skt{ततस्त्वपाकं}}{\lemma{\skt{ततस्}}\Afootnote{N H; \skt{गतस्}\ A}}\skt{ ॥}\pend

\pstart\versepar \ledsidenote{\passref{2.11.16}}\edtext{\skt{सास्रा\-व\-मु\-न्न\-ह्यति}}{\Afootnote{N H; \skt{सास्रा\-व\-मु\-न्न\-ह्यति}\ A}}\skt{ }\edtext{\skt{मान्स\-पि\-ण्ड\-म्मां\-सा\-ङ्कु\-रै\-रा\-चि\-त\-माशु}}{\lemma{\skt{मान्सपिण्डम्}}\Afootnote{A H; \skt{मांसपिण्म्}\ N}}\skt{ }\edtext{\skt{घोरं}}{\Afootnote{H; \skt{वृद्धिम्}\ A N}}\skt{ । }\pend
\pstart\versepar \skt{स्रवत्यजस्रं }\edtext{\skt{रुधि\-र\-म्प्र\-दु\-ष्ट\-म\-सा\-ध्य\-मे\-त\-द्रु\-धि\-रा\-त्म\-कन्तु}}{\lemma{\skt{एत\-द्रु\-धि\-रा\-त्म\-कन्तु}}\Afootnote{N H; \skt{एत\-द्रु\-धि\-रा\-त्मकं स्यात् ॥}\ A}}\skt{ ॥}\pend

\pstart\versepar \ledsidenote{\passref{2.11.17}}\edtext{\edtext{\skt{रक्त\-क्ष\-यो\-प\-हू\-त\-पी\-डि\-त\-त्वा\-त्पा\-ण्डु\-र्व्भ\-वे\-त्सोऽ\-र्व्वु\-द\-पी\-डि\-तस्तु}}{\lemma{\skt{रक्त\-क्ष\-यो\-प\-हू\-त\-पी\-डि\-त\-त्वा\-त्पा\-ण्डु\-र्व्भ\-वेत्सो}}\Afootnote{H; \skt{रक्त\-क्ष\-यो\-प\-द्र\-व\-पी\-डि\-त\-त्वा\-त्पा\-ण्डु\-र्भ\-वेत्सो}\ A; \skt{रक्त\-क्ष\-प\-द्ग\-त\-पि\-डि\-त\-त्वा\-त्पा\-ण्डू\-र्भ\-वेत्सो}\ N}}}{\lemma{\skt{ऽर्व्वु\-द\-पी\-डि\-तस्तु}}\Afootnote{H; \skt{ऽर्बु\-द\-पी\-डि\-तस्तु ।}\ A; \skt{वुर्दपीडितस्तु ॥}\ N}}\skt{ । }\pend
\pstart\versepar \skt{}\edtext{\edtext{\skt{मुष्टि\-प्र\-हा\-रा\-दि\-भि\-र\-र्द्दिते}}{\lemma{\skt{मुष्टि\-प्र\-हा\-रा\-दि\-भिर्}}\Afootnote{A H; \skt{मुष्ठि\-प्र\-हा\-रा\-दि\-भिर्}\ N}}\skt{ ङ्गे}}{\lemma{\skt{अर्द्दिते ङ्गे}}\Afootnote{H; \skt{अर्दितेऽङ्गे}\ A; \skt{अर्दितंगे}\ N}}\skt{ }\edtext{\edtext{\skt{मान्स\-म्प्र\-दु\-ष्ट\-ञ्ज\-न\-यत्तु}}{\lemma{\skt{मान्सम्}}\Afootnote{A H; \skt{मासं}\ N}}}{\lemma{\skt{जनयत्तु}}\Afootnote{H; \skt{प्रकरोति}\ A; \skt{जनयेत्तु}\ N}}\skt{ }\edtext{\skt{शोफम्}}{\Afootnote{A H; \skt{शोफा ॥}\ N}}\skt{ ॥}\pend

\pstart\versepar \ledsidenote{\passref{2.11.18}}\edtext{\skt{अवेदनं}}{\Afootnote{A H; \skt{अवेदन}\ N}}\skt{ }\edtext{\edtext{\edtext{\skt{स्निग्ध\-म\-न\-न्य\-व\-र्ण्ण\-म\-पा\-क\-म\-श्मो\-प\-म\-म\-प्र\-चाल्यं}}{\lemma{\skt{अपाकम्}}\Afootnote{A H; \skt{अवाकम्}\ N}}}{\lemma{\skt{अश्मोपमम्}}\Afootnote{A H; \skt{अश्मोपम}\ N}}}{\lemma{\skt{अप्रचाल्यं}}\Afootnote{A H; \skt{मप्रचाल्य ।}\ N}}\skt{ । }\pend
\pstart\versepar \skt{प्रदुष्टमान्सस्य नरस्य }\edtext{\edtext{\skt{गाढ\-मे\-त\-द्भ\-वे\-त्मा\-न्स\-प\-रा\-य\-णस्य}}{\lemma{\skt{गाढम्}}\Afootnote{N H; \skt{बाढम्}\ A}}}{\lemma{\skt{एत\-द्भ\-वे\-त्मा\-न्स\-प\-रा\-य\-णस्य}}\Afootnote{H; \skt{एत\-द्भ\-वे\-न्मां\-स\-प\-रा\-य\-णस्य ॥}\ A; \skt{एत\-द्भ\-वे\-त्मां\-स\-रा\-य\-नस्य ॥}\ N}}\skt{ ॥}\pend

\pstart\versepar \ledsidenote{\passref{2.11.19}}\edtext{\skt{मांसार्व्वुदं}}{\Afootnote{N H; \skt{मांसार्बुदं}\ A}}\skt{ त्वेत\-द\-सा\-ध्य\-मुक्तं साध्येष्वपीमानि }\edtext{\skt{विवर्ज्जयीत}}{\Afootnote{N H; \skt{विवर्जयेत्तु ।}\ A}}\skt{ । }\pend
\pstart\versepar \skt{सम्प्र\-स्रु\-त\-म्म\-र्म्मणि यच्च जातं }\edtext{\skt{स्रोतस्सु वा}}{\Afootnote{A H; \skt{स्रोतस्सु}\ N}}\skt{ यच्च भवेदचाल्यम् ॥}\pend

\pstart\versepar \ledsidenote{\passref{2.11.20}}\edtext{\edtext{\skt{यज्जा\-य\-तेऽ\-न्य\-त्खलु}}{\lemma{\skt{यज्जायते}}\Afootnote{A; \skt{यज्जायते}\ N; \skt{> यज्जायते}\ H}}}{\lemma{\skt{ऽन्यत्खलु}}\Afootnote{A H; \skt{न्यद्खलु}\ N}}\skt{ पूर्व्वजाते }\edtext{\edtext{\skt{ज्ञेय\-न्त\-द\-ध्य\-र्व्वु\-द\-म\-र्बु\-दज्ञैः}}{\lemma{\skt{तदध्यर्व्वुदम्}}\Afootnote{H; \skt{तदध्यर्बुदम्}\ A; \skt{तदन्यर्वुदम्}\ N}}}{\lemma{\skt{अर्बुदज्ञैः}}\Afootnote{A H; \skt{अर्वुदज्ञैः ।}\ N}}\skt{ । }\pend
\pstart\versepar \skt{यद्द्वन्द्वजातं युगपत्क्रमाद्वा }\edtext{\skt{द्विरर्व्वुदन्तच्च}}{\lemma{\skt{द्विरर्व्वुदन्}}\Afootnote{N H; \skt{द्विरर्बुदं}\ A}}\skt{ }\edtext{\skt{भवेदसाध्यम्}}{\Afootnote{A H; \skt{भवेदसाध्य ॥}\ N}}\skt{ ॥}\pend

\pstart\versepar \ledsidenote{\passref{2.11.21}}\skt{न पाकमायान्ति }\edtext{\skt{कफा\-न्वि\-त\-त्वा\-न्मे\-दो\-व\-हु\-त्वाच्च}}{\Afootnote{N H; \skt{कफा\-धि\-क\-त्वा\-न्मे\-दो\-ब\-हु\-त्वाच्च}\ A}}\skt{ विशेषतस्तु । }\pend
\pstart\versepar \skt{दोष\-स्थि\-र\-त्वा\-द्ग्र\-थ\-नाच्च तेषां }\edtext{\skt{सर्व्वार्व्वुदान्येव}}{\lemma{\skt{सर्व्वार्व्वुदान्य्}}\Afootnote{N H; \skt{सर्वार्बुदान्य्}\ A}}\skt{ }\edtext{\skt{विशेतम्तु}}{\lemma{\skt{विशेतम्}}\Afootnote{H; \skt{निसर्गतस्}\ A; \skt{विषेतस्}\ N}}\skt{ ॥}\pend

\pstart\versepar \ledsidenote{\passref{2.11.22}}\skt{वातः }\edtext{\skt{कफश्चापि}}{\lemma{\skt{चापि}}\Afootnote{N H; \skt{चैव}\ A}}\skt{ गले }\edtext{\skt{प्र त्तु}}{\Afootnote{\skt{प्रवृद्धौ मन्ये तु}\ A; \skt{प्रदुष्ठो मन्ये तु}\ N; \skt{प्र… त्तु}\ H}}\skt{ }\edtext{\skt{संश्रित्य}}{\Afootnote{H; \skt{संसृत्य}\ A; \skt{संस्रित्य}\ N}}\skt{ तथैव मेदः । }\pend
\pstart\versepar \skt{कुर्व्वन्ति गण्डं क्रमशः स्वलिङ्गैः }\edtext{\edtext{\skt{सम\-न्वि\-त\-न्त\-ङ्ग\-ल\-ग\-ण्ड\-माहुः}}{\lemma{\skt{समन्वितन्तङ्}}\Afootnote{A H; \skt{समन्वितन्त}\ N}}}{\lemma{\skt{गलगण्डम्}}\Afootnote{A H; \skt{न्गलगण्डम्}\ N}}\skt{ ॥}\pend

\pstart \ledsidenote{\passref{2.11.23}}\skt{तोदान्वितः कृष्ण\-सि\-रा\-व\-नद्धः }\edtext{\skt{कृष्णो}}{\Afootnote{A H; \skt{कृष्णो}\ N}}\skt{ }\edtext{\skt{रुणो वा}}{\Afootnote{N H; \skt{ऽरुणो वा}\ A}}\skt{ }\edtext{\edtext{\skt{पवनात्मकस्तु}}{\lemma{\skt{पवनात्मकस्}}\Afootnote{A H; \skt{पणनात्मकस्}\ N}}}{\lemma{\skt{तु}}\Afootnote{N H; \skt{तु । मेदो\-न्वि\-त\-श्चो\-प\-चि\-तश्च काला\-द्भ\-वे\-द\-ति\-स्नि\-ग्ध\-त\-रोऽ\-रु\-जश्च ॥}\ A}}\skt{ ।}\pend

\pstart \ledsidenote{\passref{2.11.24}}\edtext{\skt{पारु\-ष्य\-यु\-क्ति\-श्चि\-र\-वृ\-द्ध्य\-पाको}}{\lemma{\skt{पारुष्ययुक्तिश्}}\Afootnote{H; \skt{पारुष्ययुक्तश्}\ A N}}\skt{ यदृच्छया }\edtext{\skt{चैष}}{\Afootnote{N H; \emph{om.}\ A}}\skt{ }\edtext{\edtext{\skt{भवे\-त्क\-दा\-चि\-त्पा\-क\-मि\-या\-त्क\-दा\-चित्}}{\lemma{\skt{भवे\-त्क\-दा\-चि\-त्पा\-कम्}}\Afootnote{H; \skt{पाक}\ A; \skt{भवे}\ N}}\skt{ ॥ वैरस्यमास्यस्य}}{\lemma{\skt{इयात्कदाचित् ॥ वैरस्यम्}}\Afootnote{H; \skt{मिया\-त्क\-दा\-चि\-त्वै\-र\-स्यम्}\ A; \skt{त्कदाचित् । वैरस्यम्}\ N}}\skt{ च तस्य }\edtext{\skt{जन्तोर्व्भवेत्तथा}}{\lemma{\skt{व्भवेत्तथा}}\Afootnote{H; \skt{भवेत्तथा}\ A N}}\skt{ तालु\-ग\-ल\-प्र\-शोषः ।}\pend

\pstart\versepar \ledsidenote{\passref{2.11.25}}\edtext{\skt{स्थिरः}}{\Afootnote{A H; \skt{स्थिरे}\ N}}\skt{ सवर्णो }\edtext{\skt{गुरुरुग्रकण्डू}}{\Afootnote{H; \skt{ऽल्प\-रु\-गु\-ग्र\-कण्डूः}\ A; \skt{गुतुरुग्रकण्डू}\ N}}\skt{ शीतो }\edtext{\skt{महां\-श्चा\-पि\-क\-फा\-त्म\-कस्तु}}{\lemma{\skt{महांश्}}\Afootnote{A H; \skt{महाश्}\ N}}\skt{ ॥ }\pend
\pstart\versepar \skt{}\edtext{\skt{चिरा\-भि\-वृ\-द्धि\-म्भ\-जते}}{\lemma{\skt{भजते}}\Afootnote{N H; \skt{कुरुते}\ A}}\skt{ चिराच्च }\edtext{\skt{प्रपद्यते}}{\Afootnote{N H; \skt{प्रपच्यते}\ A}}\skt{ मन्दरुजः कदाचित् ।}\pend

\pstart\versepar \ledsidenote{\passref{2.11.26}}\skt{माधुर्यमास्यस्य च }\edtext{\edtext{\skt{जन्तोर्व्भवेत्तथा}}{\lemma{\skt{जन्तोर्}}\Afootnote{H; \skt{तस्य जन्तोर्}\ A N}}}{\lemma{\skt{व्भवेत्तथा}}\Afootnote{H; \skt{भवेत्तथा}\ A N}}\skt{ तालु\-ग\-ल\-प्र\-लेपः ॥ }\pend
\pstart\versepar \skt{स्निग्धो }\edtext{\skt{गुदुः}}{\Afootnote{H; \skt{मृदुः}\ A; \skt{गुरु}\ N}}\skt{ पाण्डु\-र\-नि\-ष्ट\-गन्धो }\edtext{\skt{मेदोथितः}}{\Afootnote{N H; \skt{मेदःकृतो}\ A}}\skt{ }\edtext{\skt{पाण्डुयुतो}}{\Afootnote{H; \skt{नीरु\-ग\-था\-ति\-कण्डूः ॥}\ A; \skt{कण्डूयुतो}\ N}}\skt{ }\edtext{\edtext{\skt{रुजश्च}}{\lemma{\skt{रुजश्}}\Afootnote{N H; \emph{om.}\ A}}}{\lemma{\skt{च}}\Afootnote{N H; \emph{om.}\ A}}\skt{ ।}\pend

\pstart\versepar \ledsidenote{\passref{2.11.27}}\edtext{\edtext{\skt{प्रल\-म्व\-तेऽ\-ला\-वु\-व\-द\-ल्प\-मूलो}}{\lemma{\skt{प्रलम्वते}}\Afootnote{N H; \skt{प्रलम्बते}\ A}}}{\lemma{\skt{ऽला\-वु\-व\-द\-ल्प\-मूलो}}\Afootnote{H; \skt{ऽला\-बु\-व\-द\-ल्प\-मूलो}\ A; \skt{लावु\-व\-द\-ल्प\-मूलो}\ N}}\skt{ विवर्द्धते हीयति चात्र }\edtext{\skt{देहम्}}{\Afootnote{N H; \skt{देहा\-नु\-रू\-प\-क्षय}\ A}}\skt{ ॥ }\pend
\pstart\versepar \skt{}\edtext{\skt{स्निग्धास्यता}}{\Afootnote{N H; \skt{वृद्धियुक्तः । स्निग्धास्यता}\ A}}\skt{ तस्य भवेच्च }\edtext{\edtext{\skt{जन्तोर्ग्गले}}{\lemma{\skt{जन्तोर्}}\Afootnote{H; \skt{जन्तो}\ A N}}}{\lemma{\skt{ग्गले}}\Afootnote{H; \skt{र्गले}\ A N}}\skt{ }\edtext{\skt{च}}{\Afootnote{H; \skt{अनु}\ A; \skt{न च}\ N}}\skt{ }\edtext{\skt{शव्दङ्कुरुते}}{\lemma{\skt{शव्दङ्}}\Afootnote{N H; \skt{शब्दं}\ A}}\skt{ च }\edtext{\skt{नित्यं}}{\Afootnote{A H; \skt{नित्य ॥}\ N}}\skt{ ।}\edtext{}{\lemma{\skt{विवर्द्धते… दे}}\Afootnote{\emph{om.}\ A}}\pend

\pstart\versepar \ledsidenote{\passref{2.11.28}}\edtext{\skt{कृच्छ्रो\-च्छ्व\-स\-न्त\-म्मृ\-दु\-स\-र्व्व\-गात्रं}}{\lemma{\skt{कृच्छ्रोच्छ्वसन्तम्}}\Afootnote{H; \skt{कृच्छ्राच्छ्वसन्तं}\ A N}}\skt{ सम्वत्सराती }\edtext{\skt{तम\-रो\-च\-का\-र्त्तम्}}{\lemma{\skt{अरोचकार्त्तम्}}\Afootnote{A H; \skt{अरोचकार्त ।}\ N}}\skt{ ॥ }\pend
\pstart\versepar \skt{}\edtext{\skt{क्षीणञ्च}}{\lemma{\skt{च}}\Afootnote{A H; \skt{चैव}\ N}}\skt{ वैद्यो }\edtext{\skt{गलगण्डिनन्तु}}{\Afootnote{A H; \skt{गलगण्डितं तु}\ N}}\skt{ }\edtext{\skt{भिन्नस्वरञ्चापि}}{\lemma{\skt{चापि}}\Afootnote{N H; \skt{चैव}\ A}}\skt{ }\edtext{\skt{विवर्ज्जयीत}}{\Afootnote{H; \skt{विवर्जयेत्तु ॥}\ A; \skt{विवर्ज्जजीति ॥}\ N}}\skt{ ॥}\pend

\pstart\versepar \ledsidenote{\passref{2.11.29}}\edtext{\edtext{\skt{निव\-द्ध\-श्व\-य\-थु\-र्यस्य}}{\lemma{\skt{निव\-द्ध\-श्व\-य\-थुर्}}\Afootnote{N H; \skt{निबद्धः श्वयथु}\ A}}}{\lemma{\skt{यस्य}}\Afootnote{N H; \skt{र्यस्य}\ A}}\skt{ }\edtext{\skt{यथा}}{\Afootnote{N H; \emph{om.}\ A}}\skt{ }\edtext{\edtext{\skt{मुष्कम्प्रलम्वते}}{\lemma{\skt{मुष्कम्}}\Afootnote{H; \skt{मुष्क}\ A N}}}{\lemma{\skt{प्रलम्वते}}\Afootnote{N H; \skt{वल्लम्बते}\ A}}\skt{ । }\pend
\pstart\versepar \skt{}\edtext{\skt{महा\-नु\-वो\-ध्य\-थवा ह्रस्वस्तं}}{\lemma{\skt{अथवा ह्रस्वस्तं}}\Afootnote{H\textsuperscript{pc}; \skt{अथ ह्रस्वस्तं}\ H\textsuperscript{ac}}}\skt{ गण्डमिति निर्दिशेत् ॥ }\pend
\pstart\versepar \skt{इति ग्रन्थ्य\-प\-च्य\-र्व्वु\-द\-ग\-ल\-ग\-ण्ड\-नि\-दान एकादश ॥ ११ ॥}\edtext{}{\lemma{\skt{महानुवोध्य्… इति}}\Afootnote{\emph{om.}\ A}}\edtext{}{\lemma{\skt{महानुवोध्य्… ग्रन्}}\Afootnote{\emph{om.}\ N}}\edtext{}{\lemma{\skt{निर्दिशेत् ॥ इति ग्रन्थ्य\-प\-च्य\-र्व्वु\-द\-ग\-ल\-ग\-ण्ड\-नि\-दान एकादश}}\Afootnote{\skt{निर्दिशेत् ॥}\ H}}\pend

\endnumbering
\beginnumbering
\eledsection*{Adhyāya 12}
\pstart \ledsidenote{\passref{2.12.1}}\skt{अथातो }\edtext{\edtext{\skt{वृद्ध्यु\-प\-दं\-श\-श्ली\-प\-दा\-ना\-न्नि\-दा\-न\-म्व्या\-ख्या\-स्यामः}}{\lemma{\skt{वृद्ध्यु\-प\-दं\-श\-श्ली\-प\-दा\-नान्}}\Afootnote{A H; \skt{वृद्ध्यु\-प\-दं\-स\-श्ली\-प\-दा\-नान्}\ N}}}{\lemma{\skt{व्याख्यास्यामः}}\Afootnote{A H; \skt{व्याव्याख्यास्यामः ॥}\ N}}\skt{ ॥}\pend

\pstart \ledsidenote{\passref{2.12.3}}\edtext{\skt{वात\-पि\-त्त\-श्ले\-ष्म\-शो\-णि\-त\-मे\-दो\-मू\-त्रा\-न्त्र\-नि\-मित्ताः}}{\Afootnote{A H; \skt{वात\-पि\-त्त\-श्ले\-ष्म\-शो\-णि\-त\-मे\-दो\-मू\-त्रा\-न्त\-नि\-मि\-त्तास्}\ N}}\skt{ सप्त वृद्धयो भवन्ति । }\edtext{\skt{तेषु}}{\Afootnote{N H; \skt{तासां}\ A}}\skt{ }\edtext{\skt{मूत्रान्त्रनिमित्तौ}}{\Afootnote{N H; \skt{मूत्रान्त्रनिमित्ते}\ A}}\skt{ }\edtext{\skt{वातसमुत्थौ}}{\Afootnote{N H; \skt{वृद्धी वातसमुत्थे}\ A}}\skt{ । }\edtext{\skt{केव\-ल\-मु\-त्प\-त्ति\-हे\-तु\-रन्यः}}{\lemma{\skt{अन्यः}}\Afootnote{N H; \skt{अन्यतमः ॥}\ A}}\skt{ ॥}\pend

\pstart \ledsidenote{\passref{2.12.4}}\skt{अधः }\edtext{\edtext{\skt{कुपि\-तोऽ\-न्य\-त\-मस्तु}}{\lemma{\skt{कुपितो}}\Afootnote{H; \skt{प्रकुपितो}\ A N}}}{\lemma{\skt{ऽन्यतमस्तु}}\Afootnote{H; \skt{ऽन्यतमो हि}\ A; \skt{न्यतमस्तु}\ N}}\skt{ }\edtext{\edtext{\skt{दोष\-फ\-ल\-को\-श\-यो\-र्व्वा\-त\-वा\-हि\-नी\-ध\-म\-नी\-म\-भि\-प्र\-पद्य}}{\lemma{\skt{दोष\-फ\-ल\-को\-श\-योर्}}\Afootnote{H; \skt{दोषः फलकोश}\ A; \skt{दोषः । फलकोशयोर्}\ N}}}{\lemma{\skt{व्वात\-वा\-हि\-नी\-ध\-म\-नीम्}}\Afootnote{H; \skt{वाहिनीर्}\ A; \skt{वात\-वा\-हि\-णी\-ध\-म\-नीम्}\ N}}\skt{ }\edtext{\edtext{\skt{फल\-को\-ष\-यो\-र्व्वृ\-द्धि\-मु\-प\-ज\-न\-यति}}{\lemma{\skt{फलकोषयोर्}}\Afootnote{N H; \skt{धमनीः फलकोषयोर्}\ A}}\skt{ ॥ तम्वृ\-द्धि\-मि\-त्या\-च\-क्षते}}{\lemma{\skt{उपजनयति ॥ तम्}}\Afootnote{N H; \skt{जनयति तां}\ A}}\skt{ ॥}\pend

\pstart \ledsidenote{\passref{2.12.6}}\edtext{\edtext{\edtext{\edtext{\edtext{\skt{तत्रा\-नि\-ल\-पू\-र्ण्ण\-ब\-स्ति\-मि\-वा\-त\-त\-म्पु\-रु\-ष\-म\-नि\-मि\-त्त\-रु\-ज\-म्वा\-त\-वृ\-द्धि\-मि\-त्या\-च\-क्षते}}{\lemma{\skt{तत्रा\-नि\-ल\-पू\-र्ण्ण\-ब\-स्तिम्}}\Afootnote{N H; \skt{तत्रा\-नि\-ल\-प\-रि\-पूर्णां बस्तिम्}\ A}}}{\lemma{\skt{इवाततम्}}\Afootnote{N H; \skt{इवाततां}\ A}}}{\lemma{\skt{पुरुषम्}}\Afootnote{N H; \skt{परुषाम्}\ A}}}{\lemma{\skt{अनि\-मि\-त्त\-रु\-जम्}}\Afootnote{N H; \skt{अनि\-मि\-त्ता\-नि\-ल\-रुजां}\ A}}}{\lemma{\skt{इत्य्}}\Afootnote{N H; \emph{om.}\ A}}\skt{ । }\edtext{\edtext{\skt{पक्वो\-दु\-म्ब\-र\-स\-ङ्का\-श\-माशु}}{\lemma{\skt{पक्वो\-दु\-म्ब\-र\-स\-ङ्का\-शम्}}\Afootnote{H; \skt{पक्वो\-दु\-म्ब\-र\-स\-ङ्काशां}\ A; \skt{पक्वो\-दु\-म्ब\-र\-संका}\ N}}}{\lemma{\skt{आशु}}\Afootnote{H; \skt{ज्वर\-दा\-हो\-ष्म\-वतीं चाशु}\ A; \skt{समाशु}\ N}}\skt{ }\edtext{\edtext{\skt{समु\-त्था\-न\-म्पि\-त्त\-वृद्धिं}}{\lemma{\skt{समुत्थानम्}}\Afootnote{N H; \skt{समुत्थान}\ A}}}{\lemma{\skt{पित्तवृद्धिं}}\Afootnote{N H; \skt{पाकां पित्तवृद्धिं;}\ A}}\skt{ ॥ }\edtext{\skt{कठिनं}}{\Afootnote{N H; \skt{कठि}\ A}}\skt{ }\edtext{\edtext{\skt{स्निग्ध\-म\-ल्प\-वे\-दनं}}{\lemma{\skt{स्निग्धम्}}\Afootnote{H; \emph{om.}\ A; \skt{स्निग्ध}\ N}}}{\lemma{\skt{अल्पवेदनं}}\Afootnote{H; \skt{नामल्पवेदनां शीतां कण्}\ A; \skt{कण्डु\-म\-ल्प\-वे\-दनं}\ N}}\skt{ }\edtext{\skt{श्लेष्मवृद्धिं}}{\Afootnote{N H; \skt{डूमतीं श्लेष्मवृद्धिं;}\ A}}\skt{ ॥ }\edtext{\skt{कृष्णस्फोटावृतं}}{\Afootnote{H; \skt{कृष्णस्फोटावृतां}\ A; \skt{कृष्णास्फोटावृतं}\ N}}\skt{ }\edtext{\skt{पित्तवृद्धिलिङ्गं}}{\Afootnote{H; \skt{पित्तवृद्धिलिङ्गां}\ A; \skt{पित्त\-वृ\-इ\-द्धि\-लिङ्गं}\ N}}\skt{ रक्तवृद्धिं ॥ }\edtext{\skt{मृदुस्निग्धं}}{\Afootnote{N H; \skt{मृदुस्निग्धां}\ A}}\skt{ }\edtext{\edtext{\edtext{\skt{सक\-ण्डु\-र\-ल्प\-वे\-द\-न\-न्ता\-ल\-फ\-ल\-प्र\-का\-श\-म्मे\-दो\-वृद्धिं}}{\lemma{\skt{सकण्डुर्}}\Afootnote{N H; \skt{कण्डूमतीम्}\ A}}}{\lemma{\skt{अल्पवेदनन्}}\Afootnote{N H; \skt{अल्पवेदनां}\ A}}}{\lemma{\skt{ताल\-फ\-ल\-प्र\-का\-शम्}}\Afootnote{N H; \skt{ताल\-फ\-ल\-प्र\-काशां}\ A}}\skt{ ॥ }\edtext{\skt{मूत्र\-स\-न्धा\-र\-ण\-शी\-लस्य}}{\Afootnote{A H; \skt{मूत्र\-स\-न्धा\-र\-ण\-सी\-लस्य}\ N}}\skt{ }\edtext{\skt{मूत्र\-वृ\-द्धि\-र्ब्भ\-वति ॥ स}}{\lemma{\skt{ब्भवति ॥ स}}\Afootnote{N H; \skt{भवति, सा}\ A}}\skt{ }\edtext{\skt{गच्छतोऽम्बुपूर्ण्णा}}{\lemma{\skt{ऽम्बुपूर्ण्णा}}\Afootnote{A H; \skt{बुपूर्ण्ण}\ N}}\skt{ दृतिरिव }\edtext{\skt{क्षुभ्यते}}{\Afootnote{N H; \skt{क्षुभ्यति}\ A}}\skt{ }\edtext{\skt{मूत्र\-कृ\-च्छ्र\-वे\-द\-ना\-वन्तं}}{\Afootnote{H; \skt{मूत्र\-कृ\-च्छ्र\-वे\-दनां वृषणयोः श्वयथुं कोश\-यो\-श्चा\-पा\-द\-यति, तां}\ A; \skt{मूत्र\-कृ\-च्छ्र\-वे\-द\-ना\-वंत्तं}\ N}}\skt{ }\edtext{\edtext{\edtext{\skt{मूत्र\-वृ\-द्धि\-मि\-त्या\-च\-क्षते}}{\lemma{\skt{मूत्रवृद्धिम्}}\Afootnote{A H; \skt{मूत्रवृद्धिर्}\ N}}}{\lemma{\skt{इत्य्}}\Afootnote{N H; \skt{विद्या}\ A}}}{\lemma{\skt{आचक्षते}}\Afootnote{N H; \skt{त्; भारह}\ A}}\skt{ ॥ }\edtext{\skt{तत्र}}{\Afootnote{N H; \skt{रण}\ A}}\skt{ }\edtext{\edtext{\skt{बल\-व\-द्वि\-ग्र\-हा\-दि\-भि\-र्व्वि\-शे\-षै\-र्व्वायुः}}{\lemma{\skt{बल\-व\-द्वि\-ग्र\-हा\-दि\-भिर्}}\Afootnote{N H; \skt{बल\-व\-द्वि\-ग्र\-ह\-वृ\-क्ष\-प्र\-प\-त\-ना\-दि\-भिर्}\ A}}}{\lemma{\skt{व्विशेषैर्}}\Afootnote{N H; \skt{आया\-स\-वि\-शे\-षैर्}\ A}}\skt{ }\edtext{\skt{प्रकुपितो}}{\Afootnote{N H; \skt{अभिप्रवृद्धः प्रकुपितश्}\ A}}\skt{ }\edtext{\skt{न्त्रस्य}}{\Afootnote{N H; \skt{च}\ A}}\skt{ }\edtext{\skt{स्थूलस्य}}{\Afootnote{N H; \skt{स्थूलान्त्र}\ A}}\skt{ }\edtext{\skt{चेत\-र\-स्यै\-क\-द\-श\-म्वि\-गु\-ण\-मा\-दा\-याधो}}{\lemma{\skt{चेत\-र\-स्यै\-क\-द\-शम्}}\Afootnote{H; \skt{स्येतरस्य चैकदेशं}\ A; \skt{चेत\-र\-स्यै\-क\-दे\-शम्}\ N}}\skt{ }\edtext{\edtext{\skt{वंक्ष\-ण\-स\-न्धि\-मा\-गम्य}}{\lemma{\skt{वंक्षणसन्धिम्}}\Afootnote{N H; \skt{गत्वा वङ्क्षणसन्धिम्}\ A}}}{\lemma{\skt{आगम्य}}\Afootnote{N H; \skt{उपेत्य ग्रन्थिरूपेण स्थित्वाऽप्र}\ A}}\skt{ }\edtext{\skt{कालान्तरेण}}{\Afootnote{N H; \skt{तिक्रियमाणे च कालान्तरेण}\ A}}\skt{ }\edtext{\edtext{\skt{मुष्क\-को\-ष\-मु\-पैति}}{\lemma{\skt{मुष्ककोषम्}}\Afootnote{H; \skt{फलकोशं प्रविश्य मुष्कशोफम्}\ A; \skt{मुश्ककोशम्}\ N}}}{\lemma{\skt{उपैति}}\Afootnote{N H; \skt{आपादयति,}\ A}}\skt{ }\edtext{\skt{आध्मा\-त\-ब\-स्ति\-रि\-वा\-ततः}}{\lemma{\skt{आध्मातबस्तिर्}}\Afootnote{N H; \skt{आध्मातो बस्तिर्}\ A}}\skt{ }\edtext{\skt{प्रदीर्घः}}{\Afootnote{N H; \skt{प्रदीर्घः स}\ A}}\skt{ शोफो भवति । }\edtext{\edtext{\edtext{\edtext{\skt{सश\-ब्द\-म\-नु\-पी\-डि\-त\-श्चो\-र्द्ध्व\-मु\-त्प\-तति}}{\lemma{\skt{सशब्दम्}}\Afootnote{H; \skt{सशब्द}\ A; \skt{ससब्दम्}\ N}}}{\lemma{\skt{अनुपीडितश्}}\Afootnote{N H; \skt{मवपीडित}\ A}}}{\lemma{\skt{चोर्द्ध्वम्}}\Afootnote{N H; \skt{श्चोर्ध्वम्}\ A}}}{\lemma{\skt{उत्पतति}}\Afootnote{N H; \skt{उपैति,}\ A}}\skt{ विमुक्तश्च }\edtext{\skt{पुनराध्माति । }\edtext{\skt{तम\-न्त्र\-वृ\-द्धि\-म\-सा\-ध्य\-मि\-त्या\-च\-क्षते}}{\lemma{\skt{असाध्यम्}}\Afootnote{N H; \skt{असाध्याम्}\ A}}}{\lemma{\skt{आध्माति । तम्}}\Afootnote{N H; \skt{आध्मायते, ताम्}\ A}}\skt{ ॥}\pend

\pstart \ledsidenote{\passref{2.12.7}}\edtext{\skt{तत्र}}{\Afootnote{N H; \skt{तत्रा\-ति\-मै\-थु\-ना\-द\-ति\-ब्र\-ह्म\-च\-र्याद्वा तथाऽति}\ A}}\skt{ }\edtext{\edtext{\skt{ब्रह्म\-चा\-रि\-णी\-न्दी\-र्घ\-क\-र्क्क\-श\-रोमां}}{\lemma{\skt{ब्रह्मचारिणीन्}}\Afootnote{A H; \skt{ब्रह्मचारि}\ N}}}{\lemma{\skt{दीर्घ\-क\-र्क्क\-श\-रोमां}}\Afootnote{H; \skt{चिरोत्सृष्टां रजस्वलां दीर्घरोमां कर्कशरोमां}\ A; \skt{णी दीर्घकर्कशरो}\ N}}\skt{ }\edtext{\edtext{\edtext{\skt{योनि\-रो\-गो\-प\-सृ\-ष्टा\-म\-ल्प\-द्वा\-रा\-म\-प्रि\-या\-म\-का\-मा\-म\-चौ\-क्ष\-स\-लि\-ल\-प्र\-क्षा\-लि\-त\-यो\-निम्वा}}{\lemma{\skt{योनि\-रो\-गो\-प\-सृ\-ष्टाम्}}\Afootnote{H; \skt{सङ्कीर्णरोमां निगूढरोमाम्}\ A; \skt{माणो\-नि\-रो\-गो\-प\-सृ\-ष्टाम्}\ N}}}{\lemma{\skt{अल्पद्वाराम्}}\Afootnote{N H; \skt{अल्पद्वारां महाद्वाराम्}\ A}}}{\lemma{\skt{अचौ\-क्ष\-स\-लि\-ल\-प्र\-क्षा\-लि\-त\-यो\-निम्वा}}\Afootnote{H; \skt{अचौ\-क्ष\-स\-लि\-ल\-प्र\-क्षा\-लि\-त\-यो\-नि\-म\-प्र\-क्षा\-लि\-त\-योनिं योनि\-रो\-गो\-प\-सृष्टां स्वभावतो वा दुष्टयोनिं वियोनिं वा}\ A; \skt{अचो\-क्ष\-स\-लि\-ल\-प्र\-क्षा\-लि\-त\-यो\-निम्}\ N}}\skt{ }\edtext{\skt{नारीं }\edtext{\skt{योऽधिगच्छेत्तथा}}{\lemma{\skt{ऽधिगच्छेत्तथा}}\Afootnote{H; \skt{त्यर्थ\-मु\-प\-से\-व\-मा\-नस्य तथा}\ A; \skt{वा नारी यो धिगच्छेत्तथा}\ N}}}{\lemma{\skt{नारीं यो}}\Afootnote{H; \skt{नारीम}\ A; \emph{om.}\ N}}\skt{ }\edtext{\skt{कर\-ज\-द\-श\-न\-स\-वि\-ष\-शू\-क\-नि\-पा\-त\-ना\-द्ध\-स्ता\-भि\-घा\-ता\-च्च\-तु\-ष्प\-द\-ग\-म\-ना\-द\-चौ\-क्ष\-स\-लि\-ल\-प्र\-क्षा\-ल\-नाद्वा प्रकुपिता मेढ्रगतास्तु}}{\Afootnote{H; \skt{कर\-ज\-द\-श\-न\-वि\-ष\-शू\-क\-नि\-पा\-त\-ना\-द्ब\-न्ध\-ना\-द्ध\-स्ता\-भि\-घा\-ता\-च्च\-तु\-ष्प\-दी\-ग\-म\-ना\-द\-चौ\-क्ष\-स\-लि\-ल\-प्र\-क्षा\-ल\-ना\-द\-व\-पी\-ड\-ना\-च्छ्रु\-क्र\-वे\-ग\-वि\-घा\-र\-णा\-न्मै\-थु\-नान्ते वा+अ\-प्र\-क्षा\-ल\-ना\-दि\-भि\-र्मे\-ढ्र\-मा\-गम्य प्रकुपिता}\ A; \skt{कर\-ज\-द\-श\-न\-स\-वि\-ष\-शू\-क\-नि\-पा\-त\-ना\-द्ध\-स्ता\-भि\-घा\-ता\-च्च\-तु\-ष्प\-द\-ग\-म\-ना\-द\-चौ\-क्ष\-स\-लि\-ल\-प्र\-क्षा\-ल\-नाद्वा प्रकुपितास्तु}\ N}}\skt{ }\edtext{\skt{दोषाः}}{\Afootnote{A N H\textsuperscript{pc}; \skt{दोषा}\ H\textsuperscript{ac}}}\skt{ क्षते }\edtext{\skt{क्षते वा}}{\Afootnote{N H; \skt{ऽक्षते वा}\ A}}\skt{ }\edtext{\skt{श्वय\-थु\-मु\-प\-ज\-न\-यन्ति}}{\lemma{\skt{श्वयथुम्}}\Afootnote{A H; \skt{स्वयथुम्}\ N}}\skt{ । }\edtext{\skt{तमु\-प\-दं\-श\-मि\-त्या\-च\-क्षते}}{\lemma{\skt{उपदंशम्}}\Afootnote{A H; \skt{उपदंसम्}\ N}}\skt{ ॥}\pend

\pstart \ledsidenote{\passref{2.12.8}}\skt{स }\edtext{\skt{पञ्च\-वि\-ध\-स्त्रि\-भि\-र्द्दोषैः}}{\lemma{\skt{त्रिभिर्}}\Afootnote{A H; \skt{त्रिभि}\ N}}\skt{ }\edtext{\skt{पृथ\-क्स\-म\-स्तै\-र\-सृ\-जश्च}}{\lemma{\skt{असृजश्च}}\Afootnote{N H; \skt{असृजा चेति ॥}\ A}}\skt{ ।}\pend

\pstart \ledsidenote{\passref{2.12.9}}\skt{तत्र वातिके }\edtext{\skt{त्वक्परिमोटनं}}{\Afootnote{N; \skt{पारुष्यं त्वक्परिपुटनं}\ A; \skt{त्वक्परि पा टनं}\ H\textsuperscript{ac}; \skt{त्वक्परि मो टनं}\ H\textsuperscript{pc}}}\skt{ स्तब्धमेढ्रता परुषशोफता }\edtext{\skt{विवि\-ध\-वे\-द\-ना\-प्रा\-दु\-र्भा\-वश्च}}{\Afootnote{N H; \skt{विविधाश्च वातवेदनाः;}\ A}}\skt{ ॥ पैत्तिके }\edtext{\skt{श्वय\-थु\-रु\-दु\-म्ब\-र\-प्र\-ति\-काशो}}{\Afootnote{H; \skt{ज्वरः श्वयथुः पक्वो\-डु\-म्ब\-र\-स\-ङ्का\-शस्}\ A; \skt{स्वय\-थु\-रु\-दु\-म्ब\-र\-प्र\-ति\-काशो}\ N}}\skt{ }\edtext{\skt{ज्वरदाहयुक्तः}}{\Afootnote{N H; \skt{तीव्रदाहः क्षिप्रपाकः पित्तवेदनाश्च;}\ A}}\skt{ ॥ श्लैष्मिके }\edtext{\skt{श्वयथुः}}{\Afootnote{A H; \skt{स्वयथुः}\ N}}\skt{ }\edtext{\skt{कण्डूमान्कठिनः}}{\Afootnote{A H; \skt{कण्डून्मां कठिनः}\ N}}\skt{ }\edtext{\skt{स्निग्धश्च ॥ रक्तजे कृष्ण\-स्फो\-ट\-प्रा\-दु\-र्ब्भा\-वोऽ\-त्य\-र्थ\-म\-सृ\-क्प्र\-वृत्तिः}}{\lemma{\skt{च ॥ रक्तजे कृष्ण\-स्फो\-ट\-प्रा\-दु\-र्ब्भावो}}\Afootnote{H; \skt{श्लेष्मवेदनाश्च; रक्तजे कृष्ण\-स्फो\-ट\-प्रा\-दु\-र्भावो}\ A; \skt{च रक्तजे मेढ्रगतो कृष्ण\-स्फो\-ट\-प्रा\-दु\-र्ब्भावो}\ N}}\skt{ }\edtext{\skt{पित्तलिङ्गानि}}{\Afootnote{A H; \skt{पित्तालिङ्गानि}\ N}}\skt{ }\edtext{\skt{च}}{\Afootnote{N H; \skt{अत्यर्थं ज्वरदाहौ शोषश्च, याप्यश्चैव कदाचि}\ A}}\skt{ ॥ }\edtext{\skt{सर्व्वजे}}{\Afootnote{N H; \skt{त्; सर्वजे}\ A}}\skt{ सर्व्व\-लि\-ङ्ग\-द\-र्श\-न\-म\-व\-द\-र\-णञ्च }\edtext{\skt{शेफसः}}{\Afootnote{A H; \skt{सेफसः ॥}\ N}}\skt{ । }\edtext{\skt{क्रिमिप्रादुर्ब्भावो}}{\Afootnote{N\textsuperscript{pc} H; \skt{कृमिप्रादुर्भावो}\ A; \skt{क्रिमि दा दुर्भावो}\ N\textsuperscript{ac}}}\skt{ मरणञ्चेति ॥}\pend

\pstart \ledsidenote{\passref{2.12.10}}\edtext{\skt{प्रकुपितास्तु}}{\Afootnote{N H; \skt{कुपितास्तु}\ A}}\skt{ }\edtext{\edtext{\skt{दोषा\-स्त्व\-धः\-प्र\-पन्ना}}{\lemma{\skt{दोषास्}}\Afootnote{N H; \skt{दोषा}\ A}}}{\lemma{\skt{त्वधःप्रपन्ना}}\Afootnote{N H; \skt{वात\-पि\-त्त\-श्ले\-ष्मा\-णोऽ\-धः\-प्र\-पन्ना}\ A}}\skt{ }\edtext{\edtext{\skt{वंक्ष\-णो\-रु\-जं\-घा\-स्व\-व\-ति\-ष्ठन्ते}}{\lemma{\skt{वंक्ष\-णो\-रु\-जं\-घास्व्}}\Afootnote{N H; \skt{वङ्क्ष\-णो\-रु\-जा\-नु\-ज\-ङ्घास्व्}\ A}}}{\lemma{\skt{अवतिष्ठन्ते}}\Afootnote{N H; \skt{अवतिष्ठ}\ A}}\skt{ । }\edtext{\skt{ततः}}{\Afootnote{H; \skt{मानाः}\ A; \skt{तत्र}\ N}}\skt{ कालान्तरेण }\edtext{\skt{पादमाश्रित्य}}{\lemma{\skt{आश्रित्य}}\Afootnote{N H; \skt{आश्रित्य्}\ A}}\skt{ }\edtext{\skt{शनैः}}{\Afootnote{N H; \skt{अ}\ A}}\skt{ शनैः }\edtext{\skt{शोफ\-मु\-प\-ज\-न\-यन्ति}}{\lemma{\skt{उपजनयन्ति}}\Afootnote{N H; \skt{जनयन्ति,}\ A}}\skt{ । }\edtext{\skt{तच्छ्ली\-प\-द\-मि\-त्या\-च\-क्षते}}{\lemma{\skt{तच्छ्लीपदम्}}\Afootnote{N H; \skt{तं श्लीपदम्}\ A}}\skt{ ॥ }\edtext{\skt{तत्तु}}{\Afootnote{H; \skt{त}\ A; \skt{तं तु}\ N}}\skt{ त्रिवि\-ध\-म्वा\-त\-पि\-त्त\-क\-फ\-नि\-मि\-त्त\-मिति ॥}\pend

\pstart \ledsidenote{\passref{2.12.11}}\skt{तत्र वातजं }\edtext{\edtext{\skt{कृष्णम्परुषं}}{\lemma{\skt{कृष्णम्}}\Afootnote{N H; \skt{खरं कृष्णं}\ A}}}{\lemma{\skt{परुषं}}\Afootnote{N H; \skt{परुष}\ A}}\skt{ }\edtext{\edtext{\edtext{\skt{खर\-म\-नि\-मि\-त्त\-रु\-ज\-म\-न\-ल्प\-वे\-द\-न\-म्प\-रि\-स्फु\-टति}}{\lemma{\skt{खरम्}}\Afootnote{N H; \skt{म}\ A}}}{\lemma{\skt{अनि\-मि\-त्त\-रु\-जम्}}\Afootnote{N H; \skt{निमि\-त्ता\-नि\-ल\-रुजं}\ A}}}{\lemma{\skt{अन\-ल्प\-वे\-द\-नम्}}\Afootnote{H; \emph{om.}\ A; \skt{अनल्पवेदन}\ N}}\skt{ च बहुशः ॥ पित्तजन्तु }\edtext{\skt{पीतावभासं}}{\Afootnote{N H; \skt{पीता\-व\-भा\-स\-मी\-षन्}\ A}}\skt{ }\edtext{\skt{मृदु\-ज्व\-र\-दा\-ह\-प्र\-दञ्च}}{\Afootnote{N H; \skt{मृदु ज्वरदाहप्रायं च;}\ A}}\skt{ ॥ श्लेष्मजन्तु }\edtext{\skt{स्निग्धं}}{\Afootnote{N H; \skt{श्वेतं}\ A}}\skt{ }\edtext{\skt{स्वेतावभासं}}{\Afootnote{N H; \skt{स्निग्धावभासं}\ A}}\skt{ }\edtext{\edtext{\skt{मवे\-द\-न\-म\-ति\-म\-हान्तं}}{\lemma{\skt{मवेदनम्}}\Afootnote{H; \skt{मन्दवेदानं}\ A; \skt{अवेदनम्}\ N}}}{\lemma{\skt{अतिमहान्तं}}\Afootnote{H; \skt{भारिकं महा}\ A; \skt{अतिमहान्त्र}\ N}}\skt{ }\edtext{\skt{ग्रन्थि\-क\-ण्ट\-कै\-रु\-प\-चि\-तञ्च}}{\lemma{\skt{ग्रन्थिकण्टकैर्}}\Afootnote{N H; \skt{ग्रन्थिकं कण्टकैर्}\ A}}\skt{ ॥}\pend

\pstart \ledsidenote{\passref{2.12.12}}\skt{तत्र }\edtext{\edtext{\skt{सम्व\-त्स\-रा\-ती\-त\-म\-ति\-म\-हन्तं}}{\lemma{\skt{सम्व\-त्स\-रा\-ती\-तम्}}\Afootnote{N H; \skt{संवत्सरातीत}\ A}}}{\lemma{\skt{अतिमहन्तं}}\Afootnote{N H; \skt{मतिमह}\ A}}\skt{ }\edtext{\edtext{\skt{वल्मी\-क\-जा\-त\-म\-ति\-प्र\-स्रु\-त\-मिति}}{\lemma{\skt{वल्मीकजातम्}}\Afootnote{N H; \skt{द्वल्मीकजातं}\ A}}}{\lemma{\skt{अतिप्रस्रुतम्}}\Afootnote{N H; \skt{प्रसृतम्}\ A}}\skt{ वर्जनीयानि ॥}\pend

\pstart\versepar \ledsidenote{\passref{2.12.13}}\edtext{\skt{त्रीण्यप्येतानि}}{\lemma{\skt{त्रीण्य्}}\Afootnote{N H; \skt{भवन्ति चात्र त्रीण्य्}\ A}}\skt{ }\edtext{\skt{जानी\-या\-च्छ्ली\-प\-दानि}}{\Afootnote{A N; \skt{जानीयाच्छ्लीपदी नि}\ H\textsuperscript{ac}; \skt{जानीयाच्छ्लीपदा नि}\ H\textsuperscript{pc}}}\skt{ कफोच्छ्रयात् । }\pend
\pstart\versepar \skt{गुरुत्वञ्च महत्त्वञ्च यस्मान्नास्ति विना कफात् ॥}\pend

\pstart\versepar \ledsidenote{\passref{2.12.14}}\edtext{\skt{पुरा\-णो\-द\-क\-भू\-यिष्ठाः}}{\Afootnote{A H; \skt{पुरा\-णो\-द\-क\-भू\-यिष्ठा}\ N}}\skt{ सर्व्वर्त्तुषु च शीतलाः । }\pend
\pstart\versepar \skt{ये देशास्तेषु जायन्ते श्लीपदानि विशेषतः ॥}\pend

\pstart\versepar \ledsidenote{\passref{2.12.15}}\skt{पाद\-व\-द्ध\-स्त\-यो\-श्चापि श्लीपदञ्जायते नृणाम् । }\pend
\pstart\versepar \skt{}\edtext{\edtext{\skt{कर्णा\-क्षि\-ना\-सा\-स्वपि}}{\lemma{\skt{कर्णाक्षिनासास्व्}}\Afootnote{N H; \skt{कर्णाक्षिनासिकौ}\ A}}\skt{ च}}{\lemma{\skt{अपि च}}\Afootnote{N H; \skt{ष्ठेषु}\ A}}\skt{ केचिदिच्छन्ति }\edtext{\skt{तद्विद}}{\Afootnote{N H; \skt{तद्विदः ॥}\ A}}\skt{ }\edtext{\skt{इति}}{\Afootnote{N H; \emph{om.}\ A}}\skt{ ॥}\pend

\pstart \ledsidenote{\passref{2.12.trailer}}\edtext{\edtext{\skt{वृद्ध्यु\-प\-द\-ङ्श\-श्ली\-प\-दा\-ना\-न्नि\-दाने}}{\lemma{\skt{वृद्ध्यु\-प\-द\-ङ्श\-श्ली\-प\-दा\-नान्}}\Afootnote{H; \skt{इति सुश्रु\-त\-सं\-हि\-तायां निदानस्थाने वृद्ध्यु\-प\-दं\-श\-श्ली\-प\-द\-निदा}\ A; \skt{नि ॥ १२ ॥ ० ॥}\ N}}}{\lemma{\skt{निदाने}}\Afootnote{H; \skt{नं नाम}\ A; \emph{om.}\ N}}\skt{ }\edtext{\skt{द्वादश}}{\Afootnote{H; \skt{द्वादशोऽध्यायः ॥१२ ॥}\ A; \emph{om.}\ N}}\skt{ ॥ ॥}\pend

\endnumbering
\beginnumbering
\eledsection*{Adhyāya 13}
\pstart \ledsidenote{\passref{2.13.1}}\skt{अथातः }\edtext{\skt{क्षुद्र\-रो\-ग\-नि\-दानं}}{\Afootnote{N H; \skt{क्षुद्ररोगाणां निदानं}\ A}}\skt{ व्याख्यास्यामः ॥}\pend

\pstart \ledsidenote{\passref{2.13.3}}\edtext{\skt{सप्त\-च\-त्वा\-रिं\-श\-त्क्षु\-द्र\-रोगा}}{\Afootnote{N H; \skt{समासेन चतु\-श्च\-त्वा\-रिं\-श\-त्क्षु\-द्र\-रोगा}\ A}}\skt{ भवन्ति ॥ तद्यथा ॥ अजगल्लिका । यवप्रख्या । }\edtext{\skt{अम्वालजी}}{\Afootnote{H; \skt{अन्धालजी}\ A; \skt{अन्धाजलजी ।}\ N}}\skt{ । विवृता । कच्छपिका । }\edtext{\skt{वल्मीका}}{\Afootnote{N H; \skt{वल्मीअकं}\ A}}\skt{ । }\edtext{\skt{इन्द्रविद्धा}}{\Afootnote{N H; \skt{इन्द्रवृद्धा}\ A}}\skt{ । }\edtext{\skt{गर्द्दभिका}}{\Afootnote{N H; \emph{om.}\ A}}\skt{ । पनसिका । }\edtext{\skt{पाषाणगर्द्दभः}}{\Afootnote{H; \skt{पाषाणगर्दभः जालगर्दभः}\ A; \skt{पाषाणगर्दभ ॥ जालगर्दभः ।}\ N}}\skt{ । }\edtext{\skt{इचिवल्लिका}}{\Afootnote{H; \emph{om.}\ A; \skt{इरिवेण्णिका ।}\ N}}\skt{ । }\edtext{\skt{कक्षा}}{\Afootnote{A H; \skt{कक्ष्या ।}\ N}}\skt{ । }\edtext{\skt{गन्धनामा}}{\Afootnote{N H; \emph{om.}\ A}}\skt{ । }\edtext{\skt{विस्फोटकं}}{\Afootnote{N H; \skt{विस्फोटकः}\ A}}\skt{ । अग्निरोहिणी । }\edtext{\skt{चिप्प}}{\Afootnote{H; \skt{चिप्पं}\ A; \skt{चि}\ N}}\skt{ । }\edtext{\skt{अनुशयी}}{\Afootnote{H; \skt{कुनखः अनुशयी}\ A; \skt{प्पा । अनुशयी ।}\ N}}\skt{ । विदारिका । }\edtext{\skt{शर्क्करा}}{\Afootnote{N H; \skt{श}\ A}}\skt{ । }\edtext{\skt{शर्करार्व्वुदं}}{\Afootnote{H; \skt{र्करार्बुदं}\ A; \skt{सर्क्करार्बुदं ।}\ N}}\skt{ । पामा । विचर्च्चिका । रकसा । }\edtext{\skt{पाददारी}}{\Afootnote{N H; \skt{पाददारिका}\ A}}\skt{ । कदरं । }\edtext{\skt{अलसकः}}{\Afootnote{N H; \skt{अलसेन्द्र}\ A}}\skt{ । }\edtext{\skt{रुजा}}{\Afootnote{N H; \skt{लुप्तौ}\ A}}\skt{ }\edtext{\skt{दारूणकं}}{\Afootnote{H; \skt{दारुणकः}\ A; \skt{दारुणाक ।}\ N}}\skt{ । }\edtext{\skt{अरूषिका}}{\Afootnote{H; \skt{अरुंषिका}\ A; \skt{अरुषिका ।}\ N}}\skt{ । पलितं । }\edtext{\skt{मयूरिका}}{\Afootnote{H; \skt{मसूरिका}\ A; \skt{मशूरिका ।}\ N}}\skt{ । }\edtext{\skt{युवानपिटका}}{\Afootnote{N H; \skt{यौवनपिडका}\ A}}\skt{ । }\edtext{\skt{पद्मिनीकण्टका}}{\Afootnote{H; \skt{पद्मिनीकण्टकः}\ A; \skt{पद्मिनीकण्ठका ।}\ N}}\skt{ । }\edtext{\skt{जतुमणि}}{\Afootnote{N H; \skt{जतुमणिः}\ A}}\skt{ । }\edtext{\skt{मसकं}}{\Afootnote{H; \skt{मशकः}\ A; \skt{म}\ N}}\skt{ । }\edtext{\skt{तिलका}}{\Afootnote{H; \skt{चर्मकीलः तिलका}\ A; \skt{शक । तिलका ।}\ N}}\skt{ । }\edtext{\skt{नच्छ}}{\Afootnote{N H; \skt{लकः न्यच्छं}\ A}}\skt{ । }\edtext{\skt{व्यंगः}}{\Afootnote{A H; \skt{व्यङ्ग ।}\ N}}\skt{ । }\edtext{\skt{नीलिका}}{\Afootnote{H; \emph{om.}\ A; \skt{नालिका ।}\ N}}\skt{ । }\edtext{\skt{परिवर्त्तिका}}{\Afootnote{N H; \skt{परिवर्तिका अवपाटिका}\ A}}\skt{ । }\edtext{\skt{निरुद्धप्रकाशं}}{\Afootnote{H; \skt{निरुद्धप्रकशः}\ A; \skt{निरुद्धप्रकशं ।}\ N}}\skt{ । }\edtext{\skt{अवपाटिका}}{\Afootnote{H; \emph{om.}\ A; \skt{अवपाठिका ।}\ N}}\skt{ । }\edtext{\skt{सन्निरुद्धगुदं}}{\Afootnote{N H; \skt{संनिरुद्धगुदः}\ A}}\skt{ । अहिपूतनं । }\edtext{\skt{वृषणकच्छूः}}{\Afootnote{H; \skt{वृषणकच्छ्रः}\ A; \skt{वृषणकच्छृ}\ N}}\skt{ । }\edtext{\skt{गुदभ्रन्सश्चेति}}{\lemma{\skt{गुदभ्रन्सश्}}\Afootnote{H; \skt{गुदभ्रंशश्}\ A; \skt{गुदं भ्रमश्}\ N}}\skt{ ॥}\pend

\pstart\versepar \ledsidenote{\passref{2.13.4}}\edtext{\skt{स्निग्धाः}}{\Afootnote{H; \skt{स्निग्धा}\ A; \skt{श्निग्धाः}\ N}}\skt{ }\edtext{\skt{सवर्ण्णा}}{\Afootnote{A H; \skt{सवर्ण्ण}\ N}}\skt{ ग्रथिता }\edtext{\skt{नीरुजा}}{\Afootnote{N H; \skt{नीरुजाअ}\ A}}\skt{ }\edtext{\skt{मुद्गसन्निताः}}{\Afootnote{H; \skt{मुद्रसन्निभा ।}\ A; \skt{मुद्गसम्मिता ।}\ N}}\skt{ । }\pend
\pstart\versepar \skt{कफवातोत्थिता ज्ञेया }\edtext{\edtext{\skt{वाला\-ना\-म\-ज\-ग\-ल्लिका}}{\lemma{\skt{वालानाम्}}\Afootnote{H; \skt{बालाना}\ A N}}}{\lemma{\skt{अजगल्लिका}}\Afootnote{H; \skt{मजगल्लिका ॥}\ A N}}\skt{ ॥}\pend

\pstart\versepar \ledsidenote{\passref{2.13.5}}\skt{यवाकारा सुकठिना ग्रथिता }\edtext{\skt{मांससंश्रिता}}{\Afootnote{A H; \skt{मांससंश्रिताः ।}\ N}}\skt{ । }\pend
\pstart\versepar \skt{}\edtext{\skt{पिडका}}{\Afootnote{A H; \skt{पिटका}\ N}}\skt{ श्लेष्मवाताभ्यां यवप्रख्येति }\edtext{\skt{चोच्यते}}{\Afootnote{N H; \skt{सोच्यते ॥}\ A}}\skt{ ॥}\pend

\pstart\versepar \ledsidenote{\passref{2.13.6}}\edtext{\edtext{\edtext{\skt{घना\-ल्प\-व\-क्रा\-म्पि\-ड\-का\-मु\-न्न\-ता\-म्प\-रि\-म\-ण्डलां}}{\lemma{\skt{घनाल्पवक्राम्}}\Afootnote{H; \skt{घनमवक्त्रां}\ A; \skt{घनाल्पवक्त्रम्}\ N}}}{\lemma{\skt{पिडकाम्}}\Afootnote{H; \skt{पिडका}\ A; \skt{पिटका}\ N}}}{\lemma{\skt{उन्नताम्}}\Afootnote{H; \skt{मुन्नतां}\ A N}}\skt{ । }\pend
\pstart\versepar \skt{}\edtext{\edtext{\edtext{\skt{अम्वा\-ल\-जी\-म\-ल्प\-पू\-या\-न्ता\-म्वि\-द्या\-त्क\-फ\-वा\-तजां}}{\lemma{\skt{अम्वालजीम्}}\Afootnote{H; \skt{अन्धालजी}\ A N}}}{\lemma{\skt{अल्पपूयान्ताम्}}\Afootnote{H; \skt{मल्पपूयां तां}\ A; \skt{मल्पपूया ताम्}\ N}}}{\lemma{\skt{विद्या\-त्क\-फ\-वा\-तजां}}\Afootnote{A H; \skt{विद्या\-त्क\-फ\-वा\-तजा ॥}\ N}}\skt{ ॥}\pend

\pstart\versepar \ledsidenote{\passref{2.13.7}}\edtext{\edtext{\skt{विवृ\-ता\-भ्या\-म्म\-हा\-दा\-हा\-म्प\-क्वो\-दु\-म्व\-र\-स\-न्निभां}}{\lemma{\skt{विवृताभ्याम्}}\Afootnote{H; \skt{विवृतास्यां}\ A N}}}{\lemma{\skt{पक्वो\-दु\-म्व\-र\-स\-न्निभां}}\Afootnote{H; \skt{पक्वो\-दु\-म्ब\-र\-स\-न्नि\-भाम्}\ A N}}\skt{ । }\pend
\pstart\versepar \skt{}\edtext{\edtext{\edtext{\edtext{\edtext{\skt{परि\-म\-ण्ड\-ला\-म्पि\-त्त\-कृ\-ता\-म्वि\-वृ\-ता\-न्ना\-स\-ता\-म्विदुः}}{\lemma{\skt{परिमण्डलाम्}}\Afootnote{N H; \skt{विवृतामिति तां विद्यात्पित्तोत्थां परिमण्डलाम् ॥}\ A}}}{\lemma{\skt{पित्तकृताम्}}\Afootnote{N H; \emph{om.}\ A}}}{\lemma{\skt{विवृतान्}}\Afootnote{N H; \emph{om.}\ A}}}{\lemma{\skt{नासताम्}}\Afootnote{H; \emph{om.}\ A; \skt{नामताम्}\ N}}}{\lemma{\skt{विदुः}}\Afootnote{N H; \emph{om.}\ A}}\skt{ ॥}\pend

\pstart\versepar \ledsidenote{\passref{2.13.8}}\edtext{\skt{ग्रथिता}}{\Afootnote{H; \skt{ग्रथिताः}\ A N}}\skt{ पञ्च वा षड्वा }\edtext{\skt{दारूणाः}}{\Afootnote{H; \skt{दारुणाः}\ A N}}\skt{ कच्छपोन्नताः । }\pend
\pstart\versepar \skt{}\edtext{\skt{कफा\-नि\-ला\-भ्या\-म्पि\-डका}}{\lemma{\skt{पिडका}}\Afootnote{A H; \skt{पिटका}\ N}}\skt{ ज्ञेया कच्छपिका }\edtext{\skt{वुधैः}}{\Afootnote{H; \skt{बुधैः ॥}\ A N}}\skt{ ॥}\pend

\pstart\versepar \ledsidenote{\passref{2.13.9}}\skt{पाणिपादतले सन्धौ }\edtext{\edtext{\skt{ग्रीवा\-या\-मू\-र्ध्व\-ज\-त्रुणि}}{\lemma{\skt{ग्रीवायाम्}}\Afootnote{N H; \skt{ग्रीवाया}\ A}}}{\lemma{\skt{ऊर्ध्वजत्रुणि}}\Afootnote{H; \skt{मूर्ध्वजत्रुणि}\ A N}}\skt{ । }\pend
\pstart\versepar \skt{}\edtext{\edtext{\skt{ग्रन्थि\-र्व्व\-ल्मी\-क\-व\-द्यस्य}}{\lemma{\skt{ग्रन्थिर्}}\Afootnote{H; \skt{ग्रन्थि}\ A N}}}{\lemma{\skt{व्वल्मीकवद्यस्य}}\Afootnote{H; \skt{र्वल्मीकवद्यस्तु}\ A; \skt{र्वल्मीकवद्यस्य}\ N}}\skt{ शनैः समुपचीयते ॥}\pend

\pstart\versepar \ledsidenote{\passref{2.13.10}}\edtext{\skt{तोदः}}{\Afootnote{H; \skt{तोद}\ A N}}\skt{ }\edtext{\skt{क्लेदपरीदाहः}}{\Afootnote{N H; \skt{क्लेदपरीदाह}\ A}}\skt{ कण्डू\-म\-द्भि\-र्म्मु\-खै\-र्व्वृतः । }\pend
\pstart\versepar \skt{}\edtext{\skt{व्याधिर्व्वल्मीक}}{\lemma{\skt{व्याधिर्}}\Afootnote{A H; \skt{व्याधि}\ N}}\skt{ इत्येष }\edtext{\skt{विज्ञेयः}}{\Afootnote{H; \emph{om.}\ A; \skt{विज्ञायः}\ N}}\skt{ }\edtext{\skt{कफपैत्तिकः}}{\Afootnote{N H; \skt{कफ\-पि\-त्ता\-नि\-लि\-द्भवः ॥}\ A}}\skt{ ॥}\pend

\pstart\versepar \ledsidenote{\passref{2.13.11}}\edtext{\skt{पद्म\-क\-र्ण्णि\-क\-व\-त्मध्ये}}{\Afootnote{H; \skt{पद्मा\-पु\-ष्क\-र\-व\-न्मध्ये}\ A; \skt{पद्म\-क\-र्ण्णि\-क\-व\-न्मध्ये}\ N}}\skt{ }\edtext{\skt{पिडकाभिः}}{\Afootnote{A H; \skt{पिटकाभिः}\ N}}\skt{ समाचितां । }\pend
\pstart\versepar \skt{}\edtext{\skt{इन्द्रा}}{\Afootnote{H; \skt{इन्द्र}\ A N}}\skt{ }\edtext{\edtext{\skt{विद्धा\-म्वि\-जा\-नी\-या\-द्वा\-त\-पि\-त्तो\-त्थि\-ता\-म्भि\-षक्}}{\lemma{\skt{विद्धाम्}}\Afootnote{H; \skt{वृद्धां}\ A; \skt{विद्वा}\ N}}}{\lemma{\skt{विजा\-नी\-या\-द्वा\-त\-पि\-त्तो\-त्थि\-ताम्}}\Afootnote{N H; \skt{तु तां विद्या\-द्वा\-त\-पि\-त्तो\-त्थितां}\ A}}\skt{ ॥}\pend

\pstart\versepar \ledsidenote{\passref{2.13.12}}\edtext{\edtext{\skt{मण्ड\-ल\-म्वृ\-त्त\-मु\-त्सन्नं}}{\lemma{\skt{वृत्तम्}}\Afootnote{H; \skt{वृत्त}\ A N}}}{\lemma{\skt{उत्सन्नं}}\Afootnote{H; \skt{मुत्सन्नं}\ A N}}\skt{ }\edtext{\skt{सर\-क्त\-म्पि\-ड\-का\-चितं}}{\lemma{\skt{पिडकाचितं}}\Afootnote{A H; \skt{पिटकाचितं ।}\ N}}\skt{ । }\pend
\pstart\versepar \skt{}\edtext{\edtext{\skt{रुजा\-क\-री\-ङ्ग\-र्द्द\-भिकां}}{\lemma{\skt{रुजाकरीङ्}}\Afootnote{A H; \skt{रुजाकरी}\ N}}\skt{ }\edtext{\skt{ताम्व\-दे\-द्वा\-त\-पि\-त्ततः}}{\lemma{\skt{वदे\-द्वा\-त\-पि\-त्ततः}}\Afootnote{H; \skt{विद्या\-द्वा\-त\-पि\-त्त\-जाम् ।}\ A; \skt{वदेद्वातपित्तः ।}\ N}}}{\lemma{\skt{गर्द्दभिकां ताम्}}\Afootnote{N H; \skt{गदर्भिकां तां}\ A}}\skt{ ॥ }\pend
\pstart\versepar \skt{}\edtext{\skt{कर्ण्णस्याभ्यन्तरे}}{\Afootnote{H; \skt{कर्णौ परि समन्ताद्वा}\ A; \skt{कर्ण्णस्याभ्यान्तरे}\ N}}\skt{ }\edtext{\edtext{\edtext{\skt{जाता\-म्पि\-ड\-का\-मु\-ग्र\-वे\-दनां}}{\lemma{\skt{जाताम्}}\Afootnote{H; \skt{पृष्ठे वा}\ A; \skt{जाता}\ N}}}{\lemma{\skt{पिडकाम्}}\Afootnote{H; \skt{पिडको}\ A; \skt{पिट्का}\ N}}}{\lemma{\skt{उग्रवेदनां}}\Afootnote{H; \emph{om.}\ A; \skt{मुग्रवेदनां ।}\ N}}\skt{ । }\pend
\pstart\versepar \skt{}\edtext{\edtext{\skt{स्थिरा\-म्पे\-न\-सि\-का\-न्तान्तु}}{\lemma{\skt{स्थिराम्}}\Afootnote{N H; \emph{om.}\ A}}}{\lemma{\skt{पेनसिकान्तान्तु}}\Afootnote{H; \skt{ग्ररुक् । शालू\-क\-व\-त्प\-न\-सिकां तां}\ A; \skt{परि\-न\-सि\-का\-न्तान्तु}\ N}}\skt{ }\edtext{\skt{विद्या\-द्वा\-त\-क\-फो\-त्थि\-ताम्}}{\Afootnote{H; \skt{विद्या\-च्छ्ले\-ष्म\-वा\-त\-जाम् ॥}\ A; \skt{विद्या\-द्वा\-त\-क\-फो\-त्थि\-तान् ॥}\ N}}\skt{ ॥}\pend

\pstart\versepar \ledsidenote{\passref{2.13.13}}\edtext{\skt{हन्वाः}}{\Afootnote{N H; \skt{हनु}\ A}}\skt{ }\edtext{\skt{सन्धिसमुद्भूतं}}{\Afootnote{N H; \skt{सन्धौ समुद्भूतं}\ A}}\skt{ }\edtext{\edtext{\skt{शोफमल्परुजं}}{\lemma{\skt{शोफम्}}\Afootnote{H; \skt{शोफ}\ A N}}}{\lemma{\skt{अल्परुजं}}\Afootnote{H; \skt{मल्परुजं}\ A N}}\skt{ स्थिरं । }\pend
\pstart\versepar \skt{}\edtext{\skt{पाषा\-ण\-ग\-र्द्द\-भ\-मिति}}{\lemma{\skt{इति}}\Afootnote{N H; \skt{विद्या}\ A}}\skt{ }\edtext{\skt{व्रूया\-त्त\-त्क\-फ\-पि\-त्ततः}}{\Afootnote{H; \skt{द्बला\-स\-प\-व\-ना\-त्म\-कम् ॥}\ A; \skt{ब्रूया\-त्क\-फ\-पि\-त्ततः ॥}\ N}}\skt{ ॥}\pend

\pstart\versepar \ledsidenote{\passref{2.13.14}}\skt{विस\-र्प्प\-व\-त्स\-र्प्पति यः }\edtext{\skt{शोफ\-स्त\-नु\-र\-पा\-क\-वान्}}{\Afootnote{H; \emph{om.}\ A; \skt{सोफ\-म्त\-तु\-र्पा\-क्वान् ।}\ N}}\skt{ । }\pend
\pstart\versepar \skt{}\edtext{\skt{दाहज्वरकरः}}{\Afootnote{N H; \skt{दाह\-ज्व\-र\-क\-र\-स्तनुः । अपाकः}\ A}}\skt{ }\edtext{\skt{पित्तात्सज्ञेयो}}{\Afootnote{N H; \skt{श्वयथुः पित्तात्स ज्ञेयो}\ A}}\skt{ जालगर्द्दभः ॥}\pend

\pstart\versepar \ledsidenote{\passref{2.13.15}}\edtext{\edtext{\edtext{\edtext{\skt{पिड\-का\-मु\-त्त\-मा\-ङ्ग\-स्था\-म्वृ\-त्ता\-मु\-ग्र\-रुजां}}{\lemma{\skt{पिडकाम्}}\Afootnote{H; \skt{पिडिका}\ A; \skt{पिटका}\ N}}}{\lemma{\skt{उत्तमाङ्गस्थाम्}}\Afootnote{H; \skt{मुत्तमाङ्गस्थां}\ A N}}}{\lemma{\skt{वृत्ताम्}}\Afootnote{H; \skt{वृत्ता}\ A N}}}{\lemma{\skt{उग्ररुजां}}\Afootnote{H; \skt{मुग्ररुजा}\ A N}}\skt{ ज्वरां । }\pend
\pstart\versepar \skt{}\edtext{\skt{सर्व्वात्मकां}}{\Afootnote{A H; \skt{सर्वात्मिकां}\ N}}\skt{ सर्व्वलिङ्गां जानी\-या\-दि\-रि\-वे\-ल्लि\-काम् ॥}\pend

\pstart\versepar \ledsidenote{\passref{2.13.16}}\edtext{\skt{वाहुकन्द्यां}}{\Afootnote{H; \skt{बाहुपार्श्वां}\ A; \skt{बाहुकक्ष्यां}\ N}}\skt{ }\edtext{\skt{सपार्श्वषु}}{\Afootnote{H; \skt{सकक्षासु}\ A; \skt{सपार्श्वेषु}\ N}}\skt{ }\edtext{\skt{कृष्णां}}{\Afootnote{N H; \skt{कृष्ण}\ A}}\skt{ स्फोटां }\edtext{\skt{सवेदनां}}{\Afootnote{A H; \skt{सवेनां ।}\ N}}\skt{ । }\pend
\pstart\versepar \skt{पित्त\-प्र\-को\-प\-सं\-भूतां }\edtext{\skt{कक्ष्येतिभिषगा}}{\Afootnote{H; \skt{कक्षामिति विनि}\ A; \skt{कति भिषगा}\ N}}\skt{ }\edtext{\skt{दिशेत्}}{\Afootnote{H; \skt{र्दिशेत् ॥}\ A; \skt{दिसेत् ॥}\ N}}\skt{ ॥}\pend

\pstart\versepar \ledsidenote{\passref{2.13.17}}\edtext{\edtext{\edtext{\skt{एका\-मे\-ता\-दृ\-शी\-न्दृष्टां}}{\lemma{\skt{एकाम्}}\Afootnote{H; \skt{एका}\ A N}}}{\lemma{\skt{एतादृशीन्}}\Afootnote{H; \skt{मेवंविधां}\ A; \skt{मेतान्दृशी}\ N}}}{\lemma{\skt{दृष्टां}}\Afootnote{H; \skt{दृष्ट्वा}\ A N}}\skt{ }\edtext{\skt{पिटिकां}}{\Afootnote{A H; \skt{पिटकां}\ N}}\skt{ स्फोटसन्निभां । }\pend
\pstart\versepar \skt{}\edtext{\skt{त्वग्ग\-ता\-म्पि\-त्त\-को\-पेन}}{\lemma{\skt{त्वग्गताम्}}\Afootnote{A H; \skt{त्वग्तां}\ N}}\skt{ }\edtext{\skt{गन्ध\-ना\-मा\-म्वि\-नि\-र्द्दे\-शेत्}}{\lemma{\skt{विनिर्द्देशेत्}}\Afootnote{H; \skt{प्रचक्षते ॥}\ A; \skt{विनिदिशेत् ॥}\ N}}\skt{ ॥}\pend

\pstart\versepar \ledsidenote{\passref{2.13.18}}\skt{अग्निदग्धनिभाः स्फोटाः सज्वराः पित्तरक्ततः । }\pend
\pstart\versepar \skt{}\edtext{\skt{चि}}{\Afootnote{H; \skt{क्वचि}\ A; \skt{त्वचि}\ N}}\skt{ }\edtext{\skt{सर्व्वत्र वा}}{\Afootnote{N H; \skt{त्सर्वत्र वा}\ A}}\skt{ देहे }\edtext{\skt{विस्फोटक}}{\Afootnote{N H; \skt{स्मृता विस्फोटका}\ A}}\skt{ इति }\edtext{\skt{स्मृतः}}{\Afootnote{N H; \emph{om.}\ A}}\skt{ ॥}\pend

\pstart\versepar \ledsidenote{\passref{2.13.19}}\edtext{\skt{कक्षभागेषु ये}}{\Afootnote{H; \skt{कक्षाभागेषु ये}\ A; \skt{कक्षकारोषु ये}\ N}}\skt{ स्फोटा जायन्ते }\edtext{\skt{मान्सदारूणा}}{\Afootnote{H; \skt{मांसदारु(र)णाः ।}\ A; \skt{मांसदारुणाः ।}\ N}}\skt{ । }\pend
\pstart\versepar \skt{}\edtext{\skt{अन्त\-र्द्दा\-ह\-ज्व\-र\-करा}}{\Afootnote{A H; \skt{अन्त\-र्दा\-ह\-ज्व\-रा\-करा}\ N}}\skt{ }\edtext{\skt{दीप्त\-पा\-व\-क\-स\-न्निभाः}}{\Afootnote{A H; \skt{दीप्त\-पा\-व\-क\-स\-न्निभा ॥}\ N}}\skt{ ॥}\pend

\pstart\versepar \ledsidenote{\passref{2.13.20}}\skt{सप्ताहाद्वा }\edtext{\skt{दशाहाद्वा}}{\Afootnote{H; \skt{दशाहाद्वा}\ A N}}\skt{ }\edtext{\skt{पक्षाद्वा}}{\Afootnote{H; \skt{पक्षाद्वा}\ A N}}\skt{ घ्नन्ति मानवं । }\pend
\pstart\versepar \skt{}\edtext{\edtext{\skt{ताम\-ग्नि\-रो\-हि\-णी\-म्वि\-द्या\-द\-साध्यां}}{\lemma{\skt{अग्निरोहिणीम्}}\Afootnote{H; \skt{मग्निरोहिणीं}\ A; \skt{मग्निरोहिणी}\ N}}}{\lemma{\skt{विद्यादसाध्यां}}\Afootnote{H; \skt{विद्यादसाध्यां}\ A N}}\skt{ सन्निपाततः ॥}\edtext{}{\lemma{\skt{मानवं । ताम्}}\Afootnote{H; \skt{मानवम्ता}\ A N}}\pend

\pstart\versepar \ledsidenote{\passref{2.13.21}}\edtext{\edtext{\skt{नख\-मा\-न्स\-म\-धि\-ष्ठाय}}{\lemma{\skt{नखमान्सम्}}\Afootnote{H; \skt{नखमांस}\ A N}}}{\lemma{\skt{अधिष्ठाय}}\Afootnote{H; \skt{मधिष्ठाय}\ A N}}\skt{ }\edtext{\skt{वातः}}{\Afootnote{H; \skt{पित्तं}\ A; \skt{वात}\ N}}\skt{ }\edtext{\skt{पित्तञ्देहिनां}}{\Afootnote{H; \skt{वातश्च वेदनाम् ।}\ A; \skt{पित्तञ्च देहिनां ।}\ N}}\skt{ । }\pend
\pstart\versepar \skt{}\edtext{\skt{कुर्या तान्दाहपाकौ}}{\lemma{\skt{कुर्या तान्}}\Afootnote{H; \skt{करोति}\ A; \skt{कुर्याता}\ N}}\skt{ च तं व्याधि\-ञ्चि\-प्प\-मा\-दि\-शेत् ॥}\pend

\pstart \ledsidenote{\passref{2.13.22}}\edtext{\edtext{\skt{तदे\-वा\-ल्प\-त\-त\-रै\-र्दोषैः}}{\lemma{\skt{तदे\-वा\-ल्प\-त\-त\-रैर्}}\Afootnote{H; \skt{तदे\-वा\-क्ष\-त\-रो\-गाख्यं तथो\-प\-न\-ख\-मि\-त्यपि । अभिघातात्प्र}\ A; \skt{तदे\-वा\-ल्प\-त\-रै\-र्दो\-षै\-र्कु\-न\-ख\-म्पा\-रुषं खर वा}\ N}}}{\lemma{\skt{दोषैः}}\Afootnote{H; \skt{दुष्टो}\ A; \emph{om.}\ N}}\skt{ }\edtext{\edtext{\skt{कुनखम्परूषं}}{\lemma{\skt{कुनखम्}}\Afootnote{H; \skt{यो नखो}\ A; \emph{om.}\ N}}}{\lemma{\skt{परूषं}}\Afootnote{H; \skt{रूक्षोऽसितः}\ A; \emph{om.}\ N}}\skt{ }\edtext{\skt{खरम्}}{\Afootnote{H; \skt{खरः ॥}\ A; \emph{om.}\ N}}\skt{ ॥}\pend

\pstart \ledsidenote{\passref{2.13.23}}\edtext{\edtext{\skt{गम्भी\-रा\-म\-ल्प\-सं\-रम्भां}}{\lemma{\skt{गम्भीराम्}}\Afootnote{N H; \skt{भवेत्तं कुनखं विद्या\-त्कु\-ली\-न\-मिति संज्ञितम् । गम्भीरा}\ A}}}{\lemma{\skt{अल्पसंरम्भां}}\Afootnote{H; \skt{मल्पसंरम्भां}\ A; \skt{अल्पसन्रम्भां}\ N}}\skt{ }\edtext{\skt{सव\-र्ण्णा\-मु\-प\-रि\-स्थितां}}{\lemma{\skt{सवर्ण्णाम्}}\Afootnote{A H; \skt{सवर्ङ्णाम्}\ N}}\skt{ ।}\pend

\pstart\versepar \ledsidenote{\passref{2.13.24}}\edtext{\edtext{\edtext{\skt{कफा\-द\-नु\-श\-यी\-म्य\-स्या\-म्वि\-द्या\-दन्तः}}{\lemma{\skt{कफा\-द\-नु\-श\-यीम्}}\Afootnote{H; \skt{कफा}\ A; \skt{कफादनुशयी}\ N}}}{\lemma{\skt{यस्याम्}}\Afootnote{H; \emph{om.}\ A; \skt{पङ्ख}\ N}}}{\lemma{\skt{विद्यादन्तः}}\Afootnote{N H; \skt{दन्तः}\ A}}\skt{ }\edtext{\skt{प्रपाकिनीं}}{\Afootnote{H; \skt{प्रपाकां तां}\ A; \skt{प्र}\ N}}\skt{ ॥ }\pend
\pstart\versepar \skt{}\edtext{\skt{विदा\-री\-क\-न्द\-व\-द्धृ\-ता\-न्क\-क्ष\-वं\-क्ष\-ण\-स\-न्धिषु}}{\Afootnote{H; \skt{विद्यादनुशयीं भिषक् । विदा\-री\-क\-न्द\-व\-द्वृत्तां कक्षा\-व\-ङ्क्ष\-ण\-स\-न्धिषु ॥}\ A; \skt{याकिरीं । विदा\-री\-क\-न्द\-व\-द्वृत्तां कक्षा\-वं\-क्ष\-ण\-स\-न्धिषु ॥}\ N}}\skt{ ।}\pend

\pstart\versepar \ledsidenote{\passref{2.13.25}}\edtext{\edtext{\skt{विदारिकामिति}}{\lemma{\skt{विदारिकाम्}}\Afootnote{N H; \skt{रक्तां विदारिकां}\ A}}}{\lemma{\skt{इति}}\Afootnote{N H; \emph{om.}\ A}}\skt{ }\edtext{\skt{वदेत्सरुजां}}{\Afootnote{H; \skt{विद्यात्सर्वजां}\ A; \skt{वदेद्सरुजां}\ N}}\skt{ सर्व्वलक्षणां ॥ }\pend
\pstart\versepar \skt{}\edtext{\skt{प्राप्य\-मां\-स\-सि\-रा\-म्नायु}}{\Afootnote{\skt{प्राप्य मांससिरास्नायु}\ A H; \skt{प्राप्य मांससिरास्नायुं}\ N}}\skt{ श्लेष्मा }\edtext{\skt{मेदस्तथानिलः}}{\Afootnote{\skt{मेदस्तथाऽनिलः ॥}\ A; \skt{मेदस्तथानिलः}\ N H}}\skt{ ।}\pend

\pstart\versepar \ledsidenote{\passref{2.13.26}}\edtext{\edtext{\skt{ग्रन्थिङ्कुर्व्वन्त्यसौ}}{\lemma{\skt{कुर्व्वन्त्य्}}\Afootnote{A H; \skt{कुर्वन्}\ N}}}{\lemma{\skt{असौ}}\Afootnote{H; \emph{om.}\ A; \skt{त्य सौ}\ N}}\skt{ भिन्नो }\edtext{\skt{मधु\-स\-र्प्पि\-र्व्व\-सा\-निभं}}{\Afootnote{N H; \skt{ऽसौ मधु\-स\-र्पि\-र्व\-सा\-नि\-मम् ।}\ A}}\skt{ ॥ }\pend
\pstart\versepar \skt{}\edtext{\edtext{\skt{करो\-त\-स्रा\-व\-म\-नि\-ल\-स्तत्र}}{\lemma{\skt{करोतस्रावम्}}\Afootnote{H; \skt{स्रवत्यास्राव}\ A; \skt{करोत्यास्राव}\ N}}}{\lemma{\skt{अनिलस्}}\Afootnote{H; \skt{मत्यर्थं}\ A; \skt{मनिलस्}\ N}}\skt{ वृद्धिङ्गतः }\edtext{\skt{पुनः}}{\Afootnote{N H; \skt{ऽनिलः ॥}\ A}}\skt{ ॥}\pend

\pstart\versepar \ledsidenote{\passref{2.13.27}}\edtext{\skt{मान्स\-म्वि\-शो\-ग्र\-थितां}}{\lemma{\skt{विशोग्रथितां}}\Afootnote{H; \skt{विशोष्य ग्रथितां}\ A N}}\skt{ }\edtext{\edtext{\skt{शर्क्क\-रा\-ञ्ज\-न\-त्यतः}}{\lemma{\skt{शर्क्कराञ्जनत्य्}}\Afootnote{H; \skt{शर्करां जनये}\ A; \skt{शर्क्करां जनय}\ N}}}{\lemma{\skt{अतः}}\Afootnote{H; \skt{त्पुनः ।}\ A; \skt{न्यतः ।}\ N}}\skt{ ॥ }\pend
\pstart\versepar \skt{}\edtext{\edtext{\skt{दुर्ग\-न्धि\-क्लि\-न्न\-म\-त्य\-र्थ\-न्ना\-ना\-व\-र्ण्ण\-न्ततः}}{\lemma{\skt{दुर्ग\-न्धि\-क्लि\-न्न\-मत्य्}}\Afootnote{H; \skt{दुर्गन्धं क्लिन्नम}\ A; \skt{दुर्गन्ध क्लिन्नम}\ N}}}{\lemma{\skt{अर्थन्}}\Afootnote{H; \skt{त्य्र्थं}\ A; \skt{त्यर्थन्}\ N}}\skt{ सिराः ।}\pend

\pstart\versepar \ledsidenote{\passref{2.13.28}}\edtext{\skt{सृजन्ति}}{\Afootnote{N H; \skt{स्रवन्ति}\ A}}\skt{ }\edtext{\skt{रक्तं}}{\Afootnote{N H; \emph{om.}\ A}}\skt{ सहसा }\edtext{\edtext{\edtext{\skt{तम्वि\-द्या\-च्छ\-र्श\-र्क्क\-रा\-र्व्वु\-दम्}}{\lemma{\skt{तम्}}\Afootnote{N H; \skt{रक्तं त}\ A}}}{\lemma{\skt{विद्या\-च्छ\-र्श\-र्क्क\-रार्}}\Afootnote{H; \skt{द्विद्याच्छर्करा}\ A; \skt{विद्याच्छर्क्करा}\ N}}}{\lemma{\skt{व्वुदम्}}\Afootnote{H; \skt{र्बुदम्}\ A N}}\skt{ ॥ }\pend
\pstart\versepar \skt{पामाविचर्च्यौ कुष्ठेषु }\edtext{\skt{रकसा च}}{\Afootnote{A H; \skt{रकशा च}\ N}}\skt{ प्रकीर्तिता ॥}\pend

\pstart\versepar \ledsidenote{\passref{2.13.29}}\edtext{\skt{परि\-क्र\-म\-ण\-शी\-लस्य}}{\Afootnote{A H; \skt{परि\-क्र\-म\-ण\-सी\-लस्य}\ N}}\skt{ }\edtext{\edtext{\skt{वायु\-र\-त्य\-र्थ\-रू\-क्षयोः}}{\lemma{\skt{वायुरत्य्}}\Afootnote{H; \skt{वायुर}\ A; \skt{वायुर}\ N}}}{\lemma{\skt{अर्थरूक्षयोः}}\Afootnote{H; \skt{त्यर्थरूक्षयोः ।}\ A N}}\skt{ । }\pend
\pstart\versepar \skt{पादयोः }\edtext{\skt{कुरुते}}{\Afootnote{A H; \skt{कुरते}\ N}}\skt{ }\edtext{\edtext{\skt{दारी\-स\-रू\-जा\-स्त\-ल\-सं\-श्रिताः}}{\lemma{\skt{दारीसरूजास्}}\Afootnote{H; \skt{दारीं सरुजां}\ A; \skt{दारी सरुजाम्}\ N}}}{\lemma{\skt{तलसंश्रिताः}}\Afootnote{H; \skt{तलसंश्रितः ॥}\ A; \skt{तलसंस्रिताः ॥}\ N}}\skt{ ॥}\pend

\pstart \ledsidenote{\passref{2.13.30}}\edtext{\skt{शर्क्करोन्मथिते}}{\Afootnote{A H; \skt{सर्क्करोन्मथिते}\ N}}\skt{ पादे }\edtext{\skt{क्षत वा}}{\Afootnote{H; \skt{क्षते वा}\ A N}}\skt{ }\edtext{\skt{कण्टकादिभिः}}{\Afootnote{H; \skt{कण्टकादिभिः । मेदो\-र\-क्ता\-नु\-गै\-श्चैव दोषैर्वा जायते नृणाम् ॥}\ A; \skt{कण्ठकादिभिः ।}\ N}}\skt{ ।}\pend

\pstart \ledsidenote{\passref{2.13.31}}\edtext{\skt{ग्रन्थिं}}{\Afootnote{H; \skt{सकीलकठिनो ग्रन्थि\-र्नि\-म्न\-म\-ध्यो\-न्न\-तोऽपि वा ।}\ A; \skt{ग्रन्थि}\ N}}\skt{ }\edtext{\skt{कीलवदुत्सन्नो}}{\Afootnote{H; \skt{कोलमात्रः सरुक्स्रावी}\ A; \skt{कील वदुत्सन्नो}\ N}}\skt{ जायते }\edtext{\skt{कदरन्तु}}{\Afootnote{N H; \skt{कदरस्तु}\ A}}\skt{ }\edtext{\skt{तत}}{\Afootnote{N H; \skt{सः ॥}\ A}}\skt{ ॥}\pend

\pstart\versepar \ledsidenote{\passref{2.13.32}}\skt{क्लिन्नाङ्गुल्यन्तरौ पादौ }\edtext{\skt{कण्ड\-दा\-ह\-रु\-जा\-न्वितौ}}{\Afootnote{H; \skt{कण्डू\-दा\-ह\-रु\-ग\-न्वितौ ।}\ A; \skt{कण्डू\-दा\-ह\-रु\-जा\-न्वितौ ।}\ N}}\skt{ । }\pend
\pstart\versepar \skt{}\edtext{\skt{दुष्ट\-क\-र्द्द\-म\-सं\-स्प\-र्शा\-द\-ल\-सेति}}{\Afootnote{H; \skt{दुष्ट\-क\-र्द\-म\-सं\-स्प\-र्शा\-द\-लसं तं}\ A; \skt{दुष्ट\-क\-र्द्द\-म\-सं\-स्प\-र्शा\-द\-ल\-सेति}\ N}}\skt{ }\edtext{\skt{विभावयेत्}}{\Afootnote{N H; \skt{विनिर्दिशेत् ॥}\ A}}\skt{ ॥}\pend

\pstart\versepar \ledsidenote{\passref{2.13.33}}\edtext{\skt{रोम\-कू\-पा\-नु\-ग\-म्पि\-त्त\-म्वा\-तेन}}{\lemma{\skt{रोम\-कू\-पा\-नु\-गम्}}\Afootnote{A H; \skt{रोम\-कू\-पा\-नु\-राम्}\ N}}\skt{ सह मूर्च्छितं । }\pend
\pstart\versepar \skt{प्रच्यावयति रोमाणि ततः श्लेष्मा }\edtext{\skt{सशोणितः}}{\Afootnote{A H; \skt{सशोणितं ॥}\ N}}\skt{ ॥}\pend

\pstart\versepar \ledsidenote{\passref{2.13.34}}\edtext{\skt{रुण्द्धि}}{\Afootnote{H; \skt{रुणद्धि}\ A N}}\skt{ }\edtext{\skt{रोमकूपास्तु}}{\Afootnote{N H; \skt{रोमकूपांस्तु}\ A}}\skt{ }\edtext{\skt{ततोऽ\-न्ये\-षा\-म\-स\-म्भवः}}{\lemma{\skt{ऽन्ये\-षा\-म\-स\-म्भवः}}\Afootnote{A H; \skt{न्येषामसम्भवः ।}\ N}}\skt{ । }\pend
\pstart\versepar \skt{}\edtext{\skt{तदिन्द्रलुप्तं}}{\Afootnote{H; \skt{तदिन्द्रलुप्तं}\ A; \skt{तदिन्द्रलुप्त}\ N}}\skt{ }\edtext{\skt{खालित्यं}}{\Afootnote{A H; \skt{खालित्यं}\ N}}\skt{ }\edtext{\skt{रूजेति च}}{\Afootnote{H; \skt{रुज्येति च}\ A; \skt{रुजेति च}\ N}}\skt{ विभाव्यते ॥}\pend

\pstart\versepar \ledsidenote{\passref{2.13.35}}\skt{दारुणाकण्डुरा रूक्षा }\edtext{\edtext{\skt{केश\-भू\-मि\-म्प्र\-यो\-द्यते}}{\lemma{\skt{केशभूमिम्}}\Afootnote{H; \skt{केशभूमिः}\ A N}}}{\lemma{\skt{प्रयोद्यते}}\Afootnote{H; \skt{प्रपाट्यते ।}\ A; \skt{प्रचोद्यते ।}\ N}}\skt{ । }\pend
\pstart\versepar \skt{}\edtext{\skt{कफ\-मा\-रु\-त\-को\-पेन}}{\Afootnote{H; \skt{कफ\-वा\-त\-प्र\-को\-पेण}\ A; \skt{कफ\-मा\-रु\-त\-को\-पेन}\ N}}\skt{ }\edtext{\skt{विद्या\-द्दा\-रु\-ण\-कन्तु}}{\Afootnote{N H; \skt{विद्याद्दारुणकं तु}\ A}}\skt{ }\edtext{\skt{तत्}}{\Afootnote{N H; \skt{तम् ॥}\ A}}\skt{ ॥}\pend

\pstart\versepar \ledsidenote{\passref{2.13.36}}\edtext{\skt{अरुंषि}}{\Afootnote{H; \skt{अरूंषि}\ A N}}\skt{ }\edtext{\skt{वहुवक्राणि}}{\Afootnote{H; \skt{बहुवक्त्राणि}\ A N}}\skt{ }\edtext{\skt{वहुक्लेदानि}}{\Afootnote{H; \skt{बहुक्लेदीनि}\ A; \skt{बहुक्लेदानि}\ N}}\skt{ }\edtext{\skt{मूर्द्धन्तु}}{\Afootnote{H; \skt{मूर्धनि ।}\ A; \skt{मूर्ध्धित्र ।}\ N}}\skt{ । }\pend
\pstart\versepar \skt{}\edtext{\skt{कफा\-सृ\-क्क्रि\-त\-को\-पेन}}{\Afootnote{H; \skt{कफा\-सृ\-क्कृ\-मि\-को\-पेन}\ A; \skt{कफा\-सृ\-क्षि\-न्त\-को\-पेन}\ N}}\skt{ }\edtext{\edtext{\skt{व्रणा\-म्वि\-द्या\-द\-रू\-षि\-काम्}}{\lemma{\skt{व्रणाम्}}\Afootnote{N H; \skt{नृणां}\ A}}}{\lemma{\skt{विद्या\-द\-रू\-षि\-काम्}}\Afootnote{H; \skt{विद्या\-द\-रुं\-षि\-काम् ॥}\ A; \skt{विद्यादरूंषिकां ॥}\ N}}\skt{ ॥}\pend

\pstart\versepar \ledsidenote{\passref{2.13.37}}\edtext{\skt{क्रोध\-शो\-क\-श्र\-म\-कृतः}}{\Afootnote{A H; \skt{क्रोध\-सो\-क\-श्र\-म\-कृतः}\ N}}\skt{ शरीरोष्मा शिरोगतः । }\pend
\pstart\versepar \skt{}\edtext{\skt{पित्तञ्च}}{\Afootnote{A N\textsuperscript{pc} H; \skt{पित्तआञ्च}\ N\textsuperscript{ac}}}\skt{ }\edtext{\skt{केशान्पचति}}{\Afootnote{A H; \skt{केसां पचति}\ N}}\skt{ पलितन्तेन जायते ॥}\pend

\pstart\versepar \ledsidenote{\passref{2.13.38}}\skt{दाह\-ज्व\-र\-रु\-जा\-व\-न्त\-स्ताम्राः स्फोटाः }\edtext{\skt{सवेदनाः}}{\Afootnote{H; \skt{सपीतकाः ।}\ A; \skt{शवेद{\symfont ¦}नाः ।}\ N}}\skt{ । }\pend
\pstart\versepar \skt{}\edtext{\edtext{\skt{गात्रेष्वन्तश्च}}{\lemma{\skt{गात्रेष्व्}}\Afootnote{A H; \skt{गात्रेन्}\ N}}}{\lemma{\skt{अन्तश्च}}\Afootnote{H; \emph{om.}\ A; \skt{तश्च}\ N}}\skt{ वदने }\edtext{\skt{विज्ञेया तु}}{\Afootnote{H; \skt{चान्तर्विज्ञेयास्ता}\ A; \skt{विज्ञेयात्तु}\ N}}\skt{ }\edtext{\skt{मसूरिका}}{\Afootnote{H; \skt{मसूरिकाः ॥}\ A; \skt{मशूरिकाः ॥}\ N}}\skt{ ॥}\pend

\pstart\versepar \ledsidenote{\passref{2.13.39}}\edtext{\skt{शाल्म\-ली\-क\-ण्ट\-क\-प्रख्याः}}{\Afootnote{A H; \skt{शाल्म\-ली\-क\-ण्ठ\-क\-प्रख्याः}\ N}}\skt{ }\edtext{\skt{कफ\-मा\-रु\-त\-र\-क्तजाः}}{\Afootnote{H; \skt{कफ\-मा\-रु\-त\-शोणि}\ A; \skt{कपा\-फ\-मा\-रु\-त\-र\-क्तजाः ।}\ N}}\skt{ । }\pend
\pstart\versepar \skt{}\edtext{\skt{युवानपिडका}}{\Afootnote{H; \skt{तैः । जायन्ते पिडका}\ A; \skt{युवान पिटका}\ N}}\skt{ }\edtext{\edtext{\skt{यूनाम्विज्ञेया}}{\lemma{\skt{यूनाम्}}\Afootnote{N H; \skt{यूनं}\ A}}}{\lemma{\skt{विज्ञेया}}\Afootnote{N H; \skt{वक्त्रे या}\ A}}\skt{ }\edtext{\skt{मुखदूषणाः}}{\Afootnote{H; \skt{मुखदूषिकाः ॥}\ A; \skt{मुखैदूषणा ॥}\ N\textsuperscript{ac}; \skt{मुखदूषणा ॥}\ N\textsuperscript{pc}}}\skt{ ॥}\pend

\pstart\versepar \ledsidenote{\passref{2.13.40}}\edtext{\edtext{\edtext{\edtext{\skt{कण्ट\-कै\-रा\-चि\-त\-म्वि\-द्या\-न्म\-ण्ड\-ल\-म्पाण्डु}}{\lemma{\skt{कण्टकैर्}}\Afootnote{H; \skt{कण्टकै}\ A; \skt{कण्ठकैर्}\ N}}}{\lemma{\skt{आचितम्}}\Afootnote{N H; \skt{राचितं}\ A}}}{\lemma{\skt{विद्यान्मण्डलम्}}\Afootnote{N H; \skt{वृत्तं कण्डूम}\ A}}}{\lemma{\skt{पाण्डु}}\Afootnote{N H; \skt{त्पाण्डु}\ A}}\skt{ }\edtext{\skt{कण्डुरं}}{\Afootnote{H; \skt{मण्डलम् ।}\ A; \skt{कण्डूरं ।}\ N}}\skt{ । }\pend
\pstart\versepar \skt{}\edtext{\edtext{\edtext{\skt{पद्मि\-नी\-क\-ण्ट\-क\-प्र\-ख्यै\-स्त\-दा\-ख्य\-ङ्क\-फ\-बा\-त\-जम्}}{\lemma{\skt{पद्मि\-नी\-क\-ण्ट\-क\-प्र\-ख्यैस्}}\Afootnote{A H; \skt{पद्मि\-नी\-क\-ण्ठ\-क\-प्र\-ख्यैस्}\ N}}}{\lemma{\skt{तदाख्यङ्}}\Afootnote{A H; \skt{तदास्य}\ N}}}{\lemma{\skt{कफबातजम्}}\Afootnote{H; \skt{कफवातजम्}\ A N}}\skt{ ॥}\pend

\pstart\versepar \ledsidenote{\passref{2.13.41}}\edtext{\edtext{\edtext{\skt{सम\-मु\-त्स\-न्न\-म\-रु\-ज\-म्म\-ण्ड\-ल\-ङ्क\-फ\-र\-क्तजं}}{\lemma{\skt{समम्}}\Afootnote{N H; \skt{नीरुजं सम}\ A}}}{\lemma{\skt{उत्सन्नम्}}\Afootnote{N H; \skt{मुत्सन्नं}\ A}}}{\lemma{\skt{अरुजम्}}\Afootnote{N H; \emph{om.}\ A}}\skt{ । }\pend
\pstart\versepar \skt{सहजं }\edtext{\skt{लक्ष्म}}{\Afootnote{H; \skt{रक्त}\ A; \skt{लक्ष्य}\ N}}\skt{ }\edtext{\skt{चैकेषां}}{\Afootnote{N H; \skt{मीषच्च}\ A}}\skt{ }\edtext{\skt{लक्ष्यो}}{\Afootnote{H; \skt{श्लक्ष्णं}\ A; \skt{लक्ष्या}\ N}}\skt{ }\edtext{\skt{जतुमणीति}}{\Afootnote{H; \skt{जतुमणिं वि}\ A; \skt{जतुमनीति}\ N}}\skt{ }\edtext{\skt{सा}}{\Afootnote{N H; \skt{दुः ॥}\ A}}\skt{ ॥}\pend

\pstart\versepar \ledsidenote{\passref{2.13.42}}\skt{अवेदनं स्थिरञ्चैव }\edtext{\skt{यस्मिन्गात्रे}}{\Afootnote{H; \skt{यस्य गात्रे}\ A; \skt{यस्मिं गात्रे}\ N}}\skt{ }\edtext{\skt{प्रदृश्यते}}{\Afootnote{H; \skt{षु दृश्यते ।}\ A N}}\skt{ । }\pend
\pstart\versepar \skt{}\edtext{\edtext{\edtext{\edtext{\skt{माष\-व\-त्कृ\-ष्ण\-मु\-त्स\-न्न\-म\-नि\-ला\-त्म\-ष\-क\-मा\-दि\-शेत्}}{\lemma{\skt{माषवत्कृष्णम्}}\Afootnote{H; \skt{माषवत्कृष्ण}\ A N}}}{\lemma{\skt{उत्सन्नम्}}\Afootnote{H; \skt{मुत्सन्न}\ A N}}}{\lemma{\skt{अनि\-ला\-त्म\-ष\-कम्}}\Afootnote{H; \skt{मनिलान्मषकं}\ A; \skt{मनि\-ला\-न्म\-श\-कम्}\ N}}}{\lemma{\skt{आदिशेत्}}\Afootnote{H; \skt{वदेत् ॥}\ A; \skt{आदिसेत् ॥}\ N}}\skt{ ॥}\pend

\pstart\versepar \ledsidenote{\passref{2.13.43}}\skt{कृष्णानि तिलमात्राणि नीरुजानि }\edtext{\skt{समानि च}}{\Afootnote{A H; \skt{समाणि च ।}\ N}}\skt{ । }\pend
\pstart\versepar \skt{}\edtext{\edtext{\skt{वात\-पि\-त्त\-ल\-वो\-त्से\-ता\-म्वि\-ध्या\-त्ति\-का\-ल\-काम्}}{\lemma{\skt{वात\-पि\-त्त\-ल\-वो\-त्से\-ताम्}}\Afootnote{H; \skt{वात\-पि\-त्त\-क\-फो\-च्छो\-षात्ता}\ A; \skt{वात\-पि\-त्त\-ल\-छो\-त्से\-धा\-न्तास्}\ N}}}{\lemma{\skt{विध्या\-त्ति\-का\-ल\-काम्}}\Afootnote{H; \skt{न्विद्या\-त्ति\-ल\-का\-ल\-कान् ॥}\ A; \skt{विर्या\-त्ति\-ल\-का\-म्भि\-षक् ॥}\ N}}\skt{ ॥}\pend

\pstart\versepar \ledsidenote{\passref{2.13.44}}\edtext{\skt{महद्वा}}{\Afootnote{N H; \skt{मण्डलं मह}\ A}}\skt{ }\edtext{\skt{यदि}}{\Afootnote{N H; \emph{om.}\ A}}\skt{ }\edtext{\edtext{\skt{वात्यल्पं}}{\lemma{\skt{वात्य्}}\Afootnote{H; \emph{om.}\ A; \skt{वात्य}\ N}}}{\lemma{\skt{अल्पं}}\Afootnote{H; \skt{दल्पं}\ A; \skt{न्यल्पं}\ N}}\skt{ }\edtext{\skt{श्यावम्वा}}{\Afootnote{H; \skt{वा}\ A; \skt{स्पाचं वा}\ N}}\skt{ }\edtext{\skt{यदि वा}}{\Afootnote{A H; \skt{यं तिवा}\ N}}\skt{ }\edtext{\skt{सितं}}{\Afootnote{A H; \skt{मितं ।}\ N}}\skt{ । }\pend
\pstart\versepar \skt{}\edtext{\edtext{\skt{नीरुजम्मण्डलं}}{\lemma{\skt{नीरुजम्}}\Afootnote{H; \skt{सहजं नीरुजं}\ A; \skt{निरुजं}\ N}}}{\lemma{\skt{मण्डलं}}\Afootnote{N H; \emph{om.}\ A}}\skt{ }\edtext{\skt{गात्रे}}{\Afootnote{A H; \skt{गात्र}\ N}}\skt{ }\edtext{\skt{नच्छ\-मि\-त्य\-भि\-धी\-यते}}{\lemma{\skt{नच्छम्}}\Afootnote{N H; \skt{न्यच्छम्}\ A}}\skt{ ॥}\pend

\pstart \ledsidenote{\passref{2.13.45}}\edtext{\skt{क्रोधा\-या\-स\-प्र\-कु\-पितो}}{\Afootnote{H; \skt{समु\-त्था\-न\-नि\-दा\-नाभ्यां चर्मकीलं प्रकीर्तितम् । क्रोधा\-या\-स\-प्र\-कु\-पितो}\ A; \skt{क्रोधायासः प्रकुपितो}\ N}}\skt{ वायुः पित्तेन संयुतः ।}\pend

\pstart\versepar \ledsidenote{\passref{2.13.46}}\edtext{\edtext{\skt{मुखमागम्य}}{\lemma{\skt{मुखम्}}\Afootnote{H; \skt{सहसा मुख}\ A; \skt{मुख}\ N}}}{\lemma{\skt{आगम्य}}\Afootnote{H; \skt{मागत्य}\ A; \skt{मागम्य}\ N}}\skt{ }\edtext{\skt{सहसा}}{\Afootnote{N H; \emph{om.}\ A}}\skt{ }\edtext{\edtext{\skt{मण्ड\-ल\-म्वि\-सृ\-ज\-त्यतः}}{\lemma{\skt{विसृजत्य्}}\Afootnote{A H; \skt{विसृज}\ N}}}{\lemma{\skt{अतः}}\Afootnote{A H; \skt{त्यतः ॥}\ N}}\skt{ ॥ }\pend
\pstart\versepar \skt{}\edtext{\skt{निरुजन्तनुकं}}{\lemma{\skt{निरुजन्}}\Afootnote{H; \skt{नीरुजं}\ A N}}\skt{ }\edtext{\skt{श्यावं}}{\Afootnote{A H; \skt{श्याव}\ N}}\skt{ मुखे व्यङ्ग\-न्त\-मा\-दि\-शेत् ।}\pend

\pstart\versepar \ledsidenote{\passref{2.13.47}}\edtext{\edtext{\skt{कृष्णमेवङ्गुणं}}{\lemma{\skt{कृष्णम्}}\Afootnote{N H; \skt{कृष्ण}\ A}}}{\lemma{\skt{एवङ्}}\Afootnote{N H; \skt{मेवं}\ A}}\skt{ गात्रे मुखे वा }\edtext{\skt{नीलिकाम्विदुः}}{\lemma{\skt{नीलिकाम्}}\Afootnote{A N\textsuperscript{pc} H; \skt{नीइलिकाम्}\ N\textsuperscript{ac}}}\skt{ ॥ }\pend
\pstart\versepar \skt{}\edtext{\skt{मर्द्द\-ना\-त्पी\-ड\-ना\-द्चापि}}{\Afootnote{H; \skt{मर्द\-ना\-त्पी\-ड\-ना\-च्चाति}\ A; \skt{मर्द्दना पीडानाञ्चापि}\ N}}\skt{ }\edtext{\skt{तथै\-वा\-प्य\-भि\-घा\-ततः}}{\Afootnote{A H; \skt{तथै\-वा\-प्य\-भि\-घा\-ततः ।}\ N}}\skt{ । }\pend
\pstart\versepar \skt{मेढ्रचर्म्म यदा वायुर्भजते }\edtext{\skt{सर्व्वतश्चरन्}}{\lemma{\skt{चरन्}}\Afootnote{N H; \skt{चरः ॥}\ A}}\skt{ ॥}\pend

\pstart \ledsidenote{\passref{2.13.48ab}}\skt{तदा }\edtext{\skt{वातोपसंसृष्टं}}{\Afootnote{N H; \skt{वातोपसृष्टं तु}\ A}}\skt{ चर्म्म }\edtext{\skt{प्रतिनिवर्त्तये}}{\Afootnote{H; \skt{प्रतिनिवर्तते ।}\ A; \skt{प्रति\-नि\-व\-र्त्त\-येत् ।}\ N}}\skt{ ।}\pend

\pstart \ledsidenote{\passref{2.13.49ab}}\edtext{\skt{स वेद\-न\-न्स\-दा\-हश्च}}{\Afootnote{H; \skt{सवेदनः सदाहश्च}\ A; \skt{सवेदनं सदाहञ्च}\ N}}\skt{ }\edtext{\skt{पाकं}}{\Afootnote{N H; \skt{पाकं च}\ A}}\skt{ }\edtext{\skt{व्रजति}}{\Afootnote{A H; \skt{प्रजति}\ N}}\skt{ }\edtext{\skt{चासकृत्}}{\Afootnote{N H; \skt{क्वचित् ।}\ A}}\skt{ ॥}\pend

\pstart \ledsidenote{\passref{2.13.48cd}}\edtext{\edtext{\skt{मणे\-र\-ध\-स्ता\-त्को\-शस्तु}}{\lemma{\skt{मणेर्}}\Afootnote{N H; \skt{मणे}\ A}}}{\lemma{\skt{अधस्तात्कोशस्तु}}\Afootnote{N H; \skt{रधस्तात्कोशश्च}\ A}}\skt{ }\edtext{\skt{ग्रन्थिरूपेण}}{\Afootnote{A H; \skt{ग्रन्थिरूपपेण}\ N}}\skt{ }\edtext{\skt{लम्वते}}{\Afootnote{H; \skt{लम्बते}\ A N}}\skt{ ।}\pend

\pstart \ledsidenote{\passref{2.13.49cd}}\edtext{\skt{परिवर्तिकेति }\edtext{\edtext{\skt{ताम्वि\-द्या\-त्स\-रु\-जा\-म्वा\-त\-सं\-भ\-वाम्}}{\lemma{\skt{विद्यात्सरुजाम्}}\Afootnote{H; \emph{om.}\ A; \skt{विद्याद्सरुजाम्}\ N}}}{\lemma{\skt{वातसंभवाम्}}\Afootnote{H; \emph{om.}\ A; \skt{वातसंभवास् ॥}\ N}}}{\lemma{\skt{परिवर्तिकेति ताम्}}\Afootnote{N H; \skt{मारु\-ता\-ग\-न्तु\-सं\-भूतां विद्यात्तां परिवर्तिकाम् ॥}\ A}}\skt{ ॥}\pend

\pstart\versepar \ledsidenote{\passref{2.13.50}}\edtext{\skt{सकण्डूः}}{\Afootnote{A H; \skt{सण्डूकण्डूः}\ N}}\skt{ कठिना }\edtext{\skt{चापि}}{\Afootnote{A H; \skt{वापि}\ N}}\skt{ सैव }\edtext{\skt{श्लेष्मसमुत्थिता}}{\Afootnote{A H; \skt{स्लेष्मसमुत्थिता ।}\ N}}\skt{ । }\pend
\pstart\versepar \skt{अल्पीयः }\edtext{\skt{खां}}{\Afootnote{A H; \skt{खाखा}\ N}}\skt{ }\edtext{\skt{यदा}}{\Afootnote{A H; \skt{न्यदा}\ N}}\skt{ }\edtext{\edtext{\skt{हर्षो\-द्व\-ला\-ङ्ग\-च्छे\-न्स्त्रि\-य\-न्नरः}}{\lemma{\skt{हर्षोद्वलाङ्}}\Afootnote{H; \skt{हर्षाद्बालां}\ A; \skt{हर्षाद्बला}\ N}}}{\lemma{\skt{गच्छेन्स्त्रियन्}}\Afootnote{H; \skt{गच्छेत्स्त्रियं}\ A; \skt{द्गच्छेत्यन्}\ N}}\skt{ ॥}\pend

\pstart \ledsidenote{\passref{2.13.51}}\edtext{\skt{हस्ता\-भि\-घा\-ता\-दपि वा}}{\Afootnote{H; \skt{हस्ता\-भि\-घा\-ता\-द\-थवा}\ A; \skt{हस्ताभिघातापि वा}\ N}}\skt{ }\edtext{\skt{चर्म्माण्युद्वर्त्तते}}{\Afootnote{H; \skt{चर्मण्युद्वर्तिते}\ A; \skt{चर्मण्युद्वर्त्तते}\ N}}\skt{ }\edtext{\skt{वलात्}}{\Afootnote{H; \skt{बलात् । मर्द\-ना\-त्पी\-ड\-ना\-द्वाऽपि शुक्र\-वे\-ग\-वि\-घा\-ततः ॥}\ A; \skt{वा बलात् ।}\ N}}\skt{ ।}\pend

\pstart\versepar \ledsidenote{\passref{2.13.52}}\edtext{\skt{यस्यावपाते}}{\Afootnote{H; \skt{यस्यावपाठ्यते}\ A; \skt{यस्यावपाट्यते}\ N}}\skt{ चर्म्म\-ता\-म्वि\-द्या\-द\-व\-पा\-टि\-काम् ॥ }\pend
\pstart\versepar \skt{}\edtext{\edtext{\skt{वातो\-प\-सृ\-ष्ट\-मे\-वेत}}{\lemma{\skt{वातोपसृष्टम्}}\Afootnote{A H; \skt{वातोपसृष्ट}\ N}}}{\lemma{\skt{एवेत}}\Afootnote{H; \skt{एवं तु}\ A; \skt{मेवैत}\ N}}\skt{ }\edtext{\skt{च्चर्म्म}}{\Afootnote{H; \skt{चर्म}\ A; \skt{च्चर्म}\ N}}\skt{ संश्रयते मणिं ।}\pend

\pstart\versepar \ledsidenote{\passref{2.13.53}}\edtext{\skt{मणि\-श्च\-र्म्मो\-प\-रु\-द्धश्च}}{\lemma{\skt{चर्म्मोपरुद्धश्च}}\Afootnote{H; \skt{चर्मोपनद्धस्तु}\ A; \skt{चर्मोपरुद्धस्तु}\ N}}\skt{ }\edtext{\skt{मूत्रस्रोतो}}{\Afootnote{A H; \skt{मूत्रसोतो}\ N}}\skt{ }\edtext{\skt{रुणद्धि तु}}{\Afootnote{N H; \skt{रुणद्धि च ।}\ A}}\skt{ ॥ }\pend
\pstart\versepar \skt{}\edtext{\edtext{\skt{निरु\-द्ध\-प्र\-का\-श\-म्वि\-द्या\-न्म\-न्द\-धारं}}{\lemma{\skt{निरु\-द्ध\-प्र\-का\-शम्}}\Afootnote{H; \skt{निरुद्धप्रकशे}\ A; \skt{निरु\-द्ध\-प्र\-क\-शम्}\ N}}}{\lemma{\skt{विद्यान्मन्दधारं}}\Afootnote{H; \skt{तस्मि\-न्म\-न्द\-घा\-रम्}\ A; \skt{विद्यात्सन्दधारं}\ N}}\skt{ }\edtext{\skt{सवेदनं}}{\Afootnote{N H; \skt{अवेदनम् ॥}\ A}}\skt{ ।}\pend

\pstart\versepar \ledsidenote{\passref{2.13.54}}\skt{मूत्रं प्रवर्तते }\edtext{\edtext{\skt{जन्तोर्म्मणिर्न्न}}{\lemma{\skt{जन्तोर्}}\Afootnote{N H; \skt{जन्तो}\ A}}\skt{ च}}{\lemma{\skt{म्मणिर्न्न च}}\Afootnote{H; \skt{र्मणिर्न च}\ A; \skt{मणिर्न च}\ N}}\skt{ }\edtext{\skt{विदार्यते}}{\Afootnote{N H; \skt{विदीर्यते ।}\ A}}\skt{ ॥ }\pend
\pstart\versepar \skt{}\edtext{\edtext{\edtext{\skt{निरु\-द्ध\-प्र\-का\-श\-म्वि\-द्या\-त्स\-रु\-ज\-म्वा\-त\-स\-म्भवं}}{\lemma{\skt{निरु\-द्ध\-प्र\-का\-शम्}}\Afootnote{H; \skt{निरुद्धप्रकशं}\ A N}}}{\lemma{\skt{विद्यात्सरुजम्}}\Afootnote{N H; \skt{विद्याद्दुरूढां}\ A}}}{\lemma{\skt{वातसम्भवं}}\Afootnote{N H; \skt{चावपाटिकाम् ॥}\ A}}\skt{ ।}\pend

\pstart\versepar \ledsidenote{\passref{2.13.55}}\edtext{\skt{वेग\-स\-न्धा\-र\-णा\-द्वा\-यु\-र्व्वि\-हतो}}{\lemma{\skt{वेग\-स\-न्धा\-र\-णा\-द्वा\-युर्}}\Afootnote{H; \skt{वेग\-सं\-घा\-र\-णा\-द्वा\-युर्}\ A; \skt{वेग\-स\-न्धा\-र\-णा\-द्वायुः}\ N}}\skt{ }\edtext{\edtext{\skt{गुदमाश्रितः}}{\lemma{\skt{गुदम्}}\Afootnote{H; \skt{गुद}\ A N}}}{\lemma{\skt{आश्रितः}}\Afootnote{H; \skt{माश्रितः ।}\ A; \skt{मास्रितः ॥}\ N}}\skt{ ॥ }\pend
\pstart\versepar \skt{निरुणद्धि }\edtext{\skt{महोत्स्रोतः}}{\Afootnote{H; \skt{महत्स्रोतः}\ A N}}\skt{ सूक्ष्मद्वारं करोति च ।}\pend

\pstart\versepar \ledsidenote{\passref{2.13.56}}\skt{मार्गस्य सौक्ष्म्यात्कृच्छ्रेण पुरीषन्तस्य गच्छति ॥ }\pend
\pstart\versepar \skt{}\edtext{\edtext{\skt{सन्नि\-रु\-द्ध\-गु\-द\-म्व्या\-धि\-मे\-त\-द्वि\-द्या\-त्सु\-दु\-स्तरं}}{\lemma{\skt{सन्नि\-रु\-द्ध\-गु\-दम्}}\Afootnote{A H; \skt{सनिरुद्धगुदं}\ N}}}{\lemma{\skt{एत\-द्वि\-द्या\-त्सु\-दु\-स्तरं}}\Afootnote{H; \skt{एनं विद्या\-त्सु\-दु\-स्त\-रम् ॥}\ A; \skt{एत विद्यात्सुदुस्तरं ।}\ N}}\skt{ ।}\pend

\pstart\versepar \ledsidenote{\passref{2.13.57}}\edtext{\skt{शकृ\-त्मू\-त्र\-स\-मा\-युक्ते}}{\Afootnote{H; \skt{शकृ\-न्मू\-त्र\-स\-मा\-युक्ते}\ A; \skt{सकृ\-न्मू\-त्र\-स\-मा\-युक्ता}\ N}}\skt{ }\edtext{\skt{धौते}}{\Afootnote{N H; \skt{ऽधौते}\ A}}\skt{ }\edtext{\skt{पाने}}{\Afootnote{N H; \skt{ऽपाने}\ A}}\skt{ }\edtext{\skt{शिशोर्भवेत्}}{\lemma{\skt{शिशोर्}}\Afootnote{A H; \skt{सिसोर्}\ N}}\skt{ ॥ }\pend
\pstart\versepar \skt{स्विन्ने }\edtext{\skt{वास्वेद्यमाने वा}}{\Afootnote{H; \skt{स्यास्नाप्यमानस्य}\ A; \skt{वा स्वेद्यय्मानो वा}\ N}}\skt{ }\edtext{\skt{कण्डू\-र\-क्त\-क\-फो\-द्भवाः}}{\Afootnote{H; \skt{कण्डू रक्तकफोद्भवा}\ A N}}\skt{ ॥}\edtext{}{\lemma{\skt{भवेत् ॥ स्विन्ने}}\Afootnote{H; \skt{भवेत्स्विन्न}\ A N}}\pend

\pstart\versepar \ledsidenote{\passref{2.13.58}}\skt{कण्डूयनात्ततः क्षिप्रं }\edtext{\skt{स्फोटास्रावश्च}}{\Afootnote{N H; \skt{स्फोटाः स्रावश्च}\ A}}\skt{ जायते ॥ }\pend
\pstart\versepar \skt{}\edtext{\skt{एकीभूतम्व्रणं}}{\lemma{\skt{व्रणं}}\Afootnote{N H; \skt{व्रणैर्}\ A}}\skt{ घोरं तम्वि\-द्या\-द\-हि\-पू\-त\-नम् ॥}\pend

\pstart\versepar \ledsidenote{\passref{2.13.59}}\skt{स्नानो\-त्सा\-द\-न\-ही\-नस्य मलो }\edtext{\skt{वृषणसंस्थितः}}{\Afootnote{H; \skt{वृष्णसंश्रितः ।}\ A; \skt{वृष्यणसंस्थितः ॥}\ N}}\skt{ । }\pend
\pstart\versepar \skt{यदा प्रक्लिद्यते स्वेदात्कण्डूं }\edtext{\skt{सञ्जनयेत्तदा}}{\Afootnote{A H; \skt{स जनयेत्तदा ॥}\ N}}\skt{ ॥}\pend

\pstart\versepar \ledsidenote{\passref{2.13.60}}\edtext{\skt{ततः}}{\Afootnote{N H; \skt{तत्र}\ A}}\skt{ कण्डूयनात्क्षिप्रं }\edtext{\skt{स्फोटास्रावश्च}}{\Afootnote{H; \skt{स्फोटाः स्रावश्च}\ A; \skt{स्फोटा स्राचश्च}\ N}}\skt{ जायते । }\pend
\pstart\versepar \skt{}\edtext{\edtext{\skt{प्राहु\-र्व्वृ\-ष\-ण\-क\-च्छ्रून्तां}}{\lemma{\skt{प्राहुर्}}\Afootnote{N H; \skt{प्राहु}\ A}}}{\lemma{\skt{व्वृषणकच्छ्रून्तां}}\Afootnote{H; \skt{र्वृषणकच्छूं तां}\ A; \skt{वृषणकच्छूस्तां}\ N}}\skt{ }\edtext{\skt{श्लेष्म\-र\-क्त\-प्र\-को\-जाम्}}{\Afootnote{H; \skt{श्लेष्म\-र\-क्त\-प्र\-को\-प\-जाम्}\ A N}}\skt{ ॥}\pend

\pstart\versepar \ledsidenote{\passref{2.13.61}}\skt{प्रवा\-ह\-णा\-ति\-सा\-रा\-भ्या\-न्नि\-र्ग\-च्छति }\edtext{\skt{गुदो}}{\Afootnote{H; \skt{गुदं}\ A; \skt{गुदा}\ N}}\skt{ }\edtext{\skt{वहिः}}{\Afootnote{H; \skt{बहिः}\ A N}}\skt{ । }\pend
\pstart\versepar \skt{}\edtext{\skt{रूक्ष\-दु\-र्व्व\-ल\-दे\-हस्य तं}}{\Afootnote{H; \skt{रूक्ष\-दु\-र्ब\-ल\-दे\-हस्य तं}\ A N}}\skt{ }\edtext{\skt{गुद\-भ्र\-न्श\-मा\-दि\-शेत्}}{\lemma{\skt{गुदभ्रन्शम्}}\Afootnote{A H; \skt{गुदं भ्रंशम्}\ N}}\skt{ ॥}\pend

\pstart\versepar \ledsidenote{\passref{2.13.62}}\skt{इत्येवं क्षुद्ररोगाणां स्थानं सम्प\-रि\-की\-र्ति\-तम् । }\pend
\pstart\versepar \skt{तदवेक्ष्य भिषक्प्राज्ञो यथा\-दो\-ष\-मु\-पा\-च\-रेत् ॥ ॥}\pend

\pstart \ledsidenote{\passref{2.13.trailer}}\edtext{\skt{इति}}{\Afootnote{H; \emph{om.}\ A; \skt{इत्येवं क्षुद्रारोगानां ज्ञानं सम्यरि}\ N}}\skt{ }\edtext{\skt{क्षुद्र\-रो\-ग\-नि\-दाने}}{\Afootnote{H; \skt{इति सुश्रु\-त\-सं\-हि\-तायां निदानस्थाने क्षुद्र\-रो\-ग\-नि\-दानं नाम}\ A; \skt{कीर्त्तितं ॥ तदवेक्ष्या भिषक्प्राज्ञो यथा दोष\-मु\-पा\-च\-रेत् ॥ निनिदानस्थाने ॥ ॥१३ ॥}\ N}}\skt{ }\edtext{\skt{त्रयोदशमः}}{\Afootnote{H; \skt{त्रयोदशोऽध्यायः ॥१३ ॥}\ A; \emph{om.}\ N}}\skt{ ॥ ० ॥}\pend

\endnumbering
\beginnumbering
\eledsection*{Adhyāya 14}
\pstart \ledsidenote{\passref{2.14.1}}\skt{अथातः }\edtext{\edtext{\skt{शूक\-दो\-ष\-नि\-दा\-न\-म्व्या\-ख्या\-स्यामः}}{\lemma{\skt{शूक\-दो\-ष\-नि\-दा\-नम्}}\Afootnote{A H; \skt{शूक\-दो\-ष\-नि\-दाना}\ N}}}{\lemma{\skt{व्याख्यास्यामः}}\Afootnote{A H; \skt{ध्यायं व्या ॥}\ N}}\skt{ ॥}\pend

\pstart \ledsidenote{\passref{2.14.3}}\edtext{\skt{लिङ्ग\-वृ\-द्धि\-मि\-च्छ\-ता\-म\-क्र\-म\-प्र\-वृ\-त्तानां}}{\Afootnote{A; \skt{लिङ्ग\-वृ\-द्धि\-मि\-इ\-च्छ\-ता\-म\-क्र\-म\-प्र\-वृ\-त्ताना}\ N\textsuperscript{ac}; \skt{लिङ्ग\-वृ\-द्धि\-मि\-च्छ\-ता\-म\-क्र\-म\-प्र\-वृ\-त्ताना}\ N\textsuperscript{pc}; \skt{लिङ्ग\-वृ\-द्धि\-मि\-च्छ\-ता\-म\-क्र\-म\-प्र\-वृ\-त्तानां}\ H}}\skt{ }\edtext{\skt{शूक\-दो\-ष\-नि\-मित्तं}}{\Afootnote{H; \skt{शूक\-दो\-ष\-नि\-मित्ता}\ A; \skt{शूक\-दो\-ष\-नि\-मित्तां}\ N}}\skt{ }\edtext{\skt{व्याधयो}}{\Afootnote{N H; \emph{om.}\ A}}\skt{ दश }\edtext{\skt{चाष्टौ}}{\Afootnote{H; \skt{चाष्टौ च}\ A N}}\skt{ }\edtext{\skt{जायन्ते}}{\Afootnote{N H; \skt{व्याधयो जायन्ते ।}\ A}}\skt{ ॥ }\edtext{\skt{तद्यथा}}{\Afootnote{A H; \skt{तद्यथा ॥}\ N}}\skt{ ॥ }\edtext{\skt{सर्षपिकाष्ठीला}}{\Afootnote{N H; \skt{सर्षपिका अष्ठीलिका}\ A}}\skt{ ग्रथितं कुम्भीका }\edtext{\skt{अलजीमृदितं}}{\Afootnote{A H; \skt{अलजी मृदित}\ N}}\skt{ }\edtext{\skt{सम्मूढपिटका}}{\Afootnote{H; \skt{संमूढपिडका}\ A; \skt{संमूढं समूढपिडका ।}\ N}}\skt{ । अवमन्थः }\edtext{\skt{पुष्करिका}}{\Afootnote{A H; \skt{पुश्करिका}\ N}}\skt{ स्पर्शहानिः । }\edtext{\skt{उत्तवः}}{\Afootnote{H; \skt{उत्तमा}\ A; \skt{उत्तव}\ N}}\skt{ शतपोनकः । }\edtext{\skt{त्वष्काकः}}{\Afootnote{H; \skt{त्वक्पाकः}\ A N}}\skt{ }\edtext{\skt{शोणितार्ब्बुदं}}{\Afootnote{A H; \skt{शोणितार्वुदः}\ N}}\skt{ }\edtext{\skt{मांसार्ब्बुदं}}{\Afootnote{A H; \skt{मांसार्वुदः}\ N}}\skt{ मांसपाको }\edtext{\edtext{\skt{विद्र\-धि\-स्ति\-ल\-का\-ल\-क\-ञ्चेति}}{\lemma{\skt{विद्रधिस्}}\Afootnote{A N; \skt{विद्रधि}\ H}}}{\lemma{\skt{तिलकालकञ्}}\Afootnote{\skt{तिलकालकश्}\ A N; \skt{स्तिलकालकञ्}\ H}}\skt{ ॥}\pend

\pstart\versepar \ledsidenote{\passref{2.14.4}}\edtext{\skt{गौर\-स\-र्ष\-प\-सं\-स्थाना}}{\Afootnote{H; \skt{गौर\-स\-र्ष\-प\-तुल्या तु}\ A; \skt{गौर\-स\-र्ष\-प\-सं\-स्थान ।}\ N}}\skt{ }\edtext{\skt{शूक\-नि\-र्व्दु\-ग्न\-हे\-तुकाः}}{\Afootnote{H; \skt{शूक\-दु\-र्भ\-ग्न\-हे\-तुका ।}\ A; \skt{शूक\-नि\-र्भु\-ग्न\-हे\-तुका ।}\ N}}\skt{ । }\pend
\pstart\versepar \skt{}\edtext{\skt{पिट\-का\-क\-फ\-र\-क्ताभ्यां}}{\Afootnote{H; \skt{पिडका कफरक्ताभ्यां}\ A N}}\skt{ ज्ञेया }\edtext{\skt{सर्षपिसा॥}}{\lemma{\skt{सर्षपिसा}}\Afootnote{H; \skt{सर्षपिका बुधैः ॥}\ A; \skt{सर्षपिका तु सा ॥}\ N}}\pend

\pstart\versepar \ledsidenote{\passref{2.14.5}}\skt{कठिना }\edtext{\skt{विषमा}}{\Afootnote{H; \skt{विषमैर्}\ A; \skt{विष}\ N}}\skt{ }\edtext{\edtext{\skt{भग्नै\-र्वा\-यु\-ना\-ष्ठी\-लिका}}{\lemma{\skt{भग्नैर्}}\Afootnote{H; \skt{अन्तैर्}\ A; \emph{om.}\ N}}}{\lemma{\skt{वायुनाष्ठीलिका}}\Afootnote{H; \skt{मारुतस्य प्रको}\ A; \skt{मा}\ N}}\skt{ भवेत् । }\pend
\pstart\versepar \skt{}\edtext{\skt{शूकैर्यत्पूरितं}}{\lemma{\skt{यत्पूरितं}}\Afootnote{H; \skt{तु विषसं}\ A}}\skt{ }\edtext{\skt{शश्वद्ग्रथितं}}{\Afootnote{H; \skt{भुग्नैः पिडकाऽष्ठीलिका भवे}\ A; \skt{भग्नै\-र्वा\-यु\-ना\-ष्ठी\-लिका भवे}\ N}}\skt{ }\edtext{\edtext{\skt{नाम तत्कफात् ॥}}{\Afootnote{\emph{om.}\ A}}}{\Afootnote{\emph{om.}\ N}}\edtext{}{\lemma{\skt{युनाष्ठीलिका… पूरितं}}\Afootnote{\emph{om.}\ N}}\edtext{}{\lemma{\skt{भवेत् । शूकैर्}}\Afootnote{H; \skt{पतः । शूकैस्}\ A}}\pend

\pstart\versepar \ledsidenote{\passref{2.14.6}}\edtext{\skt{कुम्भीका}}{\Afootnote{H; \skt{शूकैर्यत्पूरितं शश्वद्ग्रथितं तत्कफोत्थितम् । कुम्भीका}\ A; \skt{शूकैर्यत्पूरितं शश्व\-द्ग्र\-थि\-त\-न्नाम तत्कफात् ॥ कुम्भीका}\ N}}\skt{ रक्तपित्तोत्था जाव्व\-स्थि\-स\-दृ\-शा\-शु\-णत् । }\pend
\pstart\versepar \skt{}\edtext{\skt{तुल्य\-जा\-त्व\-ल\-जी\-म्वि\-द्या\-द्यथा}}{\lemma{\skt{विद्याद्यथा}}\Afootnote{H; \emph{om.}\ A N}}\skt{ }\edtext{\edtext{\skt{प्रोक्ताम्विचक्षणः}}{\lemma{\skt{प्रोक्ताम्}}\Afootnote{H; \emph{om.}\ A N}}}{\lemma{\skt{विचक्षणः}}\Afootnote{H; \emph{om.}\ A N}}\skt{ ॥}\edtext{}{\lemma{\skt{जाव्व\-स्थि\-स\-दृ\-शा\-शु\-णत् । तुल्य\-जा\-त्व\-ल\-जीम्}}\Afootnote{H; \skt{जाम्ब\-वा\-स्थि\-नि\-भाऽ\-शुभा ॥}\ A; \skt{जंव्वस्थिसदृशा शुभा ।}\ N}}\pend

\pstart \ledsidenote{\passref{2.14.7}}\edtext{\skt{मृदितं}}{\Afootnote{H; \skt{अल\-जी\-ल\-क्ष\-णै\-र्यु\-क्ता\-म\-लजीं च वितर्कयेत् । मृदितं}\ A; \skt{तुल्य\-जा\-त्व\-ल\-जी\-म्वि\-द्या\-द्य\-था\-प्रो\-क्ता\-म्वि\-च\-क्षणः ॥ मृदितम्}\ N}}\skt{ पीडितं यत्तु }\edtext{\edtext{\skt{संर\-ब्ध\-म्वा\-त\-को\-पतः}}{\lemma{\skt{संरब्धम्}}\Afootnote{A H; \skt{संरव्धम्}\ N}}}{\lemma{\skt{वातकोपतः}}\Afootnote{H; \skt{वायुकोपतः ॥}\ A; \skt{वायतकोपतः ॥}\ N}}\skt{ ।}\pend

\pstart\versepar \ledsidenote{\passref{2.14.8}}\skt{पाणिभ्यां }\edtext{\skt{भृशसम्मूढे}}{\Afootnote{A H; \skt{भृशसंमूढ}\ N}}\skt{ }\edtext{\skt{संमूढपिटका}}{\Afootnote{H; \skt{संमूढपिडका}\ A; \skt{समू\-ढा\-पि\-ड\-कम्}\ N}}\skt{ भवेत् ॥ }\pend
\pstart\versepar \skt{दीर्घा }\edtext{\skt{बह्व्यश्च}}{\Afootnote{A H; \skt{वह्वश्च}\ N}}\skt{ }\edtext{\skt{पिटका}}{\Afootnote{H; \skt{पिडका}\ A N}}\skt{ }\edtext{\skt{दीर्यन्ते}}{\Afootnote{A H; \skt{दीर्यते}\ N}}\skt{ }\edtext{\skt{मध्धतस्तु}}{\Afootnote{H; \skt{मध्यतस्तु}\ A N}}\skt{ याः ॥}\pend

\pstart \ledsidenote{\passref{2.14.9}}\edtext{\skt{सोऽवमन्थः}}{\Afootnote{A H; \skt{सोवमन्थः}\ N}}\skt{ }\edtext{\edtext{\skt{कफा\-सृ\-ग्भ्या\-म्वे\-द\-ना\-रो\-म\-ह\-र्ष\-कृत्}}{\lemma{\skt{कफासृग्भ्याम्}}\Afootnote{A H; \skt{कफासृग्भ्या}\ N}}}{\lemma{\skt{वेद\-ना\-रो\-म\-ह\-र्ष\-कृत्}}\Afootnote{\skt{वेद\-ना\-रो\-म\-ह\-र्ष\-कृत् । पित्त\-शो\-णि\-त\-सं\-भूता पिडका पिडकाचिता ॥}\ A; \skt{स्वेदना रोमहर्षकृत् ॥ पिडकाचिता या पिडका पित्त\-शो\-णि\-त\-स\-म्भवा ।}\ N; \skt{वेद\-ना\-रो\-म\-ह\-र्ष\-कृत् ॥ पिट\-का\-चि\-ता\-या\-पि\-ट\-का\-पि\-त्त\-शो\-णि\-त\-सं\-भवा ॥}\ H}}\skt{ ॥ ॥}\pend

\pstart\versepar \ledsidenote{\passref{2.14.10}}\edtext{\skt{पद्म\-क\-र्णि\-क\-सं\-स्थाना}}{\Afootnote{N H; \skt{पद्म\-पु\-ष्क\-र\-स्ं\-स्थाना}\ A}}\skt{ ज्ञेया पुष्करिकेतिसा ॥ }\pend
\pstart\versepar \skt{}\edtext{\skt{स्पर्शहानिन्तु}}{\Afootnote{N H; \skt{जन\-ये\-त्स्प\-र्श\-हानिं तु}\ A}}\skt{ }\edtext{\skt{जनयेच्छोणितं}}{\Afootnote{N H; \skt{शोणितं}\ A}}\skt{ }\edtext{\skt{शूकदूषितं}}{\Afootnote{N H; \skt{शूकदूपितम् ॥}\ A}}\skt{ ।}\pend

\pstart\versepar \ledsidenote{\passref{2.14.11}}\edtext{\skt{मुद्गमाषोपमा}}{\Afootnote{A H; \skt{मुद्गमाशोपमा}\ N}}\skt{ रक्ता }\edtext{\skt{पिटका}}{\Afootnote{H; \skt{पिडका}\ A N}}\skt{ रक्तपित्तजा । }\pend
\pstart\versepar \skt{}\edtext{\edtext{\skt{व्याधिरेवोत्तवो}}{\lemma{\skt{व्याधिर्}}\Afootnote{H; \skt{उत्त}\ A; \skt{व्याधि}\ N}}}{\lemma{\skt{एवोत्तवो}}\Afootnote{H; \skt{मैषा तु वि}\ A; \skt{रेषोत्तवो}\ N}}\skt{ }\edtext{\skt{नाम}}{\Afootnote{N H; \skt{ज्ञेया}\ A}}\skt{ }\edtext{\skt{शूका\-जी\-र्ण्ण\-नि\-मि\-त्तजः}}{\Afootnote{N H; \skt{शूका\-जी\-र्ण\-नि\-मि\-त्तजा ॥}\ A}}\skt{ ।}\pend

\pstart\versepar \ledsidenote{\passref{2.14.12}}\edtext{\skt{छिद्रै\-र\-णु\-मु\-खै\-र्ल्लि\-ङ्ग\-ञ्चितं}}{\lemma{\skt{ल्लिङ्गञ्}}\Afootnote{N H; \skt{वस्तु}\ A}}\skt{ यस्य समन्ततः ॥ }\pend
\pstart\versepar \skt{}\edtext{\skt{वातशोणितजो}}{\Afootnote{A H; \emph{om.}\ N}}\skt{ }\edtext{\skt{व्याधिः स}}{\Afootnote{H; \skt{व्याधिर्वि}\ A; \emph{om.}\ N}}\skt{ }\edtext{\skt{ज्ञेयः}}{\Afootnote{A H; \emph{om.}\ N}}\skt{ }\edtext{\skt{शतपोनकः}}{\Afootnote{A H; \emph{om.}\ N}}\skt{ ।}\pend

\pstart\versepar \ledsidenote{\passref{2.14.13}}\edtext{\skt{वातपित्तकृतो}}{\Afootnote{N H; \skt{पित्तरक्तकृतो}\ A}}\skt{ }\edtext{\edtext{\skt{ज्ञेयस्त्वक्काको}}{\lemma{\skt{ज्ञेयस्}}\Afootnote{N H; \skt{ज्ञेय}\ A}}}{\lemma{\skt{त्वक्काको}}\Afootnote{H; \skt{स्त्वक्पाको}\ A; \skt{त्व क्पको}\ N}}\skt{ ज्वरदाहकृत् । }\pend
\pstart\versepar \skt{कृष्णैः स्फोटैः }\edtext{\skt{सरभिः}}{\Afootnote{H; \skt{सरक्तैश्च}\ A; \skt{सरक्ताभिः}\ N}}\skt{ }\edtext{\skt{पिटकाभिश्च}}{\Afootnote{H; \skt{पिडकाभिश्च}\ A N}}\skt{ }\edtext{\skt{पीडितं}}{\Afootnote{A H; \skt{पीडीतं ।}\ N}}\skt{ । }\pend
\pstart\versepar \skt{यस्य }\edtext{\edtext{\skt{बस्तिरुजञ्चोग्रा}}{\lemma{\skt{बस्तिरुजञ्}}\Afootnote{H; \skt{वस्तु रुज}\ A; \skt{वस्तिरुजश्}\ N}}}{\lemma{\skt{चोग्रा}}\Afootnote{N H; \skt{श्चोग्रा}\ A}}\skt{ }\edtext{\skt{ज्ञेय\-न्त\-च्छो\-णि\-ता\-र्ब्बु\-दम्}}{\lemma{\skt{तच्छो\-णि\-ता\-र्ब्बु\-दम्}}\Afootnote{A H; \skt{तच्छोणितार्वुदं ॥}\ N}}\skt{ ॥}\edtext{}{\lemma{\skt{ज्वरदाहकृत् । कृष्णैः}}\Afootnote{N H; \skt{ज्वरदाहवान् । कृष्णैः}\ A}}\pend

\pstart\versepar \ledsidenote{\passref{2.14.14}}\skt{मांसदोषेण }\edtext{\skt{जानी\-या\-द\-र्ब्बु\-द\-म्मा\-न्स\-सं\-भवं}}{\lemma{\skt{जानी\-या\-द\-र्ब्बु\-दम्}}\Afootnote{A H; \skt{जानी\-या\-द\-र्वु\-दम्}\ N}}\skt{ । }\pend
\pstart\versepar \skt{}\edtext{\skt{शीर्यन्ते}}{\Afootnote{A H; \skt{सीर्यन्ते}\ N}}\skt{ यस्य मान्सानि यत्र सर्व्वाश्च वेदनाः ॥}\pend

\pstart\versepar \ledsidenote{\passref{2.14.15}}\skt{विद्या\-त्त\-म्मा\-न्स\-पा\-कन्तु सर्व\-दो\-ष\-कृ\-त\-म्भि\-षक् । }\pend
\pstart\versepar \skt{विद्रधिः सन्निपातेन }\edtext{\edtext{\skt{यथो\-क्त\-म\-भि\-नि\-र्द्दि\-शेत्}}{\lemma{\skt{यथोक्तम्}}\Afootnote{N H; \skt{यथोक्त}\ A}}}{\lemma{\skt{अभिनिर्द्दिशेत्}}\Afootnote{N H; \skt{मभिनिर्दिशेत् ॥}\ A}}\skt{ ॥}\edtext{}{\lemma{\skt{भिषक् । विद्रधिः}}\Afootnote{H; \skt{भिषक्विद्रधिं}\ A N}}\pend

\pstart\versepar \ledsidenote{\passref{2.14.16}}\skt{कृष्णानि }\edtext{\edtext{\skt{चित्राण्यथवा}}{\lemma{\skt{चित्राण्य्}}\Afootnote{N H; \skt{चित्रा}\ A}}}{\lemma{\skt{अथवा}}\Afootnote{N H; \skt{ण्यथवा}\ A}}\skt{ शूकानि }\edtext{\skt{सविषाणि वा}}{\Afootnote{H; \skt{सविषाणि च ।}\ A; \skt{सविषानि च ।}\ N}}\skt{ । }\pend
\pstart\versepar \skt{पातितानि }\edtext{\edtext{\skt{पचन्त्याशु}}{\lemma{\skt{पचन्त्य्}}\Afootnote{H; \skt{पचन्}\ A N}}}{\lemma{\skt{आशु}}\Afootnote{H; \skt{त्याशु}\ A N}}\skt{ मेढ्र\-न्नि\-र\-व\-शे\-षतः ॥}\pend

\pstart\versepar \ledsidenote{\passref{2.14.17}}\edtext{\skt{कृष्णानि}}{\Afootnote{N H; \skt{कालानि}\ A}}\skt{ भूत्वा मान्सानि शीर्यन्ते यस्य देहिनः । }\pend
\pstart\versepar \skt{सन्नि\-पा\-त\-स\-मु\-त्था\-न\-न्त\-म्वि\-द्या\-त्ति\-ल\-का\-ल\-कम् ॥}\pend

\pstart\versepar \ledsidenote{\passref{2.14.18}}\edtext{\skt{तत्र}}{\Afootnote{A H; \skt{तनु}\ N}}\skt{ }\edtext{\skt{मांसार्ब्बुदं}}{\Afootnote{H; \skt{मांसार्बुदः}\ A; \skt{मासार्वुर्द}\ N}}\skt{ यच्च }\edtext{\skt{मांसपाकश्च यः}}{\Afootnote{A H; \skt{मांसा पाकश्च यः}\ N}}\skt{ स्मृतः । }\pend
\pstart\versepar \skt{विद्रधिश्च न }\edtext{\skt{सिध्धन्ति ये च }\edtext{\skt{स्युस्तिलकालकाः}}{\lemma{\skt{तिलकालकाः}}\Afootnote{A H; \skt{त्तिलकालकाः ॥}\ N}}}{\lemma{\skt{सिध्धन्ति ये च स्युस्}}\Afootnote{H; \skt{सिध्यन्ति ये च स्युस्}\ A; \skt{सिध्यन्ति ये च स्यु}\ N}}\skt{ ॥}\pend

\pstart \ledsidenote{\passref{2.14.trailer}}\edtext{\skt{इति}}{\Afootnote{H; \skt{इति सुश्रु\-त\-सं\-हि\-तायां निदानस्थाने}\ A; \emph{om.}\ N}}\skt{ }\edtext{\skt{शूक\-दो\-ष\-नि\-दाने}}{\Afootnote{H; \skt{शूक\-दो\-ष\-नि\-दानं नाम}\ A; \skt{नि ॥ ॥ ओ ॥}\ N}}\skt{ }\edtext{\skt{चतुर्द्दशमः}}{\Afootnote{H; \skt{चतुर्दशोऽध्यायः ॥१४ ॥}\ A; \emph{om.}\ N}}\skt{ ॥}\pend

\endnumbering
\beginnumbering
\eledsection*{Adhyāya 15}
\pstart \ledsidenote{\passref{2.15.1}}\skt{अथातो }\edtext{\edtext{\skt{भग्न\-नि\-दा\-न\-म्व्या\-ख्या\-स्यामः}}{\lemma{\skt{भग्ननिदानम्}}\Afootnote{H; \skt{भग्नानां निदानं}\ A; \skt{भग्निनिदाना}\ N\textsuperscript{ac}; \skt{भग्ननिदाना}\ N\textsuperscript{pc}}}}{\lemma{\skt{व्याख्यास्यामः}}\Afootnote{A H; \skt{ध्यायं व्याव्याख्यास्यामः ॥}\ N}}\skt{ ॥}\pend

\pstart \ledsidenote{\passref{2.15.3}}\edtext{\skt{पत\-न\-पी\-ड\-न\-प्र\-हा\-र\-वि\-ष\-दु\-ष्ट\-प्र\-भृ\-ति\-भि\-र\-भि\-घा\-त\-वि\-शेषैः}}{\lemma{\skt{पत\-न\-पी\-ड\-न\-प्र\-हा\-र\-वि\-ष\-दु\-ष्ट\-प्र\-भृ\-ति\-भिर्}}\Afootnote{N H; \skt{पत\-न\-पी\-ड\-न\-प्र\-हा\-र\-क्षे\-प\-ण\-व्या\-ल\-मृ\-ग\-द\-श\-न\-प्र\-भृ\-ति\-भिर्}\ A}}\skt{ । }\edtext{\skt{अने\-क\-वि\-ध\-मस्थ्नां}}{\Afootnote{A H; \skt{अने\-क\-वि\-ध\-मस्थ्नां}\ N}}\skt{ भङ्ग\-मु\-प\-दि\-शन्ति ।}\pend

\pstart \ledsidenote{\passref{2.15.4}}\skt{तत्र }\edtext{\skt{भग्न\-जा\-त\-म\-ने\-क\-वि\-ध\-म\-नु\-सा\-र्य\-माणं}}{\Afootnote{H; \skt{भङ्ग(\-ग्न)\-जा\-त\-म\-ने\-क\-वि\-ध\-म\-नु\-सा\-र्य\-माणं}\ A; \skt{भग्न\-जा\-त\-म\-ने\-क\-वि\-ध\-म\-नु\-सा\-र्य\-माण}\ N}}\skt{ }\edtext{\skt{द्विवि\-ध\-मे\-वो\-प\-प\-द्यते}}{\Afootnote{A H; \skt{द्विवि\-ध\-मे\-वो\-प\-प\-द्यते}\ N}}\skt{ । }\edtext{\skt{सन्धिमुकतं}}{\Afootnote{A H; \skt{सन्धिमुक्तं}\ N}}\skt{ काण्डभग्नञ्च । तत्र षड्विधं }\edtext{\edtext{\skt{सन्धिमुक्तं}}{\lemma{\skt{सन्धिम्}}\Afootnote{N H; \skt{सन्धि}\ A}}}{\lemma{\skt{उक्तं}}\Afootnote{N H; \skt{मुक्तं}\ A}}\skt{ }\edtext{\skt{द्वादशविधं}}{\Afootnote{A H; \skt{द्वादशविध}\ N}}\skt{ काण्ड\-भ\-ग्न\-म्भ\-वति ।}\pend

\pstart \ledsidenote{\passref{2.15.5}}\skt{तत्र }\edtext{\edtext{\edtext{\edtext{\edtext{\skt{सन्धि\-मु\-क्त\-मु\-त्पि\-ष्ट\-म्वि\-श्लि\-ष्ट\-म्वि\-प\-रि\-व\-र्त्ति\-त\-म\-व\-क्षि\-प्त\-म\-ति\-क्षि\-प्त\-न्ति\-र्य\-क्क्षि\-प्त\-मिति}}{\lemma{\skt{सन्धि\-मु\-क्त\-मु\-त्पि\-ष्टम्}}\Afootnote{H; \skt{सन्धि\-मु\-क्त\-मु\-त्पिष्टं}\ A; \skt{सन्धि\-मु\-क्त\-मु\-त्पि\-ष्टम्}\ N}}}{\lemma{\skt{विप\-रि\-व\-र्त्ति\-तम्}}\Afootnote{N H; \skt{विवर्तितम्}\ A}}}{\lemma{\skt{अवक्षिप्तम्}}\Afootnote{A H; \skt{अव}\ N}}}{\lemma{\skt{अतिक्षिप्तन्}}\Afootnote{A H; \skt{क्षिप्तन्}\ N}}}{\lemma{\skt{इति}}\Afootnote{N H; \skt{इति षड्विधम् ॥}\ A}}\skt{ ।}\pend

\pstart\versepar \ledsidenote{\passref{2.15.6}}\edtext{\skt{प्रसा\-र\-णा\-कु\-ञ्च\-ना\-श\-क्ति\-रु\-ग्र\-रु\-जता}}{\Afootnote{H; \skt{तत्र प्रसा\-र\-णा\-कु\-ञ्च\-न\-वि\-व\-र्त\-ना\-क्षे\-प\-णा\-श\-क्ति\-रु\-ग्र\-रु\-जत्वं}\ A; \skt{प्रसा\-र\-णा\-कु\-ञ्च\-ना\-श\-क्ति\-रु\-ग्र\-रु\-जता}\ N}}\skt{ }\edtext{\skt{स्पर्शा\-स\-ह\-त्वा\-ञ्चेति}}{\lemma{\skt{स्पर्शासहत्वाञ्}}\Afootnote{H; \skt{स्पर्शासहत्वं}\ A N}}\skt{ । }\pend
\pstart\versepar \skt{}\edtext{\skt{सामान्यतः}}{\Afootnote{N H; \skt{सामान्यं}\ A}}\skt{ }\edtext{\edtext{\edtext{\skt{सन्धि\-मु\-क्त\-ल\-क्ष\-ण\-मुक्तं}}{\lemma{\skt{सन्धिम्}}\Afootnote{N H; \skt{सन्धि}\ A}}}{\lemma{\skt{उक्तलक्षणम्}}\Afootnote{N H; \skt{मुक्तलक्षण}\ A}}}{\lemma{\skt{उक्तं}}\Afootnote{N H; \skt{मुक्तम् ॥}\ A}}\skt{ ।}\pend

\pstart \ledsidenote{\passref{2.15.7}}\skt{वैशेषिकं तूत्पिष्टे }\edtext{\skt{संधावुभयतः}}{\Afootnote{H; \skt{सन्धावुमयतः}\ A; \skt{सन्धावुभयतः ।}\ N}}\skt{ शोफो वेद\-ना\-प्रा\-दु\-र्ब्भावो विशेषतश्च }\edtext{\skt{रात्रौ}}{\Afootnote{N H; \skt{नानाप्रकारा वेदना रात्रौ प्रादु}\ A}}\skt{ }\edtext{\skt{भवति}}{\Afootnote{N H; \skt{र्भवन्ति;}\ A}}\skt{ । }\edtext{\skt{विश्लि\-ष्टेऽ\-ल्प\-शो\-फता}}{\lemma{\skt{ऽल्पशोफता}}\Afootnote{H; \skt{ऽल्पः शोफो}\ A; \skt{ल्पशोफता}\ N}}\skt{ वेदनासातत्यं सन्धिविक्रिया च ॥ }\edtext{\skt{विपरिवर्त्तिते च}}{\Afootnote{H; \skt{विवर्तिते तु}\ A; \skt{विप\-रि\-व\-र्त्ति\-तञ्च}\ N}}\skt{ }\edtext{\skt{सन्धौ}}{\Afootnote{N H; \skt{सन्धि}\ A}}\skt{ }\edtext{\skt{पार्श्व\-ग\-म\-ना\-द्वि\-ष\-मा\-ङ्गता}}{\Afootnote{N H; \skt{पार्श्वा\-प\-ग\-म\-ना\-द्वि\-ष\-मा\-ङ्गता}\ A}}\skt{ वेदना च । अवक्षिप्ते }\edtext{\edtext{\skt{सन्धि\-वि\-श्ले\-ष\-स्ती\-व्र\-रु\-जता}}{\lemma{\skt{सन्धिविश्लेषस्}}\Afootnote{A; \skt{शन्धिविश्लेषस्}\ N; \skt{सन्धिविश्लेष}\ H}}\skt{ च}}{\lemma{\skt{तीव्ररुजता च}}\Afootnote{N; \skt{तीव्ररुजत्वं च;}\ A; \skt{स्तीव्ररुजता च ॥}\ H}}\skt{ ॥ }\edtext{\skt{अतिक्षिप्ते}}{\Afootnote{A H; \skt{अवक्षिप्ते}\ N}}\skt{ द्वयोः }\edtext{\skt{सन्ध्ध\-स्थ्नो\-र\-ति\-क्रा\-न्तता}}{\Afootnote{H; \skt{सन्ध्य\-स्थ्नो\-र\-ति\-क्रा\-न्तता}\ A; \skt{सन्ध्य\-स्थ्नो\-र\-ति\-क्रा\-न्तता}\ N}}\skt{ वेदना च । तिर्यक्क्षिप्ते }\edtext{\edtext{\skt{त्वेका\-स्थि\-पा\-र्श्व\-ग\-म\-न\-म\-त्य\-र्थ\-मिति}}{\lemma{\skt{त्वेका\-स्थि\-पा\-र्श्व\-ग\-म\-नम्}}\Afootnote{N H; \skt{त्वेका\-स्थि\-पा\-र्श्वा\-प\-ग\-म\-नम्}\ A}}}{\lemma{\skt{इति}}\Afootnote{N H; \skt{वेदना चेति ॥}\ A}}\skt{ ॥}\pend

\pstart \ledsidenote{\passref{2.15.8}}\edtext{\edtext{\skt{काण्डभग्नमत}}{\lemma{\skt{काण्डभग्नम्}}\Afootnote{N H; \skt{काण्डभग्न}\ A}}}{\lemma{\skt{अत}}\Afootnote{N H; \skt{मत}\ A}}\skt{ ऊर्द्ध्वम्वक्ष्यामः ॥ }\edtext{\edtext{\edtext{\skt{कर्क्क\-ट\-क\-म\-श्व\-क\-र्ण्ण\-ञ्चू\-र्णि\-त\-म्पि\-च्चि\-त\-म\-स्थि\-च्छ\-ल्लि\-ता\-का\-ण्ड\-भ\-ग्न\-म\-ति\-पा\-ति\-त\-म्म\-ज्जा\-नु\-गतं}}{\lemma{\skt{अस्थि\-च्छ\-ल्लि\-ता\-का\-ण्ड\-भ\-ग्नम्}}\Afootnote{H; \skt{अस्थिच्छल्लितं काण्डभग्नं}\ A; \skt{अच्छल्लिता काण्डा भग्नम्}\ N\textsuperscript{ac}; \skt{अस्थिच्छल्लिता काण्डा भग्नम्}\ N\textsuperscript{pc}}}}{\lemma{\skt{अतिपातितम्}}\Afootnote{N H; \emph{om.}\ A}}}{\lemma{\skt{मज्जानुगतं}}\Afootnote{H; \skt{मज्जा\-नु\-ग\-त\-म\-ति\-पा\-तितं}\ A; \skt{मज्जानुगत}\ N}}\skt{ वक्रञ्च्छिन्नं }\edtext{\skt{स्फुटितं}}{\Afootnote{N H; \skt{पाटितं}\ A}}\skt{ }\edtext{\edtext{\skt{स्फाटितमिति}}{\lemma{\skt{स्फाटितम्}}\Afootnote{N H; \skt{स्फुटितम्}\ A}}}{\lemma{\skt{इति}}\Afootnote{N H; \skt{इति द्वादशविधम् ॥}\ A}}\skt{ ॥}\pend

\pstart \ledsidenote{\passref{2.15.9}}\edtext{\edtext{\skt{श्वय\-थु\-बा\-हु\-ल्य\-म\-व\-पी\-ड्य\-माने}}{\lemma{\skt{श्वय\-थु\-बा\-हु\-ल्यम्}}\Afootnote{A H; \skt{स्वय\-थु\-बा\-हु\-ल्यम्}\ N}}}{\lemma{\skt{अवपीड्यमाने}}\Afootnote{N H; \skt{स्पन्द\-न\-वि\-व\-र्त\-न\-स्प\-र्शा\-स\-हि\-ष्णु\-त्व\-म\-व\-पी\-ड्य\-माने}\ A}}\skt{ शब्दः }\edtext{\skt{स्रस्ताङ्गता}}{\Afootnote{A H; \skt{श्रस्तांगता}\ N}}\skt{ }\edtext{\edtext{\skt{विवि\-ध\-वे\-द\-ना\-प्र\-वृ\-त्ति\-रिति}}{\lemma{\skt{विवि\-ध\-वे\-द\-ना\-प्र\-वृ\-त्तिर्}}\Afootnote{N H; \skt{विवि\-ध\-वे\-द\-ना\-प्रा\-दु\-र्भावः सर्वास्ववस्थासु न शर्म}\ A}}}{\lemma{\skt{इति}}\Afootnote{N H; \skt{लाभ इति}\ A}}\skt{ ॥ }\edtext{\skt{सामान्यतः}}{\Afootnote{N H; \skt{समासेन}\ A}}\skt{ काण्ड\-भ\-ग्न\-ल\-क्ष\-ण\-मुक्तं ॥}\pend

\pstart \ledsidenote{\passref{2.15.10}}\edtext{\skt{विशेषतस्तु}}{\Afootnote{N H; \skt{विशेषस्तु}\ A}}\skt{ }\edtext{\edtext{\skt{सम्मू\-ढ\-मु\-भ\-य\-तोऽस्थि}}{\lemma{\skt{सम्मूढम्}}\Afootnote{N H; \skt{संमूढ}\ A}}}{\lemma{\skt{उभयतो}}\Afootnote{N H; \skt{मुभयतो}\ A}}\skt{ मध्ये }\edtext{\skt{भग्नं}}{\Afootnote{H; \skt{भ(ल)ग्नं}\ A; \skt{भग्न}\ N}}\skt{ }\edtext{\skt{ग्रन्थिरिवोन्नतं}}{\Afootnote{H; \skt{ग्रन्थिरिवोन्नतं}\ A; \skt{गन्थिरिवोन्नं}\ N}}\skt{ }\edtext{\skt{कर्क्कटकं}}{\Afootnote{A H; \skt{कर्कटकः ॥}\ N}}\skt{ ॥ }\edtext{\skt{अश्व\-क\-र्ण्ण\-व\-दु\-द्ग\-त\-म\-श्व\-कर्ण्णं}}{\Afootnote{H; \skt{अश्व\-क\-र्ण\-व\-दु\-द्ग\-त\-म\-श्व\-क\-र्णकं,}\ A; \skt{अश्व\-क\-र्ण्ण\-व\-दु\-द्ग\-त\-म\-श्व\-कर्ण्णं}\ N}}\skt{ ॥ }\edtext{\skt{शब्द\-स्प\-र्शा\-भ्या\-ञ्चू\-र्ण्णि\-त\-म\-व\-ग\-च्छेत्}}{\lemma{\skt{शब्द\-स्प\-र्शा\-भ्या\-ञ्चू\-र्ण्णि\-तम्}}\Afootnote{H; \skt{स्पृश्यमानं शब्द\-व\-च्चू\-र्णि\-तम्}\ A; \skt{सब्दस्पर्शाभ्यां चूर्ण्णितम्}\ N}}\skt{ ॥ }\edtext{\edtext{\skt{पिच्चि\-त\-म्पृ\-थु\-ल\-ता\-ङ्ग\-त\-म\-ल्प\-शोफं}}{\lemma{\skt{पृथु\-ल\-ता\-ङ्ग\-तम्}}\Afootnote{N H; \skt{पृथुतां गतम्}\ A}}}{\lemma{\skt{अल्पशोफं}}\Afootnote{H; \skt{अनल्पशोफं}\ A N}}\skt{ ॥ }\edtext{\edtext{\edtext{\skt{पार्श्व\-तोऽ\-न्व\-स्थि\-ही\-नो\-द्ग\-तोऽ\-अ\-स्थि\-च्छ\-लि\-त\-क\-न्ना\-म\-वे\-ल्लिते}}{\lemma{\skt{पार्श्वतो}}\Afootnote{N H; \skt{पार्श्वयोर्}\ A}}}{\lemma{\skt{ऽन्व\-स्थि\-ही\-नो\-द्गतो}}\Afootnote{N H; \skt{अस्थि हीनोद्गतम्}\ A}}}{\lemma{\skt{ऽअ\-स्थि\-च्छ\-लि\-त\-क\-न्ना\-म\-वे\-ल्लिते}}\Afootnote{H; \skt{अस्थिच्छलितं, वेल्लते}\ A; \skt{स्थिच्छ\-ल्लि\-त\-क\-न्ना\-म\-वे\-ल्लिते}\ N}}\skt{ }\edtext{\skt{प्रकम्प्यमानं}}{\Afootnote{N H; \skt{प्रकम्पमानं}\ A}}\skt{ ॥ }\edtext{\edtext{\skt{काण्ड\-भ\-ग्न\-म\-स्थि\-निः\-शे\-ष\-च्छि\-न्ना\-स्थ्य\-व\-यवः}}{\lemma{\skt{काण्डभग्नम्}}\Afootnote{A H; \skt{कण्डभग्नम्}\ N}}}{\lemma{\skt{अस्थि\-निः\-शे\-ष\-च्छि\-न्ना\-स्थ्य\-व\-यवः}}\Afootnote{H; \skt{अस्थ्यवयवो}\ A; \skt{अस्थि निशे\-ष\-च्छि\-न्ना\-स्थ्य\-व\-यवः ॥}\ N}}\skt{ ॥ }\edtext{\skt{अस्थि\-म\-ध्य\-प्र\-वि\-ष्टा\-ति\-पा\-ति\-त\-संज्ञः}}{\Afootnote{H; \skt{ऽस्थि\-म\-ध्य\-म\-नु\-प्रवि}\ A; \skt{अस्थि\-म\-ध्य\-प्र\-वि\-ष्ठा\-ति\-पा\-तितः संज्ञः ॥}\ N}}\skt{ ॥ }\edtext{\edtext{\edtext{\skt{क्षत\-भ\-ग्न\-मु\-न्न\-ह्य\-मा\-न\-म\-ज्ज\-म्म\-ज्जा\-नु\-गतं}}{\lemma{\skt{क्षतभग्नम्}}\Afootnote{N H; \skt{श्य मज्जानम्}\ A}}}{\lemma{\skt{उन्नह्यमानम्}}\Afootnote{H; \skt{उन्नह्यती}\ A; \skt{उन्नह्यमान}\ N}}}{\lemma{\skt{अज्ज\-म्म\-ज्जा\-नु\-गतं}}\Afootnote{H; \skt{ति मज्जानुगतम्, अस्थि निःशे\-ष\-त\-श्छि\-न्न\-म\-ति\-पा\-ति\-तम्,}\ A; \skt{मज्जा\-म्म\-ज्जा\-नु\-गतं ।}\ N}}\skt{ ॥ }\edtext{\skt{आभुग्नमिव}}{\lemma{\skt{इव}}\Afootnote{N H; \emph{om.}\ A}}\skt{ }\edtext{\skt{यद्विमुक्तास्थि}}{\Afootnote{N H; \skt{अविमुक्तास्थि}\ A}}\skt{ }\edtext{\edtext{\skt{तद्वक्रन्नाम}}{\lemma{\skt{तद्वक्रन्}}\Afootnote{H; \skt{वक्रम्,}\ A; \skt{त वक्रन्}\ N}}}{\lemma{\skt{नाम}}\Afootnote{H; \emph{om.}\ A; \skt{नामः ॥}\ N}}\skt{ ॥ अन्य\-त\-र\-पा\-र्श्वा\-व\-शिष्टं छिन्नं ॥ }\edtext{\skt{स्फुटितम}}{\Afootnote{H; \skt{पाटितम}\ A; \skt{स्फुटितं अ}\ N}}\skt{ }\edtext{\skt{नुब\-हु\-दा\-रि\-त\-वे\-द\-ना\-वान् ॥ }\edtext{\skt{शूक\-पू\-र्ण्ण\-मि\-वा\-ध्मा\-त\-म्वि\-पु\-लै\-कदा}}{\lemma{\skt{विपुलैकदा}}\Afootnote{N H; \skt{विपुलं}\ A}}}{\lemma{\skt{नुब\-हु\-दा\-रि\-त\-वे\-द\-ना\-वान् ॥ शूकपूर्ण्णम्}}\Afootnote{H; \skt{णुब\-हु\-वि\-दा\-रितं वेदनावच्च, शूकपूर्णम्}\ A; \skt{नुब\-हु\-द\-रि\-त\-वे\-द\-ना\-वान् ॥ शूकपूर्ण्णम्}\ N}}\skt{ }\edtext{\edtext{\skt{विस्फाटितन्नाम}}{\lemma{\skt{विस्फाटितन्}}\Afootnote{H; \skt{विस्फुटितं}\ A; \skt{रि स्फाटितन्}\ N}}}{\lemma{\skt{नाम}}\Afootnote{N H; \skt{स्फुटितमिति ॥}\ A}}\skt{ ॥}\pend

\pstart \ledsidenote{\passref{2.15.11}}\skt{तेषु }\edtext{\skt{चूर्ण्णि\-तं\-च्छि\-न्ना\-ति\-पा\-ति\-त\-म\-ज्जा\-नु\-ग\-तानि}}{\Afootnote{H; \skt{चूर्णि\-त\-च्छि\-न्ना\-ति\-पा\-ति\-त\-म\-ज्जा\-नु\-ग\-तानि}\ A; \skt{चूर्ण्णतं च्छिन्ना\-ति\-पा\-ति\-त\-म\-ज्जा\-नु\-ग\-तानि}\ N}}\skt{ }\edtext{\skt{कृच्छ्रसाध्धानि}}{\Afootnote{H; \skt{कृच्छ्रसाध्यानि}\ A N}}\skt{ । }\edtext{\skt{कृशा\-ति\-वृ\-द्ध\-बा\-ला\-स\-हा\-यानां}}{\Afootnote{N H; \skt{कृश\-वृ\-द्ध\-बा\-लानां}\ A}}\skt{ }\edtext{\skt{क्षत\-क्षी\-ण\-कु\-ष्टि\-श्वा\-सिनां}}{\Afootnote{N H; \skt{क्षत\-क्षी\-ण\-कु\-ष्ठि\-श्वा\-सिनां}\ A}}\skt{ }\edtext{\skt{सव्रणाः}}{\Afootnote{N H; \emph{om.}\ A}}\skt{ }\edtext{\skt{निःस\-न्ध्यु\-प\-ग\-तानि}}{\Afootnote{H; \skt{सन्ध्युपगतं}\ A; \skt{नि सन्ध्युपगतानि}\ N}}\skt{ चेति ॥}\pend

\pstart \ledsidenote{\passref{2.15.12}}\edtext{\skt{भवति}}{\Afootnote{H; \skt{भवन्ति}\ A; \emph{om.}\ N}}\skt{ }\edtext{\skt{चात्र}}{\Afootnote{A H; \skt{भभवति चात्र ॥}\ N}}\skt{ ॥ भिन्नङ्कपालं }\edtext{\skt{कट्यां तु}}{\Afootnote{N H; \skt{कठ्यां तु}\ A}}\skt{ }\edtext{\edtext{\skt{सन्धिमुक्तं}}{\lemma{\skt{सन्धिम्}}\Afootnote{N H; \skt{सन्धि}\ A}}}{\lemma{\skt{उक्तं}}\Afootnote{N H; \skt{मुक्तं}\ A}}\skt{ तथा च्युतं । जघ\-न\-म्प्र\-ति\-पि\-ष्टञ्च }\edtext{\skt{वर्ज्जयीत}}{\Afootnote{N H; \skt{वर्जयेत्त}\ A}}\skt{ }\edtext{\skt{विचक्षणः}}{\Afootnote{N H; \skt{च्चिकित्सकः ॥}\ A}}\skt{ ।}\pend

\pstart\versepar \ledsidenote{\passref{2.15.13}}\edtext{\skt{असं\-क्लि\-ष्ट\-ङ्क\-पा\-लञ्च}}{\Afootnote{H; \skt{असंश्लिष्टं कपालं तु}\ A; \skt{अलं क्लिष्टङ्कपालञ्च}\ N}}\skt{ ललाटे चूर्ण्णितञ्च यत् ॥ }\pend
\pstart\versepar \skt{भग्नं स्तनान्तरे }\edtext{\skt{शंखं}}{\Afootnote{H; \skt{शङ्खे}\ A; \skt{सङ्ख्यम्}\ N}}\skt{ पृष्ठे मूर्द्ध्नि च वर्ज्जयेत् ॥}\pend

\pstart \ledsidenote{\passref{2.15.14cd}}\edtext{\edtext{\skt{सम्य\-क्स\-न्धि\-त\-म\-प्य\-स्थि\-दु\-र्न्नि\-क्षे\-प्त\-नि\-ब\-न्ध\-नात्}}{\lemma{\skt{सम्यक्सन्धितम्}}\Afootnote{N H; \skt{सम्यग्यमितम्}\ A}}}{\lemma{\skt{अस्थि\-दु\-र्न्नि\-क्षे\-प्त\-नि\-ब\-न्ध\-नात्}}\Afootnote{N H; \skt{अस्थि दुर्न्या\-सा\-द्दु\-र्नि\-ब\-न्ध\-नात् ॥}\ A}}\skt{ ।}\pend

\pstart \ledsidenote{\passref{2.15.15ab}}\edtext{\skt{संक्षोभाद्वापि}}{\Afootnote{H; \skt{सङ्खोभाद्वाऽपि}\ A; \skt{सक्षोभाद्वापि}\ N}}\skt{ }\edtext{\edtext{\skt{यद्ग\-च्छे\-त्वि\-क्रि\-या\-न्त\-मपि}}{\lemma{\skt{यद्ग\-च्छे\-त्वि\-क्रि\-या\-न्तम्}}\Afootnote{H; \skt{यद्गच्छेद्विक्रियां त}\ A N}}}{\lemma{\skt{अपि}}\Afootnote{H; \skt{च्च}\ A; \skt{वि}\ N}}\skt{ वर्ज्जयेत् ॥}\pend

\pstart \ledsidenote{\passref{2.15.16cd}}\skt{तरुणास्थीनि }\edtext{\skt{नाम्यन्ते}}{\Afootnote{N H; \skt{नम्यन्ते}\ A}}\skt{ भज्यन्ते नलकानि तु ।}\pend

\pstart \ledsidenote{\passref{2.15.17}}\edtext{\skt{कपा\-ला\-न्वि\-भ\-ज्यन्ते}}{\Afootnote{H; \skt{कपा\-ला\-न्वि\-भि\-द्यन्ते}\ A; \skt{कपालानि विभज्यन्ते}\ N}}\skt{ }\edtext{\skt{स्फुल्लन्ति}}{\Afootnote{N H; \skt{स्फुटन्ति}\ A}}\skt{ रुचकानि }\edtext{\skt{तु}}{\Afootnote{N H; \skt{च ॥}\ A}}\skt{ ॥}\pend

\pstart \ledsidenote{\passref{2.15.trailer}}\edtext{\skt{इति भग्ननिदाने}}{\Afootnote{H; \skt{इति सुश्रु\-त\-सं\-हि\-तायां निदानस्थाने भग्ननिदानं नाम}\ A; \skt{१५ ॥ ० ॥}\ N}}\skt{ }\edtext{\skt{पञ्चदशमः}}{\Afootnote{H; \skt{पञ्चदशोऽध्यायः ॥१५ ॥}\ A; \emph{om.}\ N}}\skt{ ॥}\pend

\endnumbering
\beginnumbering
\eledsection*{Adhyāya 16}
\pstart \ledsidenote{\passref{2.16.1}}\skt{अथातो }\edtext{\edtext{\skt{मुख\-रो\-ग\-नि\-दा\-न\-म्व्या\-ख्या\-स्यामः}}{\lemma{\skt{मुख\-रो\-ग\-नि\-दा\-नम्}}\Afootnote{H; \skt{मुखरोगाणां निदानं}\ A; \skt{मुख\-रो\-ग\-नि\-दा\-ना\-ध्यायं}\ N}}}{\lemma{\skt{व्याख्यास्यामः}}\Afootnote{A H; \skt{वाव्याव्याख्यास्यामः ॥}\ N}}\skt{ ॥}\pend

\pstart \ledsidenote{\passref{2.16.3}}\skt{मुखरोगाः पञ्च }\edtext{\edtext{\skt{षष्टिर्भवन्ति}}{\lemma{\skt{षष्टिर्}}\Afootnote{N H; \skt{षष्टि}\ A}}}{\lemma{\skt{भवन्ति}}\Afootnote{N H; \skt{र्भवन्ति}\ A}}\skt{ ॥ }\edtext{\edtext{\skt{सप्त\-स्वा\-य\-त\-नेषु}}{\lemma{\skt{सप्तस्व्}}\Afootnote{N H; \skt{सप्त}\ A}}}{\lemma{\skt{आयतनेषु}}\Afootnote{N H; \skt{स्वायतनेषु ।}\ A}}\skt{ ॥ तत्रायतनानि ओष्ठौ दन्तमूलानि दन्ता }\edtext{\skt{जिह्वा}}{\Afootnote{A H; \skt{जिजिह्वा}\ N}}\skt{ }\edtext{\skt{तालू}}{\Afootnote{H; \skt{तालु}\ A N}}\skt{ }\edtext{\skt{कण्ठ}}{\Afootnote{N H; \skt{कण्ठः,}\ A}}\skt{ सर्व्वाणि }\edtext{\skt{चेति}}{\Afootnote{A H; \skt{च्चेति}\ N}}\skt{ ॥ }\edtext{\edtext{\skt{तत्राष्टावोष्ठयोः}}{\lemma{\skt{तत्राष्टाव्}}\Afootnote{H; \skt{तत्राष्टा}\ A; \skt{तराष्टा}\ N}}\skt{ ॥पञ्चदश}}{\lemma{\skt{ओष्ठयोः ॥पञ्चदश}}\Afootnote{H; \skt{वोष्ठयोः, पञ्चदश}\ A; \skt{वोष्थयोः पञ्चदश}\ N}}\skt{ }\edtext{\skt{दन्तमूले}}{\Afootnote{H; \skt{दन्तमूलेषु,}\ A; \skt{दन्तमूल}\ N}}\skt{ ॥ अष्टौ दन्तेषु ॥ पञ्चजिह्वायां ॥ }\edtext{\skt{नवतालुनि}}{\Afootnote{A H; \skt{नव तालूनि}\ N}}\skt{ ॥ सप्तदशकण्ठे ॥ }\edtext{\edtext{\skt{त्रय\-स्स\-र्व्वे\-ष्वा\-य\-त\-नेषु}}{\lemma{\skt{सर्व्वेष्व्}}\Afootnote{H; \skt{सर्वे}\ A; \skt{सर्वेष्}\ N}}}{\lemma{\skt{आयतनेषु}}\Afootnote{H; \skt{ष्वायतनेषु ॥}\ A; \skt{आयतनेष्विति ।}\ N}}\skt{ ॥}\pend

\pstart \ledsidenote{\passref{2.16.4}}\edtext{\skt{तत्रोष्ठे}}{\Afootnote{H; \skt{तत्रौष्ठ}\ A N}}\skt{ }\edtext{\edtext{\skt{प्रको\-पा\-द्वा\-त\-पि\-त्त\-क\-फ\-स\-न्नि\-पा\-त\-र\-क्त\-मा\-न्स\-मे\-दोऽ\-भि\-घा\-त\-नि\-मित्ताः}}{\lemma{\skt{प्रको\-पा\-द्वा\-त\-पि\-त्त\-क\-फ\-स\-न्नि\-पा\-त\-र\-क्त\-मा\-न्स\-मेदो}}\Afootnote{H; \skt{प्रकोपा वात\-पि\-त्त\-श्ले\-ष्म\-स\-न्नि\-पा\-त\-र\-क्त\-मां\-स\-मेदो}\ A; \skt{प्रको\-पा\-द्वा\-त\-पि\-त्त\-क\-फ\-स\-न्नि\-पा\-त\-र\-क्त\-मां\-स\-मेदो}\ N}}}{\lemma{\skt{ऽभि\-घा\-त\-नि\-मित्ताः}}\Afootnote{H; \skt{भिघातनिमित्ताः}\ A N}}\skt{ ॥}\pend

\pstart\versepar \ledsidenote{\passref{2.16.5}}\skt{कर्क्कशौ परुषौ }\edtext{\skt{स्तब्धौ}}{\Afootnote{A H; \skt{स्त}\ N}}\skt{ }\edtext{\skt{सम्प्रा\-प्ता\-नि\-ल\-वे\-दनौ}}{\Afootnote{H; \skt{कृष्णौ तीव्ररुगन्वितौ ।}\ A; \skt{संप्रा\-प्ता\-नि\-ल\-वे\-दनु ।}\ N}}\skt{ । }\pend
\pstart\versepar \skt{दाल्येते }\edtext{\skt{परिपोट्येते}}{\Afootnote{H; \skt{परिपाठ्येते}\ A; \skt{परिपाट्येते । ह्य्}\ N}}\skt{ }\edtext{\skt{ओष्ठौ\-मा\-रु\-त\-को\-पतः}}{\Afootnote{N H; \skt{ह्योष्ठौ मारुतकोपतः ॥}\ A}}\skt{ ॥}\pend

\pstart\versepar \ledsidenote{\passref{2.16.6}}\edtext{\skt{चीयेते}}{\Afootnote{H; \skt{आचितौ}\ A; \skt{चीयते}\ N}}\skt{ }\edtext{\skt{पिटकाभिश्च}}{\Afootnote{H; \skt{पिडाकाभिस्तु}\ A; \skt{पिडकाभिश्च}\ N}}\skt{ }\edtext{\skt{सरुजाभिः}}{\Afootnote{N H; \skt{सर्ष\-पा\-कृ\-ति\-भिर्}\ A}}\skt{ }\edtext{\skt{समन्ततः}}{\Afootnote{N H; \skt{भृशम् ।}\ A}}\skt{ । }\pend
\pstart\versepar \skt{}\edtext{\skt{सदा\-ह\-पा\-क\-पि\-टकौ}}{\Afootnote{H; \skt{सदा\-ह\-पा\-क\-सं\-स्रावौ}\ A; \skt{सदा\-ह\-पा\-क\-पि\-डकौ}\ N}}\skt{ }\edtext{\skt{पीताभासौ च}}{\Afootnote{N H; \skt{नीलौ पीतौ च}\ A}}\skt{ पित्ततः ॥}\pend

\pstart\versepar \ledsidenote{\passref{2.16.7}}\edtext{\skt{सवर्ण्णाभिश्च}}{\Afootnote{H; \skt{सवर्णाभिस्तु}\ A N}}\skt{ }\edtext{\skt{चीयेते}}{\Afootnote{A H; \skt{चीयते}\ N}}\skt{ }\edtext{\skt{पिट\-का\-भि\-र\-वे\-दनौ}}{\lemma{\skt{पिटकाभिर्}}\Afootnote{H; \skt{पिडकाभिर्}\ A N}}\skt{ । }\pend
\pstart\versepar \skt{}\edtext{\skt{भवतस्तु}}{\Afootnote{N\textsuperscript{pc} H; \skt{कण्डूमन्तौ}\ A; \skt{भव स्तु}\ N\textsuperscript{ac}}}\skt{ }\edtext{\skt{कफादोष्ठौ}}{\Afootnote{N H; \skt{कफाच्छूनौ}\ A}}\skt{ पिच्छिलौ शीतलौ }\edtext{\skt{गुरू}}{\Afootnote{A H; \skt{गुरु ॥}\ N}}\skt{ ॥}\pend

\pstart\versepar \ledsidenote{\passref{2.16.8}}\skt{सकृत्कृष्णौ सकृत्पीतौ }\edtext{\edtext{\skt{सकृत्स्वेतौ}}{\lemma{\skt{सकृत्स्व्}}\Afootnote{H; \skt{सकृ}\ A N}}}{\lemma{\skt{एतौ}}\Afootnote{H; \skt{च्छ्वेतौ}\ A; \skt{च्छ्वेतौ}\ N}}\skt{ तथैव च । }\pend
\pstart\versepar \skt{सन्निपातेन }\edtext{\edtext{\skt{विज्ञे\-या\-व\-ने\-क\-पि\-ड\-का\-चितौ}}{\lemma{\skt{विज्ञेयाव्}}\Afootnote{H; \skt{विज्ञेया}\ A N}}}{\lemma{\skt{अने\-क\-पि\-ड\-का\-चितौ}}\Afootnote{H; \skt{वने\-क\-पि\-डि\-का\-चितौ ॥}\ A; \skt{वने\-क\-पि\-ड\-का\-चितौ ॥}\ N}}\skt{ ॥}\pend

\pstart\versepar \ledsidenote{\passref{2.16.9}}\edtext{\skt{खर्ज्जू\-र\-र\-क्त\-व\-र्ण्णाभिः}}{\Afootnote{H; \skt{खर्जू\-र\-फ\-ल\-व\-र्णाभिः}\ A; \skt{खर्जू\-र\-र\-ङ्ग\-व\-र्ण्णाभिः}\ N}}\skt{ }\edtext{\edtext{\skt{पिट\-का\-भि\-र्न्नि\-पी\-डितौ}}{\lemma{\skt{पिटकाभिर्}}\Afootnote{H; \skt{पिडकाभिः}\ A; \skt{पिडकाभि}\ N}}}{\lemma{\skt{न्निपीडितौ}}\Afootnote{N H; \skt{समाचितौ ।}\ A}}\skt{ । }\pend
\pstart\versepar \skt{}\edtext{\skt{रक्तो\-प\-सृ\-ष्ट\-रु\-धिरं}}{\Afootnote{H; \skt{रक्तोपसृष्टौ रुधिरं}\ A N}}\skt{ स्रवतः शोणितप्रभौ ॥}\pend

\pstart\versepar \ledsidenote{\passref{2.16.10}}\edtext{\skt{गुरू}}{\Afootnote{H; \skt{मासदुष्टौ गुरू}\ A; \skt{गुरु}\ N}}\skt{ स्थूलौ }\edtext{\skt{मांसदुष्टौ}}{\Afootnote{N H; \emph{om.}\ A}}\skt{ }\edtext{\skt{मांस\-पि\-ण्ड\-व\-दु\-द्गतौ}}{\Afootnote{N H; \skt{मांस\-पि\-ण्ड\-व\-दु\-द्गतौ ।}\ A}}\skt{ । }\pend
\pstart\versepar \skt{जन्त\-व\-श्चा\-त्र\-मू\-र्च्छन्ति }\edtext{\skt{नरस्योभयतो}}{\Afootnote{N H; \skt{सृक्कस्योभयतो}\ A}}\skt{ मुखात् ॥}\pend

\pstart \ledsidenote{\passref{2.16.11ab}}\edtext{\edtext{\skt{सर्प्पि\-र्म्म\-ण्ड\-प्र\-ती\-काशौ}}{\Afootnote{H; \skt{सर्पि\-र्म\-ण्ड\-प्र\-ती\-काशौ}\ N}}\skt{ मेदसा}}{\lemma{\skt{… मे}}\Afootnote{\emph{om.}\ A}}\skt{ }\edtext{\skt{कण्डुरौ}}{\Afootnote{H; \skt{घृतमण्डाभौ}\ A; \emph{om.}\ N}}\skt{ }\edtext{\skt{गुरू}}{\Afootnote{H; \skt{कण्डूमन्तौ स्थिरौ मृदू ॥}\ A; \skt{कण्डुरौ गुरु ॥}\ N}}\skt{ ।}\pend

\pstart \ledsidenote{\passref{2.16.12ab}}\edtext{\skt{ओष्ठौ\-प\-र्य\-व\-दी\-र्येते}}{\Afootnote{H; \skt{क्षतजाभौ विदीर्येते}\ A; \skt{ओष्ठौ पर्यवदीर्यते}\ N}}\skt{ }\edtext{\skt{पोट्येते}}{\Afootnote{H; \skt{पाठ्येते}\ A; \skt{पाट्येते}\ N}}\skt{ चाभिघाततः ॥}\pend

\pstart \ledsidenote{\passref{2.16.11cd}}\skt{अच्छं }\edtext{\skt{स्फटिकसंकाशं}}{\Afootnote{H; \skt{स्फटिकसङ्काश}\ A N}}\skt{ }\edtext{\skt{सस्रावं}}{\Afootnote{H; \skt{मास्रावं}\ A N}}\skt{ }\edtext{\skt{स्रवतो}}{\Afootnote{H; \skt{स्रवयो}\ A; \skt{स्रंवतो}\ N}}\skt{ }\edtext{\skt{भृशं}}{\Afootnote{N H; \skt{गुरू ॥}\ A}}\skt{ ।}\pend

\pstart\versepar \ledsidenote{\passref{2.16.11.add}}\edtext{\skt{दन्त\-च्छ\-दा\-व\-दी\-र्येते}}{\Afootnote{H; \skt{दन्त\-च्छ\-दा\-व\-दी\-र्यते}\ N}}\skt{ पोट्येते चासकृत्पुनः ॥ }\pend
\pstart\versepar \skt{ततो व्रणः }\edtext{\skt{सुसंरूढो}}{\Afootnote{H; \skt{सुसंरुढो}\ N}}\skt{ मृदु\-त्व\-मु\-प\-ग\-च्छति ॥}\pend

\pstart \ledsidenote{\passref{2.16.13}}\skt{दन्त\-मू\-ल\-ग\-तास्तु शीतादो }\edtext{\skt{दन्तपुप्पुटो}}{\Afootnote{H; \skt{दन्तपुप्पुटको,}\ A; \skt{दन्तपुपुटो}\ N}}\skt{ }\edtext{\skt{दन्तवेष्टः}}{\Afootnote{H; \skt{दन्तवेष्टकः}\ A N}}\skt{ }\edtext{\skt{सौशिरो}}{\Afootnote{H; \skt{शौषिरो,}\ A; \skt{शौशिरो}\ N}}\skt{ }\edtext{\skt{महासौशिरः}}{\Afootnote{H; \skt{महाशौषिरः}\ A N}}\skt{ परिदरः । उपकुशः । दन्तवैदर्ब्भः । }\edtext{\skt{खलिवर्द्धनः}}{\Afootnote{N H; \skt{वर्धनः,}\ A}}\skt{ । अधिमांसः । नाड्यः }\edtext{\skt{पञ्चेति}}{\Afootnote{A H; \skt{पञ्चति ॥}\ N}}\skt{ ॥}\pend

\pstart\versepar \ledsidenote{\passref{2.16.14}}\skt{शोणितं दन्तवेष्टेभ्यो }\edtext{\skt{यस्या\-क\-स्मा\-त्प्र\-व\-र्त्तते}}{\Afootnote{A H; \skt{यस्यो कस्मात्प्रवर्तते ।\textbackslash\{\}}\ N}}\skt{ । }\pend
\pstart\versepar \skt{}\edtext{\skt{दुर्ग्गन्धी}}{\Afootnote{A H; \skt{दुर्गन्दी}\ N}}\skt{ निसकृष्णानि प्रक्लेदीनि }\edtext{\skt{मृदूनि}}{\Afootnote{A H; \skt{भवन्ति}\ N}}\skt{ च ॥}\pend

\pstart\versepar \ledsidenote{\passref{2.16.15}}\skt{दन्तमांसानि }\edtext{\skt{पच्यन्ते}}{\Afootnote{H; \skt{शीर्यन्ते}\ A; \skt{पर्यन्ते}\ N}}\skt{ पचन्ति च परस्परं । }\pend
\pstart\versepar \skt{}\edtext{\skt{शीता\-दो\-ना\-म\-स\-व्याधिः}}{\Afootnote{H; \skt{शीतादो नाम स व्याधिः}\ A N}}\skt{ कफ\-शो\-णि\-त\-स\-म्भवः ॥}\pend

\pstart\versepar \ledsidenote{\passref{2.16.16}}\edtext{\edtext{\skt{दन्तयोस्त्रिषु}}{\lemma{\skt{दन्तयोस्}}\Afootnote{H; \skt{दन्तयो}\ A N}}\skt{ वा}}{\lemma{\skt{त्रिषु वा}}\Afootnote{H; \skt{स्त्रिषु वा}\ A; \skt{स्तृषु वा}\ N}}\skt{ यस्य श्वयथुः }\edtext{\skt{स्पर्शना}}{\Afootnote{H; \skt{सरुजो}\ A; \skt{स्पर्शनार्थ}\ N}}\skt{ }\edtext{\skt{सहः}}{\Afootnote{N H; \skt{महान् ।}\ A}}\skt{ । }\pend
\pstart\versepar \skt{}\edtext{\skt{दन्तपुप्पुटको}}{\Afootnote{A H; \skt{दन्तपुप्र्पटको}\ N}}\skt{ }\edtext{\skt{नाम}}{\Afootnote{N H; \emph{om.}\ A}}\skt{ }\edtext{\skt{सव्याधिः}}{\Afootnote{N H; \skt{ज्ञेयः}\ A}}\skt{ }\edtext{\skt{कफरक्तजः}}{\Afootnote{H; \skt{कफ\-र\-क्त\-नि\-मि\-त्तजः ॥}\ A; \skt{कफरक्तजाः ॥}\ N}}\skt{ ॥}\pend

\pstart\versepar \ledsidenote{\passref{2.16.17}}\skt{स्रवन्ति पूयरुधिरं चलादन्ता भवन्ति च । }\pend
\pstart\versepar \skt{दन्तवेष्टः सविज्ञेयो दुष्ट\-शो\-णि\-त\-स\-म्भवः ॥}\pend

\pstart \ledsidenote{\passref{2.16.18}}\edtext{\edtext{\skt{श्वय\-थू\-र्द्द\-न्त\-मू\-लेषु}}{\lemma{\skt{श्वयथूर्}}\Afootnote{H; \skt{श्वयथु}\ A; \skt{श्वयथुर्}\ N}}}{\lemma{\skt{द्दन्तमूलेषु}}\Afootnote{N H; \skt{र्दन्तमूलेषु}\ A}}\skt{ }\edtext{\skt{रुजा\-वा\-न्क\-फ\-स\-म्भवः}}{\Afootnote{N H; \skt{रुजा\-वा\-न्क\-फ\-र\-क्तजः ।}\ A}}\skt{ । लालास्रावी स विज्ञेयः }\edtext{\edtext{\skt{सौशि\-रो\-ना\-म\-ना\-मतः}}{\lemma{\skt{सौशिरोनाम्}}\Afootnote{H; \skt{कण्डूमाञ्शौषिरो}\ A; \skt{शौशिरो ना}\ N}}}{\lemma{\skt{अनामतः}}\Afootnote{H; \skt{गदः ॥}\ A; \skt{म नामतः ॥}\ N}}\skt{ ॥}\pend

\pstart \ledsidenote{\passref{2.16.19}}\skt{दन्ताश्चलन्ति }\edtext{\skt{वेष्टभ्यस्तालु}}{\lemma{\skt{वेष्टभ्यस्}}\Afootnote{H; \skt{वेष्टेभ्यस्}\ A N}}\skt{ }\edtext{\skt{चाप्यवदीर्यते}}{\lemma{\skt{अवदीर्यते}}\Afootnote{N H; \skt{अवदीर्यते । दन्तमांसानि पच्यन्ते मुखं च परिपीड्यते ॥}\ A}}\skt{ ।}\pend

\pstart\versepar \ledsidenote{\passref{2.16.20}}\edtext{\skt{यस्मिन्स}}{\Afootnote{A H; \skt{यस्मिम्स}\ N}}\skt{ सर्व्वजो }\edtext{\edtext{\skt{व्याधि\-र्म्म\-हा\-सौ\-शि\-र\-सं\-ज्ञितः}}{\lemma{\skt{व्याधिर्}}\Afootnote{H; \skt{व्याधि}\ A; \skt{वाधिर्}\ N}}}{\lemma{\skt{म्महा\-सौ\-शि\-र\-सं\-ज्ञितः}}\Afootnote{H; \skt{र्महा\-शौ\-षि\-र\-सं\-ज्ञकः ।}\ A; \skt{महा\-शौ\-सि\-र\-सं\-ज्ञेतः ॥}\ N}}\skt{ ॥ }\pend
\pstart\versepar \skt{दन्तमांसानि शीर्यन्ते }\edtext{\skt{यस्मिन्ष्ठीवति}}{\Afootnote{A H; \skt{यस्मिन्वति}\ N}}\skt{ }\edtext{\skt{चासकृत्}}{\Afootnote{N H; \skt{चाप्यसृक् ॥}\ A}}\skt{ ।}\pend

\pstart\versepar \ledsidenote{\passref{2.16.21}}\skt{पित्तासृक्कफजो }\edtext{\edtext{\skt{व्याधिर्ज्ञेयः}}{\lemma{\skt{व्याधिर्}}\Afootnote{N H; \skt{व्याधि}\ A}}}{\lemma{\skt{ज्ञेयः}}\Afootnote{H; \skt{र्ज्ञेयः}\ A; \skt{ज्ञेय या}\ N}}\skt{ परिदरो हि सः ॥ }\pend
\pstart\versepar \skt{}\edtext{\skt{वेष्टः}}{\Afootnote{H; \skt{चेष्टे}\ A; \skt{वेष्टेस्}\ N}}\skt{ }\edtext{\skt{सदाहम्पचति}}{\Afootnote{N H; \skt{षु दाहः पाकश्च}\ A}}\skt{ }\edtext{\skt{यस्य\-द\-न्ता\-श्च\-ल\-न्तिच।}}{\lemma{\skt{यस्य\-द\-न्ता\-श्च\-ल\-न्तिच}}\Afootnote{H; \skt{तेभ्यो दन्ताश्चलन्ति च ॥}\ A; \skt{यस्य दन्ताश्चलन्ति च ।}\ N}}\skt{ }\edtext{\skt{दाहेन}}{\Afootnote{N H; \emph{om.}\ A}}\skt{ }\edtext{\skt{वेष्टास्तप्ताश्च}}{\Afootnote{H; \emph{om.}\ A; \skt{वेष्टास्तप्ताश्च}\ N}}\skt{ }\edtext{\skt{शुष्यन्ते}}{\Afootnote{N H; \emph{om.}\ A}}\skt{ }\edtext{\skt{चाग्निना}}{\Afootnote{N H; \emph{om.}\ A}}\skt{ }\edtext{\skt{यथा}}{\Afootnote{N H; \emph{om.}\ A}}\skt{ ॥}\pend

\pstart \ledsidenote{\passref{2.16.22}}\edtext{\skt{अघट्टिताः}}{\Afootnote{N H; \skt{आघट्टिताः}\ A}}\skt{ }\edtext{\skt{प्रस्रवन्ति}}{\Afootnote{A H; \skt{प्तस्रवन्ति}\ N}}\skt{ }\edtext{\skt{शोणि\-त\-म्म\-न्द\-वे\-दनाः}}{\lemma{\skt{मन्दवेदनाः}}\Afootnote{H; \skt{मन्दवेदनाः । आध्मायन्ते स्रुते रक्ते मुखं पूति च जायते ॥}\ A; \skt{मन्दवेदना ।}\ N}}\skt{ ।}\pend

\pstart\versepar \ledsidenote{\passref{2.16.23}}\edtext{\skt{यस्मि\-न्सो\-प\-कु\-शो\-नाम}}{\Afootnote{H; \skt{यस्मिन्नुपकुशः स स्या}\ A; \skt{यस्मिं सोपकुशो नाम}\ N}}\skt{ }\edtext{\skt{पित्तरक्तकृतो}}{\Afootnote{N H; \skt{त्पित्तरक्तकृतो}\ A}}\skt{ गदः ॥ }\pend
\pstart\versepar \skt{घृष्टेषु दन्तमूलेषु संरम्भो जायते महान् ।}\pend

\pstart\versepar \ledsidenote{\passref{2.16.24}}\skt{भवन्ति }\edtext{\skt{दन्ताश्च}}{\Afootnote{H; \skt{च चला}\ A; \emph{om.}\ N}}\skt{ }\edtext{\skt{चलाः}}{\Afootnote{H; \skt{दन्ताः}\ A; \skt{दन्तास्चलास्}\ N}}\skt{ }\edtext{\skt{सवै\-द\-र्ब्भाऽ\-भि\-घा\-तजः}}{\Afootnote{H; \skt{स वैद\-र्भोऽ\-भि\-घा\-तजः ।}\ A; \skt{सच्वै\-द\-र्भा\-भि\-घा\-तजः ॥}\ N}}\skt{ ॥ }\pend
\pstart\versepar \skt{मारुतेनाधिको दन्तो जायते तीव्रवेदनः ।}\pend

\pstart\versepar \ledsidenote{\passref{2.16.25}}\edtext{\skt{खलि\-व\-र्द्ध\-न\-संज्ञो}}{\Afootnote{H; \skt{वर्धनः स मतो}\ A; \skt{खलि\-व\-र्ध\-न\-संज्ञो}\ N}}\skt{ }\edtext{\skt{सौ}}{\Afootnote{H; \skt{व्याधि}\ A; \skt{सौ}\ N}}\skt{ }\edtext{\edtext{\skt{जातेरुक्च}}{\lemma{\skt{जातेर्}}\Afootnote{H; \skt{र्जाते}\ A; \skt{जायते}\ N}}}{\lemma{\skt{उक्च}}\Afootnote{H; \skt{रुक्च}\ A N}}\skt{ प्रशाम्यति ॥ }\pend
\pstart\versepar \skt{हानव्ये पश्चिमे दन्ते }\edtext{\skt{महा\-शो\-फो\-म\-हा\-रुजः}}{\Afootnote{H; \skt{महाञ्छोथो महारुजः ॥}\ A; \skt{महासौफो महारुजाः ।}\ N}}\skt{ ।}\pend

\pstart\versepar \ledsidenote{\passref{2.16.26}}\skt{लालास्रावी }\edtext{\skt{कफकृतो}}{\Afootnote{A H; \skt{कफकृ}\ N}}\skt{ }\edtext{\skt{ज्ञेयः }\edtext{\skt{सोऽधिकमांसकः}}{\lemma{\skt{ऽधिकमांसकः}}\Afootnote{H; \skt{ऽधिमांसकः ।}\ A; \skt{धिकमांसकः ॥}\ N}}}{\lemma{\skt{ज्ञेयः सो}}\Afootnote{H; \skt{विज्ञेयः सो}\ A; \skt{ज्ञेयः सो}\ N}}\skt{ ॥ }\pend
\pstart\versepar \skt{दन्तमूलगता नाड्यः पञ्च ज्ञेया }\edtext{\skt{यथोदिताः}}{\Afootnote{N H; \skt{यथेरिताः ॥}\ A}}\skt{ ॥}\pend

\pstart \ledsidenote{\passref{2.16.27}}\skt{दन्तगतास्तु दालनः क्रिमिदन्त }\edtext{\skt{कोदन्त}}{\Afootnote{H; \skt{को दन्त}\ A N}}\skt{ हर्षो भञ्जनकः । }\edtext{\skt{दन्तशर्क्करा}}{\Afootnote{A H; \skt{शर्क्करा ।}\ N}}\skt{ । }\edtext{\skt{कपालिका}}{\Afootnote{A H; \skt{ककपालिकाः}\ N}}\skt{ । }\edtext{\skt{श्यावदन्तः}}{\Afootnote{H; \skt{श्यावदन्तको,}\ A; \skt{श्यावदन्ता ।}\ N}}\skt{ । }\edtext{\skt{हनुर्म्मोक्षश्चेति}}{\lemma{\skt{हनुर्}}\Afootnote{H; \skt{हनु}\ A N}}\skt{ ॥}\pend

\pstart\versepar \ledsidenote{\passref{2.16.28}}\edtext{\edtext{\skt{दार्य\-मा\-णे\-ष्वि\-व\-रुजा}}{\lemma{\skt{दार्यमाणेष्व्}}\Afootnote{H; \emph{om.}\ A; \skt{दीर्यमानेष्व्}\ N}}}{\lemma{\skt{इवरुजा}}\Afootnote{N H; \emph{om.}\ A}}\skt{ }\edtext{\skt{यस्य}}{\Afootnote{N H; \skt{दा}\ A}}\skt{ }\edtext{\skt{दन्तेषु}}{\Afootnote{H; \skt{ल्यन्ते बहुधा दन्ता यस्मिंस्तीव्ररु}\ A; \skt{दन्तेषु जायते ।}\ N}}\skt{ }\edtext{\skt{जायते}}{\Afootnote{H; \skt{गन्विताः ।}\ A; \emph{om.}\ N}}\skt{ । }\pend
\pstart\versepar \skt{}\edtext{\skt{दालनोनाम}}{\Afootnote{H; \skt{दालनः स इ}\ A; \skt{दालनो नामस्}\ N}}\skt{ }\edtext{\skt{सव्याधिः}}{\Afootnote{N H; \skt{ति ज्ञेयः}\ A}}\skt{ सदा\-ग\-ति\-नि\-मि\-त्तजः ।}\pend

\pstart\versepar \ledsidenote{\passref{2.16.29}}\edtext{\skt{कृष्णश्च्छिद्रश्च}}{\Afootnote{H; \skt{कृष्णश्च्छिद्री च}\ A; \skt{कृष्णस्छिद्रस्च}\ N}}\skt{ }\edtext{\skt{चलति}}{\Afootnote{N H; \skt{लः स्रावी}\ A}}\skt{ }\edtext{\skt{ससंरम्भो}}{\Afootnote{N H; \skt{सस्ंरम्भो}\ A}}\skt{ महारुजः । }\pend
\pstart\versepar \skt{}\edtext{\skt{अनि\-मि\-त्त\-रु\-जो\-वा\-ता\-त्स\-ज्ञेयः}}{\Afootnote{H; \skt{अनिमित्तरुजो वाताद्विज्ञेयः}\ A; \skt{अनिमित्तरुजो वातात्सङ्जेयः}\ N}}\skt{ }\edtext{\skt{क्रिमिदन्तकः}}{\Afootnote{N H; \skt{कृमिदन्तकः ॥}\ A}}\skt{ ॥}\pend

\pstart\versepar \ledsidenote{\passref{2.16.30}}\edtext{\skt{शीतमुष्णञ्च}}{\Afootnote{A H; \skt{शीतमुष्णञ्च}\ N}}\skt{ दशनाः सहन्ते स्पर्शनन्न च । }\pend
\pstart\versepar \skt{यस्य }\edtext{\skt{तद्दन्तहर्षन्तु}}{\Afootnote{N H; \skt{तं दन्तहर्षं तु}\ A}}\skt{ व्याधि\-म्वि\-द्या\-त्स\-मी\-र\-णात् ॥}\pend

\pstart\versepar \ledsidenote{\passref{2.16.31}}\edtext{\edtext{\skt{वक्त्र\-म्व\-क्र\-म्भ\-वे\-द्यस्य}}{\lemma{\skt{वक्त्रम्}}\Afootnote{A H; \skt{वक्रम्}\ N}}}{\lemma{\skt{भवेद्यस्य}}\Afootnote{H; \skt{भवेद्यस्मिन्}\ A; \skt{भवेद्यस्य}\ N}}\skt{ }\edtext{\skt{दन्तभङ्गश्च}}{\Afootnote{A H; \skt{दन्तभल्न्गश्च}\ N}}\skt{ }\edtext{\skt{जायते}}{\Afootnote{H; \skt{तीव्ररु}\ A; \skt{जजते ।}\ N}}\skt{ । }\pend
\pstart\versepar \skt{}\edtext{\skt{कफवातकृतो}}{\Afootnote{N H; \skt{क् । कफवातकृतो}\ A}}\skt{ व्याधिः स भञ्ज\-न\-क\-सं\-ज्ञितः ॥}\pend

\pstart\versepar \ledsidenote{\passref{2.16.32}}\edtext{\edtext{\skt{शर्क्करेव}}{\lemma{\skt{शर्क्कर्}}\Afootnote{N H; \skt{शर्क}\ A}}}{\lemma{\skt{एव}}\Afootnote{N H; \skt{रेव}\ A}}\skt{ स्थिरीभूतो मलो }\edtext{\skt{दन्तेषु}}{\Afootnote{N H; \skt{दनतेषु}\ A}}\skt{ यस्य वै । }\pend
\pstart\versepar \skt{सा दन्तानां }\edtext{\skt{गुणहरी}}{\Afootnote{A K H; \skt{गुणां हरी}\ N}}\skt{ विज्ञेया दन्तशर्क्करा ॥}\edtext{}{\lemma{\skt{… दन्तानां}}\Afootnote{\emph{om.}\ K}}\edtext{}{\lemma{\skt{यस्य वै । सा}}\Afootnote{A H; \skt{यस्य वैः । सा}\ N}}\pend

\pstart\versepar \ledsidenote{\passref{2.16.33}}\edtext{\edtext{\skt{कपालेष्विव}}{\lemma{\skt{कपालेष्व्}}\Afootnote{K N H; \emph{om.}\ A}}}{\lemma{\skt{इव}}\Afootnote{K N H; \emph{om.}\ A}}\skt{ }\edtext{\skt{दीर्यत्सुदन्तानां}}{\Afootnote{K N H; \skt{दलन्ति दन्}\ A}}\skt{ }\edtext{\skt{सैव}}{\Afootnote{K N H; \skt{तव}\ A}}\skt{ }\edtext{\skt{शर्क्करा}}{\Afootnote{K N H; \skt{ल्कानि यदा शर्करया सह । ज्ञेया}\ A}}\skt{ । }\pend
\pstart\versepar \skt{}\edtext{\skt{कपालिकेति}}{\Afootnote{K N H; \skt{कपालिका}\ A}}\skt{ }\edtext{\skt{पठिता}}{\Afootnote{K N H; \skt{सै}\ A}}\skt{ }\edtext{\skt{सदा}}{\Afootnote{K N H; \skt{व द}\ A}}\skt{ }\edtext{\skt{दन्तविनाशिनी}}{\Afootnote{N H; \skt{शनानां विनाशिनी ॥}\ A; \skt{दन्तविनाशनी ॥}\ K}}\skt{ ॥}\pend

\pstart\versepar \ledsidenote{\passref{2.16.34}}\edtext{\skt{असृण्मिश्रेण}}{\Afootnote{N H; \skt{योऽसृङ्निश्रेण}\ A; \skt{असृङ्भिश्रेण}\ K}}\skt{ पित्तेन }\edtext{\edtext{\skt{दग्धो\-द\-न्त\-स्त्व\-शे\-षतः}}{\lemma{\skt{दग्धोदन्तस्}}\Afootnote{H; \skt{दग्धो दन्तस्}\ A N; \skt{दग्धां दन्तस्}\ K}}}{\lemma{\skt{त्वशेषतः}}\Afootnote{A K H; \skt{त्ववशेषतः ।}\ N}}\skt{ । }\pend
\pstart\versepar \skt{}\edtext{\skt{श्याव\-ता\-न्नी\-ल\-ता\-म्वापि}}{\lemma{\skt{वापि}}\Afootnote{K N H; \skt{वाऽपि}\ A}}\skt{ गतः स श्यावदन्तकः ॥}\pend

\pstart\versepar \ledsidenote{\passref{2.16.35}}\edtext{\skt{वातेन }\edtext{\skt{तैस्तैर्ब्भावैश्च}}{\lemma{\skt{ब्भावैश्च}}\Afootnote{N H; \skt{भावैस्तु}\ A K}}}{\lemma{\skt{वातेन तैस्तैर्}}\Afootnote{A N H; \skt{वातेन तैर्स्तैर्}\ K}}\skt{ }\edtext{\skt{हनु\-स\-न्धि\-वि\-स\-ङ्गतः}}{\Afootnote{H; \skt{हनु\-स\-न्धि\-र्वि\-दं\-हतः ।}\ A; \skt{हनु\-स\-न्धि\-र्वि\-द\-ङ्गतः ।}\ K; \skt{हनु\-स\-न्धि\-र्विं\-स\-गतः ।}\ N}}\skt{ । }\pend
\pstart\versepar \skt{}\edtext{\skt{हनुर्म्मोक्ष}}{\Afootnote{N H; \skt{हनुमोक्ष}\ A K}}\skt{ इति }\edtext{\skt{ज्ञेयो}}{\Afootnote{A K H; \skt{विज्ञेयो}\ N}}\skt{ }\edtext{\skt{व्याधि\-र\-र्द्दि\-त\-ल\-क्षणः}}{\lemma{\skt{अर्द्दितलक्षणः}}\Afootnote{A K H; \skt{अर्द्धितलक्षणः ॥}\ N}}\skt{ ॥}\pend

\pstart\versepar \ledsidenote{\passref{2.16.36}}\edtext{\skt{जिह्वागतास्तु}}{\Afootnote{K N H; \skt{जिह्वागतास्तु}\ A}}\skt{ । }\pend
\pstart\versepar \skt{}\edtext{\edtext{\edtext{\edtext{\edtext{\skt{कण्ट\-का\-स्त्रि\-वि\-धा\-स्त्रि\-भि\-र्द्दो\-षै\-र\-लास}}{\lemma{\skt{कण्टकास्}}\Afootnote{K N H; \skt{कण्टका}\ A}}}{\lemma{\skt{त्रिविधास्}}\Afootnote{N H; \skt{स्त्रिविधास्}\ A; \skt{त्रिविधांस्}\ K}}}{\lemma{\skt{त्रिभिर्}}\Afootnote{A N H; \skt{त्रिभि}\ K}}}{\lemma{\skt{द्दोषैर्}}\Afootnote{A K; \skt{दोषै}\ N H}}}{\lemma{\skt{अलास}}\Afootnote{A K; \skt{लास}\ N; \skt{रलास}\ H}}\skt{ }\edtext{\skt{उपजिह्वा}}{\Afootnote{N H; \skt{उपजिह्विका}\ A K}}\skt{ चेति ॥}\pend

\pstart\versepar \ledsidenote{\passref{2.16.37}}\edtext{\skt{जिह्वानिलेन}}{\Afootnote{N H; \skt{जिह्वाऽनिलेन}\ A; \skt{जिज्वानिलेन}\ K}}\skt{ }\edtext{\skt{स्फुटिता}}{\Afootnote{A N H; \skt{स्फुटिआ}\ K}}\skt{ प्रसुप्ता भवेच्च }\edtext{\skt{शाक\-च्छ\-द\-न\-प्र\-काशा}}{\Afootnote{A K H; \skt{शाक\-च्छा\-द\-न\-प्र\-काशा ॥}\ N}}\skt{ । }\pend
\pstart\versepar \skt{}\edtext{\edtext{\skt{पित्ता\-त्स\-दा\-है\-र\-नु\-ची\-यते}}{\lemma{\skt{पित्तात्सदाहैर्}}\Afootnote{K N H; \skt{पित्तेन पीता}\ A}}\skt{ च}}{\lemma{\skt{अनुचीयते च}}\Afootnote{K H; \skt{परिदह्यते च}\ A; \skt{अणुमीयते च}\ N}}\skt{ }\edtext{\skt{दीर्घैः}}{\Afootnote{K N H; \skt{चिता}\ A}}\skt{ सर\-क्तै\-र\-पि\-क\-ण्ट\-कैश्च ॥ }\pend
\pstart\versepar \skt{कफेन गुर्व्वी }\edtext{\skt{बहला}}{\Afootnote{A H; \skt{वहला}\ K N}}\skt{ चिता च }\edtext{\skt{मान्साङ्कुरैः}}{\Afootnote{H; \skt{मांसोद्गमैः}\ A; \skt{मांसोद्गरैः}\ K; \skt{मांसोद्ः}\ N}}\skt{ }\edtext{\skt{शाल्म\-लि\-क\-ण्ट\-कौ\-भाभैः}}{\Afootnote{\skt{शाल्म\-लि\-क\-ण्ट\-काभैः}\ A K N; \skt{शाल्म\-लि\-क\-ण्ट\-कौ\-भैभैः ।}\ H}}\skt{ ।}\pend

\pstart\versepar \ledsidenote{\passref{2.16.38}}\edtext{\edtext{\skt{अधोगतो यः}}{\Afootnote{N H; \skt{जिह्वातले यः}\ A}}\skt{ श्वयथुः }\edtext{\skt{प्रगाढः }\edtext{\skt{सोऽला\-स\-सं\-ज्ञः\-क\-फ\-र\-क्त\-मूर्त्तिः}}{\lemma{\skt{ऽला\-स\-सं\-ज्ञः\-क\-फ\-र\-क्त\-मूर्त्तिः}}\Afootnote{A H; \skt{लाससंज्ञः कफरक्तमूर्तिः ॥}\ N}}}{\lemma{\skt{प्रगाढः सो}}\Afootnote{A H; \skt{प्रगाढः सो}\ N}}}{\lemma{\skt{…ऽलाससंज्ञःक}}\Afootnote{\emph{om.}\ K}}\skt{ ॥ }\pend
\pstart\versepar \skt{}\edtext{\skt{जिह्वां स तु}}{\Afootnote{A K H; \skt{जिह्वा स तु}\ N}}\skt{ }\edtext{\skt{स्तम्भयते}}{\Afootnote{K N H; \skt{स्तम्भयति}\ A}}\skt{ प्रवृद्धो }\edtext{\skt{मूले च}}{\Afootnote{K N H; \skt{मूले तु}\ A}}\skt{ जिह्वा }\edtext{\edtext{\skt{भृशमेति}}{\lemma{\skt{भृशम्}}\Afootnote{A N H; \skt{भृश}\ K}}}{\lemma{\skt{एति}}\Afootnote{A N H; \skt{मेति}\ K}}\skt{ पाकं ।}\pend

\pstart\versepar \ledsidenote{\passref{2.16.39}}\skt{जिह्वाग्ररूपः श्वयथुर्हि जिह्वाम् }\pend
\pstart\versepar \skt{}\edtext{\skt{उन्नाम्य}}{\Afootnote{K N H; \skt{उन्नम्य}\ A}}\skt{ जातः }\edtext{\skt{कफरक्तमूलः}}{\Afootnote{H; \skt{कफरक्तयोनिः । प्रसे\-क\-क\-ण्डू\-प\-रि\-दा\-ह\-युक्ता प्रक\-थ्य\-तेऽ\-सा\-वु\-प\-जि\-ह्वि\-केति ॥}\ A; \skt{कफरक्तमूलः ॥ लालाकरः कण्डुयुतः सचोषः सा तूपजिह्वा पठिता भिषग्भिः ॥}\ K; \skt{कफरक्तमूलः ॥ लाअकरः कण्डुयुतः स चोषः । सातुपजिह्वा पठिता भिषग्भिः ॥}\ N}}\skt{ ॥}\pend

\pstart \ledsidenote{\passref{2.16.39a}}\skt{लालाकरः कण्डुयुतः स चोषः सातुपजिह्वा पठिता भिषग्भिः । ।}\pend

\pstart \ledsidenote{\passref{2.16.40}}\skt{तालुगतास्तु ॥ गलशुण्डिका }\edtext{\skt{तुण्डिकेरी}}{\Afootnote{A N H; \skt{तुण्डीकेरी}\ K}}\skt{ अध्रुषः कच्छपः । }\edtext{\skt{अर्ब्बुदः}}{\Afootnote{H; \skt{अर्बुदं,}\ A; \skt{अर्वुदः}\ K; \skt{अर्वुद्धः}\ N}}\skt{ । }\edtext{\skt{मान्ससंघातः}}{\Afootnote{A K H; \skt{माससंघातः}\ N}}\skt{ । }\edtext{\skt{तालु पुप्पुटकः}}{\Afootnote{H; \skt{तालुपुप्पुटः,}\ A; \skt{तालुपुष्पुटाकः}\ K; \skt{तालुपुर्प्पुटकः}\ N}}\skt{ । तालुशोषः । तालुपाक इति ॥}\pend

\pstart \ledsidenote{\passref{2.16.41}}\edtext{\skt{कफा\-सृ\-ग्भ्या\-न्ता\-लु\-मू\-ला\-त्प्र\-वृद्धो}}{\Afootnote{H; \skt{श्लेष्मासृग्भ्यां तालु\-मू\-ला\-त्प्र\-वृद्धो}\ A; \skt{कफां सृभ्यां तालु\-मू\-ला\-त्प्र\-वृद्धा}\ K\textsuperscript{ac}; \skt{कफां सृगभ्यां तालु\-मू\-ला\-त्प्र\-वृद्धा}\ K\textsuperscript{pc}; \skt{कफो\-सृ\-ग्भ्या\-न्ता\-लु\-मू\-ला\-त्प्र\-वृद्धो}\ N}}\skt{ दीर्घः }\edtext{\skt{शोफो\-ध्मा\-त\-ब\-स्ति\-प्र\-काशः}}{\Afootnote{H; \skt{शोफो ध्मात\-ब\-स्ति\-प्र\-काशः ।}\ A; \skt{शोफो ध्मात\-व\-स्ति\-प्र\-काशः ।}\ K; \skt{शोफो ध्मांत\-व\-स्ति\-प्र\-काशः ।}\ N}}\skt{ । }\edtext{\skt{तृष्णाकासः}}{\Afootnote{K N H; \skt{तृष्णाकास}\ A}}\skt{ }\edtext{\edtext{\skt{श्वास\-कृ\-त्त\-म्व\-दन्ति}}{\lemma{\skt{श्वासकृत्तम्}}\Afootnote{K N H; \skt{श्वासकृत्सं}\ A}}}{\lemma{\skt{वदन्ति}}\Afootnote{K N H; \skt{प्रदिष्टो}\ A}}\skt{ }\edtext{\edtext{\skt{व्याधिम्वैद्या}}{\lemma{\skt{व्याधिम्}}\Afootnote{N H; \skt{व्याधिर्}\ A; \skt{व्याधि}\ K}}}{\lemma{\skt{वैद्या}}\Afootnote{K H; \skt{वैद्यैः}\ A; \skt{वैद्यः}\ N}}\skt{ }\edtext{\skt{गलशुण्डेति}}{\Afootnote{N H; \skt{कण्ठशुण्डीति}\ A; \skt{गलशुण्डति}\ K}}\skt{ }\edtext{\skt{नाम्ना}}{\Afootnote{A K H; \skt{नाम्नाः ॥}\ N}}\skt{ ॥}\pend

\pstart \ledsidenote{\passref{2.16.42}}\edtext{\skt{शोफः}}{\Afootnote{A K H; \skt{शोफाः}\ N}}\skt{ }\edtext{\skt{स्थूलस्तोददाह प्रपाकी}}{\Afootnote{K N H; \skt{स्त}\ A}}\skt{ प्रागुक्ताभ्यां तुण्डिकेरीमता तु । मृदुः }\edtext{\skt{शोफो}}{\Afootnote{N H; \skt{शोथो}\ K}}\skt{ }\edtext{\skt{लोहितोत्थो}}{\Afootnote{N H; \skt{लोहिं\-त\-स्ता\-लु\-देशे रक्ता}\ A; \skt{लोहिता}\ K}}\skt{ }\edtext{\edtext{\skt{ज्ञेयोऽध्रुषः}}{\lemma{\skt{ज्ञेयो}}\Afootnote{K N H; \skt{ज्ञेयः सो}\ A}}}{\lemma{\skt{ऽध्रुषः}}\Afootnote{A H; \skt{ध्रुषः}\ K N}}\skt{ }\edtext{\edtext{\skt{सज्वरस्तीव्र}}{\lemma{\skt{सज्वरस्}}\Afootnote{K N H; \skt{रुग्ज्वरा}\ A}}}{\lemma{\skt{तीव्र}}\Afootnote{K N H; \skt{ढ्यः ॥}\ A}}\skt{ }\edtext{\skt{रुक्च}}{\Afootnote{K H; \emph{om.}\ A; \skt{रुक् ॥}\ N}}\skt{ ॥}\edtext{}{\lemma{\skt{पाकी… शोफो}}\Afootnote{\emph{om.}\ A}}\pend

\pstart \ledsidenote{\passref{2.16.43}}\edtext{\skt{कूर्म्मोन्नतो}}{\Afootnote{N H; \skt{कूर्मोत्सन्नो}\ A; \skt{कूर्मोन्न}\ K}}\skt{ }\edtext{\skt{वेद\-नाऽ\-शी\-घ्र\-जन्मा}}{\Afootnote{H; \skt{ऽवे\-द\-नोऽ\-शी\-घ्र\-जन्मा}\ A; \skt{तोवे\-द\-नो\-श्री\-घ्र\-जन्मा}\ K; \skt{वेद\-ना\-शी\-घ्र\-जर्मा}\ N}}\skt{ }\edtext{\skt{रोगो}}{\Afootnote{K N H; \skt{ऽरक्तो}\ A}}\skt{ }\edtext{\skt{ज्ञेयः}}{\Afootnote{A K H; \skt{ज्ञेउ}\ N}}\skt{ }\edtext{\skt{कच्छपः}}{\Afootnote{A K H; \skt{अः च्छपः}\ N}}\skt{ }\edtext{\skt{श्लेष्मणा तु}}{\Afootnote{K N H; \skt{श्लेष्मणा स्या}\ A}}\skt{ ॥ }\edtext{\skt{पद्मा\-क\-र\-न्ता\-लु\-मध्ये}}{\lemma{\skt{पद्माकरन्}}\Afootnote{H; \skt{त् । पद्माकारं}\ A; \skt{पद्माकारं}\ K N}}\skt{ तु }\edtext{\skt{शोफं }\edtext{\skt{तम्वि\-द्या\-द्र\-क्ता\-र्ब्बुदं}}{\lemma{\skt{विद्याद्रक्तार्ब्बुदं}}\Afootnote{H; \skt{विद्या\-द्र\-क्ता\-द\-र्बुदं}\ A; \skt{विद्याद्रक्तार्द्वुदं}\ K; \skt{विद्याद्रक्तार्वुदं}\ N}}}{\lemma{\skt{शोफं तम्}}\Afootnote{K N H; \skt{शोफं}\ A}}\skt{ }\edtext{\skt{प्रोक्तलिंगं}}{\Afootnote{A N H; \skt{प्रोक्तलिंङ्गं ॥}\ K}}\skt{ ।}\pend

\pstart \ledsidenote{\passref{2.16.44}}\edtext{\edtext{\skt{मांस\-न्नि\-रु\-ज\-न्ता\-लु\-मध्ये}}{\lemma{\skt{मांसन्निरुजन्}}\Afootnote{K N H; \skt{दुष्टं मांसं श्लेष्मणा नीरुजं}\ A}}}{\lemma{\skt{तालुमध्ये}}\Afootnote{N H; \emph{om.}\ A; \skt{तालु}\ K}}\skt{ }\edtext{\skt{कफात्तु प्रोक्तो}}{\Afootnote{H; \skt{च ताल्वन्तः}\ A; \skt{कफान्प्रोर्वै}\ K; \skt{कफा प्रोक्तो}\ N}}\skt{ }\edtext{\edtext{\skt{वैद्यै\-र्म्मां\-स\-सं\-घात}}{\lemma{\skt{वैद्यैर्}}\Afootnote{H; \skt{स्थं}\ A; \skt{द्यै}\ K; \skt{वैदैर्}\ N}}}{\lemma{\skt{म्मांससंघात}}\Afootnote{K N H; \skt{मांससङ्घातम्}\ A}}\skt{ }\edtext{\skt{एषः}}{\Afootnote{K N H; \skt{आहुः ।}\ A}}\skt{ ॥ }\edtext{\skt{अरुक्छिरः}}{\Afootnote{H; \skt{नीरुक्स्थायी}\ A; \skt{अरुक्स्थिरः}\ K; \skt{अरुस्थिरः}\ N}}\skt{ कोलमात्रः }\edtext{\skt{कफेन}}{\Afootnote{K N H; \skt{कफात्स्या}\ A}}\skt{ }\edtext{\edtext{\skt{समे\-द\-सा\-पु\-प्पु\-ट\-स्ता\-लु\-देशे}}{\lemma{\skt{समे\-द\-सा\-पु\-प्पु\-टस्}}\Afootnote{H; \skt{न्मेदो\-यु\-क्ता\-त्पु\-प्पुट}\ A; \skt{समेदसा पुर्पुटास्}\ K; \skt{समेदसा पुर्प्पुटस्}\ N}}}{\lemma{\skt{तालुदेशे}}\Afootnote{K N H; \skt{स्तालुदेशे ॥}\ A}}\skt{ ।}\pend

\pstart \ledsidenote{\passref{2.16.45}}\edtext{\edtext{\skt{शुष्य\-त्य\-त्य\-र्थ\-न्दी\-र्यते}}{\lemma{\skt{शुष्यत्य्}}\Afootnote{K N H; \skt{शो}\ A}}}{\lemma{\skt{अत्यर्थन्}}\Afootnote{K H; \skt{सोऽत्यर्थं}\ A; \skt{अत्यर्थन्}\ N}}\skt{ चापि }\edtext{\edtext{\edtext{\skt{तालु\-श्वा\-स\-श्चै\-वो\-ग्र\-स्ता\-लु\-शो\-षोऽ\-नि\-लात्तु}}{\lemma{\skt{तालुश्वासश्}}\Afootnote{H; \skt{तालुः श्वासो}\ A K N}}}{\lemma{\skt{चैवो\-ग्र\-स्ता\-लु\-शोषो}}\Afootnote{H; \skt{वातात्तालुशोषः}\ A; \skt{चैवो\-ग्र\-स्ता\-लु\-शोषो}\ K N}}}{\lemma{\skt{ऽनिलात्तु}}\Afootnote{H; \skt{सपित्ता}\ A; \skt{निलात्तु ॥}\ K N}}\skt{ ॥ }\edtext{\skt{पित्तं}}{\Afootnote{K N H; \skt{त् । पित्तं}\ A}}\skt{ कुर्या\-त्पा\-क\-म\-त्य\-र्थ\-घोरं }\edtext{\edtext{\skt{तालु\-न्ये\-व\-न्ता\-लु\-पा\-क\-म्व\-दन्ति}}{\lemma{\skt{तालुन्य्}}\Afootnote{A N H; \skt{तालन्य्}\ K}}}{\lemma{\skt{एवन्}}\Afootnote{K N H; \skt{एनं}\ A}}\skt{ ॥}\pend

\pstart \ledsidenote{\passref{2.16.46}}\skt{कण्ठगतास्तु ॥ रोहिण्यः पञ्च । }\edtext{\edtext{\skt{कण्ठ\-शा\-लू\-क\-म\-धि\-जिह्वां}}{\lemma{\skt{कण्ठशालूकम्}}\Afootnote{A K H; \skt{कण्ठशालूक ।}\ N}}}{\lemma{\skt{अधिजिह्वां}}\Afootnote{H; \skt{अधिजिह्वो,}\ A; \skt{अधिजिह्वा}\ K N}}\skt{ }\edtext{\skt{वलयः}}{\Afootnote{H; \skt{बलयो,}\ A; \skt{वलयोर्}\ K; \skt{बलयोर्}\ N}}\skt{ । }\edtext{\skt{अलासः}}{\Afootnote{N H; \skt{बलास,}\ A; \skt{अलाम ।}\ K}}\skt{ । }\edtext{\skt{एकवृन्दः}}{\Afootnote{A N H; \skt{एक}\ K}}\skt{ । }\edtext{\skt{शतघ्नी}}{\Afootnote{N H; \skt{वृन्दः, शतघ्नी,}\ A; \skt{वृक्षः । शतघी}\ K}}\skt{ । गिलायुः । }\edtext{\skt{गलविद्रधिः}}{\Afootnote{A K H; \skt{गलविद्रपिः}\ N}}\skt{ । गलौघः । स्वरघ्नः । }\edtext{\skt{मांस\-ता\-नो\-वि\-चा\-रि\-श्चेति}}{\lemma{\skt{मांस\-ता\-नो\-वि\-चा\-रिश्}}\Afootnote{H; \skt{मांसतानो, विदारी}\ A; \skt{मांसतानो विदारिञ्}\ K; \skt{मांसतानो विदारिश्}\ N}}\skt{ ॥}\pend

\pstart \ledsidenote{\passref{2.16.47}}\edtext{\skt{गलेऽनिलः}}{\lemma{\skt{ऽनिलः}}\Afootnote{A H; \skt{निलः}\ K N}}\skt{ }\edtext{\skt{पित्तकफौ तु}}{\Afootnote{K N H; \skt{पित्तकफौ च}\ A}}\skt{ }\edtext{\skt{मूर्च्छितौ}}{\Afootnote{A K H; \skt{मूतौ}\ N}}\skt{ }\edtext{\skt{प्रदूष्य}}{\Afootnote{K N H; \skt{पृथक्स}\ A}}\skt{ }\edtext{\skt{मांसञ्च}}{\Afootnote{K N H; \skt{मस्ताश्च}\ A}}\skt{ }\edtext{\skt{तथैव}}{\Afootnote{A K H; \skt{तोथैव}\ N}}\skt{ शोणितं । }\edtext{\edtext{\edtext{\edtext{\edtext{\skt{गलो\-प\-सं\-रो\-ध\-क\-रै\-स्त\-था\-ङ्कु\-रै\-र्नि\-ह\-न्त्य\-सु\-न्व्या\-धि\-रयं}}{\lemma{\skt{गलो\-प\-सं\-रो\-ध\-क\-रैस्}}\Afootnote{K N H; \skt{प्रदूष्य मांसं गलरोधि}\ A}}}{\lemma{\skt{तथाङ्कुरैर्}}\Afootnote{N H; \skt{नोऽङ्कुरा}\ A; \skt{तथाङ्कुरौ}\ K}}}{\lemma{\skt{निहन्त्यस्}}\Afootnote{N H; \skt{न्सृजन्}\ A; \skt{हं न्य}\ K}}}{\lemma{\skt{उन्व्याधिर्}}\Afootnote{N H; \skt{ति यान्सा}\ A; \skt{स्तं व्याधिर्}\ K}}\skt{ हि}}{\lemma{\skt{अयं हि}}\Afootnote{K N H; \skt{ऽसुहरा हि}\ A}}\skt{ रोहिणी ॥}\pend

\pstart\versepar \ledsidenote{\passref{2.16.48}}\edtext{\skt{जिह्वा}}{\Afootnote{K N H; \skt{जिह्वां}\ A}}\skt{ }\edtext{\skt{सम\-न्ता\-द्भृ\-श\-वे\-द\-नान्तु}}{\Afootnote{N H; \skt{सम\-न्ता\-द्भृ\-श\-वे\-दना ये}\ A; \skt{सम\-न्ता\-द्धृ\-त\-श\-वे\-द\-नात्तु}\ K}}\skt{ }\edtext{\skt{मान्साङ्कुरैः}}{\Afootnote{K H; \skt{मांसाङ्कुराः}\ A; \skt{मासांङ्कुरैः}\ N}}\skt{ कण्ठ\-नि\-रो\-धि\-नीया । }\pend
\pstart\versepar \skt{तां }\edtext{\edtext{\skt{रोहि\-णी\-म्वा\-त\-कृ\-ता\-म्व\-दन्ति}}{\lemma{\skt{रोहिणीम्}}\Afootnote{A H; \skt{रोहिणी}\ K N}}}{\lemma{\skt{वदन्ति}}\Afootnote{K N H; \skt{वदन्}\ A}}\skt{ }\edtext{\skt{तिस्रश्च}}{\Afootnote{K N H; \skt{ति}\ A}}\skt{ }\edtext{\edtext{\skt{तास्त्वादित}}{\Afootnote{H; \skt{तास्त्वादित}\ K N}}\skt{ एव कृच्छ्रा ॥}}{\lemma{\skt{स्रश्… कृच्छ्रा}}\Afootnote{\emph{om.}\ A}}\edtext{}{\lemma{\skt{कण्ठ\-नि\-रो\-धि\-नीया । तां}}\Afootnote{H; \skt{कण्ठनिरोधिनः स्युः । तां}\ A; \skt{कण्ठं निरोधिनी या । तां}\ K; \skt{कण्ठ\-नि\-रो\-ध\-नीया । तां}\ N}}\pend

\pstart\versepar \ledsidenote{\passref{2.16.49}}\skt{क्षिप्रोद्गमा }\edtext{\edtext{\skt{क्षिप्र\-वि\-दा\-ह\-पा\-का\-स्ती\-व्र\-ज्वरा}}{\lemma{\skt{क्षिप्र\-वि\-दा\-ह\-पा\-कास्}}\Afootnote{N H; \skt{क्षिप्र\-वि\-दा\-ह\-पाका}\ A K}}}{\lemma{\skt{तीव्रज्वरा}}\Afootnote{A N H; \skt{स्तीव्रज्वरा}\ K}}\skt{ }\edtext{\skt{पित्तनिमित्तजा तु}}{\Afootnote{K N H; \skt{पित्तनिमित्तजा स्या}\ A}}\skt{ । }\pend
\pstart\versepar \skt{}\edtext{\skt{स्रोतो}}{\Afootnote{N H; \skt{त् । स्रोतो}\ A; \skt{स्रोतो}\ K}}\skt{ }\edtext{\skt{निरो\-धि\-न्य\-च\-लो\-न्नता च}}{\Afootnote{K N H; \skt{निरोधिन्यपि मन्दपाका}\ A}}\skt{ }\edtext{\skt{स्थिराङ्कुरा या}}{\Afootnote{K H; \skt{गुर्वी स्थिरा सा}\ A; \skt{स्थितांकुरा या}\ N}}\skt{ कफसंभवा }\edtext{\skt{सा}}{\Afootnote{K N H; \skt{वै ॥}\ A}}\skt{ ॥}\pend

\pstart\versepar \ledsidenote{\passref{2.16.50}}\edtext{\edtext{\skt{गम्भी\-र\-पा\-कि\-न्य\-नि\-वा\-र\-वीर्या}}{\lemma{\skt{गम्भीरपाकिन्य्}}\Afootnote{H; \skt{गम्भीरपाकाऽ}\ A; \skt{गम्भीरपाकि}\ K; \skt{गम्भीरा पाकिन्य्}\ N}}}{\lemma{\skt{अनिवारवीर्या}}\Afootnote{N H; \skt{प्रतिवारवीर्या}\ A; \skt{न्यनिवारवी}\ K}}\skt{ त्रिदोषलिङ्गा }\edtext{\skt{त्रितयोत्थितं तु}}{\Afootnote{N H; \skt{त्रयसंभवा स्या}\ A; \skt{त्रितयोक्ति तत्तु ॥}\ K}}\skt{ । }\pend
\pstart\versepar \skt{}\edtext{\skt{स्फोटैश्चिता}}{\Afootnote{H; \skt{त् । स्फोटाचिता}\ A; \skt{स्फोटौम्चिता}\ K; \skt{स्फोटचिता}\ N}}\skt{ }\edtext{\skt{पित्त\-स\-मा\-न\-लिं\-गा\-साध्या}}{\Afootnote{K N H; \skt{पित्त\-स\-मा\-न\-लि\-ङ्गाऽ\-साध्या}\ A}}\skt{ प्रदिष्टा }\edtext{\skt{रुधिरात्मिका}}{\Afootnote{K N H; \skt{रुधिरात्मिके}\ A}}\skt{ }\edtext{\skt{तु}}{\Afootnote{K N H; \skt{यम् ॥}\ A}}\skt{ ॥}\pend

\pstart\versepar \ledsidenote{\passref{2.16.51}}\skt{कोलास्थिमात्रः }\edtext{\skt{कफसंभवो यो}}{\Afootnote{K N H; \skt{कफस्ंभवो यो}\ A}}\skt{ }\edtext{\skt{ग्रन्थि\-र्ग\-ले\-क\-ण्ट\-क\-शू\-क\-भूतः}}{\lemma{\skt{गले\-क\-ण्ट\-क\-शू\-क\-भूतः}}\Afootnote{A H; \skt{गले कष्ट\-क\-शु\-क\-भूतः ।}\ K; \skt{गले कण्ठ\-क\-शू\-क\-भूतः ।}\ N}}\skt{ । }\pend
\pstart\versepar \skt{}\edtext{\skt{खरः}}{\Afootnote{A H; \skt{खर}\ K N}}\skt{ स्थिरः }\edtext{\skt{शस्त्र\-नि\-पा\-त\-सा\-ध्य\-स्त\-त्क\-ण्ट\-शा\-लू\-क\-मिति}}{\lemma{\skt{तत्क\-ण्ट\-शा\-लू\-कम्}}\Afootnote{H; \skt{तं कण्ठशालूकम्}\ A K N}}\skt{ }\edtext{\skt{ब्रुवन्ति}}{\Afootnote{A H; \skt{व्रुवन्ति ॥}\ K N}}\skt{ ॥}\pend

\pstart\versepar \ledsidenote{\passref{2.16.52}}\skt{जिह्वाग्ररूपः }\edtext{\skt{श्वयथुः}}{\Afootnote{A K H; \skt{स्वयथुः}\ N}}\skt{ }\edtext{\skt{कफात्तु}}{\Afootnote{A K H; \skt{कफां तु}\ N}}\skt{ }\edtext{\skt{जिह्वो\-प\-रि\-ष्टा\-द\-पि\-र\-क्त\-मिश्रः}}{\Afootnote{H; \skt{जिह्वा\-प्र\-ब\-न्धो\-परि रक्तमिश्रा}\ A; \skt{जिह्वो\-प\-रि\-ष्टा\-दपि रक्तमिश्रः ।}\ K; \skt{जिह्वो\-प\-रि\-ष्ठा\-दपि रक्तमिश्रः ।}\ N}}\skt{ । }\pend
\pstart\versepar \skt{}\edtext{\edtext{\skt{ज्ञेयोऽधिजिह्वा}}{\lemma{\skt{ज्ञेयो}}\Afootnote{K H; \skt{त् । ज्ञेयो}\ A; \skt{ज्ञेयो}\ N}}}{\lemma{\skt{ऽधिजिह्वा}}\Afootnote{H; \skt{ऽधिजिह्वः}\ A; \skt{धिजिह्व}\ K; \skt{धिजिह्वा}\ N}}\skt{ खलु रोग }\edtext{\skt{एष}}{\Afootnote{A N H; \emph{om.}\ K}}\skt{ }\edtext{\skt{विव\-र्ज्ज\-ये\-दा\-ग\-त\-पा\-क\-मे\-नम्}}{\lemma{\skt{एनम्}}\Afootnote{A K H; \skt{एन ॥}\ N}}\skt{ ॥}\pend

\pstart \ledsidenote{\passref{2.16.53}}\edtext{\skt{बलास}}{\Afootnote{A H; \skt{वलास}\ K N}}\skt{ }\edtext{\edtext{\skt{एवा\-य\-त\-मु\-न्न\-तञ्च}}{\lemma{\skt{एवायतम्}}\Afootnote{K N H; \skt{एवायत}\ A}}}{\lemma{\skt{उन्नतञ्च}}\Afootnote{K N H; \skt{मुन्नतं च}\ A}}\skt{ शोफं }\edtext{\edtext{\skt{करो\-त्य\-न्न\-ग\-ति\-न्नि\-वा\-र्य।तं}}{\lemma{\skt{करोत्य्}}\Afootnote{K N H; \skt{करो}\ A}}}{\lemma{\skt{अन्नगतिन्}}\Afootnote{K N H; \skt{त्यन्नगतिं}\ A}}\skt{ }\edtext{\skt{सर्व्व\-थै\-वा\-प्र\-ति\-वा\-र\-वी\-र्य\-म्वि\-व\-र्ज्ज\-नी\-य\-म्व\-ल\-य\-म्व\-दन्ति}}{\lemma{\skt{विवर्ज्जनीयम्}}\Afootnote{A N H; \skt{विवर्जनीय}\ K}}\skt{ ॥}\pend

\pstart\versepar \ledsidenote{\passref{2.16.54}}\skt{गले तु शोफङ्कुरुतः प्रवृद्धौ श्लेष्मानिलौ श्वास\-रु\-जो\-प\-पन्नं । }\pend
\pstart\versepar \skt{}\edtext{\edtext{\edtext{\edtext{\skt{मर्म्म\-च्छि\-द\-न्दु\-स्त\-र\-मे\-त\-दा\-हु\-र\-ला\-स\-सं\-ज्ञ\-न्नि\-पु\-णा\-धि\-का\-रम्}}{\lemma{\skt{दुस्तरम्}}\Afootnote{A N H; \skt{दुस्तर}\ K}}}{\lemma{\skt{एतदाहुर्}}\Afootnote{A N H; \skt{मेतदाहुर्}\ K}}}{\lemma{\skt{अलाससंज्ञन्}}\Afootnote{K H; \skt{बलासस्ंज्ञं}\ A; \skt{अलाससंज्ञ}\ N}}}{\lemma{\skt{निपु\-णा\-धि\-का\-रम्}}\Afootnote{H; \skt{निपुणा विकारम्}\ A K; \skt{निपुणाधिकार ॥}\ N}}\skt{ ॥}\pend

\pstart\versepar \ledsidenote{\passref{2.16.55}}\edtext{\skt{वृत्तोन्नतान्तः}}{\Afootnote{N H; \skt{वृत्तोन्नतो यः}\ A; \skt{वृत्तोन्नतान्नः}\ K}}\skt{ श्वयथुः सदाहः }\edtext{\edtext{\skt{सक\-क\-ण्डु\-रोऽ\-पा\-व्य\-गु\-रु\-र्म्मृ\-दुश्च}}{\lemma{\skt{सककण्डुरो}}\Afootnote{H; \skt{कण्ड्वन्वितो}\ A; \skt{स कण्डुरो}\ K; \skt{सकण्डुरो}\ N}}}{\lemma{\skt{ऽपा\-व्य\-गु\-रु\-र्म्मृ\-दुश्च}}\Afootnote{H; \skt{ऽपा\-क्य\-मृ\-दु\-र्गु\-रुश्च ।}\ A; \skt{पाक्गुरूमृदुश्च ॥}\ K; \skt{वाक्य\-गु\-रु\-र्मृ\-दुश्च ।}\ N}}\skt{ । }\pend
\pstart\versepar \skt{}\edtext{\skt{नाम्नैकवृन्दः}}{\Afootnote{A K H; \skt{नाम्वैकवृन्दः}\ N}}\skt{ }\edtext{\edtext{\skt{परिकीर्त्यतेऽसौ}}{\lemma{\skt{परिकीर्त्यते}}\Afootnote{N H; \skt{परिकीर्तितो}\ A; \skt{परिकीर्तिते}\ K}}}{\lemma{\skt{ऽसौ}}\Afootnote{A H; \skt{सौ}\ K N}}\skt{ }\edtext{\edtext{\skt{व्याधि\-र्ब्ब\-ला\-स\-क्ष\-तजः}}{\lemma{\skt{व्याधिर्}}\Afootnote{A N H; \skt{व्याधि}\ K}}}{\lemma{\skt{ब्बलासक्षतजः}}\Afootnote{H; \skt{बलासक्षतज}\ A; \skt{र्वलासः क्षतजः}\ K; \skt{वलासः क्षतजः}\ N}}\skt{ प्रसूतः ॥}\pend

\pstart\versepar \ledsidenote{\passref{2.16.56}}\edtext{\edtext{\skt{समु\-न्न\-त\-म्वृ\-त्त\-म\-म\-न्द\-दा\-ह\-न्ती\-व्र\-ज्वरं}}{\lemma{\skt{समु\-न्न\-त\-म्वृ\-त्तम्}}\Afootnote{A K H; \skt{समुन्नतं वर्त्तम्}\ N}}}{\lemma{\skt{अम\-न्द\-दा\-ह\-न्ती\-व्र\-ज्वरं}}\Afootnote{A K H; \skt{अमन्ददाहं तीव्रं ज्वरां}\ N}}\skt{ }\edtext{\edtext{\skt{वृन्द\-मु\-दा\-ह\-रन्ति}}{\lemma{\skt{वृन्दम्}}\Afootnote{K N H; \skt{वृन्द}\ A}}}{\lemma{\skt{उदाहरन्ति}}\Afootnote{K N H; \skt{मुदाहरन्ति ।}\ A}}\skt{ । }\pend
\pstart\versepar \skt{}\edtext{\skt{तच्चापि}}{\Afootnote{K N H; \skt{तं चापि}\ A}}\skt{ }\edtext{\skt{पित्त\-क्ष\-त\-ज\-प्र\-को\-पा\-ज्ज्ञेयं}}{\Afootnote{H; \skt{पित्त\-क्ष\-त\-ज\-प्र\-को\-पा\-द्विद्या}\ A; \skt{पित्त\-क्ष\-त\-ज\-प्र\-को\-पा\-द्ज्ञेयं}\ K; \skt{पित्तक्षतरुः प्रकोपाज्ज्ञेयं}\ N}}\skt{ }\edtext{\edtext{\skt{सतो\-द\-म्प\-व\-ना\-त्म\-कन्तु}}{\lemma{\skt{सतोदम्}}\Afootnote{K N H; \skt{त्सतोदं}\ A}}}{\lemma{\skt{पवनात्मकन्}}\Afootnote{K N H; \skt{पवनास्रजं}\ A}}\skt{ ॥}\pend

\pstart \ledsidenote{\passref{2.16.57}}\edtext{\skt{वर्त्तिर्घना}}{\lemma{\skt{वर्त्तिर्}}\Afootnote{A K H; \skt{वर्त्ति}\ N}}\skt{ }\edtext{\skt{कण्ठनिरोधनी तु}}{\Afootnote{H; \skt{कण्ठनिरोधिनी या}\ A; \skt{कण्ठनिरोधिनी तु}\ K; \skt{कुण्ठनिरोधिनी तु}\ N}}\skt{ }\edtext{\skt{चिता\-ति\-मा\-त्र\-म्पि\-शि\-त\-प्र\-रोहैः}}{\lemma{\skt{चितातिमात्रम्}}\Afootnote{H; \skt{चिताऽतिमात्रं}\ A; \skt{चितानिमात्रं}\ K; \skt{चिताचिमात्रं}\ N}}\skt{ । }\edtext{\skt{अने\-क\-रु\-क्प्रा\-ण\-हरी}}{\Afootnote{K N H; \skt{नाना\-रु\-जो\-च्छ्रा\-य\-करी}\ A}}\skt{ }\edtext{\skt{त्रिदोषाज्ज्ञेया}}{\Afootnote{H; \skt{त्रिदोषाज्ज्ञेया}\ A; \skt{तृदोषाज्ज्ञेया}\ K; \skt{तृदोषा ज्ञेया}\ N}}\skt{ }\edtext{\skt{शतघ्नीति}}{\Afootnote{K N H; \skt{शतघ्नीव}\ A}}\skt{ }\edtext{\skt{शतघ्निरूपा}}{\Afootnote{N H; \skt{शतघ्न्यसाध्या ॥}\ A; \skt{सतघ्निरूपा ।}\ K}}\skt{ ॥}\pend

\pstart\versepar \ledsidenote{\passref{2.16.58}}\edtext{\skt{ग्रन्थिर्ग्गले}}{\Afootnote{A H; \skt{ग्रन्थिर्गले}\ K N}}\skt{ त्वाम\-ल\-का\-स्थि\-मात्रः }\edtext{\edtext{\skt{स्थिरोऽल्परुक्यः}}{\lemma{\skt{स्थिरो}}\Afootnote{A H; \skt{स्थिरो}\ K N}}}{\lemma{\skt{ऽल्परुक्यः}}\Afootnote{H; \skt{ऽल्परुक्स्या}\ A; \skt{ल्परूग्यः}\ K; \skt{ल्परुपः}\ N}}\skt{ कफपित्तमूर्त्तिः । }\pend
\pstart\versepar \skt{यो लक्ष्यते }\edtext{\edtext{\skt{रक्त\-मि\-वा\-स्र\-वश्च}}{\lemma{\skt{रक्तम्}}\Afootnote{K N H; \skt{सक्तम्}\ A}}}{\lemma{\skt{इवास्रवश्च}}\Afootnote{H; \skt{इवाशनं च}\ A; \skt{इवाश्रवश्च}\ K; \skt{इवास्रवञ्च}\ N}}\skt{ }\edtext{\skt{सशस्रसाध्धस्तु}}{\Afootnote{H; \skt{स शस्त्रसाध्यस्तु}\ A K N}}\skt{ गिलायुसंज्ञः ॥}\edtext{}{\lemma{\skt{कफपित्तमूर्त्तिः । यो}}\Afootnote{K H; \skt{त्कफरक्तमूर्तिः । सं}\ A; \skt{कफपित्तमूर्त्ति । यो}\ N}}\pend

\pstart\versepar \ledsidenote{\passref{2.16.59}}\skt{सर्व्वङ्गलं }\edtext{\skt{व्याप्य}}{\Afootnote{A K H; \skt{प्राप्य}\ N}}\skt{ समुत्थितो यः शोफो }\edtext{\skt{रुजाः}}{\Afootnote{K N H; \skt{रुजो}\ A}}\skt{ }\edtext{\skt{सन्ति च}}{\Afootnote{K N H; \skt{यत्र च}\ A}}\skt{ }\edtext{\skt{यत्र}}{\Afootnote{N H; \skt{सन्ति}\ A; \skt{यत्र तु}\ K}}\skt{ सर्व्वा । }\pend
\pstart\versepar \skt{स }\edtext{\skt{सर्व्व\-दो\-षै\-र्ग्ग\-ल\-वि\-द्र\-धिस्तु}}{\lemma{\skt{सर्व्वदोषैर्}}\Afootnote{K H; \skt{सर्वदोषो}\ A N}}\skt{ तस्यैव तुल्यः खलु सर्व्वजस्य ॥}\edtext{}{\lemma{\skt{सर्व्वा । स}}\Afootnote{K N H; \skt{सर्वाः । स}\ A}}\pend

\pstart\versepar \ledsidenote{\passref{2.16.60}}\skt{शोफो }\edtext{\skt{महान्तन्न जलावरोधी}}{\Afootnote{H; \skt{महा\-न\-न्न\-ज\-ला\-व\-रोधी}\ A K; \skt{महानन्नज जलावरोधी}\ N}}\skt{ }\edtext{\skt{तीव्रज्वरो}}{\Afootnote{A N H; \skt{तीव्रं ज्वरो}\ K}}\skt{ }\edtext{\edtext{\skt{वायु\-ग\-ते\-र्न्नि\-हन्ता}}{\lemma{\skt{वायुगतेर्}}\Afootnote{K N H; \skt{वातगते}\ A}}}{\lemma{\skt{न्निहन्ता}}\Afootnote{H; \skt{र्निहन्ता ।}\ A; \skt{निहन्त ॥}\ K N}}\skt{ । }\pend
\pstart\versepar \skt{कफेन जातो }\edtext{\skt{रुधिरान्वितेन}}{\Afootnote{A K H; \skt{रुधिरन्वितेन}\ N}}\skt{ गले गलौघः }\edtext{\skt{पठितो}}{\Afootnote{N H; \skt{परिकीर्त्यतेऽसौ ॥}\ A; \skt{परितो निषग्निः ॥}\ K}}\skt{ }\edtext{\skt{भिषग्भिः}}{\Afootnote{N H; \emph{om.}\ A K}}\skt{ ॥}\pend

\pstart\versepar \ledsidenote{\passref{2.16.61}}\edtext{\skt{यस्ताम्यमानः}}{\Afootnote{H; \skt{योऽतिप्रताम्य}\ A; \skt{यस्ताम्य मानः}\ K; \skt{यस्म्यमानः}\ N}}\skt{ }\edtext{\skt{श्वसिति}}{\Afootnote{K N H; \skt{न्श्वसिति}\ A}}\skt{ प्रस\-क्त\-म्भि\-न्न\-स्वरः }\edtext{\skt{शुष्क\-वि\-मु\-क्त\-कण्ठः}}{\Afootnote{A K H; \skt{शुष्ककण्ठः ।}\ N}}\skt{ । }\pend
\pstart\versepar \skt{}\edtext{\edtext{\skt{कफो\-प\-दि\-ष्टे\-ष्व\-नि\-ला\-य\-तेषु}}{\lemma{\skt{कफोपदिष्टेष्व्}}\Afootnote{H; \skt{कफोपदिग्धेष्व्}\ A N; \skt{कफोपदिग्धे}\ K}}}{\lemma{\skt{अनिलायतेषु}}\Afootnote{H; \skt{अनिलायनेषु}\ A N; \skt{ष्वनिलायनेषु}\ K}}\skt{ ज्ञेयः स रोगः }\edtext{\skt{श्वसनात्स्वरघ्नः}}{\Afootnote{H; \skt{श्वसनात्स्वरघ्नः}\ A N; \skt{श्वसनान्स्वरघ्नः ॥}\ K}}\skt{ ॥}\pend

\pstart\versepar \ledsidenote{\passref{2.16.62}}\skt{प्रतानवान्यः श्वयथुः }\edtext{\skt{सुकष्टो}}{\Afootnote{A N H; \skt{सुकष्ठो}\ K}}\skt{ । }\pend
\pstart\versepar \skt{}\edtext{\skt{गलो\-प\-रो\-धा\-न्कु\-रुते}}{\Afootnote{H; \skt{गलोपरोधं कुरुते}\ A; \skt{गलोपरोधां कुरुते}\ K N}}\skt{ क्रमेण । }\pend
\pstart\versepar \skt{स }\edtext{\skt{मान्सतानेति}}{\Afootnote{K N H; \skt{मांसतानः क}\ A}}\skt{ }\edtext{\skt{बिभर्त्ति}}{\Afootnote{H; \skt{थितोऽ}\ A; \skt{विनर्त्ति}\ K; \skt{विभर्त्ति}\ N}}\skt{ }\edtext{\skt{संज्ञा\-प्रा\-ण\-प्र\-णु\-त्स\-र्व्व\-कृतो}}{\Afootnote{H; \skt{वलम्बी प्राण\-प्र\-णु\-त्स\-र्व\-कृतो}\ A; \skt{संज्ञ प्राण\-प्र\-णु\-न्स\-र्व\-कृतो}\ K; \skt{सण प्रणुत्सर्वकृतो}\ N}}\skt{ विकारः ॥}\pend

\pstart\versepar \ledsidenote{\passref{2.16.63}}\skt{सदाहतोदं श्वयथुं }\edtext{\edtext{\skt{सता\-म्र\-म\-न्त\-र्ग्गले}}{\lemma{\skt{सताम्रम्}}\Afootnote{H; \skt{सरक्तम्}\ A; \skt{स ताश्र}\ K; \skt{सतादम्}\ N}}}{\lemma{\skt{अन्तर्ग्गले}}\Afootnote{A H; \skt{मन्तु जले}\ K; \skt{अन्तमले}\ N}}\skt{ }\edtext{\skt{पूति\-वि\-शी\-र्ण्ण\-मान्सं}}{\Afootnote{A K H; \skt{पूति\-वि\-शी\-र्ण्ण\-मास ।}\ N}}\skt{ । }\pend
\pstart\versepar \skt{पित्तेन विद्याद्वदने }\edtext{\edtext{\skt{विचारिम्पार्श्वे}}{\lemma{\skt{विचारिम्}}\Afootnote{H; \skt{विदारीं}\ A; \skt{विदारिम्}\ K; \skt{विचारि}\ N}}}{\lemma{\skt{पार्श्वे}}\Afootnote{A N H; \skt{पार्श्व}\ K}}\skt{ }\edtext{\skt{विशेषात्स तु}}{\Afootnote{A N H; \skt{विशेषान्स तु}\ K}}\skt{ येन शेते ॥}\pend

\pstart \ledsidenote{\passref{2.16.64}}\skt{सर्व्वसरास्तु }\edtext{\skt{वात\-पि\-त्त\-क\-फ\-नि\-मित्ताः}}{\Afootnote{K N H; \skt{वात\-पि\-त्त\-क\-फ\-शो\-णि\-त\-नि\-मित्ताः ॥}\ A}}\skt{ ॥}\pend

\pstart\versepar \ledsidenote{\passref{2.16.65}}\skt{स्फोटैः सतोदैर्व्वदनं सम\-न्ता\-द्य\-स्या\-चितं सर्व्वसरः सवातात् । }\pend
\pstart\versepar \skt{रक्तैः सदाहैः }\edtext{\skt{पिटकैः}}{\Afootnote{N H; \skt{तनुभिः}\ A; \skt{पिटाकैः}\ K}}\skt{ }\edtext{\skt{सपी\-तै\-र्य\-स्या\-चितं}}{\lemma{\skt{सपीतैर्}}\Afootnote{A N H; \skt{सपीतै}\ K}}\skt{ चापि सपित्तकोपात् ॥}\pend

\pstart \ledsidenote{\passref{2.16.66ab}}\edtext{\skt{अवेदनैः}}{\Afootnote{K N H; \emph{om.}\ A}}\skt{ }\edtext{\skt{कण्डुयुतैः}}{\Afootnote{N H; \skt{कण्डूयुतैर्}\ A; \skt{कण्डुयुतः}\ K}}\skt{ }\edtext{\edtext{\skt{सतो\-दै\-र्य\-स्या\-चितं}}{\lemma{\skt{सतोदैर्}}\Afootnote{H; \skt{अल्परुजैः सर्वणैर्}\ A; \skt{सतादेश}\ K; \skt{सतोदै}\ N}}}{\lemma{\skt{यस्याचितं}}\Afootnote{A N H; \skt{स्याचितं}\ K}}\skt{ चापि स वै }\edtext{\skt{कफेन}}{\Afootnote{A H; \skt{कफेनेति ॥}\ K N}}\skt{ ॥}\pend

\pstart \ledsidenote{\passref{2.16.trailer}}\edtext{\skt{॥ इति मुख\-रो\-ग\-नि\-दाने षोडशमः}}{\Afootnote{H; \skt{इति सुश्रु\-त\-सं\-हि\-तायां निदानस्थाने मुख\-रो\-ग\-नि\-दानं नाम षोडशोऽध्यायः ॥१६ ॥}\ A; \skt{"नि ॥}\ K; \skt{॥ निदाने ॥ ॥ {\symfont ❈} ॥}\ N}}\skt{ ॥}\pend

\pstart\versepar \ledsidenote{\passref{2.16.66.add1}}\edtext{\skt{ग्रन्थ्थ्य\-प\-च्य\-र्बु\-द\-ङ्गण्डुं}}{\Afootnote{H; \skt{ग्रन्थ्थ्य\-प\-च्य\-र्वु\-द\-ङ्गण्डुं}\ K; \skt{ग्रन्प\-च्य\-र्वु\-द\-ङ्गण्डं}\ N}}\skt{ }\edtext{\edtext{\skt{वृद्धिन्दंशञ्च}}{\lemma{\skt{वृद्धिन्}}\Afootnote{K H; \skt{वृद्विन्}\ N}}}{\lemma{\skt{दंशञ्च}}\Afootnote{K H; \skt{दञ्च}\ N}}\skt{ }\edtext{\skt{श्लीपदम्}}{\Afootnote{K H; \skt{श्लीप}\ N}}\skt{ । }\pend
\pstart\versepar \skt{}\edtext{\skt{क्षुद्रशूकं}}{\Afootnote{N H; \skt{क्षुद्रम्शुकं}\ K}}\skt{ }\edtext{\skt{सभग्नञ्च}}{\Afootnote{N H; \skt{सभग्रञ्च}\ K}}\skt{ मुखरोगेण }\edtext{\skt{षट्स्मृतः}}{\Afootnote{H; \skt{स्मृतः ॥}\ K N}}\skt{ ॥}\pend

\pstart\versepar \ledsidenote{\passref{2.16.66.add2}}\edtext{\skt{सूत्रे}}{\Afootnote{K H; \skt{मूत्रे}\ N}}\skt{ }\edtext{\edtext{\skt{प्रदि\-ष्ट\-न्द\-श\-षट्वा}}{\lemma{\skt{प्रदिष्टन्}}\Afootnote{N H; \skt{प्रदिष्ट}\ K}}}{\lemma{\skt{दशषट्वा}}\Afootnote{K H; \skt{दश षट्व}\ N}}\skt{ }\edtext{\skt{पूर्व्वं}}{\Afootnote{K H; \skt{त्पूर्वं}\ N}}\skt{ }\edtext{\edtext{\skt{स्थन\-न्नि\-दा\-न\-म्भ\-व\-तीति}}{\lemma{\skt{स्थनन्}}\Afootnote{K H; \skt{स्थानन्}\ N}}}{\lemma{\skt{भवतीति}}\Afootnote{K H; \skt{अम्भवती इ}\ N}}\skt{ }\edtext{\skt{यच्च}}{\Afootnote{K H; \skt{त्यच्च ।}\ N}}\skt{ । }\pend
\pstart\versepar \skt{सवि\-स्त\-र\-न्त\-त्क\-थि\-त\-म्म\-येह }\edtext{\skt{स्वयम्भुना}}{\Afootnote{H; \skt{स्वयम्सुना}\ K; \skt{स्वे यस्तु ना}\ N}}\skt{ }\edtext{\skt{यत्कथितं}}{\Afootnote{K H; \skt{यन्कथितं}\ N}}\skt{ }\edtext{\skt{हिताय}}{\Afootnote{K H; \skt{हिताय ॥ {\symfont ❈} ॥}\ N}}\skt{ ॥}\pend

\pstart \ledsidenote{\passref{2.trailer}}\edtext{\skt{समा\-प्त\-न्नि\-दा\-न\-स्था\-नम्}}{\lemma{\skt{निदानस्थानम्}}\Afootnote{H; \skt{निदानस्थानम् ॥ {\symfont ❈} ॥}\ K; \skt{निदानथान ॥ ० ॥}\ N}}\skt{ ॥}\pend

\endnumbering
